% Equilibria — Chemistry Upper Sixth — Cours signé OnBuch+
% Official MINESEC syllabus. Compiler avec : tectonic CHE01-equilibria.tex
\newcommand{\DOCMATIERE}{Chemistry}
\newcommand{\DOCNIVEAU}{Upper Sixth}
\newcommand{\DOCMODULE}{Upper Sixth — Chemistry}
\newcommand{\DOCLECON}{1}
\newcommand{\DOCTITRE}{Equilibria}
% OnBuch+ — préambule « cours signés » ENRICHI (compilé avec tectonic / XeLaTeX)
% Système visuel « Pop » d'OnBuch+ : fond crème (jamais blanc), bordures encre
% épaisses, ombres « dures » décalées, titres Archivo Black en capitales.
% Version MATHÉMATIQUES : graphes (pgfplots/tikz), boîtes pédagogiques
% (définition, propriété, méthode, attention, le savais-tu, exercice résolu,
% exercice, TP, bilan), page de garde et signature OnBuch+ ancrée en bas.
%
% Macros attendues dans main.tex AVANT \input{preamble} :
%   \DOCMATIERE  \DOCNIVEAU  \DOCMODULE  \DOCLECON (numéro)  \DOCTITRE
\documentclass[11pt]{article}

\usepackage{fontspec}
\usepackage[english]{babel}
\usepackage[a4paper,margin=2cm,top=3.4cm,bottom=2.8cm,headheight=1.4cm,footskip=1.3cm]{geometry}
\usepackage{amsmath,amssymb,amsfonts}
\usepackage{mathtools} % \xmapsto, \coloneqq... souvent utilisés par les modèles
\usepackage{cancel} % simplifications barrées (bilans de forces)
% \abs{z} (module, valeur absolue) et \norm{u} (norme), utilisés par les modèles
\providecommand{\abs}{}\let\abs\relax\DeclarePairedDelimiter{\abs}{\lvert}{\rvert}
\providecommand{\norm}{}\let\norm\relax\DeclarePairedDelimiter{\norm}{\lVert}{\rVert}
\usepackage[table]{xcolor} % [table] : fournit aussi \rowcolors (alternance de lignes)
\usepackage{pagecolor}
\usepackage{tikz}
\usetikzlibrary{arrows.meta,positioning,calc,shapes.geometric,shapes.misc,shapes.arrows,decorations.pathmorphing,decorations.pathreplacing,decorations.markings,patterns,shadows,mindmap,trees,fit,backgrounds,matrix,intersections,angles,quotes}
\usepackage{pgfplots}
\usepackage[version=4]{mhchem} % équations nucléaires \ce{^{235}_{92}U}
\usepackage[european,siunitx]{circuitikz} % schémas électriques
\pgfplotsset{compat=1.18}
\usepgfplotslibrary{fillbetween,groupplots} % aires entre courbes (calcul intégral)
\usepackage{enumitem}
\usepackage{fancyhdr}
% [most] charge toutes les bibliothèques tcolorbox (listings, minted, raster,
% documentation...) et gonfle inutilement la mémoire TeX sur un document long
% (~300 boîtes) : on ne charge que ce qui est réellement utilisé.
\usepackage[skins,breakable]{tcolorbox}
\usepackage{graphicx}
\usepackage{colortbl}
\usepackage{booktabs}
\usepackage{tabularx}
\usepackage{diagbox} % case d'en-tête diagonale des tableaux à double entrée
% Colonnes extensibles centrée / gauche / droite (C, L, R), souvent utilisées par les modèles
\newcolumntype{C}{>{\centering\arraybackslash}X}
\newcolumntype{L}{>{\raggedright\arraybackslash}X}
\newcolumntype{R}{>{\raggedleft\arraybackslash}X}
\usepackage{array}
\usepackage{multicol}
\usepackage{siunitx}
% parse-numbers=false : accepte aussi \SI{2,5 \times 10^{-3}}{...}
\sisetup{locale=UK,output-decimal-marker={.},group-minimum-digits=4,parse-numbers=false}
\DeclareSIUnit\ohm{\ensuremath{\symup{\Omega}}} % U+2126 absent de la police
\DeclareSIUnit\division{div} % réglages d'oscilloscope (ms/div, V/div)
\DeclareSIUnit\tr{tr} % tours (vitesse de rotation, ex. \SI{1500}{\tr\per\minute})
\usepackage{qrcode}
% \degree : commande fréquemment utilisée par les modèles pour un angle ou une
% température en dehors de \SI{}{} (non fournie par les paquets chargés).
\providecommand{\degree}{^{\circ}}
% \qed (fin de démonstration), utilisé par les modèles sans amsthm
\providecommand{\qed}{\hfill$\square$}
% \barre : double barre des valeurs interdites dans un tableau de variations (tkz-tab)
\providecommand{\barre}{\ensuremath{\|}}
% environnement proof (amsthm non chargé) utilisé par les modèles
\ifdefined\proof\else\newenvironment{proof}[1][Proof]{\par\noindent\textit{#1.}\ }{\qed\par}\fi
\usepackage{lastpage}
\usepackage{hyperref}
\hypersetup{hidelinks}

% ============================================================
% Palette « Pop » OnBuch+ — mêmes hex que la classe P (plus_ui.dart)
% ============================================================
\definecolor{poporange}{HTML}{F59321}
\definecolor{popdark}{HTML}{E07A0C}
\definecolor{popcream}{HTML}{FFF6E8}
\definecolor{popink}{HTML}{1C1714}
\definecolor{poppurple}{HTML}{7A5AE0}
\definecolor{popgreen}{HTML}{2E9E6A}
\definecolor{poppink}{HTML}{E0457B}
\definecolor{popblue}{HTML}{2F7DE1}
\definecolor{popgold}{HTML}{C98A12}
\definecolor{popmuted}{HTML}{8A7F76}
\definecolor{popline}{HTML}{EADFD2}
\definecolor{popgray}{HTML}{8A7F76}
\colorlet{popgrey}{popgray}
% Alias tolérants (le modèle nomme parfois les couleurs différemment)
\colorlet{popwhite}{white}
\colorlet{popblack}{popink}
\colorlet{popred}{poppink}
\colorlet{popyellow}{popgold}
% Teintes claires pour fonds de tableaux / zones
\colorlet{popblueL}{popblue!10!white}
\colorlet{poporangeL}{poporange!14!white}
\colorlet{popgreenL}{popgreen!12!white}
\colorlet{poppurpleL}{poppurple!12!white}
\colorlet{poppinkL}{poppink!12!white}
\colorlet{popgoldL}{popgold!14!white}

\pagecolor{popcream}

% ============================================================
% Typographie
% ============================================================
\defaultfontfeatures{Ligatures=TeX}
\setmainfont{PlusJakartaSans-Medium.ttf}[Path=../../../fonts/,BoldFont=PlusJakartaSans-Bold.ttf]
% Maths en sans-serif (Fira Math) avec chiffres et lettres droites de Plus
% Jakarta, pour que les formules se fondent dans le texte.
\usepackage[math-style=ISO,bold-style=ISO]{unicode-math}
\setmathfont{FiraMath-Regular.otf}[Path=../../../fonts/]
\setmathfont{PlusJakartaSans-Medium.ttf}[Path=../../../fonts/,range={up/num,up/{Latin,latin},\mathup}]
% (icomma retiré : décimales avec point en anglais)
% Caractères mathématiques Unicode absents des polices de texte (Plus Jakarta,
% Archivo Black) : rendus par leur équivalent mathématique, en texte comme en
% formule — sinon ils s'affichent en carré vide (ex. « Arithmétique dans ℤ »).
\usepackage{newunicodechar}
\newunicodechar{ℝ}{\ensuremath{\mathbb{R}}}
\newunicodechar{ℤ}{\ensuremath{\mathbb{Z}}}
\newunicodechar{ℂ}{\ensuremath{\mathbb{C}}}
\newunicodechar{ℕ}{\ensuremath{\mathbb{N}}}
\newunicodechar{ℚ}{\ensuremath{\mathbb{Q}}}
\newunicodechar{𝔻}{\ensuremath{\mathbb{D}}}
\newunicodechar{α}{\ensuremath{\alpha}}
\newunicodechar{β}{\ensuremath{\beta}}
\newunicodechar{γ}{\ensuremath{\gamma}}
\newunicodechar{δ}{\ensuremath{\delta}}
\newunicodechar{ε}{\ensuremath{\varepsilon}}
\newunicodechar{θ}{\ensuremath{\theta}}
\newunicodechar{λ}{\ensuremath{\lambda}}
\newunicodechar{μ}{\ensuremath{\mu}}
\newunicodechar{σ}{\ensuremath{\sigma}}
\newunicodechar{τ}{\ensuremath{\tau}}
\newunicodechar{φ}{\ensuremath{\varphi}}
\newunicodechar{ψ}{\ensuremath{\psi}}
\newunicodechar{ω}{\ensuremath{\omega}}
\newunicodechar{Σ}{\ensuremath{\Sigma}}
% majuscules produites par \MakeUppercase dans les titres (α-aminés -> Α)
\newunicodechar{Α}{\ensuremath{\alpha}}
\newunicodechar{Β}{\ensuremath{\beta}}
\newunicodechar{Γ}{\ensuremath{\Gamma}}
\newunicodechar{Δ}{\ensuremath{\Delta}}
\newunicodechar{Ε}{\ensuremath{\varepsilon}}
\newunicodechar{Θ}{\ensuremath{\Theta}}
\newunicodechar{Λ}{\ensuremath{\Lambda}}
\newunicodechar{Μ}{\ensuremath{\mu}}
\newunicodechar{Τ}{\ensuremath{\tau}}
\newunicodechar{Φ}{\ensuremath{\Phi}}
\newunicodechar{Ψ}{\ensuremath{\Psi}}
\newunicodechar{Ω}{\ensuremath{\Omega}}
\newunicodechar{↦}{\ensuremath{\mapsto}}
\newunicodechar{∈}{\ensuremath{\in}}
\newunicodechar{∉}{\ensuremath{\notin}}
\newunicodechar{∧}{\ensuremath{\wedge}}
\newunicodechar{⇔}{\ensuremath{\Leftrightarrow}}
\newunicodechar{⟺}{\ensuremath{\Longleftrightarrow}}
\newunicodechar{⇒}{\ensuremath{\Rightarrow}}
\newunicodechar{⟹}{\ensuremath{\Longrightarrow}}
\newunicodechar{∘}{\ensuremath{\circ}}
\newunicodechar{∪}{\ensuremath{\cup}}
\newunicodechar{∩}{\ensuremath{\cap}}
\newunicodechar{≡}{\ensuremath{\equiv}}
\newunicodechar{⊂}{\ensuremath{\subset}}
\newunicodechar{⊥}{\ensuremath{\perp}}
\newunicodechar{∃}{\ensuremath{\exists}}
\newunicodechar{∀}{\ensuremath{\forall}}
\newunicodechar{∛}{\ensuremath{\sqrt[3]{\,}}}
\newunicodechar{∜}{\ensuremath{\sqrt[4]{\,}}}
\newunicodechar{⃗}{\ensuremath{{}^{\to}}}
\newunicodechar{ⁿ}{\ifmmode^{n}\else\textsuperscript{\ensuremath{n}}\fi}
\newunicodechar{ˣ}{\ifmmode^{x}\else\textsuperscript{\ensuremath{x}}\fi}
\newunicodechar{ᵇ}{\ifmmode^{b}\else\textsuperscript{\ensuremath{b}}\fi}
\newunicodechar{ʳ}{\ifmmode^{r}\else\textsuperscript{\ensuremath{r}}\fi}
\newunicodechar{ᵘ}{\ifmmode^{u}\else\textsuperscript{\ensuremath{u}}\fi}
\newunicodechar{ʸ}{\ifmmode^{y}\else\textsuperscript{\ensuremath{y}}\fi}
\newunicodechar{ᵃ}{\ifmmode^{a}\else\textsuperscript{\ensuremath{a}}\fi}
\newunicodechar{ᵉ}{\ifmmode^{e}\else\textsuperscript{\ensuremath{e}}\fi}
\newunicodechar{ᵏ}{\ifmmode^{k}\else\textsuperscript{\ensuremath{k}}\fi}
\newunicodechar{ᵖ}{\ifmmode^{p}\else\textsuperscript{\ensuremath{p}}\fi}
\newunicodechar{ᴺ}{\ifmmode^{N}\else\textsuperscript{\ensuremath{N}}\fi}
\newunicodechar{ᶜ}{\ifmmode^{c}\else\textsuperscript{\ensuremath{c}}\fi}
\newunicodechar{ʰ}{\ifmmode^{h}\else\textsuperscript{\ensuremath{h}}\fi}
\newunicodechar{ᵠ}{\ifmmode^{q}\else\textsuperscript{\ensuremath{q}}\fi}
\newunicodechar{ₙ}{\ifmmode_{n}\else\textsubscript{\ensuremath{n}}\fi}
\newunicodechar{ₐ}{\ifmmode_{a}\else\textsubscript{\ensuremath{a}}\fi}
\newunicodechar{ᵢ}{\ifmmode_{i}\else\textsubscript{\ensuremath{i}}\fi}
\newunicodechar{ᵣ}{\ifmmode_{r}\else\textsubscript{\ensuremath{r}}\fi}
\newunicodechar{ₖ}{\ifmmode_{k}\else\textsubscript{\ensuremath{k}}\fi}
\newunicodechar{ᵦ}{\ifmmode_{\beta}\else\textsubscript{\ensuremath{\beta}}\fi}
\newunicodechar{ₓ}{\ifmmode_{x}\else\textsubscript{\ensuremath{x}}\fi}
\newfontfamily\popdisplay{ArchivoBlack-Regular.ttf}[Path=../../../fonts/]
\newfontfamily\poplogo{SpaceGrotesk-Bold.ttf}[Path=../../../fonts/]
\newfontfamily\popsemi{PlusJakartaSans-SemiBold.ttf}[Path=../../../fonts/]
\newfontfamily\popxbold{PlusJakartaSans-ExtraBold.ttf}[Path=../../../fonts/]
\newfontfamily\popxboldls{PlusJakartaSans-ExtraBold.ttf}[Path=../../../fonts/,LetterSpace=3.0]

\setlist[enumerate]{leftmargin=1.7em,itemsep=5pt,topsep=4pt}
\setlist[itemize]{leftmargin=1.7em,itemsep=4pt,topsep=4pt,label={\color{poporange}\textbullet}}
\setlength{\parindent}{0pt}
\setlength{\parskip}{5pt}
\renewcommand{\arraystretch}{1.25}

% pgfplots : style Pop par défaut
\pgfplotsset{
  every axis/.append style={axis line style={popink,line width=1pt},tick style={popink},
    grid style={popline,line width=0.5pt},label style={font=\small},tick label style={font=\footnotesize}},
  popcurve/.style={poporange,line width=1.8pt,no markers,smooth},
}
% Même style disponible dans un \draw TikZ simple (hors pgfplots)
\tikzset{popcurve/.style={poporange,line width=1.8pt}}
\tikzset{popgrid/.style={popline,very thin}}
\colorlet{popcurve}{poporange} % le nom du style est parfois utilisé comme couleur

% ============================================================
% Pill / squiggle
% ============================================================
\newcommand{\poppill}[3]{%
  \tikz[baseline=-0.55ex]\node[draw=popink,line width=1.2pt,rounded corners=2.6pt,%
    fill=#2,inner xsep=6.5pt,inner ysep=3pt]{%
    \popxboldls\fontsize{7}{7}\selectfont\color{#3}\MakeUppercase{#1}};%
}
\newcommand{\popsquiggle}[1]{%
  \tikz[baseline=-0.4ex]{
    \draw[#1,line width=2.6pt,line cap=round]
      plot [smooth] coordinates {(0,0.09) (0.34,0.01) (0.68,0.21) (1.02,0.01) (1.36,0.10)};
  }%
}
% Mot-clé surligné (marqueur orange sous le texte)
\newcommand{\cle}[1]{{\popxbold\color{popdark}#1}}

% ============================================================
% En-tête / pied
% ============================================================
\pagestyle{fancy}
\fancyhf{}
\renewcommand{\headrulewidth}{0pt}
\renewcommand{\footrulewidth}{0pt}
\lhead{\raisebox{-0.15ex}{\poplogo\fontsize{13}{13}\selectfont\color{popink}OnBuch\color{poporange}+}}
\rhead{\poppill{\DOCMATIERE\ \textbullet\ \DOCNIVEAU\ \textbullet\ Chapter \DOCLECON}{white}{popink}}
\lfoot{\popxbold\fontsize{7.6}{7.6}\selectfont\color{popmuted}\MakeUppercase{Signed course OnBuch+}\ \textbullet\ {\popsemi\DOCTITRE}}
\rfoot{\tikz[baseline=-0.6ex]\node[rounded corners=8.5pt,fill=popink,minimum height=17pt,minimum width=17pt,inner xsep=5pt,inner sep=0]{\popxbold\fontsize{8}{8}\selectfont\color{popcream}\thepage\,/\,\pageref*{LastPage}};}

% ============================================================
% Titres
% ============================================================
\newcommand{\chaptitre}[2]{%
  \begin{tcolorbox}[enhanced,colback=poporange,colframe=popink,boxrule=2.6pt,arc=11pt,
      drop shadow=popink,left=15pt,right=15pt,top=12pt,bottom=11pt,before skip=3pt,after skip=18pt]
    \poppill{#2}{popink}{popcream}\par\vspace{9pt}
    {\raggedright\hyphenpenalty=10000\exhyphenpenalty=10000\popdisplay\fontsize{22}{24}\selectfont\color{popcream}\MakeUppercase{#1}\par}
  \end{tcolorbox}
}
% Section numérotée : pastille orange avec le numéro + titre Archivo
\newcounter{popsec}
\newcommand{\coursec}[1]{%
  \stepcounter{popsec}\setcounter{popsub}{0}%
  \par\vspace{16pt}\noindent
  \begin{minipage}{\linewidth}
  \tikz[baseline=-0.7ex]\node[draw=popink,line width=1.6pt,fill=poporange,rounded corners=4pt,minimum size=22pt,inner sep=0]{\popdisplay\fontsize{12}{12}\selectfont\color{popcream}\thepopsec};%
  \hspace{8pt}{\popdisplay\fontsize{15}{17}\selectfont\color{popink}\MakeUppercase{#1}}\par
  \vspace{4pt}\noindent{\color{poporange}\rule{36pt}{3.6pt}}
  \end{minipage}\par\nopagebreak\vspace{8pt}
}
\newcounter{popsub}
\newcommand{\courssub}[1]{%
  \stepcounter{popsub}%
  \par\vspace{9pt}\noindent{\popxbold\fontsize{11.5}{13}\selectfont\color{popink}{\color{poporange}\thepopsec.\thepopsub}\hspace{6pt}#1}\par\nopagebreak\vspace{4pt}
}

% ============================================================
% Boîtes pédagogiques — même recette : fond blanc, bordure épaisse colorée,
% ombre dure encre, badge plein. Titre optionnel [#1].
% ============================================================
\tcbset{popbase/.style={breakable,enhanced,colback=white,boxrule=2.3pt,arc=9pt,
  shadow={3pt}{-3pt}{0pt}{popink},left=14pt,right=14pt,top=11pt,bottom=11pt,before skip=11pt,after skip=15pt,
  fontupper=\popsemi\fontsize{10.6}{14.5}\selectfont\color{popink},
  fontlower=\popsemi\fontsize{10.6}{14.5}\selectfont\color{popink}}}
% Coche dessinée (le glyphe ✓ n'existe pas dans les polices du document)
\newcommand{\popcheck}[1]{\tikz[baseline=-0.5ex]\draw[#1,line width=1.6pt,line cap=round,line join=round](-0.13,0.0)--(-0.04,-0.09)--(0.13,0.1);}
\providecommand{\coche}{\popcheck{popgreen}}
\providecommand{\resultat}[1]{\par\smallskip\noindent\fcolorbox{popgreen}{popgreenL}{\parbox{\dimexpr\linewidth-2\fboxsep-2\fboxrule\relax}{\textbf{Result.} #1}}\par\smallskip}
\newcommand{\popboxhead}[3]{\poppill{#1}{#2}{white}\ifx\relax#3\relax\else\hspace{6pt}{\popxbold\fontsize{10.6}{12}\selectfont\color{popink}#3}\fi\par\vspace{7pt}}

\newtcolorbox{aretenir}[1][]{popbase,colframe=popgreen,before upper={\popboxhead{Key points}{popgreen}{#1}}}
\newtcolorbox{exemplebox}[1][]{popbase,colframe=poporange,before upper={\popboxhead{Exemple}{poporange}{#1}}}
\newtcolorbox{definition}[1][]{popbase,colframe=popblue,before upper={\popboxhead{Definition}{popblue}{#1}}}
\newtcolorbox{propriete}[1][]{popbase,colframe=poppurple,before upper={\popboxhead{Property}{poppurple}{#1}}}
\newtcolorbox{methode}[1][]{popbase,colframe=popgold,before upper={\popboxhead{Method}{popgold}{#1}}}
\newtcolorbox{attention}[1][]{popbase,colframe=poppink,before upper={\popboxhead{Warning}{poppink}{#1}}}
\newtcolorbox{experience}[1][]{popbase,colframe=popblue,colback=popblueL,before upper={\popboxhead{Experiment}{popblue}{#1}}}
% « Le savais-tu ? » : carte violette pleine (contexte, culture, Cameroun)
\newtcolorbox{savaistu}[1][]{popbase,colback=poppurple,colframe=popink,
  fontupper=\popsemi\fontsize{10.4}{14.2}\selectfont\color{white},
  before upper={\poppill{Did you know?}{white}{poppurple}\ifx\relax#1\relax\else\hspace{6pt}{\popxbold\color{white}#1}\fi\par\vspace{7pt}}}
% Exercice résolu : énoncé en haut, \tcblower puis la solution sur fond crème
\newcounter{popexo}\newcounter{popexr}
\newtcolorbox{exoresolu}[1][]{popbase,colframe=poporange,colbacklower=popcream,
  segmentation style={popink,line width=1.2pt,dashed},
  before upper={\stepcounter{popexr}\popboxhead{Worked example \thepopexr}{poporange}{#1}},
  before lower={\poppill{Solution}{popink}{popcream}\par\vspace{6pt}}}
% Exercice (non résolu en place) — la correction est dans la partie Corrigés
\newtcolorbox{exercice}[1][]{popbase,colframe=popink,
  before upper={\stepcounter{popexo}\popboxhead{Exercise \thepopexo}{popink}{#1}}}
% Corrigé
\newtcolorbox{corrige}[1][]{popbase,colframe=popgreen,colback=popgreenL,
  before upper={\popboxhead{Solution}{popgreen}{#1}}}
% Objectifs / prérequis / bilan
\newtcolorbox{objectifs}{popbase,colframe=popink,colback=poporangeL,
  before upper={\popboxhead{Chapter objectives}{popink}{}}}
\newtcolorbox{prerequis}{popbase,colframe=popink,colback=popblueL,
  before upper={\popboxhead{Prerequisites}{popblue}{}}}
\newtcolorbox{bilan}{popbase,colframe=popink,colback=white,boxrule=2.8pt,
  before upper={\popboxhead{Fiche bilan}{poporange}{L'essentiel à connaître pour l'examen}}}
% Figure encadrée avec légende
\newtcolorbox{popfigure}[1][]{enhanced,breakable=false,colback=white,colframe=popline,boxrule=1.4pt,arc=8pt,
  left=10pt,right=10pt,top=10pt,bottom=8pt,before skip=10pt,after skip=12pt,halign=center,
  lower separated=false,halign lower=center,
  fontlower=\popsemi\fontsize{9}{11.5}\selectfont\color{popmuted},
  #1}
\newcommand{\legende}[1]{\par\vspace{6pt}{\popsemi\fontsize{9}{11.5}\selectfont\color{popmuted}#1}}

% ============================================================
% Signature OnBuch+
% ============================================================
\newcommand{\poplogoinline}[1]{{\poplogo\fontsize{#1}{#1}\selectfont\color{popink}OnBuch\color{poporange}+}}
\newcommand{\signatureonbuch}{%
  \begin{tcolorbox}[enhanced,colback=popink,colframe=popink,boxrule=2.2pt,arc=10pt,
      drop shadow=poporange,left=14pt,right=14pt,top=10pt,bottom=10pt,before skip=0pt,after skip=0pt]
    \tikz[baseline=-0.7ex]\node[circle,fill=popgreen,draw=popcream,line width=1.3pt,minimum size=22pt,inner sep=0]{\popcheck{white}};%
    \hspace{9pt}\begin{minipage}[c]{0.62\linewidth}
      {\popxbold\fontsize{12}{13}\selectfont\color{popcream}Signed\ \textbullet\ {\poplogo OnBuch\color{poporange}+}}\par\vspace{2pt}
      {\raggedright\popsemi\fontsize{8}{10.5}\selectfont\color{popline}\MakeUppercase{OnBuch+ teaching team}\\\MakeUppercase{Follows the official MINESEC syllabus}\par}
    \end{minipage}\hfill
    {\poplogo\fontsize{22}{22}\selectfont\color{popcream}OnBuch\color{poporange}+}
  \end{tcolorbox}%
}

% ============================================================
% Page de garde
% #1 titre · #2 matière · #3 niveau · #4 module · #5 n° leçon · #6 durée ·
% #7 description / objectifs (texte ou itemize)
% ============================================================
\newcommand{\pagegarde}[7]{%
  \thispagestyle{empty}
  \newgeometry{a4paper,margin=1.7cm,top=1.6cm,bottom=1.4cm}%
  \begin{tikzpicture}[remember picture,overlay]
    \fill[poporange!18] ([xshift=-3.2cm,yshift=-3.6cm]current page.north east) circle (4.6cm);
    \fill[poppurple!14] ([xshift=2.4cm,yshift=7.6cm]current page.south west) circle (3.8cm);
  \end{tikzpicture}%
  {\poplogo\fontsize{30}{30}\selectfont\color{popink}OnBuch\color{poporange}+}\hfill
  \raisebox{0.6ex}{\poppill{Signed course \textbullet\ Premium}{popink}{popcream}}\par
  \vspace{1pt}\popsquiggle{poppurple}\par
  \vspace{8pt}
  \begin{tcolorbox}[enhanced,colback=poporange,colframe=popink,boxrule=2.8pt,arc=13pt,
      drop shadow=popink,left=18pt,right=18pt,top=12pt,bottom=13pt,after skip=10pt]
    \poppill{#2 \textbullet\ #3}{popink}{popcream}\hspace{5pt}\poppill{#4}{white}{popink}\par\vspace{8pt}
    {\popdisplay\fontsize{14}{14}\selectfont\color{popink}CHAPTER #5}\par\vspace{4pt}
    {\raggedright\hyphenpenalty=10000\exhyphenpenalty=10000\popdisplay\fontsize{25}{27}\selectfont\color{popcream}\MakeUppercase{#1}\par}
    \vspace{8pt}
    {\popxbold\fontsize{9.5}{9.5}\selectfont\color{popink}\MakeUppercase{Suggested time: #6}}
  \end{tcolorbox}
  \begin{tcolorbox}[enhanced,colback=white,colframe=popink,boxrule=2.2pt,arc=10pt,
      drop shadow=popink,left=16pt,right=16pt,top=10pt,bottom=10pt,after skip=8pt]
    \poppill{In this chapter}{poppurple}{white}\par\vspace{5pt}
    \popsemi\fontsize{9.3}{12.6}\selectfont\color{popink}#7
  \end{tcolorbox}
  \noindent\begin{minipage}[t]{0.487\linewidth}
    \begin{tcolorbox}[enhanced,colback=popgreen,colframe=popink,boxrule=2.2pt,arc=10pt,
        drop shadow=popink,left=11pt,right=11pt,top=10pt,bottom=10pt,height=3.1cm,valign=center]
      \tcbox[colback=white,colframe=popink,boxrule=1.3pt,arc=5pt,boxsep=0pt,nobeforeafter,
        left=3pt,right=3pt,top=3pt,bottom=3pt,valign=center]{\qrcode[height=1.9cm]{https://play.google.com/store/apps/details?id=com.luvvix.onbuch&hl=en}}%
      \hspace{8pt}\begin{minipage}[c]{0.5\linewidth}
        {\popdisplay\fontsize{11}{12.5}\selectfont\color{white}\MakeUppercase{Download OnBuch}}\par\vspace{4pt}
        {\raggedright\popsemi\fontsize{8.4}{11}\selectfont\color{white}Free \textbullet\ Google Play\\AI tutor, past papers, results\par}
      \end{minipage}
    \end{tcolorbox}
  \end{minipage}\hfill
  \begin{minipage}[t]{0.487\linewidth}
    \begin{tcolorbox}[enhanced,colback=poppurple,colframe=popink,boxrule=2.2pt,arc=10pt,
        drop shadow=popink,left=11pt,right=11pt,top=10pt,bottom=10pt,height=3.1cm,valign=center]
      \tikz[baseline=-0.7ex]\node[circle,fill=white,draw=popink,line width=1.3pt,minimum size=1.6cm,inner sep=0]{\popdisplay\fontsize{24}{24}\selectfont\color{poppurple}+};%
      \hspace{8pt}\begin{minipage}[c]{0.6\linewidth}
        {\popdisplay\fontsize{11}{12.5}\selectfont\color{white}\MakeUppercase{Go OnBuch+}}\par\vspace{4pt}
        {\raggedright\popsemi\fontsize{8.4}{11}\selectfont\color{white}All signed courses, unlimited AI tutor, PDF sheets — OnBuch+ tab in the app\par}
      \end{minipage}
    \end{tcolorbox}
  \end{minipage}
  \vspace{8pt}
  \popcontenu
  \vfill
  \signatureonbuch
  \restoregeometry
  \newpage
}

% Rangée « Ce cours contient » (page de garde)
\newcommand{\poptile}[3]{% #1 couleur #2 gros texte #3 légende
  \begin{tcolorbox}[enhanced,colback=#1,colframe=popink,boxrule=2pt,arc=9pt,shadow={2.5pt}{-2.5pt}{0pt}{popink},
      left=6pt,right=6pt,top=9pt,bottom=9pt,halign=center,valign=center,height=1.75cm,width=\linewidth,nobeforeafter]
    {\popdisplay\fontsize{12.5}{13}\selectfont\color{white}\MakeUppercase{#2}}\par\vspace{3pt}
    {\centering\popsemi\fontsize{7.8}{9.5}\selectfont\color{white}#3\par}
  \end{tcolorbox}}
\newcommand{\popcontenu}{%
  \par\noindent{\popxboldls\fontsize{7.5}{7.5}\selectfont\color{popmuted}THIS SIGNED COURSE CONTAINS}\par\vspace{7pt}
  \noindent\begin{minipage}[t]{0.235\linewidth}\poptile{popblue}{Course}{complete, illustrated, step by step}\end{minipage}\hfill
  \begin{minipage}[t]{0.235\linewidth}\poptile{poporange}{Methods}{and worked examples}\end{minipage}\hfill
  \begin{minipage}[t]{0.235\linewidth}\poptile{poppink}{Exercises}{exam style + solutions}\end{minipage}\hfill
  \begin{minipage}[t]{0.235\linewidth}\poptile{popgreen}{Summary sheet}{the essentials to remember}\end{minipage}\par}

% Sommaire
\newtcolorbox{sommaire}{enhanced,colback=white,colframe=popink,boxrule=2.2pt,arc=10pt,
  drop shadow=popink,left=18pt,right=18pt,top=13pt,bottom=13pt,before skip=6pt,after skip=20pt,
  before upper={\poppill{Contents}{poppurple}{white}\par\vspace{9pt}\popsemi\fontsize{10.6}{14.5}\selectfont\color{popink}}}

% Signature de fin de cours — poussée tout en bas de la dernière page
\newcommand{\findecours}{%
  \par\vfill\nopagebreak
  \begin{center}\popsquiggle{poporange}\end{center}\vspace{4pt}
  \signatureonbuch
}
\AtBeginDocument{\def\llbracket{\mathopen{[\mkern-3.5mu[}}\def\rrbracket{\mathclose{]\mkern-3.5mu]}}} % intervalles d'entiers (glyphe ⟦ absent de la police)
% Glyphes absents de Fira Math : redéfinis à partir de glyphes disponibles
\AtBeginDocument{%
  \def\setminus{\mathbin{\backslash}}\let\smallsetminus\setminus
  \def\vdots{\mathord{\vbox{\baselineskip3.2pt\lineskiplimit0pt\kern4pt\hbox{.}\hbox{.}\hbox{.}}}}%
  \def\ddots{\mathinner{\mkern1mu\raise6.5pt\hbox{.}\mkern2mu\raise3.6pt\hbox{.}\mkern2mu\raise0.7pt\hbox{.}\mkern1mu}}%
  \def\triangle{\mathord{\scalebox{0.9}{$\symup{\Delta}$}}}%
  \def\nabla{\mathord{\raisebox{\depth}{\scalebox{1}[-1]{$\symup{\Delta}$}}}}%
  \def\checkmark{\popcheck{popdark}}%
  \def\star{\mathbin{\ast}}\def\bigstar{\mathord{\ast}}%
  \def\diamond{\mathbin{\scalebox{0.6}{$\Diamond$}}}%
  \def\bigcup{\mathop{\mathchoice{\raisebox{-0.3em}{\scalebox{1.7}{$\cup$}}}{\scalebox{1.25}{$\cup$}}{\cup}{\cup}}\displaylimits}%
  \def\bigcap{\mathop{\mathchoice{\raisebox{-0.3em}{\scalebox{1.7}{$\cap$}}}{\scalebox{1.25}{$\cap$}}{\cap}{\cap}}\displaylimits}%
  \def\lesseqqgtr{\mathrel{\lessgtr}}\def\bigoplus{\mathop{\oplus}\displaylimits}%
}
\providecommand{\ssi}{if and only if} % abréviation fréquente
\AtBeginDocument{\providecommand{\wideparen}{\overparen}} % arc de cercle

%% FIGFIX-BEGIN v5 — correctifs globaux des figures (docs/onbuch-generation/tools/figfix)
\usepackage{adjustbox}\usepackage{etoolbox}\tcbuselibrary{raster}
% toute figure TikZ/pgfplots plus large que la ligne est réduite à la largeur disponible
\BeforeBeginEnvironment{tikzpicture}{\begin{adjustbox}{max width=\linewidth}}
\AfterEndEnvironment{tikzpicture}{\end{adjustbox}}
% légendes pgfplots lisibles par-dessus les courbes
\pgfplotsset{every axis/.append style={legend style={fill=white,fill opacity=0.92,text opacity=1,draw=black!15,inner sep=3pt}}}
% étiquettes d'arêtes (syntaxe edge["..."]) avec fond blanc
\tikzset{every edge quotes/.append style={fill=white,inner sep=1.2pt}}
%% FIGFIX-END
\begin{document}

\pagegarde{Equilibria}{Chemistry}{Upper Sixth}{Upper Sixth — Chemistry}{1}{33 periods}{This chapter studies chemical equilibria in all their forms: reversible reactions and the equilibrium constant, redox equilibria and electrode potentials, and acid-base equilibria with pH calculations, titrations, buffers and solubility. It equips the learner with the quantitative tools required to predict, control and exploit chemical systems at the GCE Advanced Level.
\vspace{4pt}
\begin{multicols}{2}\small
\begin{itemize}
\item By the end of this chapter I can explain the dynamic nature of chemical equilibrium and state the conditions for its attainment.
\item By the end of this chapter I can apply Le Chatelier's principle to predict the effect of concentration, pressure, temperature and a catalyst on an equilibrium position.
\item By the end of this chapter I can write expressions for Kc and Kp, deduce their units and perform related calculations, including for heterogeneous equilibria.
\item By the end of this chapter I can relate the equilibrium constant to temperature and enthalpy change and interpret its magnitude.
\item By the end of this chapter I can define redox systems, balance redox equations and describe the different types of half-cells.
\item By the end of this chapter I can measure and use standard electrode potentials to calculate cell emf and predict the feasibility of redox reactions.
\end{itemize}\end{multicols}}

\begin{sommaire}
\begin{multicols}{2}
\begin{enumerate}
\item Reversibility and Dynamic Equilibrium
\item Equilibrium Position and Le Chatelier's Principle
\item The Equilibrium Law: Kc and Kp
\item Heterogeneous Equilibria and the Effect of Temperature on K
\item Redox Systems and Balancing Redox Equations
\item Types of Half-Cells and Standard Electrode Potentials
\item Electrochemical Cells, Cell emf and Uses of E\^{}θ
\item Acid-Base Theories, pH and the Ionic Product of Water
\item pH of Strong and Weak Acids and Bases, and of Mixtures
\item Acid-Base Indicators and Titration Curves
\item Buffer Solutions
\item Salt Hydrolysis, Solubility Product and Common Ion Effect
\item Integration activity
\item Methods and worked examples
\item Exercises
\item Solutions to the exercises
\item Summary sheet
\end{enumerate}
\end{multicols}
\end{sommaire}

\begin{objectifs}
\begin{itemize}
\item By the end of this chapter I can explain the dynamic nature of chemical equilibrium and state the conditions for its attainment.
\item By the end of this chapter I can apply Le Chatelier's principle to predict the effect of concentration, pressure, temperature and a catalyst on an equilibrium position.
\item By the end of this chapter I can write expressions for Kc and Kp, deduce their units and perform related calculations, including for heterogeneous equilibria.
\item By the end of this chapter I can relate the equilibrium constant to temperature and enthalpy change and interpret its magnitude.
\item By the end of this chapter I can define redox systems, balance redox equations and describe the different types of half-cells.
\item By the end of this chapter I can measure and use standard electrode potentials to calculate cell emf and predict the feasibility of redox reactions.
\item By the end of this chapter I can define acids and bases according to Arrhenius, Brønsted-Lowry and Lewis and calculate the pH of strong and weak acid/base solutions and their mixtures.
\item By the end of this chapter I can interpret titration curves, select suitable indicators and explain the functioning of buffer solutions with pH calculations.
\item By the end of this chapter I can explain salt hydrolysis and calculate solubility products, including the common ion effect.
\end{itemize}
\end{objectifs}

\begin{prerequis}
\begin{itemize}
\item Chemical equations, stoichiometry and the mole concept (Lower Sixth)
\item Amount of substance, molar concentration and gas laws (Lower Sixth)
\item Enthalpy changes of reactions and Hess's law (Lower Sixth)
\item Oxidation numbers and simple redox reactions (Lower Sixth)
\item Rates of reaction and collision theory (Lower Sixth)
\end{itemize}
\end{prerequis}

\newpage

\coursec{Reversibility and Dynamic Equilibrium}

In a Douala brewery, the production of beer relies on fermentation reactions that never go to completion — the chemist must constantly push the equilibrium towards the desired product. In the same city, a car battery delivers current because of redox reactions occurring at its electrodes, and a nurse at the Laquintinie Hospital checks the pH of a patient's blood, which is kept almost constant by buffer systems. Why do some reactions stop before all reactants are consumed? How can we predict and control the yield of a reaction, the voltage of a cell, or the pH of a solution? This chapter provides the quantitative tools to answer these questions.

Until now, most of the reactions you have met — burning magnesium, neutralising an acid with a base — were treated as if they went to completion: the reactants are entirely converted into products. In reality, a very large number of chemical reactions are \cle{incomplete}: no matter how long we wait, some reactant always remains. This is not a failure of the chemist; it is a fundamental property of matter. Such reactions are said to be \cle{reversible}, and the state they finally reach is called a state of \cle{chemical equilibrium}. Understanding equilibrium is the key to controlling industrial processes (the Haber synthesis of ammonia), to explaining how batteries work, and to calculating the pH of blood, of palm wine, or of the water of the Wouri river.

\courssub{The Concept of Reversibility}

\textbf{Observation 1: water in a closed vessel.}\par
Pour some water into a bottle, close it tightly, and leave it on the laboratory bench. The water level does not change, even after weeks. Yet water is volatile: molecules are constantly escaping from the liquid surface (evaporation). Why does the level stay constant? Because the vessel is \emph{closed}: the water vapour cannot escape, and vapour molecules continually return to the liquid (condensation). After some time, evaporation and condensation occur at exactly the same rate, and the amount of liquid water remains constant. We write:
\[
\ce{H2O(l) <=> H2O(g)}
\]

\textbf{Observation 2: a saturated salt solution.}\par
Add sodium chloride to water and stir until no more dissolves: some solid remains at the bottom. The solution is \emph{saturated}. If we could watch the ions, we would see that \ce{Na+} and \ce{Cl-} ions are still leaving the crystal (dissolution) while other ions from the solution are redepositing onto it (crystallisation), at equal rates:
\[
\ce{NaCl(s) <=> Na+(aq) + Cl-(aq)}
\]
The mass of undissolved solid stays constant, but the solid is \emph{not} the same collection of ions from one moment to the next.

\textbf{Observation 3: the Haber process.}\par
Industrially, ammonia is synthesised from nitrogen and hydrogen:
\[
\ce{N2(g) + 3H2(g) <=> 2NH3(g)} \qquad \Delta H^{\theta} = -92\ \mathrm{kJ\,mol^{-1}}
\]
Even under high pressure and moderate temperature, the conversion is never complete: the mixture leaving the reactor always contains unreacted \ce{N2} and \ce{H2}, which are recycled. The reaction is reversible — ammonia also decomposes back into nitrogen and hydrogen.

\textbf{Observation 4: esterification.}\par
Mix ethanoic acid with ethanol in the presence of a little concentrated sulphuric acid and warm. The fruity-smelling ester ethyl ethanoate forms, but the reaction stops when roughly two-thirds of the acid has reacted, because the ester simultaneously reacts with water to regenerate the acid and the alcohol:
\[
\ce{CH3COOH(l) + C2H5OH(l) <=> CH3COOC2H5(l) + H2O(l)}
\]

\begin{definition}[Reversible reaction]
A \cle{reversible reaction} is a chemical reaction that can proceed in \textbf{both directions}: the products formed can react together to regenerate the original reactants. It is represented by a \cle{double arrow} $\rightleftharpoons$ (or \ce{<=>}) between reactants and products. The reaction going from reactants to products is the \cle{forward reaction}; the reaction going from products back to reactants is the \cle{reverse reaction}.
\end{definition}

\begin{attention}[The double arrow is not decoration]
Writing \ce{A + B <=> C + D} is a statement: it claims that \emph{both} \ce{A + B -> C + D} and \ce{C + D -> A + B} occur. Never use the double arrow for a reaction that goes essentially to completion (e.g.\ the combustion of methane), and never use a single arrow for a genuinely reversible reaction such as esterification. At the GCE, marks are lost for the wrong arrow.
\end{attention}

\courssub{Static Equilibrium versus Dynamic Equilibrium}

A book lying on a table is in equilibrium: the downward pull of gravity is exactly balanced by the upward reaction of the table, so the book does not move. Nothing is happening — this is a \cle{static equilibrium}. Many students wrongly imagine chemical equilibrium in the same way: everything has stopped. Nothing could be further from the truth.

\begin{definition}[Dynamic equilibrium]
A system is in a state of \cle{dynamic equilibrium} when the forward and reverse reactions continue to occur at the microscopic (molecular) level, but at \textbf{equal rates}, so that the macroscopic properties of the system (concentrations, pressure, colour, mass of each phase) remain \textbf{constant} with time.
\end{definition}

The word \emph{dynamic} is essential: the molecules never stop reacting. A useful image is that of two escalators in a busy shop in Bamenda, one going up and one going down, carrying the same number of people per minute. The number of people on each floor stays constant, yet people are continuously moving between floors. Similarly, in the closed bottle of water, molecules evaporate and condense at equal rates: the water level is constant, but the traffic of molecules never ceases.

\begin{attention}[Equilibrium does NOT mean the reaction has stopped]
At equilibrium, the reaction has \textbf{not} stopped. Forward and reverse reactions both continue indefinitely at equal rates. What has stopped is the \emph{net change}: the rate of formation of products exactly equals their rate of consumption. Saying "the reaction stops at equilibrium" is a classic error that costs marks.
\end{attention}

\begin{attention}[Constant does not mean equal]
At equilibrium the concentrations of reactants and products are \textbf{constant}, but they are generally \textbf{not equal}. In the esterification of ethanoic acid, for instance, the equilibrium mixture contains appreciable amounts of all four species, in proportions fixed by the equilibrium constant — not in equal amounts. Constant $\neq$ equal.
\end{attention}

\courssub{Approach to Equilibrium: Rates and Concentrations}

Consider the general reversible reaction \ce{A + B <=> C + D}, starting from pure \ce{A} and \ce{B} mixed together.

At the moment of mixing ($t = 0$), the concentrations of \ce{C} and \ce{D} are zero, so the reverse reaction cannot yet occur: the reverse rate is zero. The forward rate is at its maximum because the concentrations of \ce{A} and \ce{B} are highest. As time passes:

\begin{itemize}
\item \ce{A} and \ce{B} are consumed, so their concentrations fall and the \textbf{forward rate decreases};
\item \ce{C} and \ce{D} accumulate, so the \textbf{reverse rate increases} from zero;
\item eventually the two rates become \textbf{equal}: from that instant, every molecule of \ce{C} formed is matched by one destroyed, and all concentrations remain constant — dynamic equilibrium is established.
\end{itemize}

\begin{popfigure}
\begin{center}
\begin{tikzpicture}
\begin{axis}[
  width=0.85\linewidth, height=6cm,
  xlabel={time}, ylabel={rate of reaction},
  xmin=0, xmax=10, ymin=0, ymax=1.2,
  xtick=\empty, ytick={0.6},
  yticklabels={$r_{\mathrm{eq}}$},
  axis lines=left,
  legend style={at={(0.97,0.97)}, anchor=north east, font=\small},
]
\addplot[popblue, very thick, domain=0:10, samples=100] {0.6 + 0.5*exp(-0.6*x)};
\addlegendentry{forward rate}
\addplot[poporange, very thick, domain=0:10, samples=100] {0.6 - 0.6*exp(-0.6*x)};
\addlegendentry{reverse rate}
\addplot[popmuted, dashed, domain=0:10] {0.6};
\draw[popmuted, dashed] (axis cs:6.9,0) -- (axis cs:6.9,0.6);
\node[popdark, font=\small, anchor=south] at (axis cs:6.9,0.62) {equilibrium};
\end{axis}
\end{tikzpicture}
\end{center}
\legende{As a reversible reaction proceeds in a closed system, the forward rate falls and the reverse rate rises until they become equal: dynamic equilibrium.}
\end{popfigure}

The corresponding concentration–time graph shows the concentrations of reactants falling and those of products rising, then all levelling off to constant (but generally different) values.

\begin{popfigure}
\begin{center}
\begin{tikzpicture}
\begin{axis}[
  width=0.85\linewidth, height=6cm,
  xlabel={time}, ylabel={concentration},
  xmin=0, xmax=10, ymin=0, ymax=1.2,
  xtick=\empty, ytick=\empty,
  axis lines=left,
  legend style={at={(0.97,0.5)}, anchor=east, font=\small},
]
\addplot[popblue, very thick, domain=0:10, samples=100] {0.35 + 0.75*exp(-0.55*x)};
\addlegendentry{reactants \ce{A}, \ce{B}}
\addplot[popgreen, very thick, domain=0:10, samples=100] {0.75 - 0.75*exp(-0.55*x)};
\addlegendentry{products \ce{C}, \ce{D}}
\draw[popmuted, dashed] (axis cs:6.5,0) -- (axis cs:6.5,1.05);
\node[popdark, font=\small, anchor=west] at (axis cs:6.6,1.05) {equilibrium reached};
\end{axis}
\end{tikzpicture}
\end{center}
\legende{Concentrations change rapidly at first, then level off at constant values once equilibrium is attained. Note that the equilibrium concentrations are constant but not equal.}
\end{popfigure}

\begin{popfigure}
\begin{center}
\begin{tikzpicture}[scale=1.0]
% vessel
\draw[popdark, very thick] (0,0) -- (0,3.4) -- (4,3.4) -- (4,0) -- cycle;
% liquid
\fill[popblueL] (0.06,0.06) rectangle (3.94,1.3);
\draw[popblue, thick] (0.06,1.3) -- (3.94,1.3);
\node[popdark, font=\small] at (2,0.7) {liquid water};
\node[popdark, font=\small] at (2,2.9) {water vapour};
% vapour molecules
\foreach \x/\y in {0.6/1.8, 1.4/2.4, 2.2/1.9, 3.0/2.5, 3.5/1.7, 1.0/2.9, 2.7/2.9}{
  \fill[popblue] (\x,\y) circle (0.06);}
% arrows evaporation
\draw[->, poporange, very thick] (1.2,1.35) .. controls (1.1,1.8) .. (1.3,2.2);
\draw[->, poporange, very thick] (3.2,1.35) .. controls (3.1,1.8) .. (3.3,2.2);
\node[poporange, font=\small, anchor=west] at (3.45,2.0) {evaporation};
% arrows condensation
\draw[->, popgreen, very thick] (2.0,2.2) .. controls (2.2,1.8) .. (2.0,1.35);
\node[popgreen, font=\small, anchor=east] at (1.85,2.0) {condensation};
% stopper
\draw[popdark, thick, fill=popgold] (1.6,3.4) rectangle (2.4,3.7);
\node[popdark, font=\small, anchor=west] at (2.5,3.55) {closed system};
\end{tikzpicture}
\end{center}
\legende{Liquid–vapour equilibrium of water in a closed vessel: molecules evaporate and condense at equal rates, so the liquid level stays constant.}
\end{popfigure}

\courssub{Characteristics and Conditions of Dynamic Equilibrium}

\begin{propriete}[Characteristics of a system at dynamic equilibrium]
A chemical system at dynamic equilibrium has the following characteristics:
\begin{enumerate}
\item The \textbf{forward and reverse reactions occur at equal rates}.
\item The \textbf{macroscopic properties are constant}: concentrations, total pressure (for gases), colour intensity, density and mass of each phase do not change with time.
\item The equilibrium can be approached \textbf{from either direction} (starting from reactants or from products); the same equilibrium state is reached under the same conditions.
\item The equilibrium is \textbf{dynamic}: molecular-level activity continues indefinitely.
\end{enumerate}
No proof of these points is required at the GCE, but you must be able to state and explain them.
\end{propriete}

\begin{propriete}[Conditions required for dynamic equilibrium]
Dynamic equilibrium can only be attained if \textbf{all} of the following conditions hold:
\begin{enumerate}
\item The reaction must be \textbf{reversible}.
\item The system must be \textbf{closed}: no matter (reactant, product or solvent vapour) may enter or leave, otherwise a species would be removed and the reverse reaction could not balance the forward one.
\item The \textbf{temperature must be constant}, because the rates of the forward and reverse reactions, and the position of equilibrium itself, depend on temperature.
\item The system must be given \textbf{sufficient time} for the forward and reverse rates to become equal (a catalyst shortens this time but does not change the final state).
\end{enumerate}
\end{propriete}

\begin{exemplebox}[Why an open beaker never reaches equilibrium]
If the bottle of water is left \emph{open}, the water vapour escapes into the air and is carried away. Condensation can no longer balance evaporation, the reverse process is starved of vapour, and the water level slowly falls until the bottle is dry. Equilibrium is impossible because the system is not closed. This is exactly why the esterification mixture is heated under \emph{reflux}: the condenser returns the vapours to the flask, keeping the system closed so that equilibrium can be established.
\end{exemplebox}

\begin{exoresolu}[Identifying systems at dynamic equilibrium]
For each situation, state whether a dynamic equilibrium exists and justify your answer.
\begin{enumerate}
\item Ice and liquid water together in a perfectly insulated, sealed flask at $0\ ^{\circ}\mathrm{C}$.
\item Steam escaping from a boiling open saucepan.
\item A sealed flask containing \ce{N2}, \ce{H2} and \ce{NH3} whose total pressure no longer changes with time.
\end{enumerate}
\tcblower
\textbf{Solution.}
\begin{enumerate}
\item \textbf{Yes.} The flask is closed (insulated and sealed), the temperature is constant at $0\ ^{\circ}\mathrm{C}$, and the process \ce{H2O(s) <=> H2O(l)} is reversible. Ice melts and water freezes at equal rates: the masses of ice and liquid remain constant. All conditions for dynamic equilibrium are satisfied.
\item \textbf{No.} The saucepan is an \emph{open} system: steam escapes into the kitchen and cannot condense back into the pan. The reverse process cannot balance the forward one, so no equilibrium is established — eventually all the water boils away.
\item \textbf{Yes.} The flask is sealed (closed system) and the reaction \ce{N2(g) + 3H2(g) <=> 2NH3(g)} is reversible. The total pressure is a macroscopic property; since it is constant with time, the amounts of gas no longer change, meaning the forward and reverse rates are equal: dynamic equilibrium has been reached. (Note: constant pressure is the \emph{evidence} of equilibrium here, not the cause.)
\end{enumerate}
\end{exoresolu}

\begin{methode}[Testing whether a system is at equilibrium]
To decide whether a reversible system has reached dynamic equilibrium:
\begin{enumerate}
\item Check that the system is \textbf{closed} and at \textbf{constant temperature}.
\item Identify a \textbf{macroscopic property} (concentration, pressure, colour, mass of a phase).
\item Measure it at \textbf{two different times}. If the value is unchanged, and the conditions above hold, the system is at dynamic equilibrium.
\item Remember: equal rates at the molecular level, constant properties at the macroscopic level — never "the reaction has stopped".
\end{enumerate}
\end{methode}

\begin{aretenir}[Key points of this section]
\begin{itemize}
\item A \cle{reversible reaction} proceeds in both directions and is written with a double arrow: \ce{A + B <=> C + D}.
\item At \cle{dynamic equilibrium}, forward and reverse rates are \textbf{equal}, so macroscopic properties are \textbf{constant} — but the molecules never stop reacting.
\item Constant concentrations do \textbf{not} mean equal concentrations.
\item Equilibrium requires a \textbf{reversible reaction} in a \textbf{closed system} at \textbf{constant temperature}, given sufficient time.
\item Equilibrium can be reached from either direction; the same state is obtained under the same conditions.
\end{itemize}
\end{aretenir}

\begin{savaistu}[Equilibrium all around us]
The fizz in a bottle of \emph{Top} or \emph{Guinness} kept in a Douala drinks depot is an equilibrium: carbon dioxide is dissolved in the drink under pressure, and \ce{CO2(g) <=> CO2(aq)} is established in the sealed bottle. When the cap is removed, the system is opened, \ce{CO2} escapes, and the equilibrium is destroyed — the drink slowly goes flat. The same principle explains why a sealed bottle keeps its sparkle for months.
\end{savaistu}

\begin{exercice}[Checking your understanding]
\begin{enumerate}
\item Define a reversible reaction and give one chemical example with a balanced equation.
\item Explain, in terms of molecular behaviour, why the level of water in a tightly closed bottle remains constant.
\item A student says: "At equilibrium, the concentrations of reactants and products are equal and all reaction has ceased." Identify and correct the two mistakes in this statement.
\item State the four conditions necessary for a system to attain dynamic equilibrium.
\end{enumerate}
\end{exercice}

\begin{corrige}[Exercise 1]
\begin{enumerate}
\item A reversible reaction is one that can proceed in both directions: products can regenerate the reactants. Example: \ce{CH3COOH(l) + C2H5OH(l) <=> CH3COOC2H5(l) + H2O(l)}.
\item Water molecules continually evaporate from the liquid surface, but because the bottle is closed the vapour cannot escape; vapour molecules condense back at the same rate. Equal rates of evaporation and condensation keep the amount of liquid — and hence the level — constant.
\item Mistake 1: concentrations at equilibrium are \emph{constant}, not necessarily \emph{equal}. Mistake 2: the reactions have \emph{not} ceased; forward and reverse reactions continue at equal rates (the equilibrium is dynamic).
\item The reaction must be reversible; the system must be closed; the temperature must be constant; sufficient time must be allowed for the forward and reverse rates to become equal.
\end{enumerate}
\end{corrige}

\coursec{Equilibrium Position and Le Chatelier's Principle}

In the previous section we established that a reversible reaction in a closed system does not stop: it reaches a \cle{dynamic equilibrium} in which the forward and reverse reactions continue at equal rates, so that the macroscopic composition of the mixture no longer changes. But ``reaching equilibrium'' does not mean ``half the reactants have been converted''. Some equilibria contain almost entirely products; others contain almost entirely reactants. In this section we learn to describe \emph{where} an equilibrium lies, and --- far more importantly for the chemist and the industrial engineer --- how to \emph{push} it in the direction we want.

\courssub{The Equilibrium Position}

\begin{definition}[Equilibrium position]
The \cle{equilibrium position} of a reversible reaction describes the relative proportions of reactants and products present in the equilibrium mixture.
\begin{itemize}
\item If, at equilibrium, the mixture contains \textbf{mostly products}, we say the equilibrium position \emph{lies to the \cle{right}} (products are favoured).
\item If, at equilibrium, the mixture contains \textbf{mostly reactants}, we say the equilibrium position \emph{lies to the \cle{left}} (reactants are favoured).
\end{itemize}
\end{definition}

Consider two familiar examples. The esterification of ethanol by ethanoic acid,
\[
\ce{CH3COOH(l) + C2H5OH(l) <=> CH3COOC2H5(l) + H2O(l)},
\]
settles at an equilibrium containing appreciable amounts of \emph{both} reactants and products --- the position lies roughly in the middle. By contrast, the reaction of hydrogen with oxygen to form water goes essentially to completion under suitable conditions: the equilibrium lies far to the right. The decomposition of limestone, $\ce{CaCO3(s) <=> CaO(s) + CO2(g)}$, lies far to the left at room temperature, which is why the limestone hills around Buea do not spontaneously crumble into quicklime.

\begin{attention}[Equilibrium position is not rate]
A common confusion: the equilibrium position tells us \emph{how far} a reaction goes, not \emph{how fast} it gets there. A reaction may have an equilibrium lying far to the right yet be immeasurably slow (a mixture of $\ce{H2}$ and $\ce{O2}$ at room temperature). Position is a question of \textbf{energetics and balance of rates at equilibrium}; speed is a question of \textbf{kinetics}. Keep the two ideas strictly separate.
\end{attention}

The crucial experimental fact is this: \textbf{the equilibrium position of a given reaction is not fixed}. It depends on the conditions --- concentrations, pressure (for gases) and temperature. If we change the conditions, the composition of the equilibrium mixture changes: we say the equilibrium position has \emph{shifted}. The whole art of this section is to predict the direction of the shift.

\courssub{Le Chatelier's Principle}

\begin{experience}[A thought experiment with the esterification equilibrium]
Suppose we have the esterification mixture above at equilibrium, and we now pour in a large excess of cheap ethanol. What happens? Immediately after the addition, the forward reaction $\ce{CH3COOH + C2H5OH -> CH3COOC2H5 + H2O}$ speeds up (more ethanol molecules means more fruitful collisions), while the reverse rate is initially unchanged. The forward reaction now outruns the reverse one, so ester and water accumulate and acid is consumed, until a \emph{new} equilibrium is established --- one containing \textbf{more ester} than before. The system has responded to the disturbance (extra ethanol) by \emph{removing} some of that ethanol. It has, in a sense, ``fought back'' against the change we imposed.
\end{experience}

This behaviour is completely general, and was summarised by the French chemist Henri Le Chatelier in 1884.

\begin{propriete}[Le Chatelier's Principle]
\textbf{Statement (examinable):} If a constraint (a change of concentration, pressure or temperature) is imposed on a system at dynamic equilibrium, the system adjusts itself so as to \cle{oppose} (partially counteract) that constraint, and a new equilibrium position is established.

Equivalently: the equilibrium shifts in the direction that \emph{removes} added substance, \emph{produces} removed substance, \emph{reduces} an imposed pressure increase, or \emph{absorbs} added heat.
\end{propriete}

\begin{attention}[The system only \emph{opposes} the change --- it never cancels it]
Le Chatelier's principle says the system shifts so as to \emph{reduce} the effect of the disturbance, \textbf{not} to eliminate it. If you add ethanol, the new equilibrium contains less ethanol than immediately after your addition, but \emph{more} than before the addition. The opposition is always partial.
\end{attention}

\begin{methode}[Predicting the shift of an equilibrium]
\begin{enumerate}
\item \textbf{Identify the change} imposed on the system (addition/removal of a species, change of pressure, change of temperature).
\item \textbf{Ask: which direction opposes the change?} (Which direction consumes the added species? Which direction produces more gas moles to relieve a pressure decrease? Which direction absorbs heat?)
\item \textbf{State the shift} (to the right or to the left) and \textbf{state the consequence} on the yield of the desired product.
\item \textbf{Check the special cases:} no gas moles change? A pure solid or liquid involved? A catalyst added? --- these may cause \emph{no} shift.
\end{enumerate}
\end{methode}

\courssub{Effect of Changing Concentration}

\begin{propriete}[Concentration changes]
If the \textbf{concentration of a reactant is increased} (or a product removed), the equilibrium shifts \textbf{to the right}, consuming some of the added reactant and increasing the yield of products. If the \textbf{concentration of a reactant is decreased} (or a product added), the equilibrium shifts \textbf{to the left}.
\end{propriete}

The intuition is kinetic, as in the esterification example: adding a reactant raises the forward rate above the reverse rate, so net forward reaction occurs until the two rates are equal again. Removing a product (for instance distilling off the ester as it forms) continually lowers the reverse rate, so the forward reaction keeps winning and the yield rises --- a standard trick in the laboratory preparation of esters.

\begin{exoresolu}[Shifting the equilibrium in the chromate/dichromate system]
The orange dichromate ion and the yellow chromate ion coexist in the equilibrium
\[
\ce{Cr2O7^{2-}(aq) + H2O(l) <=> 2CrO4^{2-}(aq) + 2H+(aq)} \qquad \text{(orange} \rightleftharpoons \text{yellow)}
\]
Predict and explain the colour change observed when (a) dilute sodium hydroxide is added; (b) dilute hydrochloric acid is then added in excess.
\tcblower
\textbf{(a) Adding \ce{NaOH}.} The $\ce{OH-}$ ions neutralise $\ce{H+}$ ions, so $[\ce{H+}]$ \emph{decreases}. By Le Chatelier's principle, the system opposes this removal by shifting in the direction that \emph{produces} $\ce{H+}$ --- that is, \textbf{to the right}. More yellow $\ce{CrO4^{2-}}$ forms: the solution turns \textbf{yellow}.

\textbf{(b) Adding excess \ce{HCl}.} Now $[\ce{H+}]$ \emph{increases}. The system opposes this by consuming $\ce{H+}$, shifting \textbf{to the left}. Orange $\ce{Cr2O7^{2-}}$ re-forms: the solution turns \textbf{orange}.

This reversible colour change is a classic GCE qualitative observation and a beautiful visual proof that equilibria really do shift.
\end{exoresolu}

\begin{attention}[Adding a pure solid or pure liquid causes no shift]
The ``concentration'' of a pure solid or a pure liquid is constant (it is its density in molar terms, which does not change). Adding more \emph{lumps} of $\ce{CaCO3}$ to the equilibrium $\ce{CaCO3(s) <=> CaO(s) + CO2(g)}$ does \textbf{not} shift the position at all --- there is simply more solid sitting there. Only changes in the amounts of \textbf{gases or dissolved species} matter. Never write ``adding more solid shifts the equilibrium'' in an examination.
\end{attention}

\courssub{Effect of Changing Pressure (Gaseous Systems)}

Pressure only matters when gases are involved, and it acts through the \emph{number of moles of gas} on each side of the equation. When the total pressure of an equilibrium mixture is increased (by compressing the mixture into a smaller volume), the system can relieve the constraint by shifting towards the side with \textbf{fewer moles of gas} --- fewer particles means fewer collisions with the walls, hence lower pressure.

\begin{propriete}[Pressure changes]
For a gaseous equilibrium, \textbf{increasing the total pressure} (by compression) shifts the equilibrium towards the side with the \cle{fewer moles of gas}; \textbf{decreasing the pressure} shifts it towards the side with more moles of gas. If both sides contain the \emph{same} number of moles of gas, a pressure change causes \textbf{no shift}.
\end{propriete}

\begin{exemplebox}[Applying the rule]
\begin{itemize}
\item $\ce{N2(g) + 3H2(g) <=> 2NH3(g)}$: 4 moles of gas on the left, 2 on the right. Increasing the pressure shifts the equilibrium \textbf{to the right} --- more ammonia. This is why the Haber process uses high pressure.
\item $\ce{N2O4(g) <=> 2NO2(g)}$: 1 mole versus 2 moles. Compression shifts the equilibrium \textbf{to the left}; the brown colour of $\ce{NO2}$ pales.
\item $\ce{H2(g) + I2(g) <=> 2HI(g)}$: 2 moles on each side. Changing the pressure causes \textbf{no shift} (although the rates change).
\end{itemize}
\end{exemplebox}

\begin{attention}[Two pressure traps]
\begin{enumerate}
\item Pressure changes only shift equilibria involving gases with \textbf{unequal numbers of moles of gas} on the two sides. Solids and liquids are essentially incompressible and do not count.
\item Adding an \textbf{inert gas} (e.g.\ argon) \emph{at constant volume} raises the total pressure but changes none of the partial pressures of the reacting gases: there is \textbf{no shift}. Only compression (change of volume) or changing the amount of a reacting gas matters.
\end{enumerate}
\end{attention}

\courssub{Effect of Changing Temperature}

Temperature is the most subtle constraint, because it is the only one of the three that changes the \emph{value of the equilibrium constant} (as we shall see in Section 4). The rule, however, is easy to apply once the sign of the enthalpy change is known.

\begin{propriete}[Temperature changes]
For an \textbf{exothermic} forward reaction ($\Delta H < 0$), \textbf{raising} the temperature shifts the equilibrium \textbf{to the left} (the endothermic, heat-absorbing direction is favoured); \textbf{lowering} the temperature shifts it \textbf{to the right}. For an \textbf{endothermic} forward reaction ($\Delta H > 0$), the reverse holds: raising the temperature shifts the equilibrium \textbf{to the right}.
\end{propriete}

The logic follows directly from Le Chatelier: heat added to the system is opposed by the reaction direction that \emph{absorbs} heat, i.e.\ the endothermic direction. A useful memory trick is to treat heat as a chemical species: for $\ce{N2 + 3H2 <=> 2NH3}$, $\Delta H = -92\ \mathrm{kJ\,mol^{-1}}$, write ``$\ce{N2 + 3H2 <=> 2NH3 + heat}$''. Adding heat then obviously pushes the equilibrium left, exactly like adding a product.

\begin{exemplebox}[Heating the brown gas]
A sealed tube containing the equilibrium $\ce{N2O4(g) <=> 2NO2(g)}$, $\Delta H = +57\ \mathrm{kJ\,mol^{-1}}$ (endothermic to the right), darkens when warmed and pales when cooled: heating favours the endothermic production of brown $\ce{NO2}$; cooling favours the exothermic reverse reaction.
\end{exemplebox}

\courssub{Effect of a Catalyst}

\begin{propriete}[A catalyst does not shift the equilibrium]
A catalyst speeds up the \textbf{forward and reverse reactions equally} (by providing an alternative pathway of lower activation energy for both). Consequently:
\begin{itemize}
\item the \textbf{equilibrium position is unchanged} --- the same equilibrium mixture is obtained, with or without catalyst;
\item the only effect is that equilibrium is \cle{reached faster}.
\end{itemize}
(The proof is not examinable; the statement and its application are.)
\end{propriete}

This is why catalysts are industrially precious even though they cannot improve the equilibrium yield: a process that reaches 30\% conversion in one hour instead of one month is the difference between profit and bankruptcy.

\courssub{Summary Diagram and Graphical View}

\begin{popfigure}
\begin{center}
\begin{tikzpicture}[
  box/.style={draw=popblue, fill=popblueL, rounded corners, text width=3.6cm, align=center, font=\small, minimum height=1.1cm},
  eff/.style={draw=poporange, fill=white, rounded corners, text width=4.6cm, align=center, font=\small, minimum height=1.1cm},
  arr/.style={-{Stealth[length=2.5mm]}, thick, popdark}]
\node[box] (c1) at (0,3.6) {Add a reactant (or remove a product)};
\node[eff] (e1) at (6.4,3.6) {Shift to the \textbf{right}: yield of products increases};
\node[box] (c2) at (0,2.2) {Compress a gaseous mixture};
\node[eff] (e2) at (6.4,2.2) {Shift towards the side with \textbf{fewer moles of gas}};
\node[box] (c3) at (0,0.8) {Raise the temperature};
\node[eff] (e3) at (6.4,0.8) {Shift in the \textbf{endothermic} direction ($K$ changes)};
\node[box] (c4) at (0,-0.6) {Add a catalyst};
\node[eff] (e4) at (6.4,-0.6) {\textbf{No shift}: equilibrium reached faster};
\node[box] (c5) at (0,-2.0) {Add a pure solid or liquid};
\node[eff] (e5) at (6.4,-2.0) {\textbf{No shift}: its ``concentration'' is constant};
\draw[arr] (c1) -- (e1);
\draw[arr] (c2) -- (e2);
\draw[arr] (c3) -- (e3);
\draw[arr] (c4) -- (e4);
\draw[arr] (c5) -- (e5);
\end{tikzpicture}
\end{center}
\legende{How a system at equilibrium responds to each type of constraint (Le Chatelier's principle).}
\end{popfigure}

\begin{popfigure}
\begin{center}
\begin{tikzpicture}
\begin{axis}[
  width=11cm, height=6.5cm,
  xlabel={time}, ylabel={concentration},
  xmin=0, xmax=10, ymin=0, ymax=3.2,
  xtick=\empty, ytick=\empty,
  axis lines=left, thick,
  legend style={at={(0.98,0.45)}, anchor=east, font=\small, draw=popline}]
\addplot[popblue, very thick, smooth] coordinates {(0,1.0) (2,1.0) (2.01,2.6) (3,2.2) (4,1.9) (5,1.75) (6,1.7) (8,1.7) (10,1.7)};
\addplot[poporange, very thick, smooth] coordinates {(0,2.0) (2,2.0) (2.01,2.0) (3,2.4) (4,2.7) (5,2.85) (6,2.9) (8,2.9) (10,2.9)};
\legend{reactant, product}
\draw[dashed, popmuted] (axis cs:2,0) -- (axis cs:2,3.1);
\node[popmuted, font=\small, align=center] at (axis cs:2,3.05) [above] {reactant added};
\draw[dashed, popmuted] (axis cs:6,0) -- (axis cs:6,3.1);
\node[popmuted, font=\small, align=center] at (axis cs:6.1,3.05) [above right] {new equilibrium};
\end{axis}
\end{tikzpicture}
\end{center}
\legende{Concentration--time graph: at the dashed line a reactant is suddenly added (instantaneous jump of its concentration); the system then shifts right until a new equilibrium is established, with more product and slightly more reactant than originally.}
\end{popfigure}

Note carefully the shape of this graph: the added reactant's concentration \emph{jumps} vertically at the moment of addition, then falls back --- but it settles \emph{above} its original value. This is the graphical signature of Le Chatelier's ``partial opposition''.

\courssub{Industrial Applications: the Haber and Contact Processes}

\begin{savaistu}[Feeding the world from thin air]
Ammonia, $\ce{NH3}$, is the starting material for nitrogen fertilisers without which modern agriculture --- including the maize and cocoa farms of Cameroon --- could not feed today's population. It is manufactured from nitrogen (from air) and hydrogen (from natural gas) by the \cle{Haber process}, developed by Fritz Haber and Carl Bosch in the early twentieth century. Choosing the operating conditions is a masterclass in applied Le Chatelier.
\end{savaistu}

The equilibrium is
\[
\ce{N2(g) + 3H2(g) <=> 2NH3(g)} \qquad \Delta H^{\ominus} = -92\ \mathrm{kJ\,mol^{-1}} \text{ (exothermic).}
\]

Let us apply our method to each condition:

\textbf{Pressure.} There are 4 moles of gas on the left and 2 on the right, so \emph{high pressure} shifts the equilibrium to the right and increases the yield of ammonia. High pressure also speeds up the reaction. Industry uses about $200$--$250\ \mathrm{atm}$. Why not $1000\ \mathrm{atm}$? Because compressing gases and building vessels that withstand enormous pressures is extremely expensive and dangerous: the extra yield does not repay the extra cost.

\textbf{Temperature.} The forward reaction is exothermic, so \emph{low temperature} favours a high yield of ammonia. But at low temperature the reaction is hopelessly slow --- a high yield is useless if it takes years to obtain. Industry therefore uses a \emph{compromise} temperature of about $400$--$450\ ^{\circ}\mathrm{C}$: the yield per pass is modest, but it is obtained quickly.

\textbf{Catalyst.} Finely divided \emph{iron} (with promoters) is used. It cannot improve the equilibrium yield, but it allows the compromise temperature to deliver that yield at an acceptable rate.

\textbf{Removing the product.} The gas mixture leaving the converter is cooled; ammonia liquefies and is removed. By Le Chatelier's principle, removing $\ce{NH3}$ pulls the equilibrium to the right, and the unreacted $\ce{N2}$ and $\ce{H2}$ are \emph{recycled} through the converter, pushing the overall conversion towards $98\ \%$.

\begin{popfigure}
\begin{center}
\begin{tikzpicture}[
  stage/.style={draw=popblue, fill=popblueL, rounded corners, text width=3.0cm, align=center, font=\small, minimum height=1.2cm},
  arr/.style={-{Stealth[length=2.5mm]}, very thick, poporange},
  lab/.style={font=\scriptsize, align=center, text=popdark}]
\node[stage, align=center] (mix) at (0,0){Compressor\\ $\approx 200$--$250\ \mathrm{atm}$};
\node[stage, align=center] (conv) at (4.6,0){Converter\\ iron catalyst\\ $400$--$450\ ^{\circ}\mathrm{C}$};
\node[stage, align=center] (cool) at (9.2,0){Cooler / condenser\\ $\ce{NH3}$ liquefies};
\node[stage, fill=popgreenL, draw=popgreen] (prod) at (9.2,-2.6) {Liquid $\ce{NH3}$ collected};
\draw[arr] (mix) -- node[lab, above]{$\ce{N2} + 3\ce{H2}$} (conv);
\draw[arr] (conv) -- node[lab, above]{$\ce{NH3} +$ unreacted gases} (cool);
\draw[arr] (cool) -- (prod);
\draw[arr, popgreen] (cool.north) .. controls (6,2.2) and (3,2.2) .. node[lab, above]{unreacted $\ce{N2}$, $\ce{H2}$ recycled} (mix.north);
\end{tikzpicture}
\end{center}
\legende{Schematic of the Haber process. High pressure and product removal favour the yield; the moderate temperature and iron catalyst favour the rate; recycling pushes the overall conversion close to completion.}
\end{popfigure}

\begin{exemplebox}[The Contact process for sulphuric acid]
The key step in manufacturing sulphuric acid (used in fertiliser and battery production) is
\[
\ce{2SO2(g) + O2(g) <=> 2SO3(g)} \qquad \Delta H^{\ominus} = -196\ \mathrm{kJ\,mol^{-1}}.
\]
By the same reasoning: 3 moles of gas become 2, so high pressure favours $\ce{SO3}$ --- but the yield is already so good that industry uses only about $2\ \mathrm{atm}$, saving the cost of compression. A compromise temperature of about $450\ ^{\circ}\mathrm{C}$ and a vanadium(V) oxide catalyst, $\ce{V2O5}$, give a fast reaction with an excellent yield.
\end{exemplebox}

\begin{aretenir}[Key points]
\begin{itemize}
\item The \textbf{equilibrium position} describes whether reactants or products predominate at equilibrium (``left'' or ``right'').
\item \textbf{Le Chatelier's principle:} a system at equilibrium subjected to a change adjusts so as to \emph{oppose} the change --- but only partially.
\item Adding a reactant (or removing a product) shifts the equilibrium \textbf{right}; the reverse shifts it \textbf{left}. Pure solids and liquids have no effect.
\item Compressing a gaseous equilibrium shifts it towards the side with \textbf{fewer moles of gas}; equal moles of gas means no shift; inert gas at constant volume means no shift.
\item Raising the temperature favours the \textbf{endothermic} direction --- and is the only change that alters the value of $K$.
\item A \textbf{catalyst} speeds up both directions equally: equilibrium is reached faster but the position (and the yield) is unchanged.
\item Industry chooses \textbf{compromise conditions}: the Haber process ($\approx 450\ ^{\circ}\mathrm{C}$, $\approx 200\ \mathrm{atm}$, iron catalyst, product removed and recycled) balances yield against rate and cost.
\end{itemize}
\end{aretenir}

\begin{exercice}[Predicting shifts]
For each equilibrium below, state the effect on the yield of the indicated product of (i) increasing the pressure by compression, (ii) raising the temperature, (iii) adding a catalyst. Justify each answer.
\begin{enumerate}
\item $\ce{N2(g) + 3H2(g) <=> 2NH3(g)}$, $\Delta H = -92\ \mathrm{kJ\,mol^{-1}}$ (consider $\ce{NH3}$ as the product).
\item $\ce{CaCO3(s) <=> CaO(s) + CO2(g)}$, $\Delta H = +178\ \mathrm{kJ\,mol^{-1}}$ (consider $\ce{CO2}$ as the product). What is the effect of adding more $\ce{CaCO3(s)}$?
\item $\ce{H2(g) + I2(g) <=> 2HI(g)}$, $\Delta H = +53\ \mathrm{kJ\,mol^{-1}}$ (consider $\ce{HI}$ as the product).
\end{enumerate}
\end{exercice}

\begin{corrige}[Exercise 1]
\textbf{1.} (i) 4 moles of gas $\to$ 2 moles: compression shifts right, yield of $\ce{NH3}$ \emph{increases}. (ii) Forward reaction exothermic: raising $T$ shifts left (endothermic direction), yield \emph{decreases}. (iii) No change in yield; equilibrium reached faster.

\textbf{2.} (i) 0 moles of gas on the left, 1 on the right: compression shifts left, yield of $\ce{CO2}$ \emph{decreases}. (ii) Forward reaction endothermic: raising $T$ shifts right, yield \emph{increases} (this is why lime kilns are very hot). (iii) No change in yield. Adding more $\ce{CaCO3(s)}$: \emph{no effect} --- it is a pure solid, whose ``concentration'' is constant.

\textbf{3.} (i) 2 moles of gas on each side: compression causes \emph{no shift}, yield unchanged. (ii) Forward reaction endothermic: raising $T$ shifts right, yield of $\ce{HI}$ \emph{increases}. (iii) No change in yield; equilibrium reached faster.
\end{corrige}

\coursec{The Equilibrium Law: Kc and Kp}

Le Chatelier's principle, which we studied in the previous section, is a wonderful \emph{qualitative} tool: it tells us the \emph{direction} in which an equilibrium shifts when we disturb it. But an examiner, an industrial chemist at a fertiliser plant, or a laboratory technician in Douala needs more than a direction --- they need \emph{numbers}. How much product will actually be present when equilibrium is reached? Is a reaction worth running industrially? To answer such questions we need a \emph{quantitative} description of equilibrium. That description is the \cle{equilibrium law}, and the number at its heart is the \cle{equilibrium constant}.

\courssub{From Observation to Law: The Equilibrium Constant $K_c$}

\begin{experience}[The observation behind the law]
Consider the reversible reaction between hydrogen and iodine at a fixed temperature:
\[
\ce{H2(g) + I2(g) <=> 2HI(g)}
\]
If we start several sealed tubes with \emph{different} initial amounts of \ce{H2}, \ce{I2} and \ce{HI}, allow each to reach equilibrium at the same temperature, and then analyse the mixtures, we find that the equilibrium concentrations differ from tube to tube. Yet the quantity
\[
\frac{[\ce{HI}]^2}{[\ce{H2}][\ce{I2}]}
\]
comes out to the \emph{same numerical value} in every tube. This remarkable regularity, established experimentally in the nineteenth century by Guldberg and Waage (the ``law of mass action''), is the basis of the equilibrium law: at a given temperature, the equilibrium composition is not arbitrary --- it is fixed by a constant.
\end{experience}

\begin{definition}[The equilibrium constant $K_c$]
For a reversible reaction at equilibrium at a given temperature,
\[
\ce{aA + bB <=> cC + dD}
\]
the \cle{equilibrium constant} in terms of concentration is defined by
\[
K_c=\frac{[\ce{C}]^c\,[\ce{D}]^d}{[\ce{A}]^a\,[\ce{B}]^b}
\]
where $[\ce{X}]$ denotes the \textbf{equilibrium} concentration of species \ce{X} in $\mathrm{mol\,dm^{-3}}$, and the powers $a,b,c,d$ are the stoichiometric coefficients of the balanced equation. The value of $K_c$ is constant at a fixed temperature.
\end{definition}

\begin{aretenir}[Writing a $K_c$ expression correctly]
\begin{itemize}
\item \textbf{Products} go on top, \textbf{reactants} on the bottom.
\item Each concentration is raised to the power of its \textbf{coefficient} in the balanced equation.
\item Only species whose concentration can vary appear: gases and species in solution. (Pure solids and pure liquids are omitted --- see the section on heterogeneous equilibria.)
\item The expression depends on the way the equation is written: for $\ce{2SO2 + O2 <=> 2SO3}$, $K_c=\dfrac{[\ce{SO3}]^2}{[\ce{SO2}]^2[\ce{O2}]}$; for the reverse reaction the constant is $1/K_c$, and halving all coefficients gives $\sqrt{K_c}$.
\end{itemize}
\end{aretenir}

\begin{attention}[Equilibrium concentrations, not initial concentrations!]
The single most common error at GCE is to substitute \emph{initial} concentrations (or initial amounts) into the $K_c$ expression. The law holds \textbf{only at equilibrium}. If you are given initial amounts, you must first work out the equilibrium composition (the ICE method below) before applying the expression.
\end{attention}

\courssub{Units of $K_c$}

$K_c$ is \emph{not} automatically dimensionless. Its units are obtained by substituting $\mathrm{mol\,dm^{-3}}$ for every concentration and simplifying. For a general reaction, if $\Delta n = (c+d)-(a+b)$ is the change in the number of moles of (concentration-dependent) species, then
\[
\text{units of }K_c = \left(\mathrm{mol\,dm^{-3}}\right)^{\Delta n}.
\]

\begin{exemplebox}[Deducing units of $K_c$]
\begin{itemize}
\item $\ce{H2(g) + I2(g) <=> 2HI(g)}$: \; $K_c=\dfrac{[\ce{HI}]^2}{[\ce{H2}][\ce{I2}]}$ has units $\dfrac{(\mathrm{mol\,dm^{-3}})^2}{(\mathrm{mol\,dm^{-3}})(\mathrm{mol\,dm^{-3}})}$ = \textbf{no units} ($\Delta n = 0$).
\item $\ce{N2(g) + 3H2(g) <=> 2NH3(g)}$: \; units $=\dfrac{(\mathrm{mol\,dm^{-3}})^2}{(\mathrm{mol\,dm^{-3}})(\mathrm{mol\,dm^{-3}})^3}=(\mathrm{mol\,dm^{-3}})^{-2}=\mathbf{mol^{-2}\,dm^{6}}$.
\item $\ce{PCl5(g) <=> PCl3(g) + Cl2(g)}$: \; units $=\mathrm{mol\,dm^{-3}}$ ($\Delta n = +1$).
\end{itemize}
\end{exemplebox}

\begin{attention}[$K_c$ and $K_p$ generally have units]
Never write ``$K_c$ has no units'' as a reflex. Deduce the units \emph{each time} from the expression and the balanced equation. Examiners award a mark for correct units, and a wrong or missing unit loses it.
\end{attention}

\courssub{Calculating $K_c$: The ICE Table Method}

Most GCE problems give \emph{initial} amounts and \emph{one piece of equilibrium information} (an amount, a concentration, or a percentage dissociation). The systematic tool is the \cle{ICE table}: \textbf{I}nitial, \textbf{C}hange, \textbf{E}quilibrium.

\begin{methode}[The ICE table procedure]
\begin{enumerate}
\item Write the balanced equation and the expression for $K_c$.
\item \textbf{Initial row}: enter the initial moles (or concentrations) of each species. Products are usually zero initially.
\item \textbf{Change row}: let the extent of reaction be $x$. Reactants change by $-a x$, $-b x$, \dots; products by $+c x$, $+d x$, \dots, according to the coefficients.
\item \textbf{Equilibrium row}: add the two rows algebraically.
\item Use the given equilibrium information to solve for $x$.
\item Convert equilibrium \emph{moles} to equilibrium \emph{concentrations} by dividing by the volume $V$ (in $\mathrm{dm^3}$).
\item Substitute into the $K_c$ expression and state the units.
\end{enumerate}
\end{methode}

\begin{popfigure}
\begin{center}
\begin{tikzpicture}[x=2.9cm,y=0.85cm]
\draw[popline] (-1.55,3.5) rectangle (3.55,-0.5);
\foreach \x in {-1.55,-0.05,1.15,2.35,3.55}{\draw[popline] (\x,3.5) -- (\x,-0.5);}
\foreach \y in {2.5,1.5,0.5}{\draw[popline] (-1.55,\y) -- (3.55,\y);}
\fill[popblueL] (-1.55,2.5) rectangle (3.55,3.5);
\node[popdark] at (-0.8,3.0) {\small\textbf{Stage}};
\node[popdark] at (0.55,3.0) {\small\textbf{A}};
\node[popdark] at (1.75,3.0) {\small\textbf{B}};
\node[popdark] at (2.95,3.0) {\small\textbf{C}};
\node[popmuted] at (-0.8,2.0) {\small Initial (mol)};
\node[popmuted] at (-0.8,1.0) {\small Change (mol)};
\node[popmuted] at (-0.8,0.0) {\small Equilibrium};
\node[popink] at (0.55,2.0) {\small $n_0(\mathrm{A})$};
\node[popink] at (1.75,2.0) {\small $n_0(\mathrm{B})$};
\node[popink] at (2.95,2.0) {\small $0$};
\node[poporange] at (0.55,1.0) {\small $-ax$};
\node[poporange] at (1.75,1.0) {\small $-bx$};
\node[poporange] at (2.95,1.0) {\small $+cx$};
\node[popgreen] at (0.55,0.0) {\small $n_0(\mathrm{A})-ax$};
\node[popgreen] at (1.75,0.0) {\small $n_0(\mathrm{B})-bx$};
\node[popgreen] at (2.95,0.0) {\small $cx$};
\node[popmuted, align=center, text width=3.2cm] at (4.9,1.5) {\small then divide by $V$ to get $[\,]_{\mathrm{eq}}$};
\end{tikzpicture}
\legende{ICE table template for $\ce{aA + bB <=> cC}$: the equilibrium row is initial $+$ change; concentrations follow by dividing by the volume.}
\end{center}
\end{popfigure}

\begin{exoresolu}[Calculating $K_c$ for the esterification of ethanol]
Statement. \SI{1.0}{mol} of ethanoic acid and \SI{1.0}{mol} of ethanol are mixed with a catalyst and allowed to reach equilibrium at $298\ \mathrm{K}$:
\[
\ce{CH3COOH(l) + C2H5OH(l) <=> CH3COOC2H5(l) + H2O(l)}
\]
At equilibrium, \SI{0.33}{mol} of ethanoic acid remains. Calculate $K_c$.
\tcblower
\textbf{Solution.} Build the ICE table in moles (all species are in the same liquid mixture of volume $V$):
\begin{center}
\begin{tabular}{|l|c|c|c|c|}
\hline
\rowcolor{popblueL} & \ce{CH3COOH} & \ce{C2H5OH} & \ce{CH3COOC2H5} & \ce{H2O} \\
\hline
Initial / mol & $1.0$ & $1.0$ & $0$ & $0$ \\
\hline
Change / mol & $-x$ & $-x$ & $+x$ & $+x$ \\
\hline
Equilibrium / mol & $1.0-x$ & $1.0-x$ & $x$ & $x$ \\
\hline
\end{tabular}
\end{center}
Given $1.0-x=0.33$, so $x=0.67\ \mathrm{mol}$. Since all coefficients are 1, every concentration is (moles)$/V$ and $V$ cancels:
\[
K_c=\frac{[\ce{CH3COOC2H5}][\ce{H2O}]}{[\ce{CH3COOH}][\ce{C2H5OH}]}
=\frac{(0.67/V)(0.67/V)}{(0.33/V)(0.33/V)}
=\frac{0.67^2}{0.33^2}\approx 4.1.
\]
Here $\Delta n = 0$, so $K_c$ has \textbf{no units}. Note how the unknown volume cancelled --- a useful check.
\end{exoresolu}

\begin{exoresolu}[Calculating $K_c$ when the volume does not cancel]
Statement. $\SI{2.00}{mol}$ of \ce{PCl5} vapour is introduced into a $\SI{1.00}{dm^3}$ vessel and heated. At equilibrium, $\SI{0.80}{mol}$ of \ce{Cl2} is present:
\[
\ce{PCl5(g) <=> PCl3(g) + Cl2(g)}
\]
Calculate $K_c$ at this temperature.
\tcblower
\textbf{Solution.} ICE table in moles: start $(2.00, 0, 0)$; change $(-x, +x, +x)$; equilibrium $(2.00-x,\ x,\ x)$. Given $x = 0.80\ \mathrm{mol}$, the equilibrium amounts are $n_{\ce{PCl5}}=1.20$, $n_{\ce{PCl3}}=n_{\ce{Cl2}}=0.80\ \mathrm{mol}$. Since $V = 1.00\ \mathrm{dm^3}$, the concentrations are numerically equal to the moles:
\[
K_c=\frac{[\ce{PCl3}][\ce{Cl2}]}{[\ce{PCl5}]}=\frac{0.80\times0.80}{1.20}=0.53\ \mathrm{mol\,dm^{-3}}.
\]
Units: $\dfrac{(\mathrm{mol\,dm^{-3}})^2}{\mathrm{mol\,dm^{-3}}}=\mathrm{mol\,dm^{-3}}$ ($\Delta n = +1$). Here $\Delta n \neq 0$, so the volume does \emph{not} cancel --- always check.
\end{exoresolu}

\begin{exoresolu}[Finding an equilibrium concentration from $K_c$]
Statement. For $\ce{N2(g) + 3H2(g) <=> 2NH3(g)}$ at a certain temperature, $K_c = 0.50\ \mathrm{mol^{-2}\,dm^{6}}$. At equilibrium in a $\SI{2.0}{dm^3}$ vessel there are $\SI{0.10}{mol}$ of \ce{N2} and $\SI{0.30}{mol}$ of \ce{H2}. Calculate the equilibrium concentration of \ce{NH3}.
\tcblower
\textbf{Solution.} Equilibrium concentrations: $[\ce{N2}]=0.10/2.0=0.050\ \mathrm{mol\,dm^{-3}}$; $[\ce{H2}]=0.30/2.0=0.15\ \mathrm{mol\,dm^{-3}}$. Rearranging the equilibrium law:
\[
[\ce{NH3}]^2 = K_c\,[\ce{N2}][\ce{H2}]^3 = 0.50\times 0.050\times (0.15)^3 = 8.4\times10^{-5}\ \mathrm{mol^2\,dm^{-6}}
\]
\[
[\ce{NH3}] = \sqrt{8.4\times10^{-5}} = 9.2\times10^{-3}\ \mathrm{mol\,dm^{-3}}.
\]
\end{exoresolu}

\courssub{Gaseous Equilibria: Partial Pressure and $K_p$}

For reactions in the gas phase it is often more convenient to measure \emph{pressures} than concentrations.

\begin{definition}[Partial pressure]
In a mixture of gases, the \cle{partial pressure} $p_A$ of gas $A$ is the pressure that $A$ would exert if it alone occupied the whole volume at the same temperature. It is given by
\[
p_A = x_A \times P_{\text{total}}
\]
where $x_A=\dfrac{n_A}{n_{\text{total}}}$ is the \cle{mole fraction} of $A$. The sum of all partial pressures equals the total pressure (Dalton's law). Pressures are measured in pascal ($\mathrm{Pa}$, with $\mathrm{kPa}$ and $\mathrm{atm}$ also used).
\end{definition}

\begin{definition}[The equilibrium constant $K_p$]
For the gaseous equilibrium $\ce{aA(g) + bB(g) <=> cC(g) + dD(g)}$,
\[
K_p=\frac{p_C^{\,c}\; p_D^{\,d}}{p_A^{\,a}\; p_B^{\,b}}
\]
where each $p$ is the \textbf{equilibrium} partial pressure. The units of $K_p$ are $(\text{pressure unit})^{\Delta n}$, e.g.\ $\mathrm{Pa^{-2}}$ for the Haber equilibrium, deduced from the expression each time.
\end{definition}

\begin{exoresolu}[Calculating $K_p$ for the dissociation of \ce{PCl5}]
Statement. $\SI{1.00}{mol}$ of \ce{PCl5} is placed in a vessel and heated. At equilibrium, $60\%$ has dissociated and the total pressure is $\SI{200}{kPa}$:
\[
\ce{PCl5(g) <=> PCl3(g) + Cl2(g)}
\]
Calculate $K_p$.
\tcblower
\textbf{Solution.} ICE in moles: start $(1.00, 0, 0)$; change $(-0.60, +0.60, +0.60)$; equilibrium $(0.40, 0.60, 0.60)$, so $n_{\text{total}}=1.60\ \mathrm{mol}$.
\[
p_{\ce{PCl5}}=\tfrac{0.40}{1.60}\times200=50\ \mathrm{kPa},\quad
p_{\ce{PCl3}}=p_{\ce{Cl2}}=\tfrac{0.60}{1.60}\times200=75\ \mathrm{kPa}.
\]
\[
K_p=\frac{p_{\ce{PCl3}}\,p_{\ce{Cl2}}}{p_{\ce{PCl5}}}=\frac{75\times75}{50}=112.5\ \mathrm{kPa}\approx 1.1\times10^{2}\ \mathrm{kPa}.
\]
Units: $\dfrac{\mathrm{kPa}\times\mathrm{kPa}}{\mathrm{kPa}}=\mathrm{kPa}$ ($\Delta n = +1$).
\end{exoresolu}

\begin{attention}[Mole fractions need the \emph{total} number of moles]
When dissociation occurs, the total number of moles of gas \emph{changes} (here from $1.00$ to $1.60$ mol). A frequent error is to divide by the \emph{initial} number of moles instead of the \emph{equilibrium total}. Always compute $n_{\text{total}}$ from the equilibrium row of the ICE table.
\end{attention}

\begin{propriete}[Relationship between $K_p$ and $K_c$ (GCE-useful extra)]
Using $p_i = \dfrac{n_i RT}{V}=[i]\,RT$ for each gas in the $K_p$ expression gives
\[
K_p = K_c\,(RT)^{\Delta n}
\]
where $\Delta n$ is the change in the number of moles of \emph{gas} and $R=8.314\ \mathrm{J\,K^{-1}\,mol^{-1}}$ with $T$ in kelvin (pressures then in Pa, concentrations in $\mathrm{mol\,m^{-3}}$ --- take care with units). When $\Delta n = 0$, $K_p = K_c$. This relation is not required for routine GCE calculations but explains why both constants describe the same equilibrium. (Proof not examinable.)
\end{propriete}

\courssub{What the Magnitude of $K$ Tells Us}

The size of the equilibrium constant is a direct measure of \emph{where the equilibrium position lies}:
\begin{itemize}
\item $K \gg 1$ (e.g.\ $10^{10}$): products dominate at equilibrium; the reaction goes essentially to \textbf{completion}.
\item $K \ll 1$ (e.g.\ $10^{-10}$): reactants dominate; the reaction \textbf{hardly proceeds}.
\item $K \approx 1$: appreciable amounts of \textbf{both} reactants and products are present.
\end{itemize}

\begin{popfigure}
\begin{center}
\begin{tikzpicture}
\begin{axis}[
ybar, bar width=0.55cm, width=11.5cm, height=6.2cm,
ymin=0, ymax=6,
ylabel={$\log_{10} K$ (approx.)},
symbolic x coords={Haber, Ester, HI, Contact, Combustion},
xtick=data, x tick label style={font=\small, align=center},
y tick label style={font=\small},
nodes near coords, every node near coord/.append style={font=\small, color=popdark},
axis line style={popline}, tick style={popline},
]
\addplot[fill=popblue, draw=popdark] coordinates {(Haber,0.1) (Ester,0.6) (HI,1.7) (Contact,2) (Combustion,5)};
\end{axis}
\end{tikzpicture}
\legende{Illustrative magnitudes of $K$ (log scale) for familiar reactions at typical working temperatures: combustion of fuels ($K\gg1$, complete) versus the Haber process (small $K$, poor yield without optimisation). Values are indicative only --- $K$ varies with temperature.}
\end{center}
\end{popfigure}

\begin{attention}[$K$ says nothing about rate]
A huge $K$ does \emph{not} mean a fast reaction. The combustion of hydrogen has an enormous $K$, yet a hydrogen--oxygen mixture at room temperature shows no reaction without a spark, because the \emph{activation energy} is high. $K$ describes the equilibrium position; kinetics describes how quickly it is reached.
\end{attention}

\courssub{Factors Affecting the Equilibrium Constant}

In the previous section we saw that changing concentration, pressure or adding a catalyst \emph{shifts} the equilibrium position. Does the constant itself change?

\begin{propriete}[Only temperature changes $K$]
The value of an equilibrium constant for a given reaction depends \textbf{only on temperature}:
\begin{itemize}
\item Changing \textbf{concentration}: the system shifts to restore the ratio, but $K$ is unchanged.
\item Changing \textbf{pressure}: the position may shift (if $\Delta n \neq 0$), but $K$ is unchanged.
\item Adding a \textbf{catalyst}: equilibrium is reached faster, but the position and $K$ are unchanged.
\item Changing \textbf{temperature}: $K$ itself changes --- it increases with temperature for an \textbf{endothermic} forward reaction and decreases for an \textbf{exothermic} one.
\end{itemize}
\end{propriete}

\begin{aretenir}[Key points of this section]
\begin{itemize}
\item $K_c=\dfrac{[\ce{C}]^c[\ce{D}]^d}{[\ce{A}]^a[\ce{B}]^b}$ and $K_p=\dfrac{p_C^{c}p_D^{d}}{p_A^{a}p_B^{b}}$, using \emph{equilibrium} values only.
\item Deduce the units of $K_c$ and $K_p$ from the expression every time: $(\mathrm{mol\,dm^{-3}})^{\Delta n}$ and $(\text{pressure})^{\Delta n}$.
\item Use the ICE table to convert initial amounts into equilibrium amounts before substituting.
\item $p_A = x_A P_{\text{total}}$; $K_p = K_c(RT)^{\Delta n}$.
\item $K\gg1$: products favoured; $K\ll1$: reactants favoured.
\item Only a change in \textbf{temperature} alters the value of $K$.
\end{itemize}
\end{aretenir}

\begin{exercice}[Practice]
\begin{enumerate}
\item At equilibrium, a $\SI{1.0}{dm^3}$ flask contains $\SI{0.20}{mol}$ \ce{SO3}, $\SI{0.40}{mol}$ \ce{SO2} and $\SI{0.10}{mol}$ \ce{O2} for $\ce{2SO2(g) + O2(g) <=> 2SO3(g)}$. Calculate $K_c$, stating its units.
\item For $\ce{N2O4(g) <=> 2NO2(g)}$, starting from $\SI{1.0}{mol}$ of pure \ce{N2O4}, the mixture at equilibrium contains $\SI{0.40}{mol}$ of \ce{NO2} at a total pressure of $\SI{150}{kPa}$. Calculate $K_p$.
\item State, with reasons, the effect (if any) on the value of $K_c$ for the Haber equilibrium of: (a) doubling the pressure; (b) adding iron catalyst; (c) raising the temperature (forward reaction exothermic).
\end{enumerate}
\end{exercice}

\begin{corrige}[Exercise 1]
\begin{enumerate}
\item $[\ce{SO3}]=0.20$, $[\ce{SO2}]=0.40$, $[\ce{O2}]=0.10\ \mathrm{mol\,dm^{-3}}$.
$K_c=\dfrac{(0.20)^2}{(0.40)^2(0.10)}=\dfrac{0.040}{0.016}=2.5\ \mathrm{mol^{-1}\,dm^{3}}$ ($\Delta n = 2-3=-1$).
\item Change: $+0.40$ mol \ce{NO2} means $-0.20$ mol \ce{N2O4}. Equilibrium: $n_{\ce{N2O4}}=0.80$, $n_{\ce{NO2}}=0.40$, total $1.20$ mol. $p_{\ce{N2O4}}=\tfrac{0.80}{1.20}\times150=100\ \mathrm{kPa}$; $p_{\ce{NO2}}=50\ \mathrm{kPa}$. $K_p=\dfrac{50^2}{100}=25\ \mathrm{kPa}$.
\item (a) No change in $K_c$ (position shifts, constant does not). (b) No change (catalyst speeds attainment only). (c) $K_c$ \emph{decreases}, because raising temperature favours the endothermic (reverse) direction.
\end{enumerate}
\end{corrige}

\coursec{Heterogeneous Equilibria and the Effect of Temperature on K}

In the previous section we wrote the equilibrium law for reactions in which every species is a gas or every species is in aqueous solution. But many equilibria that matter in industry and in everyday life involve more than one \emph{phase}: think of limestone being roasted in a lime kiln, where a solid (\ce{CaCO3}) produces another solid (\ce{CaO}) and a gas (\ce{CO2}). How do we write $K_c$ and $K_p$ when solids or pure liquids are present? And we have used Le Chatelier's principle to \emph{predict} how temperature shifts an equilibrium; in this section we go further and ask: what does temperature do to the \emph{value} of $K$ itself?

\courssub{Homogeneous and Heterogeneous Equilibria}

\begin{definition}[Homogeneous and heterogeneous equilibria]
A \cle{homogeneous equilibrium} is an equilibrium in which all reactants and products are in the \textbf{same phase} (all gaseous, or all in the same solution). A \cle{heterogeneous equilibrium} is an equilibrium in which the reactants and products are present in \textbf{two or more different phases} (for example solids and gases, or a pure liquid and gases).
\end{definition}

The equilibrium $\ce{N2(g) + 3H2(g) <=> 2NH3(g)}$ is homogeneous: everything is gaseous. By contrast, the decomposition of calcium carbonate is heterogeneous:
\[
\ce{CaCO3(s) <=> CaO(s) + CO2(g)}
\]
because two solids and one gas coexist.

\begin{experience}[Observing a heterogeneous equilibrium]
Place a few grams of crushed limestone, \ce{CaCO3}, in a strong, sealed container fitted with a pressure gauge, and heat it strongly. The solid begins to decompose: $\ce{CaCO3(s) -> CaO(s) + CO2(g)}$. The gauge shows the pressure of \ce{CO2} rising, then becoming \textbf{constant} even though plenty of \ce{CaCO3} and \ce{CaO} remain. The system has reached equilibrium: the forward and reverse reactions continue, but the pressure of \ce{CO2} no longer changes. Remarkably, the equilibrium pressure of \ce{CO2} is the \emph{same} whether we start with one lump or ten lumps of limestone.
\end{experience}

This last observation is the key to the whole section: the amount of solid present does not affect the equilibrium at all.

\courssub{Why Pure Solids and Pure Liquids Are Omitted from $K_c$ and $K_p$}

\begin{propriete}[Rule for pure solids and pure liquids]
The \textbf{concentrations of pure solids and pure liquids are constant} (they equal the density divided by the molar mass, which does not change as more of the substance is added). They are therefore absorbed into the value of $K$ and \textbf{do not appear} in the expressions for $K_c$ or $K_p$. Only \textbf{gases} and \textbf{species in solution} are written.
\end{propriete}

\textbf{Justification.} Concentration means ``amount per unit volume of the phase''. For a pure solid such as \ce{CaCO3}, $[\ce{CaCO3}] = \dfrac{\text{density}}{M} = \dfrac{2.7\ \mathrm{g\,cm^{-3}}}{100\ \mathrm{g\,mol^{-1}}}$, a fixed number at a given temperature. Doubling the mass of solid doubles both the amount and the volume of the solid phase, so its concentration is unchanged. Since it never varies, we fold it into the constant. The same argument holds for a pure liquid such as water when it is the solvent or a pure liquid reactant.

\begin{attention}[Do not write solids or pure liquids in $K$]
A very common GCE error is to write $K_c = \dfrac{[\ce{CaO}][\ce{CO2}]}{[\ce{CaCO3}]}$. This is \textbf{wrong}. The correct expression is $K_c = [\ce{CO2}]$. Remember: \cle{solids and pure liquids are invisible} in equilibrium expressions.
\end{attention}

\courssub{Writing $K_c$ and $K_p$ for Heterogeneous Systems}

\begin{exemplebox}[The three classic examples]
\begin{itemize}
\item $\ce{CaCO3(s) <=> CaO(s) + CO2(g)}$:\quad $K_c = [\ce{CO2}]$ \quad and \quad $K_p = p(\ce{CO2})$.
\item $\ce{NH4Cl(s) <=> NH3(g) + HCl(g)}$:\quad $K_c = [\ce{NH3}][\ce{HCl}]$ \quad and \quad $K_p = p(\ce{NH3})\,p(\ce{HCl})$.
\item $\ce{C(s) + H2O(g) <=> CO(g) + H2(g)}$:\quad $K_c = \dfrac{[\ce{CO}][\ce{H2}]}{[\ce{H2O}]}$ \quad and \quad $K_p = \dfrac{p(\ce{CO})\,p(\ce{H2})}{p(\ce{H2O})}$.
\end{itemize}
In each case the solid is simply left out.
\end{exemplebox}

\begin{methode}[Writing $K_p$ for a decomposition producing a single gas]
\begin{enumerate}
\item Write the balanced equation with state symbols.
\item Identify the \textbf{only gaseous species}; ignore all solids and pure liquids.
\item Write $K_p$ as the partial pressure of that gas raised to its stoichiometric coefficient.
\item If the gas is the only species in the gas phase, its partial pressure \emph{equals the total measured pressure} of the system at equilibrium: $K_p = p_{\text{total}}$.
\item Give units: for $K_p = p(\ce{CO2})$ the units are those of pressure (e.g.\ atm or Pa).
\end{enumerate}
\end{methode}

\begin{popfigure}
\begin{tikzpicture}[scale=1.0]
% container
\draw[very thick, popink] (0,0) -- (0,3) -- (4,3) -- (4,0);
\draw[very thick, popink] (0,0) -- (4,0);
% solids at bottom
\fill[popblueL] (0.15,0) rectangle (3.85,0.55);
\draw[popink] (0.15,0) rectangle (3.85,0.55);
\node[popdark] at (2,0.28) {\small \ce{CaCO3}(s) and \ce{CaO}(s)};
% gas molecules
\foreach \x/\y in {0.5/1.2, 1.3/2.1, 2.1/1.5, 2.9/2.4, 3.5/1.3, 0.9/2.6, 1.8/2.7, 2.6/1.9, 3.3/2.0}{
  \draw[poporange, fill=poporange!40] (\x,\y) circle (0.07);}
\node[poporange, align=center] at (2,1.05) {\small \ce{CO2}(g)};
% equilibrium arrows
\draw[<->, thick, poppurple] (4.15,0.3) -- (4.75,0.3);
% pressure gauge
\draw[very thick, popink] (4,2.4) -- (5.0,2.4);
\draw[very thick, popink, fill=white] (5.6,2.4) circle (0.6);
\draw[->, thick, poppink] (5.6,2.4) -- (5.95,2.85);
\node[popdark] at (5.6,1.55) {\small pressure gauge};
\node[popdark] at (5.6,2.4) {\tiny $p$};
% heat
\foreach \x in {0.8,1.6,2.4,3.2}{
  \draw[thick, poppink] (\x,-0.55) .. controls (\x+0.1,-0.35) and (\x-0.1,-0.2) .. (\x,-0.05);}
\node[popdark] at (2,-0.85) {\small strong heating};
\end{tikzpicture}
\legende{Thermal decomposition of \ce{CaCO3} in a closed container. At equilibrium the gauge reads a constant pressure: $K_p = p(\ce{CO2})$, whatever the amounts of the two solids.}
\end{popfigure}

\begin{exoresolu}[Decomposition of ammonium chloride]
Solid ammonium chloride is heated in an evacuated \SI{5.0}{dm^3} vessel at $600\ \mathrm{K}$:
\[
\ce{NH4Cl(s) <=> NH3(g) + HCl(g)}
\]
At equilibrium the total pressure is $1.0\ \mathrm{atm}$. Calculate $K_p$ at this temperature, stating its units.
\tcblower
\textbf{Step 1 — Write $K_p$.} The solid is omitted:
\[
K_p = p(\ce{NH3})\,p(\ce{HCl})
\]
\textbf{Step 2 — Find the partial pressures.} The two gases are produced in a $1:1$ ratio and are the only gases present, so
\[
p(\ce{NH3}) = p(\ce{HCl}) = \tfrac{1}{2}\times 1.0 = 0.50\ \mathrm{atm}.
\]
\textbf{Step 3 — Compute.}
\[
K_p = 0.50 \times 0.50 = 0.25\ \mathrm{atm^2}.
\]
The units are (atm)$^2$ because two pressures are multiplied. Note that the volume of the vessel and the mass of \ce{NH4Cl} are irrelevant, provided some solid remains at equilibrium.
\end{exoresolu}

\begin{exoresolu}[From $K_p$ to $K_c$ for a heterogeneous system]
For the equilibrium $\ce{CaCO3(s) <=> CaO(s) + CO2(g)}$ at $1100\ \mathrm{K}$, $K_p = 0.40\ \mathrm{atm}$. Calculate $K_c$ at the same temperature, in $\mathrm{mol\,dm^{-3}}$. (Take $R = 0.082\ \mathrm{dm^3\,atm\,K^{-1}\,mol^{-1}}$.)
\tcblower
\textbf{Step 1 — Relate concentration and pressure.} For the only gaseous species,
\[
[\ce{CO2}] = \frac{n}{V} = \frac{p(\ce{CO2})}{RT}.
\]
\textbf{Step 2 — Substitute.}
\[
K_c = [\ce{CO2}] = \frac{0.40}{0.082 \times 1100} = \frac{0.40}{90.2} = 4.4\times 10^{-3}\ \mathrm{mol\,dm^{-3}}.
\]
This illustrates the general relation $K_p = K_c(RT)^{\Delta n_g}$, where here $\Delta n_g = 1$ (one mole of gas produced, none consumed).
\end{exoresolu}

\begin{savaistu}[Lime production in Cameroon]
Quicklime, \ce{CaO}, is produced by roasting limestone in kilns: $\ce{CaCO3(s) <=> CaO(s) + CO2(g)}$, an endothermic process. Because $K_p = p(\ce{CO2})$, the decomposition only proceeds significantly when the temperature is high enough for the equilibrium pressure of \ce{CO2} to exceed the pressure of \ce{CO2} in the kiln atmosphere. Kiln operators also \emph{remove} the \ce{CO2} continuously (it is swept away by the draught), so the reverse reaction cannot occur and the limestone decomposes completely. This is Le Chatelier's principle applied on an industrial scale — and the same chemistry underlies cement works such as those at Figuil in the North Region.
\end{savaistu}

\courssub{The Effect of Temperature on the Value of $K$}

We already know that changing concentration or pressure shifts the equilibrium \emph{position} but leaves $K$ unchanged. Temperature is different: it is the \textbf{only factor that changes the value of the equilibrium constant}.

\begin{propriete}[Temperature and $K$]
Raising the temperature \textbf{increases} $K$ for an \textbf{endothermic} reaction ($\Delta H > 0$) and \textbf{decreases} $K$ for an \textbf{exothermic} reaction ($\Delta H < 0$). Lowering the temperature has the opposite effect in each case. (The qualitative reasoning via Le Chatelier is examinable; the van 't Hoff equation is not.)
\end{propriete}

\textbf{Qualitative justification (Le Chatelier).} Treat heat as if it were a ``species'' in the equation. For an endothermic forward reaction, heat is effectively a \emph{reactant}: adding heat (raising $T$) pushes the equilibrium to the right, so more products form and $K = \dfrac{\text{products}}{\text{reactants}}$ grows. For an exothermic forward reaction, heat is a \emph{product}: raising $T$ pushes the equilibrium to the left, and $K$ falls.

\begin{exemplebox}[Two contrasting cases]
\begin{itemize}
\item Endothermic: $\ce{CaCO3(s) <=> CaO(s) + CO2(g)}$, $\Delta H > 0$. Heating the kiln \emph{increases} $K_p = p(\ce{CO2})$: decomposition is favoured at high temperature.
\item Exothermic: $\ce{N2(g) + 3H2(g) <=> 2NH3(g)}$, $\Delta H = -92\ \mathrm{kJ\,mol^{-1}}$. Raising the temperature \emph{decreases} $K_p$: the yield of ammonia falls. This is why the Haber process uses a \emph{compromise} temperature of about $450\ ^\circ\mathrm{C}$ — high enough for a reasonable rate, low enough for a reasonable yield.
\end{itemize}
\end{exemplebox}

\begin{popfigure}
\begin{center}
\begin{tikzpicture}
\begin{axis}[
  width=11cm, height=7cm,
  xlabel={Temperature $T$ / K},
  ylabel={Equilibrium constant $K$},
  xmin=300, xmax=900, ymin=0, ymax=10,
  xtick={300,450,600,750,900},
  ytick=\empty,
  axis lines=left,
  legend style={at={(0.03,0.97)}, anchor=north west, draw=popline},
  every axis plot/.append style={very thick},
]
\addplot[domain=300:900, samples=100, poppink] {8*exp(-(x-300)/350)};
\addlegendentry{Exothermic ($\Delta H<0$): $K$ falls as $T$ rises}
\addplot[domain=300:900, samples=100, popblue] {0.4*exp((x-300)/300)};
\addlegendentry{Endothermic ($\Delta H>0$): $K$ rises as $T$ rises}
\end{axis}
\end{tikzpicture}
\end{center}
\legende{Variation of $K$ with temperature. For an exothermic reaction $K$ decreases as $T$ increases; for an endothermic reaction $K$ increases as $T$ increases.}
\end{popfigure}

\begin{attention}[Position versus value: a subtle but crucial distinction]
Do not confuse the effect of temperature on the \emph{equilibrium position} with its effect on the \emph{value of $K$} — the two are perfectly consistent. When $K$ increases, the position shifts towards the products; when $K$ decreases, it shifts towards the reactants. Temperature is special: it changes $K$ itself. Concentration, pressure and catalysts never change $K$; they only change the position (or, for a catalyst, merely the \emph{speed} at which equilibrium is reached). Saying ``the catalyst increases $K$'' is a serious error that costs marks.
\end{attention}

\courssub{The van 't Hoff Relationship (Extension)}

\begin{savaistu}[A quantitative link between $T$, $\Delta H$ and $K$]
Beyond the GCE syllabus, the quantitative link is the \cle{van 't Hoff equation}:
\[
\ln K = -\frac{\Delta H}{R}\cdot\frac{1}{T} + \text{constant}
\]
A graph of $\ln K$ against $\dfrac{1}{T}$ is a \textbf{straight line} whose gradient is $-\dfrac{\Delta H}{R}$. For an endothermic reaction ($\Delta H > 0$) the gradient is negative; for an exothermic reaction it is positive. Chemists use this plot to measure enthalpy changes experimentally. You only need to know it \emph{qualitatively}: it confirms that $K$ depends exponentially on temperature, which is why even modest heating can transform the yield of a kiln or a reactor.
\end{savaistu}

\begin{exoresolu}[Predicting the effect of heating]
For the water-gas reaction $\ce{C(s) + H2O(g) <=> CO(g) + H2(g)}$, $\Delta H = +131\ \mathrm{kJ\,mol^{-1}}$.
\begin{enumerate}
\item Write expressions for $K_c$ and $K_p$.
\item State and explain the effect of raising the temperature on (i) the value of $K_p$ and (ii) the yield of hydrogen.
\end{enumerate}
\tcblower
\textbf{1.} The carbon is a pure solid, so it is omitted:
\[
K_c = \frac{[\ce{CO}][\ce{H2}]}{[\ce{H2O}]}, \qquad
K_p = \frac{p(\ce{CO})\,p(\ce{H2})}{p(\ce{H2O})}.
\]
\textbf{2.} The forward reaction is endothermic ($\Delta H > 0$). By Le Chatelier's principle, raising the temperature shifts the equilibrium in the direction that absorbs heat, i.e.\ to the right. Consequently (i) $K_p$ \textbf{increases} (more products relative to reactants at equilibrium) and (ii) the equilibrium yield of hydrogen \textbf{increases}. This is why industrial hydrogen production by this route is carried out at high temperature.
\end{exoresolu}

\begin{exercice}[Limestone in a sealed vessel]
A sample of \ce{CaCO3} is placed in a sealed, evacuated vessel and heated to $1100\ \mathrm{K}$. At equilibrium the pressure in the vessel is $0.40\ \mathrm{atm}$.
\begin{enumerate}
\item Write the expression for $K_p$ and calculate its value, with units.
\item A second, larger lump of \ce{CaCO3} is used under identical conditions. State, with a reason, the new equilibrium pressure.
\item The decomposition is endothermic. State and explain the effect on $K_p$ of heating the vessel to $1200\ \mathrm{K}$.
\end{enumerate}
\end{exercice}

\begin{corrige}[Exercise 1]
\textbf{1.} $\ce{CaCO3(s) <=> CaO(s) + CO2(g)}$; solids are omitted, so $K_p = p(\ce{CO2})$. Since \ce{CO2} is the only gas, $p(\ce{CO2}) = 0.40\ \mathrm{atm}$, hence $K_p = 0.40\ \mathrm{atm}$.

\textbf{2.} The pressure is still $0.40\ \mathrm{atm}$. The concentration (and hence the contribution) of a pure solid is constant, so the equilibrium pressure of \ce{CO2} depends only on temperature, not on the amount of solid — provided at least some of each solid remains.

\textbf{3.} $K_p$ increases. The forward reaction is endothermic, so by Le Chatelier's principle raising the temperature shifts the equilibrium to the right, increasing $p(\ce{CO2})$ and therefore $K_p$.
\end{corrige}

\begin{exercice}[Steam over hot carbon]
Steam is passed over red-hot carbon in a closed reactor at $1000\ \mathrm{K}$:
\[
\ce{C(s) + H2O(g) <=> CO(g) + H2(g)} \qquad \Delta H = +131\ \mathrm{kJ\,mol^{-1}}
\]
\begin{enumerate}
\item Write the expression for $K_p$.
\item At equilibrium the partial pressures are $p(\ce{H2O}) = 0.20\ \mathrm{atm}$, $p(\ce{CO}) = 0.80\ \mathrm{atm}$ and $p(\ce{H2}) = 0.80\ \mathrm{atm}$. Calculate $K_p$, with units.
\item State two changes that would increase the equilibrium yield of hydrogen, justifying each.
\end{enumerate}
\end{exercice}

\begin{corrige}[Exercise 2]
\textbf{1.} Carbon is a pure solid and is omitted: $K_p = \dfrac{p(\ce{CO})\,p(\ce{H2})}{p(\ce{H2O})}$.

\textbf{2.} $K_p = \dfrac{0.80 \times 0.80}{0.20} = 3.2\ \mathrm{atm}$. The units are atm because there are two pressure terms in the numerator and one in the denominator ($\Delta n_g = +1$).

\textbf{3.} (i) \textbf{Raise the temperature}: the forward reaction is endothermic, so by Le Chatelier's principle the equilibrium shifts to the right and $K_p$ increases. (ii) \textbf{Decrease the total pressure} (or increase the volume): there are two moles of gas on the right and one on the left, so the equilibrium shifts to the side with more moles of gas, i.e.\ towards the products. (Removing \ce{CO} or \ce{H2} as they form would also work.)
\end{corrige}

\begin{aretenir}[Key points of this section]
\begin{itemize}
\item A heterogeneous equilibrium involves species in more than one phase.
\item \textbf{Pure solids and pure liquids never appear} in $K_c$ or $K_p$; their concentrations are constant.
\item For $\ce{CaCO3(s) <=> CaO(s) + CO2(g)}$: $K_p = p(\ce{CO2})$ — the basis of lime manufacture.
\item $K_p$ and $K_c$ remain linked by $K_p = K_c(RT)^{\Delta n_g}$, counting only gaseous species in $\Delta n_g$.
\item Temperature is the \textbf{only} factor that changes the value of $K$: raising $T$ increases $K$ for an endothermic reaction and decreases it for an exothermic one.
\item The van 't Hoff equation, $\ln K = -\dfrac{\Delta H}{R}\cdot\dfrac{1}{T} + \text{constant}$, makes this quantitative (extension, qualitative use only).
\end{itemize}
\end{aretenir}

\coursec{Redox Systems and Balancing Redox Equations}

In the previous section we saw that the position of a heterogeneous equilibrium, and indeed the value of any equilibrium constant, depends on temperature through the enthalpy change of the reaction. We treated equilibrium as a competition between forward and backward processes. Many of the most important reactions in chemistry — the rusting of iron roofing sheets in Douala, the working of a car battery, the extraction of metals from their ores — are equilibria of a special kind: they involve a \cle{transfer of electrons} from one species to another. Such reactions are called \cle{redox reactions}, and this section builds the language and the techniques needed to describe, balance and (in the next sections) quantify them.

\courssub{Oxidation, Reduction and Redox Agents}

\textbf{Intuition from a simple experiment.} When a strip of zinc metal is dipped into a blue solution of copper(II) sulphate, the blue colour fades and a reddish-brown deposit of copper forms on the zinc. Analysis shows that zinc has gone into solution as \ce{Zn^{2+}} ions while \ce{Cu^{2+}} ions have become copper metal:
\[ \ce{Zn(s) + Cu^{2+}(aq) -> Zn^{2+}(aq) + Cu(s)} \]
Each zinc atom has \emph{lost} two electrons; each copper(II) ion has \emph{gained} two electrons. The two processes cannot occur separately: the electrons lost by zinc are exactly those gained by copper. This coupling of electron loss and electron gain is the essence of every redox reaction.

\begin{definition}[Oxidation and Reduction]
\cle{Oxidation} is the \textbf{loss of electrons} by a species; equivalently, it is an \textbf{increase in oxidation number}. \cle{Reduction} is the \textbf{gain of electrons} by a species; equivalently, it is a \textbf{decrease in oxidation number}. The mnemonic \textbf{OILRIG} summarises this: \emph{Oxidation Is Loss, Reduction Is Gain} (of electrons).
\end{definition}

\begin{definition}[Oxidising and Reducing Agents]
An \cle{oxidising agent} (oxidant) is a species that \textbf{accepts electrons} and is therefore itself \emph{reduced} during the reaction. A \cle{reducing agent} (reductant) is a species that \textbf{donates electrons} and is therefore itself \emph{oxidised}.
\end{definition}

In the zinc–copper reaction above, \ce{Zn} is the reducing agent (it is oxidised to \ce{Zn^{2+}}) and \ce{Cu^{2+}} is the oxidising agent (it is reduced to \ce{Cu}).

\begin{attention}[The agent is named after what it does to the other species]
Students frequently confuse the labels. The \emph{oxidising agent} causes oxidation but is itself \emph{reduced}. Always ask: "Which species gains electrons?" — that species is reduced and is the oxidising agent.
\end{attention}

\courssub{Oxidation Numbers: the Book-keeping Tool}

To decide whether oxidation or reduction has occurred, we assign an \cle{oxidation number} (oxidation state) to each atom, using the conventional rules:
\begin{itemize}
\item an uncombined element has oxidation number $0$;
\item a monatomic ion has an oxidation number equal to its charge;
\item oxygen is usually $-2$ (except peroxides, $-1$, and \ce{OF2}, $+2$);
\item hydrogen is usually $+1$ (except metal hydrides such as \ce{NaH}, $-1$);
\item the sum of oxidation numbers equals the overall charge of the species.
\end{itemize}

\begin{exemplebox}[Oxidation number of Mn in \ce{MnO4^-}]
Let the oxidation number of Mn be $x$. Then $x + 4(-2) = -1$, so $x = +7$. Manganese is in the $+7$ state in the permanganate ion. When \ce{MnO4^-} is converted to \ce{Mn^{2+}} (oxidation number $+2$), manganese has been \emph{reduced} — its oxidation number has decreased by 5, corresponding to a gain of 5 electrons.
\end{exemplebox}

\courssub{Redox Systems: Half-Cells, Cells, Anode and Cathode}

\textbf{Separating the two processes.} In the zinc–copper reaction, electrons pass directly from zinc to \ce{Cu^{2+}}. But suppose we physically separate the two processes: place a zinc rod in a solution of \ce{Zn^{2+}} ions in one beaker, and a copper rod in a solution of \ce{Cu^{2+}} ions in another, then connect the rods by a wire and the solutions by a salt bridge. Electrons now flow \emph{through the wire} from zinc to copper, and we can make them do useful work. Each beaker is a \cle{half-cell}; the combination is an \cle{electrochemical cell}.

\begin{definition}[Half-cell and electrochemical cell]
A \cle{half-cell} consists of an \cle{electrode} (an electronic conductor, often a metal or inert graphite/platinum) in contact with an electrolyte containing the oxidised and reduced forms of a redox couple. An \cle{electrochemical cell} is formed by coupling two half-cells so that electrons can flow through an external circuit while ions complete the circuit internally (through a salt bridge or porous partition).
\end{definition}

\begin{definition}[Anode and cathode]
The \cle{anode} is the electrode at which \textbf{oxidation} occurs; the \cle{cathode} is the electrode at which \textbf{reduction} occurs. Mnemonic: \emph{AN OX, RED CAT} — \textbf{AN}ode = \textbf{OX}idation; \textbf{RED}uction at the \textbf{CAT}hode.
\end{definition}

\begin{attention}[The anode is always the site of oxidation]
This definition is universal: it holds for galvanic (voltaic) cells \emph{and} for electrolytic cells. In a galvanic cell the anode is the \emph{negative} terminal, whereas in electrolysis the anode is connected to the \emph{positive} terminal of the supply — but in both cases oxidation occurs at the anode. Never define the anode as "the positive electrode".
\end{attention}

\begin{popfigure}
\begin{center}
\begin{tikzpicture}[scale=1.0, every node/.style={font=\small}]
% beaker
\draw[thick, popblue] (-2.2,0) -- (-2.2,2.6) (-2.2,0) -- (2.2,0) (2.2,0) -- (2.2,2.6);
% solution
\fill[popblueL, opacity=0.6] (-2.15,0.05) rectangle (2.15,1.7);
\draw[popblue] (-2.15,1.7) -- (2.15,1.7);
\node[popdark] at (0,0.7) {solution of \ce{M^{n+}} ions};
% electrode
\fill[popmuted] (-0.35,1.2) rectangle (0.35,3.6);
\node[popdark, align=center] at (0,4.0) {metal electrode \ce{M}};
% ions
\node[popblue] at (-1.4,1.1) {\ce{M^{n+}}};
\node[popblue] at (1.3,0.5) {\ce{M^{n+}}};
\node[popblue] at (1.5,1.35) {\ce{M^{n+}}};
% electron arrow
\draw[->, very thick, poporange] (0.35,3.1) .. controls (1.6,3.6) .. (2.9,3.1);
\node[poporange, align=center] at (1.7,3.9) {$n e^-$ (oxidation)};
\draw[->, very thick, popgreen] (2.9,2.55) .. controls (1.6,2.1) .. (0.35,2.55);
\node[popgreen, align=center] at (1.7,1.95) {$n e^-$ (reduction)};
% equilibrium label
\node[popdark, align=center] at (0,-0.6) {\ensuremath{\mathrm{M^{n^{+}}(aq) + ne^{-} \rightleftharpoons  M(s)}}};
\end{tikzpicture}
\end{center}
\legende{A generic metal/metal-ion half-cell: the electrode and the ions in solution are in dynamic redox equilibrium, \ensuremath{\mathrm{M^{n^{+}} + ne^{-} \rightleftharpoons  M}}.}
\end{popfigure}

\courssub{Redox Potential and Redox Equilibria}

When a metal electrode is dipped into a solution of its ions, an equilibrium is rapidly established:
\[ \ce{M^{n+}(aq) + n e^- <=> M(s)} \]
Some metals (like zinc) "push" electrons onto the electrode strongly — the equilibrium lies towards electron release; others (like copper) hold their electrons more tightly. A \cle{potential difference} therefore develops between the electrode and the solution. This leads to the key concept:

\begin{definition}[Redox (electrode) potential]
The \cle{redox potential} (electrode potential) of a half-cell is a measure of the \textbf{tendency of the oxidised form of a redox couple to gain electrons} (be reduced). The more positive the redox potential, the greater the tendency of the couple \ensuremath{\mathrm{Ox + ne^{-} \rightleftharpoons  Red}} to proceed to the right, i.e.\ the stronger the oxidising agent \ce{Ox} and the weaker the reducing agent \ce{Red}.
\end{definition}

Every redox reaction can be regarded as the \textbf{combination of two half-equations}, one for each couple. For the zinc–copper reaction:
\begin{align*}
\text{oxidation:} \quad & \ce{Zn(s) -> Zn^{2+}(aq) + 2e^-} \\
\text{reduction:} \quad & \ce{Cu^{2+}(aq) + 2e^- -> Cu(s)} \\
\text{overall:} \quad & \ce{Zn(s) + Cu^{2+}(aq) -> Zn^{2+}(aq) + Cu(s)}
\end{align*}
The electrons cancel because the number lost equals the number gained — a point we shall exploit systematically when balancing equations.

\courssub{Writing Half-Equations for Simple Systems}

A \cle{half-equation} shows the oxidised and reduced forms of one couple, balanced for atoms and for charge using electrons. Common examples you must be able to write quickly:
\begin{align*}
\ce{Zn^{2+}(aq) + 2e^- &<=> Zn(s)} &
\ce{Fe^{3+}(aq) + e^- &<=> Fe^{2+}(aq)} \\
\ce{Cl2(g) + 2e^- &<=> 2Cl^-(aq)} &
\ce{I2(aq) + 2e^- &<=> 2I^-(aq)} \\
\ce{2H^+(aq) + 2e^- &<=> H2(g)} &
\ce{Cu^{2+}(aq) + 2e^- &<=> Cu(s)}
\end{align*}
For oxoanions in acidic solution, water and \ce{H+} must be included, e.g.\
\[ \ce{MnO4^-(aq) + 8H^+(aq) + 5e^- <=> Mn^{2+}(aq) + 4H2O(l)} \]
\[ \ce{Cr2O7^{2-}(aq) + 14H^+(aq) + 6e^- <=> 2Cr^{3+}(aq) + 7H2O(l)} \]

\courssub{Balancing Redox Equations: the Ion-Electron (Half-Equation) Method}

Many redox equations are far too complicated to balance by inspection. The \cle{ion-electron method} (valid in acidic medium) is systematic and always works.

\begin{methode}[Balancing a redox equation in acidic medium]
\begin{enumerate}
\item \textbf{Identify} the species oxidised and reduced (use oxidation numbers) and write the two skeleton half-equations.
\item \textbf{Balance the atoms} other than O and H in each half-equation.
\item \textbf{Balance oxygen} by adding \ce{H2O} molecules to the side deficient in O.
\item \textbf{Balance hydrogen} by adding \ce{H+} ions to the side deficient in H.
\item \textbf{Balance charge} by adding electrons, $e^-$, to the more positive side.
\item \textbf{Multiply} each half-equation so that the numbers of electrons are equal, then \textbf{add} the half-equations.
\item \textbf{Cancel} electrons, \ce{H+} and \ce{H2O} appearing on both sides; check atoms \emph{and} charges balance.
\end{enumerate}
\end{methode}

\begin{exoresolu}[Balancing \ce{MnO4^- + Fe^{2+} -> Mn^{2+} + Fe^{3+}} in acidic solution]
Statement: Acidified potassium permanganate oxidises iron(II) ions to iron(III) ions, the permanganate being reduced to manganese(II). Construct the balanced overall ionic equation.
\tcblower
\textbf{Step 1–5, reduction half-equation.} Mn goes from $+7$ in \ce{MnO4^-} to $+2$ in \ce{Mn^{2+}} (gain of $5e^-$). Balance O with \ce{H2O}, H with \ce{H+}:
\[ \ce{MnO4^-(aq) + 8H^+(aq) + 5e^- -> Mn^{2+}(aq) + 4H2O(l)} \]
Check charge: left $-1 + 8 - 5 = +2$; right $+2$. \checkmark

\textbf{Step 1–5, oxidation half-equation.} Fe goes from $+2$ to $+3$ (loss of $1e^-$):
\[ \ce{Fe^{2+}(aq) -> Fe^{3+}(aq) + e^-} \]

\textbf{Step 6.} Multiply the iron half-equation by 5 and add:
\[ \ce{MnO4^- + 8H^+ + 5e^- + 5Fe^{2+} -> Mn^{2+} + 4H2O + 5Fe^{3+} + 5e^-} \]

\textbf{Step 7.} Cancel the $5e^-$:
\[ \boxed{\ce{MnO4^-(aq) + 8H^+(aq) + 5Fe^{2+}(aq) -> Mn^{2+}(aq) + 5Fe^{3+}(aq) + 4H2O(l)}} \]
Final check — atoms: Mn 1/1, O 4/4, H 8/8, Fe 5/5; charge: left $-1+8+10 = +17$, right $+2+15 = +17$. \checkmark
\end{exoresolu}

\begin{popfigure}
\begin{center}
\begin{tikzpicture}[node distance=0.55cm, every node/.style={font=\small},
box/.style={draw=popblue, thick, rounded corners, fill=popblueL, text width=4.6cm, align=center, inner sep=4pt},
exbox/.style={draw=poporange, thick, rounded corners, fill=white, text width=5.4cm, align=center, inner sep=4pt},
arr/.style={->, thick, popdark}]
\node[box] (s1) {\textbf{1.} Write skeleton half-equations};
\node[box, below=of s1] (s2) {\textbf{2.} Balance atoms other than O, H};
\node[box, below=of s2] (s3) {\textbf{3.} Balance O with \ce{H2O}};
\node[box, below=of s3] (s4) {\textbf{4.} Balance H with \ce{H+}};
\node[box, below=of s4] (s5) {\textbf{5.} Balance charge with $e^-$};
\node[box, below=of s5] (s6) {\textbf{6.} Equalise $e^-$, add, cancel};
\node[exbox, right=1.1cm of s1, align=center] (e1){\ce{MnO4^- -> Mn^{2+}}\\ \ce{Fe^{2+} -> Fe^{3+}}};
\node[exbox, right=1.1cm of s3] (e3) {\ce{MnO4^- -> Mn^{2+} + 4H2O}};
\node[exbox, right=1.1cm of s4] (e4) {\ce{MnO4^- + 8H+ -> Mn^{2+} + 4H2O}};
\node[exbox, right=1.1cm of s5, align=center] (e5){\ce{MnO4^- + 8H+ + 5e- -> Mn^{2+} + 4H2O}\\ \ce{Fe^{2+} -> Fe^{3+} + e-}};
\node[exbox, right=1.1cm of s6, align=center] (e6){$\times 5$ on Fe line, add:\\ \ce{MnO4^- + 8H+ + 5Fe^{2+} -> Mn^{2+} + 5Fe^{3+} + 4H2O}};
\draw[arr] (s1) -- (s2); \draw[arr] (s2) -- (s3); \draw[arr] (s3) -- (s4);
\draw[arr] (s4) -- (s5); \draw[arr] (s5) -- (s6);
\draw[arr, dashed, popmuted] (s1) -- (e1); \draw[arr, dashed, popmuted] (s3) -- (e3);
\draw[arr, dashed, popmuted] (s4) -- (e4); \draw[arr, dashed, popmuted] (s5) -- (e5);
\draw[arr, dashed, popmuted] (s6) -- (e6);
\end{tikzpicture}
\end{center}
\legende{Flowchart of the ion-electron method, applied to the \ce{MnO4^-}/\ce{Fe^{2+}} reaction in acidic medium.}
\end{popfigure}

\begin{attention}[No electrons in the final equation]
Electrons are intermediates of the method, not reactants or products. If $e^-$ appears in your final overall equation, you have forgotten to multiply and cancel — a very common GCE error. Also check that the \emph{total charge} is the same on both sides, not only the atoms.
\end{attention}

\textbf{A note on alkaline medium.} If a reaction occurs in alkaline solution, first balance it exactly as for acidic medium, then add to \emph{both} sides as many \ce{OH^-} ions as there are \ce{H+} ions; each \ce{H+ + OH^-} pair becomes \ce{H2O}. Cancel any water molecules that then appear on both sides. For example, the oxidation of \ce{Cr^{3+}} to \ce{CrO4^{2-}} by hydrogen peroxide in alkaline solution gives
\[ \ce{2Cr^{3+}(aq) + 3H2O2(aq) + 10OH^-(aq) -> 2CrO4^{2-}(aq) + 8H2O(l)} \]
At GCE Advanced Level, acidic-medium balancing is by far the most frequently examined; the alkaline conversion is a useful extension.

\courssub{Balancing Using Oxidation Numbers}

An alternative, equally acceptable, method uses oxidation numbers directly: the \textbf{total increase in oxidation number must equal the total decrease}.

\begin{exoresolu}[Balancing \ce{Cu + HNO3 -> Cu(NO3)2 + NO + H2O}]
Statement: Dilute nitric acid oxidises copper metal, being itself reduced to nitrogen monoxide. Balance the equation.
\tcblower
Oxidation number changes: Cu: $0 \to +2$, an increase of 2. N (in the \ce{HNO3} that is reduced): $+5 \to +2$ in \ce{NO}, a decrease of 3. To equalise, take 3 Cu (total increase $3 \times 2 = 6$) and 2 N reduced (total decrease $2 \times 3 = 6$):
\[ \ce{3Cu + 2HNO3 -> 3Cu(NO3)2 + 2NO + H2O} \]
Now balance the remaining nitrate: $3\,\ce{Cu(NO3)2}$ needs $6$ more \ce{NO3^-}, i.e.\ 6 further \ce{HNO3}, giving 8 \ce{HNO3} in total; then balance H and O with water:
\[ \boxed{\ce{3Cu(s) + 8HNO3(aq) -> 3Cu(NO3)2(aq) + 2NO(g) + 4H2O(l)}} \]
Check: Cu 3/3, H 8/8, N 8/8, O 24/24. \checkmark
\end{exoresolu}

\courssub{Disproportionation: a Useful Extra}

\begin{definition}[Disproportionation]
\cle{Disproportionation} is a redox reaction in which the \textbf{same element, in the same species, is simultaneously oxidised and reduced} — its oxidation number both rises and falls.
\end{definition}

Classic examples include the reaction of chlorine with cold dilute alkali,
\[ \ce{Cl2(aq) + 2OH^-(aq) -> Cl^-(aq) + ClO^-(aq) + H2O(l)} \]
in which chlorine (oxidation number $0$) is reduced to \ce{Cl^-} ($-1$) and oxidised to \ce{ClO^-} ($+1$), and the decomposition of hydrogen peroxide,
\[ \ce{2H2O2(aq) -> 2H2O(l) + O2(g)} \]
in which oxygen ($-1$) goes to $-2$ (in water) and $0$ (in \ce{O2}). Disproportionation questions appear regularly at GCE: the key skill is recognising that one reactant plays the roles of \emph{both} oxidising and reducing agent.

\begin{savaistu}[Redox in everyday Cameroonian life]
The lead–acid battery that starts taxis and motorbikes across Cameroon works on the redox couples \ce{PbO2/PbSO4} and \ce{Pb/PbSO4}; the rusting of corrugated iron roofs in the humid coastal climate is the slow oxidation of iron by atmospheric oxygen; and the iodine value of palm oil, important in the agro-food industry, is measured by a redox titration with iodine. Redox chemistry is never far away.
\end{savaistu}

\begin{aretenir}[Key points of this section]
\begin{itemize}
\item Oxidation = loss of electrons / increase in oxidation number; reduction = gain of electrons / decrease in oxidation number (OILRIG).
\item The oxidising agent is reduced; the reducing agent is oxidised.
\item A half-cell is an electrode in contact with a redox couple, \ensuremath{\mathrm{Ox + ne^{-} \rightleftharpoons  Red}}; two half-cells make an electrochemical cell.
\item Oxidation always occurs at the anode, reduction at the cathode — in every type of cell.
\item The redox potential measures the tendency of a couple to gain electrons.
\item Balance redox equations by half-equations (atoms, then O with \ce{H2O}, then H with \ce{H+}, then charge with $e^-$) or by oxidation numbers; electrons must cancel in the final equation.
\end{itemize}
\end{aretenir}

\begin{exercice}[Practice]
\begin{enumerate}
\item Assign oxidation numbers to S in \ce{SO2}, \ce{SO3^{2-}} and \ce{SO4^{2-}}, and state whether the conversion \ce{SO2 -> SO4^{2-}} is an oxidation or a reduction.
\item Write balanced half-equations for: (a) \ce{Br2} reduced to \ce{Br^-}; (b) \ce{Sn^{2+}} oxidised to \ce{Sn^{4+}}.
\item Balance, in acidic solution: \ce{Cr2O7^{2-}(aq) + Fe^{2+}(aq) -> Cr^{3+}(aq) + Fe^{3+}(aq)}.
\item Show that the reaction \ce{2Cu^+(aq) -> Cu(s) + Cu^{2+}(aq)} is a disproportionation.
\end{enumerate}
\end{exercice}

\begin{corrige}[Exercise 1]
\textbf{1.} S in \ce{SO2}: $x + 2(-2) = 0 \Rightarrow x = +4$. In \ce{SO3^{2-}}: $x - 6 = -2 \Rightarrow x = +4$. In \ce{SO4^{2-}}: $x - 8 = -2 \Rightarrow x = +6$. The change $+4 \to +6$ is an increase, so \ce{SO2} is \emph{oxidised}.

\textbf{2.} (a) \ce{Br2(aq) + 2e^- -> 2Br^-(aq)}; (b) \ce{Sn^{2+}(aq) -> Sn^{4+}(aq) + 2e^-}.

\textbf{3.} Reduction: \ce{Cr2O7^{2-} + 14H+ + 6e^- -> 2Cr^{3+} + 7H2O}. Oxidation: \ce{Fe^{2+} -> Fe^{3+} + e^-}, multiplied by 6. Overall:
\[ \ce{Cr2O7^{2-}(aq) + 14H^+(aq) + 6Fe^{2+}(aq) -> 2Cr^{3+}(aq) + 6Fe^{3+}(aq) + 7H2O(l)} \]
Check charge: left $-2 + 14 + 12 = +24$; right $+6 + 18 = +24$. \checkmark

\textbf{4.} In \ce{Cu^+}, copper has oxidation number $+1$. One \ce{Cu^+} is reduced to \ce{Cu} ($0$) while the other is oxidised to \ce{Cu^{2+}} ($+2$): the same element in the same species is both oxidised and reduced — a disproportionation.
\end{corrige}

\coursec{Types of Half-Cells and Standard Electrode Potentials}

In the previous section we learnt to recognise redox systems, to write half-equations and to balance full redox equations. We saw, for example, that when a strip of zinc is dipped into a solution of copper(II) sulphate, electrons flow spontaneously from zinc atoms to copper(II) ions. A natural question now arises: \emph{how strongly} does each redox couple ``pull'' or ``push'' electrons? Can we attach a number to each couple that tells us how powerful an oxidising or reducing agent it is? The answer is yes, and that number is the \cle{standard electrode potential}, $E^{\theta}$. To measure it, we must first understand how a redox couple is turned into a physical device called a \cle{half-cell}, and we must agree on a universal reference against which all half-cells are compared. That is the business of this section.

\courssub{From a Redox Couple to a Half-Cell}

\begin{definition}[Half-cell]
A \cle{half-cell} is an assembly consisting of an \cle{electrode} (an electronic conductor) in contact with an \cle{electrolyte} containing the two species of a redox couple (the oxidised form and the reduced form). A half-cell is the physical realisation of one half-equation, for example $\ce{Cu^{2+}(aq) + 2e- <=> Cu(s)}$.
\end{definition}

A half-cell alone cannot drive a sustained current: electrons must have somewhere to go. Only when two half-cells are connected (through a wire and a salt bridge) does a complete electrochemical cell exist, as we shall study in the next section. Here we focus on the structure of the half-cell itself. Depending on the physical nature of the two members of the redox couple, three types of half-cell are distinguished, exactly as listed in the syllabus.

\courssub{The Three Types of Half-Cells}

\textbf{Type 1: Metal / metal-ion systems.}\par
This is the simplest case: the reduced form of the couple is a solid metal, which is itself an excellent conductor of electricity. The metal strip serves \emph{both} as the reduced species and as the electrode. It is simply dipped into a solution of its own ions. Examples:
\begin{itemize}
\item $\ce{Zn^{2+}(aq) + 2e- <=> Zn(s)}$: a zinc rod in zinc sulphate solution;
\item $\ce{Cu^{2+}(aq) + 2e- <=> Cu(s)}$: a copper strip in copper(II) sulphate solution;
\item $\ce{Ag+(aq) + e- <=> Ag(s)}$: a silver wire in silver nitrate solution.
\end{itemize}

\textbf{Type 2: Non-metal / non-metal-ion systems.}\par
Here neither member of the couple is a metallic conductor. A common example is the chlorine / chloride couple $\ce{Cl2(g) + 2e- <=> 2Cl-(aq)}$: chlorine is a gas and chloride ions are in solution; there is no solid metal to carry electrons into or out of the system. The solution is to bubble the gas over an \cle{inert electrode}, usually platinum, which dips into the solution containing the ions. The hydrogen half-cell $\ce{2H+(aq) + 2e- <=> H2(g)}$ (often written $\ce{H+(aq) + e- <=> 1/2 H2(g)}$) is built in exactly the same way, with hydrogen gas bubbled over platinum. Other examples include the oxygen / hydroxide couple $\ce{O2(g) + 2H2O(l) + 4e- <=> 4OH-(aq)}$.

\begin{attention}[Why platinum?]
Platinum is used as an \cle{inert electrode} whenever the half-cell contains no metallic conductor. Platinum is chosen because it conducts electricity, is chemically unreactive (it does not take part in the redox reaction), and its surface catalyses the electron transfer between the dissolved/gaseous species and the external circuit. A very common exam mistake is to say that ``platinum is one of the reactants'' --- it is not; it appears in \emph{no} half-equation.
\end{attention}

\textbf{Type 3: Ion--ion systems.}\par
Here both members of the couple are ions in the \emph{same solution}, for example $\ce{Fe^{3+}(aq) + e- <=> Fe^{2+}(aq)}$ or $\ce{MnO4-(aq) + 8H+(aq) + 5e- <=> Mn^{2+}(aq) + 4H2O(l)}$. Again there is no metal to act as electrode, so a platinum wire is dipped into a solution containing \emph{both} ions of the couple. Electrons are exchanged between the ions and the external circuit through the inert platinum surface.

\begin{popfigure}
\begin{tikzpicture}[scale=0.92, every node/.style={font=\small}]
% ---- Type 1 : Cu/Cu2+ ----
\draw[thick] (-0.9,0) -- (-0.9,-2.4) -- (0.9,-2.4) -- (0.9,0);
\fill[popblueL] (-0.86,-2.36) rectangle (0.86,-0.4);
\draw[fill=popgold] (-0.12,-2.3) rectangle (0.12,0.6);
\node at (0,0.85) {\textbf{Cu}};
\node[align=center] at (0,-1.5) {\ce{Cu^{2+}(aq)}};
\node[align=center] at (0,-2.9) {Metal / metal-ion\\ \ce{Cu^{2+} + 2e- <=> Cu}};
% ---- Type 2 : Cl2/Cl- ----
\begin{scope}[xshift=4.2cm]
\draw[thick] (-0.9,0) -- (-0.9,-2.4) -- (0.9,-2.4) -- (0.9,0);
\fill[popblueL] (-0.86,-2.36) rectangle (0.86,-0.4);
\draw[fill=gray!60] (-0.1,-2.3) rectangle (0.1,0.6);
\node at (0,0.85) {\textbf{Pt}};
\draw[thick] (-0.55,0.5) -- (-0.55,-1.9);
\draw[thick,->] (-0.55,0.8) -- (-0.55,0.45);
\node at (-0.55,1.05) {\ce{Cl2(g)}};
\foreach \y in {-1.8,-1.5,-1.2} {\fill[popgreen] (-0.45,\y) circle (0.045);}
\node[align=center] at (0.45,-1.5) {\ce{Cl-(aq)}};
\node[align=center] at (0,-2.9) {Non-metal / ion\\ \ce{Cl2 + 2e- <=> 2Cl-}};
\end{scope}
% ---- Type 3 : Fe3+/Fe2+ ----
\begin{scope}[xshift=8.4cm]
\draw[thick] (-0.9,0) -- (-0.9,-2.4) -- (0.9,-2.4) -- (0.9,0);
\fill[popblueL] (-0.86,-2.36) rectangle (0.86,-0.4);
\draw[fill=gray!60] (-0.1,-2.3) rectangle (0.1,0.6);
\node at (0,0.85) {\textbf{Pt}};
\node[align=center] at (0,-1.5) {\ce{Fe^{3+}} + \ce{Fe^{2+}}};
\node[align=center] at (0,-2.9) {Ion--ion\\ \ce{Fe^{3+} + e- <=> Fe^{2+}}};
\end{scope}
\end{tikzpicture}
\legende{The three types of half-cells. Platinum (Pt) is used as an inert electrode whenever the couple contains no metallic conductor.}
\end{popfigure}

\begin{aretenir}[Choosing the electrode]
If the reduced form of the couple is a \textbf{metal}, that metal is the electrode. If not (gas, or two dissolved ions), use an \textbf{inert platinum electrode}. For ion--ion half-cells, the solution must contain \emph{both} ions of the couple.
\end{aretenir}

\courssub{The Need for a Reference: the Standard Hydrogen Electrode (SHE)}

There is no instrument that measures the ``absolute'' tendency of a single half-cell to gain electrons. What a voltmeter \emph{can} measure is the \emph{difference} of potential between two half-cells. Scientists therefore agreed to compare every half-cell with one universal reference, exactly as altitudes on Earth are measured relative to sea level. The chosen ``sea level'' of electrochemistry is the \cle{standard hydrogen electrode} (SHE).

\begin{definition}[Standard hydrogen electrode]
The \cle{standard hydrogen electrode} consists of a platinum electrode coated with finely divided platinum (``platinum black''), over which pure hydrogen gas at a pressure of $1\ \mathrm{atm}$ ($101\ \mathrm{kPa}$) is bubbled, immersed in a solution of hydrogen ions (e.g.\ dilute $\ce{HCl}$) of concentration $1\ \mathrm{mol\,dm^{-3}}$, at $298\ \mathrm{K}$. Its half-equation is
\[ \ce{2H+(aq) + 2e- <=> H2(g)} \qquad E^{\theta} = 0.00\ \mathrm{V} \]
By international agreement, its standard electrode potential is assigned the value \textbf{exactly zero} at all temperatures.
\end{definition}

\begin{popfigure}
\begin{tikzpicture}[scale=1.0, every node/.style={font=\small}]
% beaker
\draw[thick] (-1.6,0.2) -- (-1.6,-3.0) -- (1.6,-3.0) -- (1.6,0.2);
\fill[popblueL] (-1.56,-2.96) rectangle (1.56,-0.6);
\node[align=center] at (0,-2.55) {\ce{H+(aq)}, $1\ \mathrm{mol\,dm^{-3}}$};
% glass bell / gas inlet
\draw[thick] (-0.85,1.6) -- (-0.85,-2.2);
\draw[thick] (0.85,1.6) -- (0.85,-2.2);
\draw[thick] (-0.85,1.6) -- (0.85,1.6);
% H2 inlet tube
\draw[thick] (-0.35,2.4) -- (-0.35,-1.9);
\draw[thick,->] (-0.35,2.75) -- (-0.35,2.35);
\node[align=center] at (-1.7,2.6) {\ce{H2(g)}, $1\ \mathrm{atm}$};
% bubbles
\foreach \p in {(-0.2,-1.7),(0.15,-1.4),(-0.1,-1.1),(0.2,-1.8)} {\fill[popgreen] \p circle (0.05);}
% Pt electrode
\draw[fill=gray!60] (0.35,-2.2) rectangle (0.55,0.9);
\node[align=center] at (1.35,0.6) {Pt foil coated\\ with Pt black};
\draw[thick] (0.45,0.9) -- (0.45,2.4);
\fill (0.45,2.4) circle (0.06);
\node[right] at (0.55,2.4) {to external circuit};
% H2 outlet
\draw[thick,->] (0.85,1.9) -- (1.5,1.9);
\node[right] at (1.5,1.9) {\ce{H2} out};
\node[align=center] at (0,-3.5) {temperature $298\ \mathrm{K}$};
\end{tikzpicture}
\legende{The standard hydrogen electrode (SHE): hydrogen at $1\ \mathrm{atm}$ bubbled over platinised platinum in $\ce{H+(aq)}$ at $1\ \mathrm{mol\,dm^{-3}}$ and $298\ \mathrm{K}$.}
\end{popfigure}

\begin{propriete}[Standard conditions]
Standard electrode potentials are measured under \cle{standard conditions}:
\begin{itemize}
\item temperature: $298\ \mathrm{K}$ ($25\ ^{\circ}\mathrm{C}$);
\item concentration of every aqueous species: $1\ \mathrm{mol\,dm^{-3}}$;
\item pressure of any gas: $1\ \mathrm{atm}$ ($101\ \mathrm{kPa}$);
\item any solid present must be pure.
\end{itemize}
(These conditions must be known and quoted precisely; the proof is not examinable.)
\end{propriete}

\courssub{Standard Electrode Potential \texorpdfstring{$E^{\theta}$}{E-theta}}

\begin{definition}[Standard electrode potential]
The \cle{standard electrode potential} $E^{\theta}$ of a half-cell is the e.m.f.\ of a cell formed by coupling that half-cell, under standard conditions, to a standard hydrogen electrode. It is measured with a high-resistance voltmeter so that (almost) no current flows. By convention, $E^{\theta}$ is always quoted for the \textbf{reduction} half-equation written as $\text{oxidised form} + n\ce{e-} \to \text{reduced form}$.
\end{definition}

\begin{attention}[$E^{\theta}$ is a reduction potential]
$E^{\theta}$ values are \emph{always} quoted for the reduction half-equation, with electrons on the \textbf{left}: $\ce{Zn^{2+} + 2e- -> Zn}$, $E^{\theta} = -0.76\ \mathrm{V}$. Never write the oxidation equation ($\ce{Zn -> Zn^{2+} + 2e-}$) and attach the same value to it. Also, $E^{\theta}$ is \emph{not} multiplied when the half-equation is multiplied: for $\ce{2Zn^{2+} + 4e- -> 2Zn}$ the potential is still $-0.76\ \mathrm{V}$, because $E^{\theta}$ is an intensive quantity (a voltage, not an amount of energy).
\end{attention}

\begin{methode}[Measuring $E^{\theta}$ of a half-cell against the SHE]
\begin{enumerate}
\item Set up the half-cell under standard conditions ($1\ \mathrm{mol\,dm^{-3}}$ solutions, $298\ \mathrm{K}$, $1\ \mathrm{atm}$ for gases).
\item Connect it to a standard hydrogen electrode with a salt bridge (e.g.\ filter paper soaked in saturated $\ce{KNO3}$ solution).
\item Join the two electrodes through a high-resistance voltmeter.
\item Read the e.m.f.: its \textbf{magnitude} is $\lvert E^{\theta} \rvert$ of the half-cell.
\item Note the \textbf{polarity}: if the electrode under test is the negative terminal (electrons flow from it to the SHE), its $E^{\theta}$ is \textbf{negative}; if it is the positive terminal, its $E^{\theta}$ is \textbf{positive}.
\end{enumerate}
\end{methode}

\begin{exemplebox}[Sign convention in practice]
Coupling a standard $\ce{Zn^{2+}/Zn}$ half-cell to the SHE gives a voltmeter reading of $0.76\ \mathrm{V}$, and the zinc electrode is found to be the \emph{negative} terminal (electrons flow from zinc to hydrogen electrode, where $\ce{2H+ + 2e- -> H2}$ occurs). Hence $E^{\theta}(\ce{Zn^{2+}/Zn}) = -0.76\ \mathrm{V}$. With a $\ce{Cu^{2+}/Cu}$ half-cell, the reading is $0.34\ \mathrm{V}$ but copper is the \emph{positive} terminal, so $E^{\theta}(\ce{Cu^{2+}/Cu}) = +0.34\ \mathrm{V}$. A negative $E^{\theta}$ means the couple is a better \emph{reducing} system than $\ce{H+/H2}$; a positive $E^{\theta}$ means it is a better \emph{oxidising} system.
\end{exemplebox}

\courssub{Other Reference Electrodes}

The SHE is the \emph{defining} reference, but it is awkward in daily laboratory work: it needs a supply of pure hydrogen gas, a precise pressure control, and the platinum surface is easily poisoned by impurities. In practice, chemists use robust secondary reference electrodes whose potentials against the SHE are accurately known.

\textbf{The calomel electrode.}\par
``Calomel'' is the old name for mercury(I) chloride, $\ce{Hg2Cl2}$. The electrode consists of mercury in contact with a paste of mercury and solid $\ce{Hg2Cl2}$, immersed in a potassium chloride solution of known concentration. The half-equation is
\[ \ce{Hg2Cl2(s) + 2e- <=> 2Hg(l) + 2Cl-(aq)} \]
With saturated KCl (the ``saturated calomel electrode'', SCE) its potential is $+0.24\ \mathrm{V}$ versus the SHE at $298\ \mathrm{K}$. It is compact, stable and easy to use, which is why it is the reference built into many pH meters.

\begin{popfigure}
\begin{tikzpicture}[scale=1.0, every node/.style={font=\small}]
% outer tube
\draw[thick] (-1.0,2.2) -- (-1.0,-2.6);
\draw[thick] (1.0,2.2) -- (1.0,-2.6);
\draw[thick] (-1.0,-2.6) .. controls (-0.6,-3.0) and (0.6,-3.0) .. (1.0,-2.6);
% KCl solution
\fill[popblueL] (-0.96,-1.6) rectangle (0.96,1.6);
\node[align=center] at (0,1.1) {saturated KCl\\ solution};
% inner tube with Hg / calomel
\draw[thick] (-0.35,2.6) -- (-0.35,-1.2);
\draw[thick] (0.35,2.6) -- (0.35,-1.2);
\fill[popgold] (-0.31,-1.16) rectangle (0.31,-0.55);
\node[align=center, right] at (1.05,-0.85) {paste of Hg +\\ \ce{Hg2Cl2} (calomel)};
\fill[gray!70] (-0.31,-0.51) rectangle (0.31,0.1);
\node[align=center, right] at (1.05,-0.2) {liquid mercury, Hg};
% Pt wire
\draw[thick] (0,0.1) -- (0,3.0);
\fill (0,3.0) circle (0.06);
\node[right] at (0.1,3.0) {Pt wire to circuit};
% porous plug
\fill[poppurple] (-0.25,-2.85) rectangle (0.25,-2.55);
\node[align=center, right] at (1.05,-2.7) {porous plug (contact\\ with test solution)};
\end{tikzpicture}
\legende{The saturated calomel electrode (SCE), a practical secondary reference electrode: $E = +0.24\ \mathrm{V}$ versus the SHE at $298\ \mathrm{K}$.}
\end{popfigure}

\textbf{The silver / silver chloride electrode.}\par
A silver wire coated with a layer of solid silver chloride is dipped into a solution of chloride ions (usually KCl):
\[ \ce{AgCl(s) + e- <=> Ag(s) + Cl-(aq)} \qquad E \approx +0.22\ \mathrm{V} \text{ (saturated KCl, vs SHE)} \]
It is even smaller and more convenient than the calomel electrode and is the reference most commonly found inside modern pH and ion-selective electrodes. To use any secondary reference, one simply subtracts its known potential: if a half-cell shows $+0.52\ \mathrm{V}$ against the SCE, then $E^{\theta} = 0.52 + 0.24 = +0.76\ \mathrm{V}$ against the SHE.

\courssub{The Electrochemical Series}

Once every half-cell has been measured against the SHE, the couples can be ranked in a single table ordered by $E^{\theta}$: this is the \cle{electrochemical series}.

\begin{center}
\begin{tabularx}{\linewidth}{|l|X|c|}
\hline
\rowcolor{popblueL} Reduction half-equation & & $E^{\theta}$ / V \\
\hline
$\ce{Li+ + e- <=> Li}$ & strongest reducing agent on the right & $-3.04$ \\
\hline
$\ce{Zn^{2+} + 2e- <=> Zn}$ & & $-0.76$ \\
\hline
$\ce{Fe^{2+} + 2e- <=> Fe}$ & & $-0.44$ \\
\hline
$\ce{2H+ + 2e- <=> H2}$ & reference & $0.00$ \\
\hline
$\ce{Cu^{2+} + 2e- <=> Cu}$ & & $+0.34$ \\
\hline
$\ce{Fe^{3+} + e- <=> Fe^{2+}}$ & & $+0.77$ \\
\hline
$\ce{Cl2 + 2e- <=> 2Cl-}$ & strongest oxidising agent on the left & $+1.36$ \\
\hline
\end{tabularx}
\end{center}

The more \emph{negative} $E^{\theta}$, the more readily the reduced form (right-hand side) \emph{gives away} electrons: lithium and zinc are powerful reducing agents. The more \emph{positive} $E^{\theta}$, the more readily the oxidised form (left-hand side) \emph{accepts} electrons: chlorine is a powerful oxidising agent. This single table will allow us, in the next section, to predict the direction of redox reactions and to calculate the e.m.f.\ of any cell.

\begin{experience}[Setting up a half-cell in the laboratory]
\textbf{Objective:} Construct a standard $\ce{Cu^{2+}/Cu}$ half-cell and a $\ce{Fe^{3+}/Fe^{2+}}$ half-cell.
\textbf{Apparatus:} Two beakers ($100\ \mathrm{cm^3}$), copper strip, platinum wire, $\ce{CuSO4}$ solution ($1\ \mathrm{mol\,dm^{-3}}$), $\ce{Fe2(SO4)3}$ and $\ce{FeSO4}$ solutions (mixed to give $1\ \mathrm{mol\,dm^{-3}}$ each), connecting wires, voltmeter, salt bridge (filter paper soaked in $\ce{KNO3}$), SHE apparatus (or a secondary reference).
\textbf{Procedure:}
\begin{enumerate}
\item Polish the copper strip and platinum wire with emery paper; rinse with distilled water.
\item Place the copper strip in the $\ce{CuSO4}$ solution (Type 1 half-cell).
\item Place the platinum wire in the mixed $\ce{Fe^{3+}}/\ce{Fe^{2+}}$ solution (Type 3 half-cell).
\item Connect each half-cell to the reference electrode via a salt bridge.
\item Measure the e.m.f. with a high-resistance voltmeter. Note the polarity.
\end{enumerate}
\textbf{Observation:} The copper electrode is positive relative to SHE ($E^{\theta} = +0.34\ \mathrm{V}$). The platinum electrode in the iron solution is positive relative to SHE ($E^{\theta} = +0.77\ \mathrm{V}$).
\textbf{Conclusion:} The measured potentials confirm the positions of these couples in the electrochemical series.
\end{experience}

\begin{savaistu}[The Daniell Cell and the Telegraph]
The first practical application of electrochemical cells was the telegraph, revolutionising communication in the 19th century. The Daniell cell ($\ce{Zn|Zn^{2+}||Cu^{2+}|Cu}$) provided a steadier voltage than the earlier Voltaic pile because the porous pot (early salt bridge) prevented polarisation. In Cameroon, the principles of these cells underpin the lead-acid batteries that start our vehicles and the dry cells powering torches in rural areas. Understanding $E^{\theta}$ allows engineers to choose the best couples for high-energy batteries (like Lithium-ion, $E^{\theta} \approx -3.04\ \mathrm{V}$ for Li) or for corrosion protection (sacrificial anodes using Zn or Mg).
\end{savaistu}

\begin{exoresolu}[Identifying half-cell types and predicting polarity]
For each of the following redox couples, state the type of half-cell, the electrode material, and whether its $E^{\theta}$ is positive or negative relative to the SHE: (a) $\ce{Ag+/Ag}$ ($E^{\theta} = +0.80\ \mathrm{V}$); (b) $\ce{MnO4-/Mn^{2+}}$ ($E^{\theta} = +1.51\ \mathrm{V}$); (c) $\ce{Fe^{2+}/Fe}$ ($E^{\theta} = -0.44\ \mathrm{V}$).
\tcblower
\textbf{(a)} $\ce{Ag+/Ag}$: the reduced form is a metal, so this is a \textbf{metal/metal-ion} half-cell: a silver wire dipped in $\ce{AgNO3}$ solution ($1\ \mathrm{mol\,dm^{-3}}$). $E^{\theta} = +0.80\ \mathrm{V}$ is \textbf{positive}: silver is the positive terminal against the SHE (reduction $\ce{Ag+ + e- -> Ag}$ occurs there).\\[2pt]
\textbf{(b)} $\ce{MnO4-/Mn^{2+}}$: both species are ions in solution, so this is an \textbf{ion--ion} half-cell requiring an \textbf{inert platinum electrode} dipped in a solution containing $\ce{MnO4-}$, $\ce{Mn^{2+}}$ and $\ce{H+}$, each at $1\ \mathrm{mol\,dm^{-3}}$. $E^{\theta} = +1.51\ \mathrm{V}$ is \textbf{positive}: acidified permanganate is a very strong oxidising agent.\\[2pt]
\textbf{(c)} $\ce{Fe^{2+}/Fe}$: \textbf{metal/metal-ion} half-cell: an iron nail in $\ce{FeSO4}$ solution. $E^{\theta} = -0.44\ \mathrm{V}$ is \textbf{negative}: iron is the negative terminal against the SHE; electrons flow from iron to the hydrogen electrode, so iron is a reducing agent stronger than $\ce{H2}$.
\end{exoresolu}

\begin{exercice}[Half-cells and the SHE]
\begin{enumerate}
\item Draw a fully labelled diagram of the half-cell you would use to measure $E^{\theta}$ for the couple $\ce{Cl2/Cl-}$, stating all standard conditions.
\item A standard $\ce{Ni^{2+}/Ni}$ half-cell is connected to the SHE. The voltmeter reads $0.25\ \mathrm{V}$ and nickel is the negative terminal. Deduce $E^{\theta}(\ce{Ni^{2+}/Ni})$ and write the half-equation that occurs at each electrode.
\item A half-cell gives an e.m.f.\ of $+0.10\ \mathrm{V}$ when measured against a saturated calomel electrode ($E = +0.24\ \mathrm{V}$ vs SHE), the unknown electrode being positive. Calculate its standard electrode potential.
\end{enumerate}
\end{exercice}

\begin{corrige}[Exercise 1]
\textbf{1.} Beaker containing $\ce{Cl-(aq)}$ at $1\ \mathrm{mol\,dm^{-3}}$ (e.g.\ NaCl solution); chlorine gas bubbled at $1\ \mathrm{atm}$ over a platinum electrode; temperature $298\ \mathrm{K}$. The diagram is the middle half-cell of the figure of the three types, labelled with these conditions.\\
\textbf{2.} Nickel is the negative terminal, so $E^{\theta}(\ce{Ni^{2+}/Ni}) = -0.25\ \mathrm{V}$. At the nickel electrode (oxidation): $\ce{Ni(s) -> Ni^{2+}(aq) + 2e-}$. At the SHE (reduction): $\ce{2H+(aq) + 2e- -> H2(g)}$.\\
\textbf{3.} $E^{\theta} = E_{\text{measured}} + E_{\text{ref}} = 0.10 + 0.24 = +0.34\ \mathrm{V}$ versus the SHE.
\end{corrige}

\begin{aretenir}[Key points]
\begin{itemize}
\item Three types of half-cells: \textbf{metal/metal-ion} (the metal is the electrode), \textbf{non-metal/ion} and \textbf{ion--ion} (both need an inert \textbf{platinum} electrode).
\item The \textbf{SHE} ($\ce{2H+ + 2e- <=> H2}$, $1\ \mathrm{mol\,dm^{-3}}$ $\ce{H+}$, $\ce{H2}$ at $1\ \mathrm{atm}$, $298\ \mathrm{K}$, platinised Pt) is assigned $E^{\theta} = 0.00\ \mathrm{V}$.
\item $E^{\theta}$ of a half-cell = e.m.f.\ of the cell it forms with the SHE under standard conditions; it is always quoted for the \textbf{reduction} half-equation and is never multiplied by stoichiometric coefficients.
\item Negative $E^{\theta}$: the electrode is the negative terminal versus the SHE (strong reducing system); positive $E^{\theta}$: strong oxidising system.
\item Practical reference electrodes: \textbf{calomel} ($+0.24\ \mathrm{V}$, saturated KCl) and \textbf{silver/silver chloride} ($\approx +0.22\ \mathrm{V}$).
\item Ranking couples by $E^{\theta}$ gives the \textbf{electrochemical series}, the tool for predicting redox reactions and cell e.m.f.s.
\end{itemize}
\end{aretenir}

\coursec{Electrochemical Cells, Cell emf and Uses of $E^\theta$}

In the previous section we learnt how to measure the \cle{standard electrode potential} $E^\theta$ of any half-cell by coupling it to the standard hydrogen electrode. But why couple a half-cell to the SHE at all? Because a single half-cell can do nothing on its own: a zinc rod dipped in zinc sulphate solution simply sits there. The moment we connect \emph{two} half-cells together, something remarkable happens: electrons begin to flow through a wire, and we have built a battery. This section is about that idea: how two half-cells combine to make an \cle{electrochemical cell}, how we write it down compactly, how we calculate its voltage, and how the table of $E^\theta$ values becomes a powerful predictive tool.

\courssub{Coupling Two Half-Cells: the Electrochemical Cell}

Consider two half-cells, for example $\ce{Zn^{2+}(aq)/Zn(s)}$ and $\ce{Cu^{2+}(aq)/Cu(s)}$. Each on its own is at equilibrium:
\[
\ce{Zn^{2+}(aq) + 2e- <=> Zn(s)} \qquad E^\theta = -0.76\ \mathrm{V}
\]
\[
\ce{Cu^{2+}(aq) + 2e- <=> Cu(s)} \qquad E^\theta = +0.34\ \mathrm{V}
\]
Zinc has the more negative $E^\theta$: it is the stronger \cle{reducing agent}, so zinc atoms release electrons more readily than copper atoms. If we now connect the two metal rods by a wire, electrons flow \emph{from zinc to copper}. Zinc is oxidised ($\ce{Zn -> Zn^{2+} + 2e-}$) and copper(II) ions are reduced ($\ce{Cu^{2+} + 2e- -> Cu}$).

Two practical problems arise immediately. First, the solutions must not mix directly, otherwise $\ce{Cu^{2+}}$ ions would simply react on the zinc surface and no current would flow through the wire. Second, as the reaction proceeds, the zinc beaker would accumulate positive charge (extra $\ce{Zn^{2+}}$) and the copper beaker negative charge (loss of $\ce{Cu^{2+}}$); this charge build-up would stop the electron flow almost instantly. Both problems are solved by a \cle{salt bridge}.

\begin{definition}[Electrochemical (galvanic) cell]
An \cle{electrochemical cell} (or galvanic/voltaic cell) is a device made by coupling two half-cells so that a spontaneous redox reaction drives a flow of electrons through an external circuit. The electrode at which \textbf{oxidation} occurs is the \cle{anode} (negative pole); the electrode at which \textbf{reduction} occurs is the \cle{cathode} (positive pole).
\end{definition}

\begin{definition}[Cell emf]
The \cle{electromotive force} (emf) of a cell, $E_{\text{cell}}$, is the potential difference between its two electrodes measured when no current is drawn (e.g.\ with a high-resistance voltmeter). Under standard conditions it is the \cle{standard cell emf} $E^\theta_{\text{cell}}$, measured in volts.
\end{definition}

\courssub{The Daniell Cell}

The classic example is the \cle{Daniell cell}, built from a $\ce{Zn/Zn^{2+}}$ half-cell and a $\ce{Cu/Cu^{2+}}$ half-cell.

\begin{experience}[The Daniell cell]
Set up a zinc rod in $\SI{1.0}{mol\,dm^{-3}}$ $\ce{ZnSO4(aq)}$ and a copper rod in $\SI{1.0}{mol\,dm^{-3}}$ $\ce{CuSO4(aq)}$. Connect the solutions with a salt bridge (a U-tube containing a gel saturated with $\ce{KNO3}$ or $\ce{KCl}$, plugged with cotton wool) and connect the rods through a voltmeter.
\textbf{Observations:} the voltmeter reads about $\SI{1.10}{V}$ with copper as the positive terminal; the zinc rod slowly dissolves; the blue colour of the copper sulphate fades and pink copper deposits on the copper rod.
\textbf{Interpretation:}
\begin{align*}
\text{Anode (oxidation):}\quad & \ce{Zn(s) -> Zn^{2+}(aq) + 2e-}\\
\text{Cathode (reduction):}\quad & \ce{Cu^{2+}(aq) + 2e- -> Cu(s)}\\
\text{Overall:}\quad & \ce{Zn(s) + Cu^{2+}(aq) -> Zn^{2+}(aq) + Cu(s)}
\end{align*}
Electrons flow through the external circuit from zinc to copper; in the salt bridge, cations ($\ce{K+}$) migrate towards the copper half-cell and anions ($\ce{NO3-}$) towards the zinc half-cell, keeping both solutions electrically neutral.
\end{experience}

\begin{popfigure}
\begin{center}
\begin{tikzpicture}[scale=0.95, every node/.style={font=\small}]
% beakers
\draw[thick] (-4.2,-2.6) -- (-4.2,0.4) (-1.6,-2.6) -- (-1.6,0.4) (-4.2,-2.6) -- (-1.6,-2.6);
\draw[thick] (1.6,-2.6) -- (1.6,0.4) (4.2,-2.6) -- (4.2,0.4) (1.6,-2.6) -- (4.2,-2.6);
% solutions
\fill[popblueL] (-4.1,-2.55) rectangle (-1.7,-0.4);
\fill[popgreenL] (1.7,-2.55) rectangle (4.1,-0.4);
\node at (-2.9,-2.1) {$\ce{ZnSO4(aq)}$};
\node at (2.9,-2.1) {$\ce{CuSO4(aq)}$};
\node at (-2.9,-1.6) {$\SI{1.0}{mol\,dm^{-3}}$};
\node at (2.9,-1.6) {$\SI{1.0}{mol\,dm^{-3}}$};
% electrodes
\fill[gray!60] (-3.3,-0.3) rectangle (-3.0,1.4);
\fill[poporange!70] (3.0,-0.3) rectangle (3.3,1.4);
\node[above] at (-3.15,1.4) {$\ce{Zn}$ (anode, $-$)};
\node[above] at (3.15,1.4) {$\ce{Cu}$ (cathode, $+$)};
% wire and voltmeter
\draw[thick] (-3.15,1.4) -- (-3.15,2.3) -- (-0.7,2.3);
\draw[thick] (3.15,1.4) -- (3.15,2.3) -- (0.7,2.3);
\draw[thick] (0,2.3) circle (0.7);
\node at (0,2.3) {\textbf{V}};
\node[above] at (0,3.05) {voltmeter $\SI{1.10}{V}$};
% electron flow arrow
\draw[->, very thick, poppurple] (-2.2,2.75) -- (-0.9,2.75);
\node[poppurple, above] at (-1.55,2.75) {$e^-$ flow};
% salt bridge
\draw[line width=5pt, popgold!80] (-1.9,-0.2) .. controls (-1.2,1.3) and (1.2,1.3) .. (1.9,-0.2);
\node[popdark] at (0,1.15) {salt bridge ($\ce{KNO3}$ gel)};
% ion migrations
\draw[->, thick, popgreen] (-0.6,0.75) -- (-1.4,0.15);
\node[popgreen, left] at (-0.7,0.6) {$\ce{NO3-}$};
\draw[->, thick, popblue] (0.6,0.75) -- (1.4,0.15);
\node[popblue, right] at (0.7,0.6) {$\ce{K+}$};
% half equations
\node[align=center] at (-2.9,-3.3) {$\ce{Zn -> Zn^{2+} + 2e-}$};
\node[align=center] at (2.9,-3.3) {$\ce{Cu^{2+} + 2e- -> Cu}$};
\end{tikzpicture}
\end{center}
\legende{The Daniell cell: electrons flow through the external circuit from the zinc anode to the copper cathode; the salt bridge completes the circuit and maintains electrical neutrality.}
\end{popfigure}

\begin{attention}[The salt bridge is not optional]
Without a salt bridge (or porous pot), charge builds up in each beaker and the cell stops working within moments. In examinations you must state its two functions: \textbf{(i)} it completes the circuit, \textbf{(ii)} it maintains electrical neutrality by allowing ion migration, \emph{without} the two solutions mixing. Use $\ce{KNO3}$ or $\ce{KCl}$ because $\ce{K+}$ and $\ce{NO3-}$ (or $\ce{Cl-}$) have similar mobilities and do not react with the half-cell species.
\end{attention}

\courssub{Cell Diagrams (Cell Notation)}

Drawing beakers every time is tedious, so chemists use a shorthand called the \cle{cell diagram} or \cle{cell notation}. The rules are strict and examinable:

\begin{methode}[Writing a cell diagram and computing $E^\theta_{\text{cell}}$]
\begin{enumerate}
\item Identify the two half-cells and look up their $E^\theta$ values.
\item The half-cell with the \textbf{more negative} $E^\theta$ is the \textbf{anode} (oxidation): write it on the \textbf{left}.
\item The half-cell with the \textbf{more positive} $E^\theta$ is the \textbf{cathode} (reduction): write it on the \textbf{right}.
\item Use a single vertical line $|$ for a phase boundary (e.g.\ solid$|$solution) and a double vertical line $\|$ for the salt bridge.
\item Write species in the order: electrode $|$ ions $\|$ ions $|$ electrode. For half-cells without a solid conductor (gas electrodes, ion-ion systems), use an inert electrode (usually Pt) on the outside: $\ce{Pt|H2(g)|H+(aq)}$, $\ce{Pt|Fe^{2+}(aq),Fe^{3+}(aq)}$.
\item Compute the standard emf:
\[
E^\theta_{\text{cell}} = E^\theta_{\text{right}} - E^\theta_{\text{left}} = E^\theta_{\text{cathode}} - E^\theta_{\text{anode}} .
\]
\end{enumerate}
\end{methode}

\begin{popfigure}
\begin{center}
\begin{tikzpicture}[scale=1.0, every node/.style={font=\small}]
\node (zn) at (0,0) {$\ce{Zn(s)}$};
\node (zn2) at (2.2,0) {$\ce{Zn^{2+}(aq)}$};
\node (cu2) at (5.4,0) {$\ce{Cu^{2+}(aq)}$};
\node (cu) at (7.8,0) {$\ce{Cu(s)}$};
\node at (1.1,0) {$|$};
\node at (3.8,0) {$\|$};
\node at (6.6,0) {$|$};
% annotations
\node[poppurple, below=0.35cm of zn, align=center] {anode\\ (oxidation, $-$)};
\node[popblue, below=0.35cm of cu, align=center] {cathode\\ (reduction, $+$)};
\node[popgreen, above=0.35cm of zn2, align=center] {phase boundary\\ solid $|$ solution};
\node[poporange, above=0.35cm of cu2, align=center] {salt bridge\\ (double line)};
\draw[->, popgreen] (2.2,0.55) -- (1.15,0.12);
\draw[->, poporange] (5.4,0.55) -- (3.85,0.12);
\draw[<->, popdark] (0,-1.35) -- (7.8,-1.35);
\node[popdark, below] at (3.9,-1.35) {electrons flow left $\rightarrow$ right in the external circuit};
\end{tikzpicture}
\end{center}
\legende{Anatomy of a cell diagram: anode on the left, cathode on the right; single lines mark phase boundaries, the double line the salt bridge.}
\end{popfigure}

\begin{exoresolu}[Emf of the Daniell cell]
Given $E^\theta(\ce{Zn^{2+}/Zn}) = -0.76\ \mathrm{V}$ and $E^\theta(\ce{Cu^{2+}/Cu}) = +0.34\ \mathrm{V}$, write the cell diagram and calculate $E^\theta_{\text{cell}}$.
\tcblower
\textbf{Solution.} Zinc has the more negative $E^\theta$, so it is the anode (left); copper is the cathode (right):
\[
\ce{Zn(s)|Zn^{2+}(aq)||Cu^{2+}(aq)|Cu(s)}
\]
\[
E^\theta_{\text{cell}} = E^\theta_{\text{right}} - E^\theta_{\text{left}} = (+0.34) - (-0.76) = +1.10\ \mathrm{V}.
\]
The positive value confirms that the reaction $\ce{Zn(s) + Cu^{2+}(aq) -> Zn^{2+}(aq) + Cu(s)}$ is spontaneous under standard conditions.
\end{exoresolu}

\begin{exoresolu}[Cell diagram for a hydrogen electrode and a metal/metal-ion electrode]
Write the cell diagram and calculate $E^\theta_{\text{cell}}$ for a cell made from the standard hydrogen electrode and a $\ce{Cu^{2+}/Cu}$ half-cell. $E^\theta(\ce{Cu^{2+}/Cu}) = +0.34\ \mathrm{V}$.
\tcblower
\textbf{Solution.} SHE has $E^\theta = 0.00\ \mathrm{V}$; copper has $+0.34\ \mathrm{V}$. SHE is more negative, so it is the anode (left). The SHE uses an inert Pt electrode:
\[
\ce{Pt(s)|H2(g, 1 atm)|H+(aq, 1 M)||Cu^{2+}(aq, 1 M)|Cu(s)}
\]
\[
E^\theta_{\text{cell}} = 0.34 - 0.00 = +0.34\ \mathrm{V}.
\]
\end{exoresolu}

\begin{exoresolu}[Cell diagram for an ion-ion half-cell]
Write the cell diagram and calculate $E^\theta_{\text{cell}}$ for a cell made from $\ce{Fe^{3+}/Fe^{2+}}$ ($E^\theta = +0.77\ \mathrm{V}$) and $\ce{Cu^{2+}/Cu}$ ($E^\theta = +0.34\ \mathrm{V}$).
\tcblower
\textbf{Solution.} Copper has the more negative $E^\theta$, so it is the anode (left). The $\ce{Fe^{3+}/Fe^{2+}}$ half-cell requires an inert Pt electrode:
\[
\ce{Cu(s)|Cu^{2+}(aq)||Fe^{3+}(aq),Fe^{2+}(aq)|Pt(s)}
\]
\[
E^\theta_{\text{cell}} = 0.77 - 0.34 = +0.43\ \mathrm{V}.
\]
\end{exoresolu}

\begin{attention}[Sign mistakes in $E^\theta_{\text{cell}}$]
The most common GCE error is to subtract the wrong way round, or to \emph{change the sign} of the anode potential because "oxidation is the reverse". \textbf{Never reverse the sign of a tabulated $E^\theta$.} All $E^\theta$ values are written as \emph{reduction} potentials; simply apply $E^\theta_{\text{cell}} = E^\theta_{\text{cathode}} - E^\theta_{\text{anode}}$ with the tabulated values as they stand.
\end{attention}

\courssub{Uses of Standard Electrode Potentials}

The table of $E^\theta$ values is one of the most useful predictive tools in A-Level chemistry.

\textbf{1. Predicting the direction of electron flow and electrode polarity.} In any cell, electrons flow through the external circuit from the half-cell with the more negative $E^\theta$ (which becomes the negative electrode, the anode) to the one with the more positive $E^\theta$ (the positive electrode, the cathode).

\textbf{2. Predicting whether a redox reaction is feasible.}

\begin{propriete}[Feasibility criterion]
A redox reaction is \cle{feasible} (spontaneous) under standard conditions if the emf of the corresponding cell is positive:
\[
E^\theta_{\text{cell}} = E^\theta_{\text{reduced species}} - E^\theta_{\text{oxidised species}} > 0 .
\]
Equivalently: a species will oxidise another if its couple has the more positive $E^\theta$. (Proof not examinable; the criterion follows from $\Delta G^\theta = -nFE^\theta_{\text{cell}}$.)
\end{propriete}

\begin{exemplebox}[Will zinc displace copper?]
Does $\ce{Zn(s)}$ react with $\ce{Cu^{2+}(aq)}$? Here $\ce{Cu^{2+}}$ is reduced ($E^\theta = +0.34\ \mathrm{V}$) and $\ce{Zn}$ is oxidised ($E^\theta = -0.76\ \mathrm{V}$):
\[
E^\theta_{\text{cell}} = 0.34 - (-0.76) = +1.10\ \mathrm{V} > 0 \quad\Rightarrow\quad \text{feasible.}
\]
Conversely, $\ce{Cu(s)}$ cannot displace $\ce{Zn^{2+}}$: $E^\theta_{\text{cell}} = -0.76 - 0.34 = -1.10\ \mathrm{V} < 0$.
\end{exemplebox}

\textbf{3. Comparing oxidising and reducing strengths.}

\begin{methode}[Ranking oxidising and reducing agents from $E^\theta$]
\begin{enumerate}
\item The \textbf{more positive} the $E^\theta$ of a couple, the \textbf{stronger the oxidising agent} is the oxidised form (left-hand side of the reduction half-equation).
\item The \textbf{more negative} the $E^\theta$, the \textbf{stronger the reducing agent} is the reduced form (right-hand side).
\item Example: since $E^\theta(\ce{Cl2/Cl-}) = +1.36\ \mathrm{V} > E^\theta(\ce{Br2/Br-}) = +1.07\ \mathrm{V}$, $\ce{Cl2}$ is a stronger oxidising agent than $\ce{Br2}$, and $\ce{Br-}$ is a stronger reducing agent than $\ce{Cl-}$.
\end{enumerate}
\end{methode}

\begin{attention}[Feasible does not mean fast]
A positive $E^\theta_{\text{cell}}$ tells us a reaction is \emph{energetically} feasible, but says nothing about its \emph{rate}. Some feasible reactions have a very high activation energy and are extremely slow in practice (e.g.\ many oxidations by $\ce{MnO4-}$ need warming). $E^\theta$ predicts \textbf{possibility, not speed}.
\end{attention}

\courssub{Factors Affecting Electrode Potential}

The tabulated values are \emph{standard} potentials, measured at $\SI{298}{K}$, with all solutions at $\SI{1.0}{mol\,dm^{-3}}$ and gases at $\SI{1}{atm}$ ($\SI{101}{kPa}$). Change these conditions and the actual electrode potential changes. Each half-cell is an equilibrium, so \cle{Le Chatelier's principle} predicts the direction of the change.

\begin{itemize}
\item \textbf{Concentration:} For $\ce{Cu^{2+}(aq) + 2e- <=> Cu(s)}$, increasing $[\ce{Cu^{2+}}]$ shifts the equilibrium to the right (Le Chatelier), so the electrode becomes a better electron acceptor: $E$ becomes \emph{more positive}. Diluting the solution makes $E$ more negative.
\item \textbf{Temperature:} Changing temperature shifts the position of the half-cell equilibrium (and changes $K$), so $E$ changes; this is why standard potentials are quoted at $\SI{298}{K}$.
\item \textbf{Pressure (gaseous electrodes):} For a half-cell involving a gas, e.g.\ $\ce{Cl2(g) + 2e- <=> 2Cl-(aq)}$, increasing the pressure of $\ce{Cl2}$ shifts the equilibrium right (Le Chatelier) and makes $E$ more positive.
\end{itemize}

The quantitative relationship is given by the \emph{Nernst equation}, which is beyond the GCE syllabus but explains exactly how $E$ varies with concentration.

\begin{attention}[Non-standard conditions can reverse predictions]
Because concentration changes the actual electrode potential, a reaction predicted unfeasible from $E^\theta$ values may occur under non-standard conditions. For example, concentrated hydrochloric acid reduces $\ce{MnO2}$ to $\ce{Mn^{2+}}$ (the laboratory preparation of chlorine) although standard potentials suggest otherwise: the very high $[\ce{Cl-}]$ and $[\ce{H+}]$ shift the half-cell potentials. Always state "\emph{under standard conditions}" when using $E^\theta$ to predict feasibility.
\end{attention}

\begin{savaistu}[Batteries around us]
Every battery is an electrochemical cell: the dry cell used in torches and radios across Cameroon is based on a $\ce{Zn}$ anode and a $\ce{MnO2}$ cathode; car batteries couple lead and lead(IV) oxide plates in sulphuric acid. The Daniell cell itself, invented in 1836, was one of the first reliable sources of electric current and powered early telegraph networks.
\end{savaistu}

\begin{exercice}[Using $E^\theta$ values]
Use the following standard electrode potentials:
\[
E^\theta(\ce{Fe^{2+}/Fe}) = -0.44\ \mathrm{V},\quad E^\theta(\ce{Cu^{2+}/Cu}) = +0.34\ \mathrm{V},\quad E^\theta(\ce{Ag+/Ag}) = +0.80\ \mathrm{V},\quad E^\theta(\ce{Cl2/Cl-}) = +1.36\ \mathrm{V}.
\]
\begin{enumerate}
\item Write the cell diagram for the cell made from the $\ce{Fe^{2+}/Fe}$ and $\ce{Ag+/Ag}$ half-cells, and calculate its standard emf.
\item Predict, with a calculation, whether iron metal will displace copper(II) ions from solution.
\item Arrange $\ce{Fe}$, $\ce{Cu}$ and $\ce{Ag}$ in order of increasing reducing power.
\item Write the cell diagram for a cell constructed from a $\ce{Cl2/Cl-}$ half-cell and a $\ce{Cu^{2+}/Cu}$ half-cell. Calculate $E^\theta_{\text{cell}}$ and state which electrode is the anode.
\item Will chlorine gas oxidise copper metal to $\ce{Cu^{2+}}$ under standard conditions? Justify with a calculation.
\end{enumerate}
\end{exercice}

\begin{corrige}[Exercise 1]
\textbf{1.} $\ce{Fe}$ has the more negative $E^\theta$, so it is the anode (left):
$\ce{Fe(s)|Fe^{2+}(aq)||Ag+(aq)|Ag(s)}$;
$E^\theta_{\text{cell}} = 0.80 - (-0.44) = +1.24\ \mathrm{V}$.

\textbf{2.} $\ce{Fe}$ oxidised, $\ce{Cu^{2+}}$ reduced: $E^\theta_{\text{cell}} = 0.34 - (-0.44) = +0.78\ \mathrm{V} > 0$, so the reaction $\ce{Fe(s) + Cu^{2+}(aq) -> Fe^{2+}(aq) + Cu(s)}$ is feasible.

\textbf{3.} Reducing power increases as $E^\theta$ becomes more negative: $\ce{Ag} < \ce{Cu} < \ce{Fe}$.

\textbf{4.} $\ce{Cu}$ has the more negative $E^\theta$ ($+0.34\ \mathrm{V}$ vs $+1.36\ \mathrm{V}$), so it is the anode (left). The chlorine half-cell requires an inert Pt electrode:
$\ce{Cu(s)|Cu^{2+}(aq)||Cl-(aq)|Cl2(g)|Pt(s)}$;
$E^\theta_{\text{cell}} = 1.36 - 0.34 = +1.02\ \mathrm{V}$.

\textbf{5.} For $\ce{Cl2}$ oxidising $\ce{Cu}$: $\ce{Cl2}$ reduced ($E^\theta = +1.36\ \mathrm{V}$), $\ce{Cu}$ oxidised ($E^\theta = +0.34\ \mathrm{V}$).
$E^\theta_{\text{cell}} = 1.36 - 0.34 = +1.02\ \mathrm{V} > 0$. Yes, the reaction is feasible.
\end{corrige}

\begin{aretenir}[Key points]
\begin{itemize}
\item An electrochemical cell couples two half-cells; oxidation occurs at the \cle{anode} (negative), reduction at the \cle{cathode} (positive); electrons flow anode $\rightarrow$ cathode through the external circuit.
\item The \cle{salt bridge} completes the circuit and keeps the half-cells electrically neutral (use $\ce{KNO3}$ or $\ce{KCl}$).
\item Cell diagram: anode on the left, cathode on the right; $|$ = phase boundary, $\|$ = salt bridge. Use $\ce{Pt}$ for half-cells without a solid conductor.
\item $E^\theta_{\text{cell}} = E^\theta_{\text{right}} - E^\theta_{\text{left}} = E^\theta_{\text{cathode}} - E^\theta_{\text{anode}}$; never change the sign of a tabulated $E^\theta$.
\item $E^\theta_{\text{cell}} > 0$ $\Rightarrow$ reaction feasible under standard conditions; but feasibility says nothing about rate, and non-standard concentrations can alter the actual potentials.
\item More positive $E^\theta$ $\Rightarrow$ stronger oxidising agent; more negative $E^\theta$ $\Rightarrow$ stronger reducing agent.
\item Electrode potential changes with concentration (Le Chatelier), temperature, and gas pressure.
\end{itemize}
\end{aretenir}

\coursec{Acid-Base Theories, pH and the Ionic Product of Water}

In the previous section we studied redox equilibria, in which the particle transferred between chemical species is the \emph{electron}. We now turn to another enormous family of equilibria in which the transferred particle is the \emph{proton}, \ce{H+}: the \cle{acid-base equilibria}. You will see that the same ideas of dynamic equilibrium, equilibrium constants and Le Chatelier's principle apply here too. Acids and bases are part of everyday Cameroonian life: the sourness of lime juice and palm wine, the slippery feel of soap, the sting of an ant bite --- all of these are acid-base chemistry.

\courssub{What is an acid? Three theories}

\textbf{The Arrhenius theory}\par
The oldest useful definition is due to the Swedish chemist Svante Arrhenius. It is based on what happens when a substance dissolves in \emph{water}.

\begin{definition}[Arrhenius acid and base]
An \cle{Arrhenius acid} is a substance which releases hydrogen ions, \ce{H+}, when dissolved in water. An \cle{Arrhenius base} is a substance which releases hydroxide ions, \ce{OH-}, when dissolved in water.
\[
\ce{HCl(g) ->[H2O] H+(aq) + Cl-(aq)} \qquad \ce{NaOH(s) ->[H2O] Na+(aq) + OH-(aq)}
\]
\end{definition}

This theory explains the behaviour of common laboratory acids (\ce{HCl}, \ce{H2SO4}, \ce{HNO3}) and alkalis (\ce{NaOH}, \ce{KOH}), but it has serious \cle{limitations}:
\begin{itemize}
\item It only works for solutions \emph{in water}, yet acid-base reactions occur in other solvents and even in the gas phase.
\item It cannot explain why \ce{NH3} is a base: ammonia contains no \ce{OH-} group to release.
\item The free proton \ce{H+} does not exist alone in water; it attaches to a water molecule to give the \cle{oxonium ion}, \ce{H3O+}.
\end{itemize}

\textbf{The Brønsted-Lowry theory}\par
In 1923, Brønsted and Lowry independently proposed a more general theory centred on \emph{proton transfer}.

\begin{definition}[Brønsted-Lowry acid and base]
A \cle{Brønsted-Lowry acid} is a \textbf{proton donor}; a \cle{Brønsted-Lowry base} is a \textbf{proton acceptor}. An acid-base reaction is a transfer of a proton, \ce{H+}, from the acid to the base.
\end{definition}

Consider hydrogen chloride dissolving in water:
\[
\ce{HCl(aq) + H2O(l) -> H3O+(aq) + Cl-(aq)}
\]
\ce{HCl} donates a proton (acid); \ce{H2O} accepts it (base). Notice that water, which Arrhenius ignored, is here an active participant. Ammonia is now easily explained as a base:
\[
\ce{NH3(aq) + H2O(l) <=> NH4+(aq) + OH-(aq)}
\]

\begin{definition}[Conjugate acid-base pairs]
A \cle{conjugate acid-base pair} consists of two species that differ by exactly one proton. When an acid donates a proton, the species left behind is its \cle{conjugate base}; when a base accepts a proton, the species formed is its \cle{conjugate acid}.
\end{definition}

For the reaction $\ce{HCl} + \ce{H2O} \rightarrow \ce{H3O+} + \ce{Cl-}$, the two conjugate pairs are \ce{HCl}/\ce{Cl-} and \ce{H3O+}/\ce{H2O}. The stronger the acid, the weaker its conjugate base: \ce{Cl-} is an extremely weak base because \ce{HCl} is a very strong acid.

\begin{popfigure}
\begin{center}
\begin{tikzpicture}[scale=1.0, every node/.style={font=\small}]
% Acid-base pair diagram
\node[font=\normalsize] (ha) at (0,0) {\ce{CH3COOH}};
\node[font=\normalsize] (b) at (3.2,0) {\ce{H2O}};
\node[font=\normalsize] (a) at (6.6,0) {\ce{CH3COO-}};
\node[font=\normalsize] (hb) at (9.8,0) {\ce{H3O+}};
\node at (4.9,0) {$\rightleftharpoons$};
% proton transfer arrows
\draw[->, thick, poporange, bend left=35] (ha.north) to node[above, text=poporange, font=\footnotesize]{\ce{H+} donated} (b.north);
\draw[->, thick, popblue, bend left=35] (hb.south) to node[below, text=popblue, font=\footnotesize]{\ce{H+} donated} (a.south);
% conjugate pair links
\draw[popgreen, thick] (0,-0.9) -- (6.6,-0.9) node[midway, below, text=popgreen, font=\footnotesize]{conjugate pair 1};
\draw[popgreen, thick] (3.2,1.0) -- (9.8,1.0) node[midway, above, text=popgreen, font=\footnotesize]{conjugate pair 2};
\node[text=poporange, font=\footnotesize] at (0,0.55) {acid 1};
\node[text=popblue, font=\footnotesize] at (3.2,0.55) {base 2};
\node[text=popblue, font=\footnotesize] at (6.6,-0.55) {base 1};
\node[text=poporange, font=\footnotesize] at (9.8,-0.55) {acid 2};
\end{tikzpicture}
\end{center}
\legende{Proton transfer in the equilibrium $\ce{CH3COOH} + \ce{H2O} \rightleftharpoons \ce{CH3COO-} + \ce{H3O+}$, showing the two conjugate acid-base pairs.}
\end{popfigure}

\textbf{The Lewis theory}\par
G. N. Lewis generalised the idea still further, to reactions where no proton is transferred at all.

\begin{definition}[Lewis acid and base]
A \cle{Lewis acid} is an \textbf{electron-pair acceptor}; a \cle{Lewis base} is an \textbf{electron-pair donor}. A Lewis acid-base reaction forms a dative (coordinate) covalent bond.
\end{definition}

\begin{exemplebox}[A Lewis acid-base reaction]
Boron trifluoride, \ce{BF3}, has only six electrons around boron --- an incomplete octet --- so it can accept an electron pair. Ammonia, \ce{NH3}, has a lone pair on nitrogen:
\[
\ensuremath{\mathrm{BF_{3} + NH_{3} \rightarrow  F_{3}B\leftarrow NH_{3}}}
\]
\ce{BF3} is the Lewis acid, \ce{NH3} the Lewis base, and the product contains a dative bond \ensuremath{\mathrm{B\leftarrow N}}. Metal ions such as \ce{Cu^2+} accepting lone pairs from water ligands are also Lewis acids.
\end{exemplebox}

\begin{aretenir}[Three theories compared]
\begin{center}
\begin{tabularx}{\linewidth}{|l|X|X|}
\hline
\rowcolor{popblueL} \textbf{Theory} & \textbf{Acid} & \textbf{Base} \\
\hline
Arrhenius & releases \ce{H+} in water & releases \ce{OH-} in water \\
\hline
Brønsted-Lowry & proton (\ce{H+}) donor & proton acceptor \\
\hline
Lewis & electron-pair acceptor & electron-pair donor \\
\hline
\end{tabularx}
\end{center}
Each theory includes the previous one as a special case: every Brønsted-Lowry base is a Lewis base, but \ce{BF3} is a Lewis acid without being a Brønsted-Lowry acid.
\end{aretenir}

\courssub{Amphoteric (amphiprotic) species}

Some species can behave as an acid \emph{or} as a base depending on the partner. They are called \cle{amphoteric} (or \cle{amphiprotic} when proton transfer is involved).

\begin{exemplebox}[Water and the hydrogencarbonate ion]
Water accepts a proton from \ce{HCl} (acting as a base) but donates a proton to \ce{NH3} (acting as an acid). Similarly:
\begin{align*}
\ce{HCO3- + H2O &<=> H2CO3 + OH-} && \text{(\ce{HCO3-} as a base)}\\
\ce{HCO3- + H2O &<=> CO3^2- + H3O+} && \text{(\ce{HCO3-} as an acid)}
\end{align*}
\end{exemplebox}

\courssub{The auto-ionisation of water and the ionic product $K_w$}

Pure water seems unreactive, yet a sensitive conductivity meter shows that it conducts electricity very slightly. This proves that a few ions exist in pure water. The explanation is the \cle{auto-ionisation} (self-ionisation) of water, a Brønsted-Lowry proton transfer between two water molecules:
\[
\ce{2H2O(l) <=> H3O+(aq) + OH-(aq)}
\]
This equilibrium lies very far to the left, but it is vitally important. Applying the equilibrium law, and treating the concentration of liquid water as constant, we absorb $[\ce{H2O}]^2$ into the constant.

\begin{definition}[Ionic product of water]
The \cle{ionic product of water} is
\[
K_w = [\ce{H3O+}][\ce{OH-}]
\]
At \SI{25}{\ensuremath{{}^{\circ}\mathrm{C}}}, $K_w = 1.0\times 10^{-14}\ \mathrm{mol^2\,dm^{-6}}$.
\end{definition}

\begin{propriete}[A universal relation]
In \textbf{any} aqueous solution --- acidic, neutral or alkaline --- at a given temperature,
\[
[\ce{H3O+}][\ce{OH-}] = K_w
\]
The two ion concentrations are linked: if one increases, the other must decrease so that the product stays constant. (Proof not examinable; it follows directly from the equilibrium law.)
\end{propriete}

In \emph{pure water}, the equation produces equal numbers of the two ions, so $[\ce{H3O+}] = [\ce{OH-}] = \sqrt{K_w} = 1.0\times 10^{-7}\ \mathrm{mol\,dm^{-3}}$ at \SI{25}{\ensuremath{{}^{\circ}\mathrm{C}}}.

\courssub{The pH scale}

Concentrations of \ce{H3O+} range over many powers of ten, which is clumsy. The Danish chemist Sørensen introduced a logarithmic scale to compress this range.

\begin{definition}[pH]
The \cle{pH} of a solution is defined by
\[
\mathrm{pH} = -\log_{10}[\ce{H3O+}] \qquad \Longleftrightarrow \qquad [\ce{H3O+}] = 10^{-\mathrm{pH}}
\]
with $[\ce{H3O+}]$ in $\mathrm{mol\,dm^{-3}}$. Similarly, $\mathrm{pOH} = -\log_{10}[\ce{OH-}]$.
\end{definition}

Taking $-\log_{10}$ of $K_w = [\ce{H3O+}][\ce{OH-}]$ gives the very useful relation, valid at \SI{25}{\ensuremath{{}^{\circ}\mathrm{C}}}:
\[
\mathrm{pH} + \mathrm{pOH} = 14
\]

\begin{methode}[{Converting between $[\ce{H3O+}]$, $[\ce{OH-}]$, pH and pOH at \SI{25}{\ensuremath{{}^{\circ}\mathrm{C}}}}]
\begin{enumerate}
\item From $[\ce{H3O+}]$ to pH: compute $\mathrm{pH} = -\log_{10}[\ce{H3O+}]$.
\item From pH to $[\ce{H3O+}]$: compute $[\ce{H3O+}] = 10^{-\mathrm{pH}}$.
\item From $[\ce{H3O+}]$ to $[\ce{OH-}]$: use $[\ce{OH-}] = K_w/[\ce{H3O+}] = 10^{-14}/[\ce{H3O+}]$.
\item From pH to pOH: use $\mathrm{pOH} = 14 - \mathrm{pH}$, then $[\ce{OH-}] = 10^{-\mathrm{pOH}}$.
\end{enumerate}
\end{methode}

\begin{exoresolu}[pH of a lime juice sample]
The oxonium ion concentration of a sample of lime juice is $[\ce{H3O+}] = 5.0\times 10^{-3}\ \mathrm{mol\,dm^{-3}}$ at \SI{25}{\ensuremath{{}^{\circ}\mathrm{C}}}. Calculate (a) the pH, (b) $[\ce{OH-}]$ in the juice.
\tcblower
\textbf{(a)} Apply the definition:
\[
\mathrm{pH} = -\log_{10}(5.0\times 10^{-3}) = -(\log_{10}5.0 - 3) = 3 - 0.70 = 2.30
\]
\textbf{(b)} Use the ionic product:
\[
[\ce{OH-}] = \frac{K_w}{[\ce{H3O+}]} = \frac{1.0\times 10^{-14}}{5.0\times 10^{-3}} = 2.0\times 10^{-12}\ \mathrm{mol\,dm^{-3}}
\]
Check: pOH $= -\log_{10}(2.0\times10^{-12}) = 11.70$, and indeed $2.30 + 11.70 = 14$. The juice is acidic, as expected since $[\ce{H3O+}] > [\ce{OH-}]$.
\end{exoresolu}

\courssub{Acidic, neutral and alkaline solutions}

\begin{aretenir}[Classifying solutions at \SI{25}{\ensuremath{{}^{\circ}\mathrm{C}}}]
\begin{center}
\begin{tabularx}{\linewidth}{|l|X|X|X|}
\hline
\rowcolor{popblueL} \textbf{Solution} & $[\ce{H3O+}]$ vs $[\ce{OH-}]$ & $[\ce{H3O+}]$ ($\mathrm{mol\,dm^{-3}}$) & \textbf{pH} \\
\hline
Acidic & $[\ce{H3O+}] > [\ce{OH-}]$ & $> 1.0\times10^{-7}$ & $< 7$ \\
\hline
Neutral & $[\ce{H3O+}] = [\ce{OH-}]$ & $= 1.0\times10^{-7}$ & $= 7$ \\
\hline
Alkaline & $[\ce{H3O+}] < [\ce{OH-}]$ & $< 1.0\times10^{-7}$ & $> 7$ \\
\hline
\end{tabularx}
\end{center}
\end{aretenir}

\begin{popfigure}
\begin{center}
\begin{tikzpicture}[scale=0.86]
% colour gradient of universal indicator
\foreach \x/\col in {0/poppink,1/poppink,2/poporange,3/poporange,4/popgold,5/popgold,6/popgreen,7/popgreen,8/popgreen,9/popblue,10/popblue,11/popblue,12/poppurple,13/poppurple,14/poppurple}
  \fill[\col!70] (\x,0) rectangle (\x+1,0.7);
\draw[thick] (0,0) rectangle (15,0.7);
\foreach \x in {0,...,14}{
  \draw[thick] (\x,0) -- (\x,-0.15);
  \node[below, font=\footnotesize] at (\x+0.5,-0.15) {\x};
}
\node[above, font=\footnotesize\bfseries] at (7.5,0.75) {pH scale with universal indicator colours (\SI{25}{\ensuremath{{}^{\circ}\mathrm{C}}})};
% examples
\node[below, font=\scriptsize, align=center, text width=2cm] at (2.5,-0.85) {lime juice\\pH $\approx 2$};
\draw[->, popdark] (2.5,-0.8) -- (2.5,-0.05);
\node[below, font=\scriptsize, align=center, text width=2cm] at (4.5,-0.85) {palm wine\\pH $\approx 4$};
\draw[->, popdark] (4.5,-0.8) -- (4.5,-0.05);
\node[below, font=\scriptsize, align=center, text width=2cm] at (7.5,-0.85) {pure water\\pH $= 7$};
\draw[->, popdark] (7.5,-0.8) -- (7.5,-0.05);
\node[above, font=\scriptsize, align=center, text width=2cm] at (8.0,1.25) {blood pH $\approx 7.4$};
\draw[->, popdark] (8.0,1.2) -- (7.9,0.75);
\node[above, font=\scriptsize, align=center, text width=2cm] at (10.5,1.25) {soap solution pH $\approx 10$};
\draw[->, popdark] (10.5,1.2) -- (10.5,0.75);
\end{tikzpicture}
\end{center}
\legende{The pH scale with universal indicator colours and familiar examples.}
\end{popfigure}

\courssub{Measuring pH}

\begin{itemize}
\item \textbf{Universal indicator paper or solution (pH paper).} A mixture of indicators whose colour changes gradually from red (strongly acidic) through green (neutral) to dark purple (strongly alkaline). The colour is matched against a printed chart. It is quick and cheap but only \emph{approximate} (to about $\pm 1$ pH unit).
\item \textbf{The pH meter.} A glass electrode sensitive to \ce{H3O+} connected to a voltmeter calibrated directly in pH units. It must first be \cle{calibrated} using standard buffer solutions of known pH. It gives \emph{accurate} readings (to $0.01$ pH unit) and is essential for following pH changes during a titration.
\end{itemize}

\courssub{Effect of temperature on $K_w$ and on the pH of pure water}

The auto-ionisation of water is \emph{endothermic}:
\[
\ce{2H2O(l) <=> H3O+(aq) + OH-(aq)} \qquad \Delta H > 0
\]
By Le Chatelier's principle, raising the temperature shifts the equilibrium to the right: more ions form, so $K_w$ \emph{increases}. For example, $K_w \approx 5.5\times 10^{-14}\ \mathrm{mol^2\,dm^{-6}}$ near \SI{40}{\ensuremath{{}^{\circ}\mathrm{C}}}. In pure water we still have $[\ce{H3O+}] = [\ce{OH-}]$, so the water is still \emph{neutral} --- but $[\ce{H3O+}] > 10^{-7}$, so the pH is \emph{less than 7}.

\begin{attention}[Neutral does not always mean pH 7]
Pure water is \textbf{always neutral} because $[\ce{H3O+}] = [\ce{OH-}]$ at every temperature. But its pH equals 7 \textbf{only at \SI{25}{\ensuremath{{}^{\circ}\mathrm{C}}}}. At \SI{40}{\ensuremath{{}^{\circ}\mathrm{C}}}, neutral water has pH $\approx 6.8$. Neutrality is defined by the \emph{equality of the two ion concentrations}, not by a fixed pH value.
\end{attention}

\begin{attention}[pH values outside 0--14]
The familiar 0--14 range applies only to moderately dilute solutions at \SI{25}{\ensuremath{{}^{\circ}\mathrm{C}}}. For a concentrated strong acid such as $10\ \mathrm{mol\,dm^{-3}}$ \ce{HCl}, $\mathrm{pH} = -\log_{10}(10) = -1$: a \emph{negative} pH. Likewise, a very concentrated strong base can have pH greater than 14. The scale does not stop at 0 and 14.
\end{attention}

\begin{exercice}[Practising pH calculations]
At \SI{25}{\ensuremath{{}^{\circ}\mathrm{C}}}, a sample of soapy water from a laundry in Bamenda has pH $= 10.5$.
\begin{enumerate}
\item Calculate $[\ce{H3O+}]$ and $[\ce{OH-}]$ in the solution.
\item State, with a reason, whether the solution is acidic, neutral or alkaline.
\item At \SI{60}{\ensuremath{{}^{\circ}\mathrm{C}}}, $K_w = 9.3\times 10^{-14}\ \mathrm{mol^2\,dm^{-6}}$. Calculate the pH of pure water at this temperature and explain why the water is still neutral.
\end{enumerate}
\end{exercice}

\begin{corrige}[Exercise 1]
\textbf{1.} $[\ce{H3O+}] = 10^{-\mathrm{pH}} = 10^{-10.5} = 3.2\times 10^{-11}\ \mathrm{mol\,dm^{-3}}$.
\[
[\ce{OH-}] = \frac{K_w}{[\ce{H3O+}]} = \frac{1.0\times 10^{-14}}{3.2\times 10^{-11}} = 3.2\times 10^{-4}\ \mathrm{mol\,dm^{-3}}
\]
\textbf{2.} Since $[\ce{OH-}] > [\ce{H3O+}]$ (and pH $> 7$), the solution is alkaline, as expected for soap.
\textbf{3.} In pure water $[\ce{H3O+}] = \sqrt{K_w} = \sqrt{9.3\times 10^{-14}} = 3.0\times 10^{-7}\ \mathrm{mol\,dm^{-3}}$, so $\mathrm{pH} = -\log_{10}(3.0\times 10^{-7}) = 6.5$. The water is still neutral because auto-ionisation produces equal concentrations of \ce{H3O+} and \ce{OH-}; neutrality depends on the equality $[\ce{H3O+}] = [\ce{OH-}]$, not on pH being 7.
\end{corrige}

\begin{savaistu}[Why blood pH matters]
Human blood is kept within the very narrow pH range 7.35--7.45 by natural \emph{buffer systems} (mainly the \ce{H2CO3}/\ce{HCO3-} pair you will study in Section 11). A shift of only 0.2 pH unit can be life-threatening. The same buffering principle is used in fermented foods and in the preparation of traditional medicines, where controlling acidity prevents the growth of harmful microorganisms.
\end{savaistu}

\begin{aretenir}[Key points of this section]
\begin{itemize}
\item Arrhenius: acids give \ce{H+}, bases give \ce{OH-} in water (limited to aqueous solutions).
\item Brønsted-Lowry: acid = proton donor, base = proton acceptor; species differing by one \ce{H+} form a conjugate pair.
\item Lewis: acid = electron-pair acceptor, base = electron-pair donor (e.g. \ce{BF3} + \ce{NH3}).
\item Water auto-ionises: $\ce{2H2O <=> H3O+ + OH-}$, with $K_w = [\ce{H3O+}][\ce{OH-}] = 1.0\times10^{-14}\ \mathrm{mol^2\,dm^{-6}}$ at \SI{25}{\ensuremath{{}^{\circ}\mathrm{C}}}.
\item $\mathrm{pH} = -\log_{10}[\ce{H3O+}]$; $\mathrm{pH} + \mathrm{pOH} = 14$ at \SI{25}{\ensuremath{{}^{\circ}\mathrm{C}}}.
\item Neutral means $[\ce{H3O+}] = [\ce{OH-}]$; pH $= 7$ only at \SI{25}{\ensuremath{{}^{\circ}\mathrm{C}}} since $K_w$ increases with temperature.
\end{itemize}
\end{aretenir}

\coursec{pH of Strong and Weak Acids and Bases, and of Mixtures}

In the previous section we defined pH, established the ionic product of water $K_w = [\ce{H3O+}][\ce{OH-}] = 1.0\times 10^{-14}$ at \SI{25}{\ensuremath{{}^{\circ}\mathrm{C}}}, and learnt how to measure pH. We now put these tools to work. The central question of this section is simple: \emph{if I dissolve a given amount of an acid or a base in water, what is the pH of the solution?} The answer depends crucially on whether the acid or base is \cle{strong} or \cle{weak} — that is, on how completely it ionises in water.

\courssub{Strong Acids and Strong Bases: Complete Ionisation}

When hydrogen chloride gas is dissolved in water, essentially \emph{every} \ce{HCl} molecule donates its proton to a water molecule:

\[ \ce{HCl(aq) + H2O(l) -> H3O+(aq) + Cl-(aq)} \]

The reaction goes to completion: no detectable \ce{HCl} molecules remain in solution. Such an acid is called \cle{strong}. Similarly, sodium hydroxide is an ionic solid that dissociates completely:

\[ \ce{NaOH(s) ->[H2O] Na+(aq) + OH-(aq)} \]

\begin{definition}[Strong and weak acids and bases]
An acid is \cle{strong} if it ionises \textbf{completely} in aqueous solution (degree of ionisation $\alpha \approx 1$); it is \cle{weak} if it ionises only \textbf{partially} ($\alpha \ll 1$), so that an equilibrium is set up between the molecules and the ions. The same distinction applies to bases: a \cle{strong base} (e.g.\ \ce{NaOH}, \ce{KOH}) dissociates completely to give \ce{OH-} ions, while a \cle{weak base} (e.g.\ \ce{NH3}) reacts only partially with water.
\end{definition}

\begin{attention}['Strong' is not 'concentrated']
\cle{Strength} describes the \emph{extent of ionisation}; \cle{concentration} describes the \emph{amount of solute per unit volume}. A \SI{10}{mol.L^{-1}} solution of ethanoic acid is \textbf{concentrated but weak}; a $10^{-4}\ \mathrm{mol\,L^{-1}}$ solution of \ce{HCl} is \textbf{dilute but strong}. Never use the two words interchangeably in the examination.
\end{attention}

Common strong acids: \ce{HCl}, \ce{HNO3}, \ce{H2SO4} (first proton fully ionised), \ce{HBr}, \ce{HClO4}. Common strong bases: \ce{NaOH}, \ce{KOH}, \ce{Ba(OH)2} (moderately soluble, but what dissolves dissociates completely).

\textbf{pH of a strong monoprotic acid.} Since ionisation is complete, each mole of acid gives one mole of \ce{H3O+}:

\[ [\ce{H3O+}] = c_a \qquad \Longrightarrow \qquad \boxed{\ \mathrm{pH} = -\log c_a\ } \]

\textbf{pH of a strong monobasic base (e.g.\ \ce{NaOH}, \ce{KOH}).} Complete dissociation gives $[\ce{OH-}] = c_b$. We then pass through $K_w$:

\[ [\ce{H3O+}] = \frac{K_w}{[\ce{OH-}]} = \frac{K_w}{c_b} \qquad \Longrightarrow \qquad \mathrm{pH} = 14 + \log c_b \quad (\text{at } \SI{25}{\ensuremath{{}^{\circ}\mathrm{C}}}) \]

For a strong diacidic base like \ce{Ba(OH)2}, $[\ce{OH-}] = 2c_b$ (solubility permitting), so $\mathrm{pH} = 14 + \log(2c_b)$.

\begin{exemplebox}[pH of strong acid and strong base solutions]
\textbf{(a)} \SI{0.05}{mol.L^{-1}} \ce{HCl}: $\mathrm{pH} = -\log(0.05) = 1.30$.

\textbf{(b)} \SI{0.02}{mol.L^{-1}} \ce{NaOH}: $[\ce{OH-}] = 0.02\ \mathrm{mol\,L^{-1}}$, so $[\ce{H3O+}] = \dfrac{1.0\times10^{-14}}{0.02} = 5.0\times10^{-13}\ \mathrm{mol\,L^{-1}}$ and $\mathrm{pH} = 12.30$.

\textbf{(c)} \SI{0.01}{mol.L^{-1}} \ce{H2SO4}: the first proton is fully released; the second proton has $K_{a2} = 1.2\times10^{-2}$, so at this concentration it is \emph{not} fully dissociated. A rigorous treatment gives $[\ce{H3O+}] \approx 0.011\ \mathrm{mol\,L^{-1}}$ and $\mathrm{pH} \approx 1.96$. (The simplified assumption ``both protons fully released'' would give $[\ce{H3O+}] = 0.02\ \mathrm{mol\,L^{-1}}$, $\mathrm{pH} = 1.70$; this overestimates the acidity.)
\end{exemplebox}

\courssub{Weak Acids and Weak Bases: Partial Ionisation, $K_a$ and $K_b$}

Ethanoic acid, the acid in vinegar and in the palm-oil kitchens of Cameroon, behaves very differently from \ce{HCl}. Dissolved in water, only a small fraction of its molecules ionise; an equilibrium is established:

\[ \ce{CH3COOH(aq) + H2O(l) <=> CH3COO-(aq) + H3O+(aq)} \]

Likewise ammonia, a weak base, reacts only partially:

\[ \ce{NH3(aq) + H2O(l) <=> NH4+(aq) + OH-(aq)} \]

\begin{popfigure}
\begin{center}
\begin{tikzpicture}[scale=0.9]
% Strong acid beaker
\draw[thick, popblue] (-0.2,0) -- (-0.2,2.6) (3.4,0) -- (3.4,2.6);
\draw[thick, popblue] (-0.2,0) -- (3.4,0);
\fill[popblueL, opacity=0.5] (-0.15,0.05) rectangle (3.35,2.2);
\foreach \x/\y in {0.5/0.5, 1.5/1.6, 2.6/0.9, 0.9/1.9, 2.2/1.9, 2.9/1.5}{
  \fill[poporange] (\x,\y) circle (0.10);
  \fill[popgreen] (\x+0.28,\y+0.12) circle (0.10);
}
\node[align=center, text width=3.2cm] at (1.6,-0.7) {\small \textbf{\ce{HCl} (strong)}\\ \small only \ce{H3O+} and \ce{Cl-} ions};
% Weak acid beaker
\begin{scope}[xshift=6.2cm]
\draw[thick, poppurple] (-0.2,0) -- (-0.2,2.6) (3.4,0) -- (3.4,2.6);
\draw[thick, poppurple] (-0.2,0) -- (3.4,0);
\fill[poppurple!12, opacity=0.6] (-0.15,0.05) rectangle (3.35,2.2);
% unionised molecules (joined pair)
\foreach \x/\y in {0.5/0.5, 1.5/1.6, 2.6/0.9, 0.9/1.9}{
  \fill[popmuted] (\x,\y) circle (0.10);
  \fill[popmuted] (\x+0.22,\y) circle (0.10);
  \draw[popmuted, thick] (\x,\y) -- (\x+0.22,\y);
}
% one ionised pair
\fill[poporange] (2.3,1.9) circle (0.10);
\fill[popgreen] (2.9,1.5) circle (0.10);
\node[align=center, text width=3.4cm] at (1.6,-0.7) {\small \textbf{\ce{CH3COOH} (weak)}\\ \small mostly molecules, few ions};
\end{scope}
\end{tikzpicture}
\end{center}
\legende{Particle picture of ionisation. Strong acid: complete ionisation, only ions present. Weak acid: partial ionisation, mostly un-ionised molecules in equilibrium with a few ions.}
\end{popfigure}

Because the ionisation of a weak acid is an \emph{equilibrium}, we can apply the equilibrium law from earlier in this chapter.

\begin{definition}[$K_a$, $pK_a$, $K_b$, $pK_b$]
For a weak acid \ce{HA(aq) + H2O(l) <=> A-(aq) + H3O+(aq)}, the \cle{acid dissociation constant} is
\[ K_a = \frac{[\ce{A-}][\ce{H3O+}]}{[\ce{HA}]} \qquad \text{and} \qquad pK_a = -\log K_a . \]
For a weak base \ce{B(aq) + H2O(l) <=> BH+(aq) + OH-(aq)}, the \cle{base dissociation constant} is
\[ K_b = \frac{[\ce{BH+}][\ce{OH-}]}{[\ce{B}]} \qquad \text{and} \qquad pK_b = -\log K_b . \]
The larger $K_a$ (the smaller $pK_a$), the stronger the weak acid. For a conjugate acid--base pair, $K_a \times K_b = K_w$, i.e.\ $pK_a + pK_b = 14$ at \SI{25}{\ensuremath{{}^{\circ}\mathrm{C}}}.
\end{definition}

\begin{exemplebox}[Typical values at \SI{25}{\ensuremath{{}^{\circ}\mathrm{C}}}]
Ethanoic acid: $K_a = 1.8\times 10^{-5}\ \mathrm{mol\,L^{-1}}$, $pK_a = 4.74$. Ammonia: $K_b = 1.8\times 10^{-5}\ \mathrm{mol\,L^{-1}}$, so for \ce{NH4+}, $pK_a = 14 - 4.74 = 9.26$.
\end{exemplebox}

\courssub{Calculating the pH of a Weak Acid}

Consider a weak acid \ce{HA} of analytical concentration $c$. Let $x$ be the concentration that ionises. The equilibrium table (in \si{mol.L^{-1}}) is:

\begin{center}
\begin{tabular}{|l|c|c|c|}
\hline
\rowcolor{popblueL} & \ce{HA} & \ce{A-} & \ce{H3O+} \\
\hline
Initial & $c$ & $0$ & $0$ \\
\hline
Equilibrium & $c - x$ & $x$ & $x$ \\
\hline
\end{tabular}
\end{center}

Hence $K_a = \dfrac{x^2}{c - x}$. For a weak acid, $x \ll c$, so $c - x \approx c$, giving

\[ K_a \approx \frac{x^2}{c} \quad\Longrightarrow\quad [\ce{H3O+}] \approx \sqrt{K_a\, c} \quad\Longrightarrow\quad \mathrm{pH} = \tfrac{1}{2}\bigl(pK_a - \log c\bigr). \]

\begin{methode}[pH of a weak acid using $K_a$]
\begin{enumerate}
\item Write the ionisation equation \ce{HA + H2O <=> A- + H3O+} and the expression for $K_a$.
\item Use the approximation $[\ce{H3O+}] = \sqrt{K_a \cdot c}$.
\item Compute $\mathrm{pH} = -\log[\ce{H3O+}]$.
\item \textbf{Check the approximation}: the degree of ionisation $\alpha = [\ce{H3O+}]/c$ must be small (less than about $5\%$). If not, solve the quadratic $x^2 + K_a x - K_a c = 0$ (the positive root is $x = \frac{-K_a + \sqrt{K_a^2 + 4K_a c}}{2}$).
\end{enumerate}
\end{methode}

\begin{attention}[{$[\ce{H3O+}] \neq c$ for a weak acid}]
Because ionisation is partial, $[\ce{H3O+}]$ is \emph{much smaller} than the acid concentration. Writing $\mathrm{pH} = -\log c$ for a weak acid is a classic examination error: it is only valid for strong monoprotic acids.
\end{attention}

\begin{exoresolu}[pH of ethanoic acid]
Calculate the pH of a \SI{0.10}{mol.L^{-1}} solution of ethanoic acid, given $K_a(\ce{CH3COOH}) = 1.8\times 10^{-5}\ \mathrm{mol\,L^{-1}}$ at \SI{25}{\ensuremath{{}^{\circ}\mathrm{C}}}.
\tcblower
\textbf{Solution.} The equilibrium is $\ce{CH3COOH + H2O <=> CH3COO- + H3O+}$.

Applying the approximation:
\[ [\ce{H3O+}] = \sqrt{K_a\, c} = \sqrt{1.8\times10^{-5}\times 0.10} = \sqrt{1.8\times10^{-6}} = 1.34\times10^{-3}\ \mathrm{mol\,L^{-1}}. \]
\[ \mathrm{pH} = -\log(1.34\times10^{-3}) = 2.87. \]
\textbf{Check:} $\alpha = \dfrac{1.34\times10^{-3}}{0.10} = 0.0134 = 1.3\% < 5\%$. The approximation is valid. (Compare: a strong acid of the same concentration would have pH $= 1.00$.)
\end{exoresolu}

\textbf{pH of a weak base.} For \ce{B + H2O <=> BH+ + OH-} with concentration $c_b$, the same reasoning gives $[\ce{OH-}] \approx \sqrt{K_b\, c_b}$; we then use $K_w$ to obtain $[\ce{H3O+}]$ and hence the pH.

\begin{exemplebox}[pH of ammonia solution]
For \SI{0.10}{mol.L^{-1}} \ce{NH3} ($K_b = 1.8\times10^{-5}\ \mathrm{mol\,L^{-1}}$):
\[ [\ce{OH-}] = \sqrt{1.8\times10^{-5}\times 0.10} = 1.34\times10^{-3}\ \mathrm{mol\,L^{-1}}, \]
\[ [\ce{H3O+}] = \frac{1.0\times10^{-14}}{1.34\times10^{-3}} = 7.5\times10^{-12}\ \mathrm{mol\,L^{-1}} \quad\Rightarrow\quad \mathrm{pH} = 11.13. \]
\end{exemplebox}

\courssub{Degree of Ionisation: Dependence on Concentration and Temperature}

The \cle{degree of ionisation} is $\alpha = \dfrac{[\ce{H3O+}]}{c} = \sqrt{\dfrac{K_a}{c}}$ (using the approximation).

\begin{propriete}[Effect of dilution and temperature on $\alpha$]
\begin{itemize}
\item \textbf{Dilution:} as $c$ decreases, $\alpha = \sqrt{K_a/c}$ \emph{increases}: a weak acid ionises more when diluted (this is Ostwald's dilution law). Note, however, that $[\ce{H3O+}] = \sqrt{K_a c}$ still \emph{decreases}, so the pH rises towards 7.
\item \textbf{Temperature:} $K_a$ depends on temperature (ionisation of most weak acids is endothermic, so $K_a$ increases slightly with temperature, by Le Chatelier's principle). $K_w$ also increases strongly with temperature, so the pH of pure water falls below 7 on heating while remaining neutral ($[\ce{H3O+}] = [\ce{OH-}]$).
\end{itemize}
\end{propriete}

\begin{popfigure}
\begin{center}
\begin{tikzpicture}
\begin{axis}[
  width=11cm, height=7cm,
  xlabel={$-\log c$ (increasing dilution $\rightarrow$)},
  ylabel={pH},
  xmin=0, xmax=6, ymin=0.8, ymax=5.2,
  xtick={0,1,2,3,4,5,6},
  legend style={at={(0.03,0.97)}, anchor=north west, font=\small},
  grid=both, grid style={popline!40},
  axis lines=left, thick,
]
% Strong acid: pH = -log c = x
\addplot[domain=0:5, samples=50, very thick, popblue] {x};
\addlegendentry{Strong acid, $\mathrm{pH}=-\log c$}
% Weak acid: pH = 0.5*(pKa + x), pKa=4.74
\addplot[domain=0:5, samples=50, very thick, poporange] {0.5*(4.74+x)};
\addlegendentry{Weak acid, $\mathrm{pH}=\tfrac12(pK_a-\log c)$}
\end{axis}
\end{tikzpicture}
\end{center}
\legende{pH versus dilution for a strong acid and a weak acid ($pK_a = 4.74$) of the same initial concentration. Both pH values rise on dilution; the weak acid is always less acidic (higher pH) and its pH changes by only 0.5 unit per tenfold dilution, against 1 unit for the strong acid.}
\end{popfigure}

\courssub{pH of Mixtures of Acids and Bases}

When an acid and a base are mixed, they \emph{neutralise} each other. The strategy is always the same: treat the neutralisation as a reaction going to completion, find what remains, then compute the pH.

\begin{methode}[pH of an acid--base mixture]
\begin{enumerate}
\item Write the neutralisation equation (ionic form: \ce{H3O+ + OH- -> 2H2O}; with a weak partner, write the molecular equation, e.g.\ \ce{CH3COOH + OH- -> CH3COO- + H2O}).
\item Calculate the moles of acid and base mixed: $n = c\,V$ (volumes in \si{L}).
\item Subtract to find the \textbf{excess moles} of the reagent in excess.
\item Divide by the \textbf{total volume} to get the excess concentration, then compute the pH (using $K_w$ if the excess is \ce{OH-}).
\item If neither is in excess (equivalence), the pH is decided by the \textbf{salt formed} (see hydrolysis, later section).
\end{enumerate}
\end{methode}

\textbf{Case 1: strong acid + strong base.} The reaction \ce{H3O+ + OH- -> 2H2O} is complete. The pH is fixed by whichever ion is in excess; at exact equivalence, pH $= 7$ (at \SI{25}{\ensuremath{{}^{\circ}\mathrm{C}}}).

\begin{exoresolu}[Strong acid in excess]
\SI{20.0}{cm^3} of \SI{0.10}{mol.L^{-1}} \ce{HCl} is mixed with \SI{15.0}{cm^3} of \SI{0.10}{mol.L^{-1}} \ce{NaOH}. Calculate the pH of the mixture.
\tcblower
\textbf{Solution.} Moles mixed:
\[ n(\ce{H3O+}) = 0.10 \times 0.0200 = 2.0\times10^{-3}\ \mathrm{mol}, \qquad n(\ce{OH-}) = 0.10\times 0.0150 = 1.5\times10^{-3}\ \mathrm{mol}. \]
Neutralisation \ce{H3O+ + OH- -> 2H2O} consumes all the \ce{OH-}; excess:
\[ n(\ce{H3O+})_{\text{excess}} = (2.0 - 1.5)\times10^{-3} = 5.0\times10^{-4}\ \mathrm{mol}. \]
Total volume $= 35.0\ \mathrm{cm^3} = 0.0350\ \mathrm{L}$, so
\[ [\ce{H3O+}] = \frac{5.0\times10^{-4}}{0.0350} = 1.43\times10^{-2}\ \mathrm{mol\,L^{-1}} \quad\Rightarrow\quad \mathrm{pH} = 1.85. \]
\end{exoresolu}

\textbf{Case 2: strong acid + weak base.} Example: \ce{HCl + NH3 -> NH4Cl}. The reaction is effectively complete. If the acid is in excess, the pH comes from the excess \ce{H3O+} as above. At equivalence, the solution contains \ce{NH4+}, a weak acid ($K_a = K_w/K_b$), so the pH is \emph{less than 7} — this is salt hydrolysis, treated in a later section.

\textbf{Case 3: weak acid + strong base.} Example: \ce{CH3COOH + NaOH -> CH3COONa + H2O}. If the base is in excess, the excess \ce{OH-} fixes the pH (the weak acid is entirely consumed). At equivalence, the solution contains \ce{CH3COO-}, a weak base, so the pH is \emph{greater than 7}. If the acid is in excess, the mixture contains a weak acid \emph{and} its salt — a \cle{buffer}, whose pH is calculated by the methods of the next section.

\begin{exemplebox}[Base in excess over a weak acid]
Mixing \SI{25.0}{cm^3} of \SI{0.10}{mol.L^{-1}} \ce{CH3COOH} with \SI{30.0}{cm^3} of \SI{0.10}{mol.L^{-1}} \ce{NaOH}: excess $n(\ce{OH-}) = (3.0 - 2.5)\times10^{-3} = 5.0\times10^{-4}$ mol in \SI{55.0}{cm^3}, giving $[\ce{OH-}] = 9.1\times10^{-3}\ \mathrm{mol\,L^{-1}}$, $[\ce{H3O+}] = 1.1\times10^{-12}\ \mathrm{mol\,L^{-1}}$, pH $= 11.96$.
\end{exemplebox}

\begin{attention}[Do not forget the total volume]
After mixing, the excess reagent is diluted in the \emph{combined} volume. Dividing the excess moles by the volume of one solution only is a very common source of error.
\end{attention}

\begin{aretenir}[Key points]
\begin{itemize}
\item Strong acids/bases ionise completely: $\mathrm{pH} = -\log c_a$ (monoprotic acid) and $\mathrm{pH} = 14 + \log c_b$ (monobasic base) at \SI{25}{\ensuremath{{}^{\circ}\mathrm{C}}}. For \ce{Ba(OH)2}, $[\ce{OH-}] = 2c_b$.
\item Weak acids/bases ionise partially: $[\ce{H3O+}] = \sqrt{K_a c}$ (acid), $[\ce{OH-}] = \sqrt{K_b c}$ (base); always check $\alpha < 5\%$.
\item $pK_a = -\log K_a$; for a conjugate pair $pK_a + pK_b = 14$ at \SI{25}{\ensuremath{{}^{\circ}\mathrm{C}}}.
\item Dilution increases the degree of ionisation of a weak acid (Ostwald's law) but raises its pH towards 7.
\item For mixtures: write the neutralisation equation, find the excess moles, divide by the total volume, then compute the pH. At equivalence the pH depends on the salt (hydrolysis).
\end{itemize}
\end{aretenir}

\begin{exercice}[Practice]
\begin{enumerate}
\item Calculate the pH of: (a) \SI{0.025}{mol.L^{-1}} \ce{HNO3}; (b) \SI{0.050}{mol.L^{-1}} \ce{KOH}.
\item Calculate the pH of \SI{0.050}{mol.L^{-1}} methanoic acid \ce{HCOOH} ($K_a = 1.8\times10^{-4}\ \mathrm{mol\,L^{-1}}$) and its degree of ionisation.
\item \SI{30.0}{cm^3} of \SI{0.10}{mol.L^{-1}} \ce{NaOH} is added to \SI{25.0}{cm^3} of \SI{0.10}{mol.L^{-1}} \ce{HCl}. Find the pH of the mixture.
\item Explain, with an equation, why the pH at equivalence of a \ce{CH3COOH}--\ce{NaOH} titration is greater than 7.
\end{enumerate}
\end{exercice}

\begin{corrige}[Exercise 1]
\textbf{1.} (a) Strong acid: $\mathrm{pH} = -\log(0.025) = 1.60$. (b) $[\ce{H3O+}] = 10^{-14}/0.050 = 2.0\times10^{-13}$, $\mathrm{pH} = 12.70$.

\textbf{2.} $[\ce{H3O+}] = \sqrt{1.8\times10^{-4}\times0.050} = 3.0\times10^{-3}\ \mathrm{mol\,L^{-1}}$; $\mathrm{pH} = 2.52$; $\alpha = 3.0\times10^{-3}/0.050 = 6\%$ (borderline; the quadratic gives pH $\approx 2.55$).

\textbf{3.} Excess $n(\ce{OH-}) = (3.0-2.5)\times10^{-3} = 5.0\times10^{-4}$ mol in \SI{55.0}{cm^3}: $[\ce{OH-}] = 9.1\times10^{-3}\ \mathrm{mol\,L^{-1}}$, $\mathrm{pH} = 11.96$.

\textbf{4.} At equivalence only \ce{CH3COONa} is present; \ce{CH3COO-} is a weak base: \ce{CH3COO- + H2O <=> CH3COOH + OH-}, producing \ce{OH-}, so pH $> 7$.
\end{corrige}

\coursec{Acid-Base Indicators and Titration Curves}

In the previous section we learnt how to calculate the pH of strong and weak acids and bases, and of their mixtures. We now put that knowledge to work in one of the most important practical activities of the Upper Sixth course: the \cle{acid-base titration}. In a titration we add a base of known concentration progressively to an acid (or vice versa) and follow the pH as the reaction proceeds. Two questions immediately arise: how do we know \emph{when} the acid has been exactly neutralised, and how does the pH vary during the addition? The answers are, respectively, \emph{acid-base indicators} and \emph{titration curves} — the two subjects of this section.

\courssub{Acid-Base Indicators and How They Work}

\begin{experience}[Observing an indicator]
Place a few drops of methyl orange in a beaker of dilute hydrochloric acid: the solution is \textbf{red}. Add dilute sodium hydroxide solution drop by drop while stirring. At first nothing seems to happen; then, quite suddenly, one extra drop turns the whole solution \textbf{yellow}. The indicator has ``signalled'' that neutralisation is complete. How can a dye ``know'' the pH?
\end{experience}

\begin{definition}[Acid-base indicator]
An \cle{acid-base indicator} is a \textbf{weak acid} (or weak base) whose \textbf{acid form and conjugate base form have different colours}. For an indicator that is a weak acid, we write the equilibrium:
\[ \ce{HIn(aq) <=> H+(aq) + In-(aq)} \qquad K_{In} = \frac{[\ce{H+}][\ce{In-}]}{[\ce{HIn}]} \]
where \ce{HIn} and \ce{In-} have different colours. We define $pK_{In} = -\log K_{In}$.
\end{definition}

Rearranging the equilibrium expression gives the key to the indicator's behaviour:
\[ [\ce{H+}] = K_{In}\,\frac{[\ce{HIn}]}{[\ce{In-}]} \qquad\Longrightarrow\qquad \mathrm{pH} = pK_{In} + \log\frac{[\ce{In-}]}{[\ce{HIn}]} \]

\begin{itemize}
\item In \textbf{acidic} solution (excess \ce{H+}), Le Chatelier's principle pushes the equilibrium to the left: the colour of \ce{HIn} dominates.
\item In \textbf{alkaline} solution, \ce{H+} is removed, the equilibrium shifts right: the colour of \ce{In-} dominates.
\item When $[\ce{HIn}] = [\ce{In-}]$, $\mathrm{pH} = pK_{In}$: this is the \cle{transition point}, where the intermediate (mixed) colour is seen.
\end{itemize}

The human eye can only detect that one form dominates when it is roughly \textbf{ten times} more concentrated than the other. Since $\log 10 = 1$, this leads to a very useful rule.

\begin{propriete}[Colour-change range of an indicator]
An indicator changes colour gradually over a pH range of approximately
\[ \mathrm{pH} = pK_{In} \pm 1 \]
Below $pK_{In}-1$ the acid colour is seen; above $pK_{In}+1$ the base colour is seen; in between, the colour changes progressively. (Proof not examinable, but you must be able to justify it using the ratio $[\ce{In-}]/[\ce{HIn}]$ varying from $1/10$ to $10$.)
\end{propriete}

\begin{center}
\begin{tabularx}{\linewidth}{|l|X|X|X|X|}
\hline
\rowcolor{popblueL} \textbf{Indicator} & \textbf{$pK_{In}$} & \textbf{pH range} & \textbf{Acid colour} & \textbf{Base colour} \\
\hline
Methyl orange & 3.7 & 3.1 -- 4.4 & red & yellow \\
\hline
Litmus & 7.0 & 5.0 -- 8.0 & red & blue \\
\hline
Phenolphthalein & 9.3 & 8.3 -- 10.0 & colourless & pink \\
\hline
\end{tabularx}
\end{center}

\begin{attention}[Litmus is a poor titration indicator]
Litmus changes over the very wide range pH 5--8, so its colour change is gradual and imprecise. It is fine for testing whether a solution is acidic or alkaline, but \textbf{never} choose it to detect the end point of a titration.
\end{attention}

\courssub{pH Changes During a Titration: Equivalence Point and End Point}

Consider adding \SI{0.100}{mol.dm^{-3}} \ce{NaOH} from a burette to \SI{25.0}{cm^3} of \SI{0.100}{mol.dm^{-3}} \ce{HCl} in a conical flask. The reaction is:
\[ \ce{HCl(aq) + NaOH(aq) -> NaCl(aq) + H2O(l)} \qquad\text{i.e.}\qquad \ce{H+(aq) + OH-(aq) -> H2O(l)} \]

Let us track the pH using the methods of the previous section.

\begin{itemize}
\item \textbf{Start (0 cm$^3$):} strong acid, $[\ce{H+}] = \SI{0.100}{mol.dm^{-3}}$, so $\mathrm{pH} = 1.00$.
\item \textbf{After 10.0 cm$^3$:} moles \ce{H+} remaining $= 0.100\times(25.0-10.0)\times10^{-3} = 1.50\times10^{-3}$ mol in a total volume of \SI{35.0}{cm^3}, giving $[\ce{H+}] = \SI{0.0429}{mol.dm^{-3}}$ and $\mathrm{pH} = 1.37$. The pH has barely moved!
\item \textbf{After 24.9 cm$^3$:} $[\ce{H+}] = \dfrac{0.100\times0.1\times10^{-3}}{49.9\times10^{-3}} = 2.0\times10^{-4}\ \mathrm{mol\,dm^{-3}}$, so $\mathrm{pH} = 3.70$.
\item \textbf{After 25.0 cm$^3$:} exactly enough \ce{OH-} has been added to react with all the \ce{H+}. Only \ce{NaCl} and water remain: $\mathrm{pH} = 7.00$.
\item \textbf{After 25.1 cm$^3$:} excess \ce{OH-} $= 0.100\times0.1\times10^{-3} = 1.0\times10^{-5}$ mol in \SI{50.1}{cm^3}, giving $[\ce{OH-}] = 2.0\times10^{-4}\ \mathrm{mol\,dm^{-3}}$, $\mathrm{pOH} = 3.70$, $\mathrm{pH} = 10.30$.
\end{itemize}

Look at the last three values: a single drop (from 24.9 to 25.1 cm$^3$) makes the pH jump from 3.7 to 10.3 — a leap of nearly 7 pH units! This near-vertical jump is the most important feature of every titration curve.

\begin{definition}[Equivalence point and end point]
The \cle{equivalence point} of a titration is the point at which the amount of titrant added is \textbf{exactly the stoichiometric amount} needed to react completely with the analyte ($n_{\text{acid}} = n_{\text{base}}$, adjusted for stoichiometric coefficients). It is a fixed, theoretical point.

The \cle{end point} is the point at which the \textbf{indicator changes colour}. It is an experimental observation. A good titration is one where the end point coincides with the equivalence point.
\end{definition}

\begin{attention}[Equivalence point $\neq$ end point]
These two words are \textbf{not} synonyms, and GCE examiners penalise candidates who confuse them. The equivalence point is stoichiometric; the end point depends on the indicator chosen. We choose the indicator precisely so that its end point falls at (or extremely close to) the equivalence point.
\end{attention}

\courssub{Titration Curves and the Choice of Indicator}

A \cle{titration curve} is a graph of pH (vertical axis) against volume of titrant added (horizontal axis). Its shape depends on the \emph{strengths} of the acid and the base.

\textbf{Strong acid titrated with strong base (SA--SB).}\par
The curve starts at very low pH ($\approx 1$), rises very slowly, then shoots up almost vertically through pH 7 at the equivalence point, and finally levels off near pH 13. The vertical region is very long (roughly pH 3 to 11). \textbf{Any indicator} whose range lies inside this window is suitable: both methyl orange and phenolphthalein work.

\textbf{Weak acid titrated with strong base (WA--SB).}\par
The curve starts higher (pH $\approx 3$ for a \SI{0.1}{mol.dm^{-3}} weak acid) and shows a gently rising ``buffer region'' before equivalence. At half-equivalence, $[\ce{HA}] = [\ce{A-}]$, so $\mathrm{pH} = pK_a$ — a very useful way of measuring $pK_a$ experimentally. At the equivalence point the solution contains only the salt \ce{NaA}, whose anion \ce{A-} is a weak base (salt hydrolysis, next section): \textbf{the pH at equivalence is above 7} (typically 8--9). The vertical region is short (roughly pH 7 to 11), so only \textbf{phenolphthalein} is suitable; methyl orange would change colour far too early.

\begin{attention}[The equivalence point of a WA--SB titration is NOT at pH 7]
A very common mistake is to write ``pH = 7 at equivalence'' for every titration. pH 7 at equivalence occurs \textbf{only} for a strong acid--strong base titration. For a weak acid--strong base titration the equivalence pH is \textbf{above 7} (basic salt); for a strong acid--weak base titration it is \textbf{below 7} (acidic salt).
\end{attention}

\textbf{Strong acid titrated with weak base (SA--WB).}\par
The curve starts at low pH; the vertical region lies roughly between pH 3 and 7, with the equivalence point \textbf{below 7} (the salt, e.g.\ \ce{NH4Cl}, is acidic). \textbf{Methyl orange} is the suitable indicator; phenolphthalein would change long after equivalence.

\textbf{Weak acid titrated with weak base (WA--WB).}\par
There is \textbf{no sharp vertical region}: the pH drifts gradually through the equivalence point. \textbf{No common indicator} changes sharply enough to locate the equivalence point, so this titration is never performed with an indicator (a pH meter must be used instead).

\begin{methode}[Choosing the correct indicator]
\begin{enumerate}
\item Identify the \textbf{strengths} of the acid and the base (strong or weak).
\item Sketch (or recall) the titration curve and locate the \textbf{vertical region} and the \textbf{equivalence point pH}.
\item Choose the indicator whose \textbf{range $pK_{In}\pm1$ falls entirely within the vertical region}.
\item Check: SA--SB $\rightarrow$ either indicator; WA--SB $\rightarrow$ phenolphthalein; SA--WB $\rightarrow$ methyl orange; WA--WB $\rightarrow$ no suitable indicator.
\end{enumerate}
\end{methode}

\begin{popfigure}
\begin{center}
\begin{tikzpicture}
\begin{axis}[
  width=13cm, height=9.5cm,
  xlabel={Volume of base or acid added / cm$^3$},
  ylabel={pH},
  xmin=0, xmax=50, ymin=0, ymax=14,
  xtick={0,10,20,25,30,40,50}, ytick={0,2,4,6,7,8,10,12,14},
  grid=both, grid style={popline, very thin},
  legend style={at={(0.03,0.97)}, anchor=north west, font=\small, draw=popline},
  axis line style={popdark},
]
% indicator ranges shaded
\fill[popgreen, opacity=0.15] (axis cs:0,8.3) rectangle (axis cs:50,10.0);
\fill[poporange, opacity=0.18] (axis cs:0,3.1) rectangle (axis cs:50,4.4);
\node[font=\scriptsize, text=popgreen!60!black] at (axis cs:44,9.15) {phenolphthalein};
\node[font=\scriptsize, text=poporange!80!black] at (axis cs:44,3.75) {methyl orange};
% SA-SB curve
\addplot[very thick, popblue, smooth] coordinates {
(0,1.00) (5,1.18) (10,1.37) (15,1.60) (20,1.95) (22,2.15) (24,2.60)
(24.5,2.95) (24.9,3.70) (25,7.00) (25.1,10.30) (25.5,11.00) (26,11.30)
(28,11.75) (30,11.96) (35,12.22) (40,12.36) (50,12.52)};
\addlegendentry{SA--SB (\ce{NaOH} added to HCl)}
% WA-SB curve
\addplot[very thick, poppurple, smooth, dashed] coordinates {
(0,2.88) (5,3.80) (10,4.15) (12.5,4.35) (15,4.57) (20,5.05) (22,5.35)
(24,5.95) (24.5,6.45) (24.9,7.30) (25,8.72) (25.1,10.10) (25.5,10.80)
(26,11.10) (28,11.60) (30,11.85) (35,12.15) (40,12.30) (50,12.48)};
\addlegendentry{WA--SB (\ce{NaOH} added to \ce{CH3COOH})}
% SA-WB curve
\addplot[very thick, poppink, smooth, dotted] coordinates {
(0,1.00) (5,1.18) (10,1.37) (15,1.60) (20,1.95) (22,2.15) (24,2.60)
(24.5,2.95) (24.9,3.70) (25,5.28) (25.1,4.40) (25.5,4.10) (26,3.95)
(28,3.70) (30,3.55) (35,3.30) (40,3.15) (50,2.95)};
\addlegendentry{SA--WB (\ce{NH3} added to HCl)}
% equivalence marker
\draw[popmuted, dashed] (axis cs:25,0) -- (axis cs:25,14);
\node[font=\scriptsize, text=popmuted, rotate=90, anchor=south] at (axis cs:25.4,0.3) {equivalence (25 cm$^3$)};
\end{axis}
\end{tikzpicture}
\end{center}
\legende{Titration curves for \SI{0.1}{mol.dm^{-3}} solutions: strong acid--strong base (solid blue), weak acid--strong base (dashed purple) and strong acid--weak base (dotted pink), with the colour-change ranges of methyl orange and phenolphthalein shaded. An indicator is suitable only if its band lies within the vertical part of the curve.}
\end{popfigure}

\begin{exoresolu}[Choosing and justifying an indicator]
\SI{25.0}{cm^3} of \SI{0.100}{mol.dm^{-3}} ethanoic acid, \ce{CH3COOH} ($K_a = 1.8\times10^{-5}\ \mathrm{mol\,dm^{-3}}$), is titrated with \SI{0.100}{mol.dm^{-3}} \ce{NaOH}. (a) Calculate the pH at the start and at half-equivalence. (b) State, with reasons, which of methyl orange or phenolphthalein is suitable.
\tcblower
\textbf{(a)} At the start, for a weak acid $[\ce{H+}] \approx \sqrt{K_a\,c} = \sqrt{1.8\times10^{-5}\times0.100} = 1.34\times10^{-3}\ \mathrm{mol\,dm^{-3}}$, so $\mathrm{pH} = 2.87$.

At half-equivalence, half the acid has been converted to \ce{CH3COO-}, so $[\ce{CH3COOH}] = [\ce{CH3COO-}]$ and
\[ \mathrm{pH} = pK_a = -\log(1.8\times10^{-5}) = 4.74. \]

\textbf{(b)} This is a \textbf{weak acid--strong base} titration: the equivalence point lies \textbf{above pH 7} (the ethanoate ion is basic) and the vertical region runs roughly from pH 7 to 11. Phenolphthalein (range 8.3--10.0) lies inside this region and is \textbf{suitable}. Methyl orange (3.1--4.4) would change colour well before equivalence and is \textbf{unsuitable}.
\end{exoresolu}

\courssub{Titrations Involving Diprotic Species: Two Equivalence Points}

Some acids can donate \textbf{two} protons (diprotic acids, e.g.\ \ce{H2SO4}, \ce{H2CO3}) and some bases can accept two protons (e.g.\ the carbonate ion \ce{CO3^2-}). When such a species is titrated with a monoprotic partner, the reaction occurs in \textbf{two successive steps}, giving \textbf{two equivalence points} and \textbf{two vertical regions} on the curve.

The classic GCE example is the titration of sodium carbonate solution with dilute hydrochloric acid:
\begin{align*}
\text{Step 1:}\quad & \ce{Na2CO3(aq) + HCl(aq) -> NaHCO3(aq) + NaCl(aq)} && \text{i.e. } \ce{CO3^2- + H+ -> HCO3-} \\
\text{Step 2:}\quad & \ce{NaHCO3(aq) + HCl(aq) -> NaCl(aq) + H2O(l) + CO2(g)} && \text{i.e. } \ce{HCO3- + H+ -> H2CO3}
\end{align*}

The first equivalence point (near pH 8.3, where the solution contains \ce{NaHCO3}) is detected by \textbf{phenolphthalein} (pink $\rightarrow$ colourless). The second equivalence point (near pH 3.9, where all carbonate has become \ce{CO2} and water) is detected by \textbf{methyl orange} (yellow $\rightarrow$ orange/red). Note that the volume of acid needed for step 2 \emph{equals} the volume for step 1, because each mole of \ce{CO3^2-} produces exactly one mole of \ce{HCO3-}, which then consumes a second mole of \ce{H+}.

\begin{popfigure}
\begin{center}
\begin{tikzpicture}
\begin{axis}[
  width=13cm, height=9cm,
  xlabel={Volume of \ce{HCl} added / cm$^3$},
  ylabel={pH},
  xmin=0, xmax=60, ymin=0, ymax=14,
  xtick={0,10,20,25,30,40,50,60}, ytick={0,2,4,6,8,10,12,14},
  grid=both, grid style={popline, very thin},
  axis line style={popdark},
]
\fill[popgreen, opacity=0.15] (axis cs:0,8.3) rectangle (axis cs:60,10.0);
\fill[poporange, opacity=0.18] (axis cs:0,3.1) rectangle (axis cs:60,4.4);
\node[font=\scriptsize, text=popgreen!60!black] at (axis cs:52,9.15) {phenolphthalein};
\node[font=\scriptsize, text=poporange!80!black] at (axis cs:52,3.75) {methyl orange};
\addplot[very thick, popblue, smooth] coordinates {
(0,11.60) (5,11.10) (10,10.70) (15,10.30) (20,9.80) (22,9.40) (24,8.90)
(24.8,8.50) (25,8.30) (25.2,8.10) (26,7.60) (28,7.10) (30,6.80) (35,6.40)
(40,6.00) (45,5.40) (48,4.70) (49.5,4.20) (50,3.90) (50.5,3.50) (52,3.00)
(55,2.60) (60,2.40)};
\draw[popmuted, dashed] (axis cs:25,0) -- (axis cs:25,14);
\draw[popmuted, dashed] (axis cs:50,0) -- (axis cs:50,14);
\node[font=\scriptsize, text=popdark, align=center] at (axis cs:25,12.6) {1st equivalence\\ (\ce{NaHCO3})};
\node[font=\scriptsize, text=popdark, align=center] at (axis cs:50,12.6) {2nd equivalence\\ (\ce{CO2} + \ce{H2O})};
\end{axis}
\end{tikzpicture}
\end{center}
\legende{Titration of \SI{25.0}{cm^3} of \SI{0.10}{mol.dm^{-3}} \ce{Na2CO3} with \SI{0.10}{mol.dm^{-3}} \ce{HCl}. Two protons are accepted in succession, giving two equivalence points of equal spacing: the first within the phenolphthalein range, the second within the methyl orange range.}
\end{popfigure}

\begin{exemplebox}[Reading the two end points]
In the titration above, phenolphthalein decolourises after \SI{25.0}{cm^3} of acid and methyl orange turns after a further \SI{25.0}{cm^3} (\SI{50.0}{cm^3} in total). The \emph{equality} of the two volumes confirms that the base is \ce{CO3^2-} (diprotic): $n(\ce{HCl})$ total $= 2\,n(\ce{Na2CO3})$. This ``double indicator'' technique is a favourite GCE practical question.
\end{exemplebox}

\begin{exoresolu}[Double-indicator calculation]
\SI{25.0}{cm^3} of a sodium carbonate solution is titrated with \SI{0.100}{mol.dm^{-3}} \ce{HCl}. Phenolphthalein decolourises after \SI{12.50}{cm^3} of acid; methyl orange then changes after a further \SI{12.50}{cm^3}. Calculate the concentration of the \ce{Na2CO3} solution.
\tcblower
The first end point corresponds to step 1 only: $\ce{CO3^2- + H+ -> HCO3-}$, so
\[ n(\ce{Na2CO3}) = n(\ce{HCl})_{\text{step 1}} = 0.100\times12.50\times10^{-3} = 1.25\times10^{-3}\ \mathrm{mol}. \]
Hence
\[ [\ce{Na2CO3}] = \frac{1.25\times10^{-3}}{25.0\times10^{-3}} = \SI{0.0500}{mol.dm^{-3}}. \]
Check with the total titre: $n(\ce{HCl})_{\text{total}} = 0.100\times25.0\times10^{-3} = 2.50\times10^{-3}$ mol $= 2\,n(\ce{Na2CO3})$, as expected for a diprotic base.
\end{exoresolu}

\courssub{Reading Titration Curves: Finding $pK_a$ at Half-Equivalence}

A titration curve is not only a way to detect equivalence — it is also an \textbf{experimental source of data}. For a weak acid titrated with a strong base, at the \textbf{half-equivalence point} exactly half of the acid \ce{HA} has been converted into its conjugate base \ce{A-}. Then:
\[ [\ce{HA}] = [\ce{A-}] \quad\Longrightarrow\quad \mathrm{pH} = pK_a + \log\frac{[\ce{A-}]}{[\ce{HA}]} = pK_a \]

\begin{methode}[Determining $pK_a$ from a titration curve (GCE skill)]
\begin{enumerate}
\item Locate the \textbf{equivalence point} (midpoint of the vertical jump) and read its volume $V_e$.
\item Halve it: at $V_e/2$, read the pH from the curve.
\item This pH \textbf{is} the $pK_a$ of the weak acid; then $K_a = 10^{-pK_a}$.
\end{enumerate}
The same logic applies to a weak base titrated with a strong acid, giving $pK_b$ (or $pK_a$ of the conjugate acid) at half-equivalence.
\end{methode}

\begin{exoresolu}[Reading $pK_a$ from a curve]
The titration of \SI{25.0}{cm^3} of a weak acid \ce{HA} with \SI{0.100}{mol.dm^{-3}} \ce{NaOH} shows a vertical jump centred at \SI{20.0}{cm^3}. The pH read at \SI{10.0}{cm^3} of added base is 4.80. Determine $pK_a$ and $K_a$ of the acid, and its initial concentration.
\tcblower
Half-equivalence occurs at $V_e/2 = \SI{10.0}{cm^3}$, so $pK_a = \mathrm{pH} = 4.80$ and
\[ K_a = 10^{-4.80} = 1.6\times10^{-5}\ \mathrm{mol\,dm^{-3}}. \]
At equivalence, $n(\ce{HA}) = n(\ce{NaOH}) = 0.100\times20.0\times10^{-3} = 2.00\times10^{-3}$ mol in \SI{25.0}{cm^3}, so $c(\ce{HA}) = \dfrac{2.00\times10^{-3}}{25.0\times10^{-3}} = \SI{0.0800}{mol.dm^{-3}}$.
\end{exoresolu}

\begin{aretenir}[Key points of this section]
\begin{itemize}
\item An indicator is a weak acid/base whose two forms have different colours; it changes over $\mathrm{pH} = pK_{In}\pm1$.
\item The \textbf{equivalence point} is stoichiometric; the \textbf{end point} is the indicator's colour change — choose the indicator so the two coincide.
\item SA--SB: equivalence at pH 7, any common indicator. WA--SB: equivalence \textbf{above} 7, phenolphthalein. SA--WB: equivalence \textbf{below} 7, methyl orange. WA--WB: \textbf{no} suitable indicator.
\item Diprotic systems (e.g.\ \ce{Na2CO3} with \ce{HCl}) show \textbf{two} equivalence points of equal spacing; phenolphthalein detects the first, methyl orange the second.
\item At \textbf{half-equivalence}, $\mathrm{pH} = pK_a$: a powerful way to read $K_a$ off a titration curve.
\end{itemize}
\end{aretenir}

\begin{exercice}[Indicator choice and curve reading]
\begin{enumerate}
\item \SI{20.0}{cm^3} of \SI{0.0500}{mol.dm^{-3}} ammonia solution is titrated with \SI{0.100}{mol.dm^{-3}} \ce{HCl}. (a) Calculate the volume of acid needed to reach equivalence. (b) State whether the pH at equivalence is above, below or equal to 7, justifying your answer with an equation. (c) Choose a suitable indicator from methyl orange and phenolphthalein.
\item A student titrates ethanoic acid with sodium hydroxide using methyl orange. Explain why the titration result is unreliable, and name the indicator that should have been used.
\item Sketch the titration curve for \SI{25}{cm^3} of \SI{0.1}{mol.dm^{-3}} \ce{Na2CO3} titrated with \SI{0.1}{mol.dm^{-3}} \ce{HCl}, marking both equivalence volumes and the two suitable indicators.
\end{enumerate}
\end{exercice}

\begin{corrige}[Exercise 1]
\textbf{1.} (a) $n(\ce{NH3}) = 0.0500\times20.0\times10^{-3} = 1.00\times10^{-3}$ mol; reaction is 1:1, so $V(\ce{HCl}) = \dfrac{1.00\times10^{-3}}{0.100} = \SI{10.0}{cm^3}$. (b) Below 7: the salt \ce{NH4Cl} formed contains \ce{NH4+}, a weak acid: $\ce{NH4+(aq) + H2O(l) <=> NH3(aq) + H3O+(aq)}$. (c) Strong acid--weak base: vertical region around pH 3--7, so \textbf{methyl orange}.

\textbf{2.} Ethanoic acid is weak and \ce{NaOH} is strong: the equivalence point is above pH 7 and the vertical region is roughly pH 7--11. Methyl orange (3.1--4.4) changes colour well before equivalence, giving a falsely small titre. \textbf{Phenolphthalein} should be used.

\textbf{3.} Curve starts near pH 11.6; first vertical region (phenolphthalein range) at \SI{25}{cm^3} (\ce{CO3^2- -> HCO3-}); buffer-like region around pH 6.3; second vertical region (methyl orange range) at \SI{50}{cm^3} (\ce{HCO3- -> H2CO3}); levels near pH 2--3 in excess acid.
\end{corrige}

\coursec{Buffer Solutions}

In the previous section we followed the pH during a titration and saw how violently it can change near the equivalence point: a single drop of sodium hydroxide can send the pH of a strong acid solution from 3 to 11. In many real situations, however, such swings are unacceptable. Your blood must stay within a very narrow pH window (about 7.35 to 7.45) or enzymes stop working; the palm-wine tappers of the South-West Region rely on a stable pH for fermentation; a shampoo that is too acidic or too alkaline damages hair. Nature and industry solve this problem with \cle{buffer solutions}, the subject of this section.

\courssub{What is a Buffer Solution?}

\textbf{An observation to start with.}\par

Take two beakers, each containing \SI{100}{cm^3} of liquid at pH 7. The first contains pure water, the second a mixture of ethanoic acid and sodium ethanoate. Add \SI{1}{cm^3} of \SI{1.0}{mol.dm^{-3}} hydrochloric acid to each. In the water, the pH crashes from 7 to about 2. In the second beaker, the pH barely moves, perhaps from 4.76 to 4.75. The second liquid has \emph{resisted} the change. It is a buffer.

\begin{definition}[Buffer solution]
A \cle{buffer solution} is a solution that \textbf{resists changes in pH} when \textbf{small amounts} of acid (\ce{H+}) or base (\ce{OH-}) are added to it, or when it is moderately diluted.
\end{definition}

\begin{attention}[A buffer does not fix the pH]
A buffer does \textbf{not} keep the pH perfectly constant. It only \emph{minimises} the change. Adding acid to a buffer always lowers the pH slightly; adding base always raises it slightly. If you keep adding acid, the buffer is eventually overwhelmed and the pH falls sharply. Examiners love the statement ``a buffer keeps the pH constant'' — it is false and costs marks.
\end{attention}

\courssub{Composition and Types of Buffers}

How can a solution fight both added acid \emph{and} added base? It must contain, in large amounts, \textbf{both} a species that can consume \ce{H+} and a species that can consume \ce{OH-}. A weak acid \ce{HA} and its conjugate base \ce{A-} do exactly this.

\begin{propriete}[Composition of a buffer]
An effective buffer contains \textbf{large, comparable amounts} of a \textbf{weak acid and its conjugate base} (or a weak base and its conjugate acid). The two species must be a conjugate pair, and neither may be strong.
\end{propriete}

There are two classic types.

\textbf{1. Acidic buffers (pH $<$ 7).} Made from a \textbf{weak acid and one of its salts} (which supplies the conjugate base). Typical example: ethanoic acid and sodium ethanoate, \ce{CH3COOH/CH3COONa}. Such a buffer can also be made by partially neutralising an excess of weak acid with sodium hydroxide: the unreacted acid and the ethanoate formed then coexist.

\textbf{2. Alkaline buffers (pH $>$ 7).} Made from a \textbf{weak base and one of its salts} (which supplies the conjugate acid). Typical example: ammonia and ammonium chloride, \ce{NH3/NH4Cl}.

\begin{attention}[A strong acid and its salt do NOT form a buffer]
\ce{HCl} and \ce{NaCl} mixed together are \textbf{not} a buffer. \ce{HCl} is fully ionised, so there is no reservoir of undissociated \ce{HA} to absorb added \ce{OH-}; and \ce{Cl-} is such a weak base that it cannot absorb added \ce{H+}. Only a \textbf{weak} acid/base pair works, because the undissociated weak species and its conjugate partner must both exist in large amounts.
\end{attention}

\courssub{How a Buffer Works}

\textbf{The acidic buffer \ce{CH3COOH/CH3COO-}.}\par

In the mixture, the weak acid is barely ionised, so we have a large reservoir of \ce{CH3COOH} molecules and (from the fully ionised salt) a large reservoir of \ce{CH3COO-} ions.

\begin{itemize}
\item \textbf{Add acid (\ce{H+}):} the ethanoate ions mop it up:
\[ \ce{CH3COO-(aq) + H+(aq) -> CH3COOH(aq)} \]
The added \ce{H+} is converted into weakly ionised ethanoic acid, so the free \ce{H+} concentration hardly changes.
\item \textbf{Add base (\ce{OH-}):} the ethanoic acid neutralises it:
\[ \ce{CH3COOH(aq) + OH-(aq) -> CH3COO-(aq) + H2O(l)} \]
The added \ce{OH-} is removed as water.
\end{itemize}

In both cases the added strong reagent is ``swallowed'' by one reservoir and converted into the other partner of the pair; only the \emph{ratio} $[\ce{CH3COO-}]/[\ce{CH3COOH}]$ shifts slightly, and we shall see that the pH depends only on this ratio.

\textbf{The alkaline buffer \ce{NH3/NH4+}.}\par

\begin{itemize}
\item \textbf{Add acid:} $\ce{NH3(aq) + H+(aq) -> NH4+(aq)}$
\item \textbf{Add base:} $\ce{NH4+(aq) + OH-(aq) -> NH3(aq) + H2O(l)}$
\end{itemize}

\begin{popfigure}
\begin{center}
\begin{tikzpicture}[font=\small]
% left reservoir
\draw[fill=popblueL, draw=popblue, rounded corners] (0,0) rectangle (3.4,2.2);
\node[align=center] at (1.7,1.6) {\textbf{Reservoir of \ce{HA}}\\ (weak acid)};
\node[align=center] at (1.7,0.6) {absorbs added \ce{OH-}\\ \ce{HA + OH- -> A- + H2O}};
% right reservoir
\draw[fill=popgreenL, draw=popgreen, rounded corners] (5.0,0) rectangle (8.4,2.2);
\node[align=center] at (6.7,1.6) {\textbf{Reservoir of \ce{A-}}\\ (conjugate base)};
\node[align=center] at (6.7,0.6) {absorbs added \ce{H+}\\ \ce{A- + H+ -> HA}};
% arrows
\draw[-{Stealth}, very thick, poporange] (4.2,3.2) -- (6.2,2.6) node[midway, above, align=center]{add \ce{H+}};
\draw[-{Stealth}, very thick, poppurple] (2.2,2.6) -- (4.2,3.2) node[midway, above, align=center]{add \ce{OH-}};
\draw[{Stealth}-{Stealth}, thick, popdark] (3.4,1.1) -- (5.0,1.1) node[midway, below]{conjugate pair};
\node[align=center] at (4.2,-0.7) {pH changes only slightly};
\end{tikzpicture}
\end{center}
\legende{An acidic buffer: two large reservoirs, \ce{HA} and \ce{A-}, neutralise added base and added acid respectively.}
\end{popfigure}

\courssub{Calculating the pH of an Acidic Buffer}

For the weak acid equilibrium $\ce{HA(aq) + H2O(l) <=> H3O+(aq) + A-(aq)}$:
\[ K_a = \frac{[\ce{H3O+}][\ce{A-}]}{[\ce{HA}]} \]
Rearranging for $[\ce{H3O+}]$:
\[ [\ce{H3O+}] = K_a \times \frac{[\ce{HA}]}{[\ce{A-}]} \]
Because the acid is weak and the concentrations of acid and salt are reasonably large (typically $> \SI{0.01}{mol.dm^{-3}}$), the equilibrium concentrations are essentially the \emph{analytical} concentrations of acid and salt mixed: $[\ce{HA}] \approx c_{\text{acid}}$ and $[\ce{A-}] \approx c_{\text{salt}}$. Taking $-\log$ of both sides gives the celebrated \cle{Henderson--Hasselbalch equation}:

\begin{propriete}[Henderson--Hasselbalch equation]
For an acidic buffer made of a weak acid \ce{HA} and its salt \ce{A-}:
\[ \mathrm{pH} = \mathrm{p}K_a + \log\frac{[\text{salt}]}{[\text{acid}]} = \mathrm{p}K_a + \log\frac{[\ce{A-}]}{[\ce{HA}]} \]
Since the two species share the same volume, concentrations may be replaced by \textbf{moles}. The derivation from $K_a$ is examinable and you should be able to reproduce it.
\end{propriete}

Note the immediate consequences: when $[\text{salt}] = [\text{acid}]$, $\mathrm{pH} = \mathrm{p}K_a$; and because pH depends on a \emph{ratio}, diluting the buffer with water changes both concentrations equally and leaves the pH (almost) unchanged — a point the GCE likes to test.

\begin{methode}[Calculating the pH of an acidic buffer]
\begin{enumerate}
\item Identify the weak acid and its conjugate base (the salt).
\item Write $K_a$ and rearrange: $[\ce{H3O+}] = K_a \dfrac{[\text{acid}]}{[\text{salt}]}$.
\item Substitute concentrations (or moles) of acid and salt.
\item Compute $[\ce{H3O+}]$, then $\mathrm{pH} = -\log[\ce{H3O+}]$ — or use Henderson--Hasselbalch directly.
\item To find the pH \emph{after} adding strong acid or base, first adjust the moles of \ce{HA} and \ce{A-} using the neutralisation equations, then reapply the formula.
\end{enumerate}
\end{methode}

\begin{exoresolu}[pH of an ethanoic acid / ethanoate buffer]
A buffer is prepared by dissolving \SI{0.10}{mol} of ethanoic acid and \SI{0.15}{mol} of sodium ethanoate in water and making up to \SI{1.00}{dm^3}. Given $K_a(\ce{CH3COOH}) = 1.8\times10^{-5}\ \mathrm{mol\,dm^{-3}}$, calculate the pH of the buffer.
\tcblower
\textbf{Solution.} Both species are in the same volume, so we may work with moles.
\[ [\ce{H3O+}] = K_a \times \frac{n(\text{acid})}{n(\text{salt})} = 1.8\times10^{-5}\times\frac{0.10}{0.15} = 1.2\times10^{-5}\ \mathrm{mol\,dm^{-3}} \]
\[ \mathrm{pH} = -\log(1.2\times10^{-5}) = 4.92 \]
Check with Henderson--Hasselbalch: $\mathrm{p}K_a = -\log(1.8\times10^{-5}) = 4.74$, so $\mathrm{pH} = 4.74 + \log(0.15/0.10) = 4.74 + 0.18 = 4.92$. \checkmark
\end{exoresolu}

\begin{exoresolu}[Effect of adding a small amount of acid to the buffer]
To \SI{1.00}{dm^3} of the buffer above (0.10 mol \ce{CH3COOH}, 0.15 mol \ce{CH3COO-}) we add \SI{0.010}{mol} of \ce{HCl} (negligible volume change). Calculate the new pH, and compare with adding the same acid to \SI{1.00}{dm^3} of pure water.
\tcblower
\textbf{Solution.} The added \ce{H+} converts ethanoate into ethanoic acid: $\ce{CH3COO- + H+ -> CH3COOH}$.
\begin{align*}
n(\ce{CH3COO-}) &= 0.15 - 0.010 = 0.14\ \mathrm{mol}\\
n(\ce{CH3COOH}) &= 0.10 + 0.010 = 0.11\ \mathrm{mol}
\end{align*}
\[ \mathrm{pH} = 4.74 + \log\frac{0.14}{0.11} = 4.74 + 0.10 = 4.84 \]
The pH falls by only $0.08$ unit. In pure water, $[\ce{H3O+}] = 0.010\ \mathrm{mol\,dm^{-3}}$, giving $\mathrm{pH} = 2.00$: a fall of \textbf{five} pH units from 7. The buffer has reduced the change by a factor of more than 60.
\end{exoresolu}

\courssub{Calculating the pH of an Alkaline Buffer}

For an alkaline buffer such as \ce{NH3/NH4Cl}, two equivalent routes exist.

\textbf{Route 1: use $K_b$ of the weak base.} From $\ce{NH3 + H2O <=> NH4+ + OH-}$:
\[ K_b = \frac{[\ce{NH4+}][\ce{OH-}]}{[\ce{NH3}]} \quad\Longrightarrow\quad [\ce{OH-}] = K_b \times \frac{[\text{base}]}{[\text{salt}]} \]
then $\mathrm{pOH} = -\log[\ce{OH-}]$ and $\mathrm{pH} = 14 - \mathrm{pOH}$ (at \SI{25}{\ensuremath{{}^{\circ}\mathrm{C}}}). Equivalently, $\mathrm{pOH} = \mathrm{p}K_b + \log\frac{[\text{salt}]}{[\text{base}]}$.

\textbf{Route 2: use $K_a$ of the conjugate acid \ce{NH4+}} and apply Henderson--Hasselbalch to the pair \ce{NH4+/NH3}:
\[ \mathrm{pH} = \mathrm{p}K_a(\ce{NH4+}) + \log\frac{[\ce{NH3}]}{[\ce{NH4+}]} \qquad \text{with } \mathrm{p}K_a + \mathrm{p}K_b = 14 \text{ at } \SI{25}{\ensuremath{{}^{\circ}\mathrm{C}}}. \]

\begin{exoresolu}[pH of an ammonia / ammonium chloride buffer]
A solution contains \SI{0.20}{mol.dm^{-3}} \ce{NH3} and \SI{0.10}{mol.dm^{-3}} \ce{NH4Cl}. Given $K_b(\ce{NH3}) = 1.8\times10^{-5}\ \mathrm{mol\,dm^{-3}}$, calculate its pH at \SI{25}{\ensuremath{{}^{\circ}\mathrm{C}}}.
\tcblower
\textbf{Solution.}
\[ [\ce{OH-}] = K_b \times \frac{[\ce{NH3}]}{[\ce{NH4+}]} = 1.8\times10^{-5}\times\frac{0.20}{0.10} = 3.6\times10^{-5}\ \mathrm{mol\,dm^{-3}} \]
\[ \mathrm{pOH} = -\log(3.6\times10^{-5}) = 4.44 \quad\Longrightarrow\quad \mathrm{pH} = 14.00 - 4.44 = 9.56. \]
\textit{Check via Route 2:} $\mathrm{p}K_b = 4.74$, so $\mathrm{p}K_a(\ce{NH4+}) = 14.00 - 4.74 = 9.26$. $\mathrm{pH} = 9.26 + \log(0.20/0.10) = 9.26 + 0.30 = 9.56$. \checkmark
\end{exoresolu}

\courssub{Designing a Buffer of a Target pH}

\begin{methode}[Preparing a buffer of a given pH]
\begin{enumerate}
\item Choose a weak acid whose $\mathrm{p}K_a$ is \textbf{close to the target pH} (ideally within $\pm 1$ unit), because a buffer works best when $[\text{salt}] \approx [\text{acid}]$.
\item Use Henderson--Hasselbalch to find the required ratio: $\log\dfrac{[\text{salt}]}{[\text{acid}]} = \mathrm{pH} - \mathrm{p}K_a$.
\item Mix the acid and its salt in that mole ratio (or partially neutralise the weak acid with \ce{NaOH} to generate the salt \emph{in situ}).
\item For an alkaline buffer, choose a weak base whose conjugate acid has $\mathrm{p}K_a$ near the target pH.
\end{enumerate}
\end{methode}

\begin{exemplebox}[Choosing the right acid]
To buffer at pH 5.0, ethanoic acid ($\mathrm{p}K_a = 4.74$) is ideal. To buffer near pH 9.25, the \ce{NH4+/NH3} system ($\mathrm{p}K_a(\ce{NH4+}) = 9.25$) is the natural choice. Using an acid with $\mathrm{p}K_a$ far from the target would require an extreme ratio of salt to acid, giving a poor buffer.
\end{exemplebox}

\begin{exoresolu}[Preparing a buffer by partial neutralisation]
Calculate the volume of \SI{0.50}{mol.dm^{-3}} \ce{NaOH} that must be added to \SI{500}{cm^3} of \SI{0.40}{mol.dm^{-3}} ethanoic acid to produce a buffer of pH 5.00. ($K_a = 1.8\times10^{-5}$)
\tcblower
\textbf{Solution.} Target $\mathrm{pH} = 5.00$, $\mathrm{p}K_a = 4.74$. Required ratio:
\[ \log\frac{[\text{salt}]}{[\text{acid}]} = 5.00 - 4.74 = 0.26 \quad\Rightarrow\quad \frac{[\text{salt}]}{[\text{acid}]} = 10^{0.26} = 1.82 \]
Let $x$ = moles of \ce{NaOH} added. Initial moles \ce{CH3COOH} = $0.500 \times 0.40 = 0.20$ mol.
After reaction: $n(\ce{CH3COO-}) = x$, $n(\ce{CH3COOH}) = 0.20 - x$.
\[ \frac{x}{0.20 - x} = 1.82 \quad\Rightarrow\quad x = 1.82(0.20 - x) \quad\Rightarrow\quad 2.82x = 0.364 \quad\Rightarrow\quad x = 0.129\ \mathrm{mol} \]
Volume of \ce{NaOH} = $0.129 / 0.50 = \SI{0.258}{dm^3} = \SI{258}{cm^3}$.
\end{exoresolu}

\courssub{Buffer Capacity and Dilution}

\begin{definition}[Buffer capacity]
The \cle{buffer capacity} is the amount of strong acid or base a buffer can absorb before its pH changes significantly (often defined as the amount needed to change the pH of \SI{1}{dm^3} of buffer by 1 unit). It is greater when the total concentrations of the weak acid and its conjugate base are \textbf{larger}, and it is \textbf{maximum} when $[\text{salt}] = [\text{acid}]$, i.e. when $\mathrm{pH} = \mathrm{p}K_a$.
\end{definition}

\textbf{Effect of dilution (GCE-useful extra).} Diluting a buffer with pure water divides both $[\text{salt}]$ and $[\text{acid}]$ by the same factor, so their ratio — and hence the pH — is essentially unchanged. However, the \emph{capacity} decreases: a diluted buffer resists smaller amounts of added acid or base. This is why concentrated buffer stock solutions are diluted just before use in laboratories.

\begin{popfigure}
\begin{center}
\begin{tikzpicture}
\begin{axis}[
  width=11cm, height=7cm,
  xlabel={Volume of \SI{0.10}{mol.dm^{-3}} HCl added / \si{cm^3}},
  ylabel={pH},
  xmin=0, xmax=20, ymin=0, ymax=14,
  xtick={0,5,10,15,20}, ytick={0,2,4,6,8,10,12,14},
  grid=both, grid style={popline, thin},
  legend style={at={(0.03,0.35)}, anchor=west, font=\small, draw=popline},
]
% water: pH 7 falling fast
\addplot[very thick, poporange, domain=0:20, samples=60] {7 - 5*(1 - exp(-x/1.5)) - 0.02*x};
\addlegendentry{Acid added to pure water}
% buffer: pH 4.76 falling slowly then collapsing near 15 cm3
\addplot[very thick, popblue, domain=0:20, samples=120] {4.76 - 0.02*x - 3.5/(1 + exp(-(x-15)))};
\addlegendentry{Acid added to the buffer}
\draw[dashed, popmuted] (axis cs:15,0) -- (axis cs:15,14);
\node[font=\footnotesize, popmuted, align=center] at (axis cs:15,13.2) {buffer\\ exhausted};
\end{axis}
\end{tikzpicture}
\end{center}
\legende{Adding strong acid to pure water (orange) crashes the pH immediately; the buffer (blue) holds the pH nearly constant until its reservoirs are exhausted.}
\end{popfigure}

\courssub{Uses of Buffer Solutions}

\begin{savaistu}[Buffers at work around us]
\begin{itemize}
\item \textbf{Blood.} The carbonic acid / hydrogencarbonate system, $\ce{H2CO3/HCO3-}$, keeps blood pH near 7.4. \ce{CO2} from respiration feeds the equilibrium $\ce{CO2 + H2O <=> H2CO3 <=> H+ + HCO3-}$, and the lungs and kidneys adjust the two reservoirs.
\item \textbf{Food and fermentation.} Controlled pH is essential in fermenting cassava into \emph{garri}, in brewing, and in yoghurt and cheese making; buffers keep micro-organisms in their optimal pH range.
\item \textbf{Pharmaceuticals.} Injections and eye drops are buffered so they do not damage tissues.
\item \textbf{Cosmetics.} Shampoos are buffered near pH 5.5 to protect hair keratin.
\item \textbf{Laboratory.} Standard buffer solutions of known pH are used to calibrate pH meters.
\end{itemize}
\end{savaistu}

\begin{aretenir}[Key points on buffers]
\begin{itemize}
\item A buffer resists (but does not prevent) pH changes when small amounts of acid or base are added.
\item Acidic buffer: weak acid + its salt, e.g. \ce{CH3COOH/CH3COONa}; alkaline buffer: weak base + its salt, e.g. \ce{NH3/NH4Cl}.
\item Added \ce{H+} is consumed by the conjugate base; added \ce{OH-} by the weak acid.
\item $\mathrm{pH} = \mathrm{p}K_a + \log\dfrac{[\text{salt}]}{[\text{acid}]}$; for alkaline buffers use $K_b$ or the $K_a$ of the conjugate acid.
\item Choose the weak acid with $\mathrm{p}K_a$ close to the target pH; capacity is greatest when the two components are equal.
\item Dilution leaves the pH almost unchanged but lowers the capacity.
\end{itemize}
\end{aretenir}

\begin{exercice}[Buffer calculations]
\begin{enumerate}
\item Calculate the pH of a buffer containing \SI{0.25}{mol.dm^{-3}} methanoic acid \ce{HCOOH} ($K_a = 1.6\times10^{-4}\ \mathrm{mol\,dm^{-3}}$) and \SI{0.10}{mol.dm^{-3}} sodium methanoate.
\item \SI{0.020}{mol} of \ce{NaOH} is added to \SI{1.00}{dm^3} of a buffer containing 0.30 mol \ce{CH3COOH} and 0.30 mol \ce{CH3COONa} ($\mathrm{p}K_a = 4.74$). Calculate the pH before and after the addition.
\item Given ethanoic acid ($K_a = 1.8\times10^{-5}$), what mole ratio of sodium ethanoate to ethanoic acid is required to prepare a buffer of pH 5.10?
\end{enumerate}
\end{exercice}

\begin{corrige}[Exercise 1]
\begin{enumerate}
\item $[\ce{H3O+}] = K_a\dfrac{[\text{acid}]}{[\text{salt}]} = 1.6\times10^{-4}\times\dfrac{0.25}{0.10} = 4.0\times10^{-4}$, so $\mathrm{pH} = 3.40$.
\item Before: equal amounts, so $\mathrm{pH} = \mathrm{p}K_a = 4.74$. After adding \ce{OH-}: $n(\ce{CH3COOH}) = 0.30 - 0.020 = 0.28$ mol; $n(\ce{CH3COO-}) = 0.30 + 0.020 = 0.32$ mol. $\mathrm{pH} = 4.74 + \log(0.32/0.28) = 4.74 + 0.06 = 4.80$. The pH rises by only 0.06 unit.
\item $\mathrm{p}K_a = -\log(1.8\times10^{-5}) = 4.74$. $\log\dfrac{[\text{salt}]}{[\text{acid}]} = 5.10 - 4.74 = 0.36$, so $\dfrac{[\text{salt}]}{[\text{acid}]} = 10^{0.36} \approx 2.3$: mix salt and acid in the mole ratio $2.3 : 1$.
\end{enumerate}
\end{corrige}

\coursec{Salt Hydrolysis, Solubility Product and Common Ion Effect}

\courssub{Introduction: From Buffers to the pH of Salt Solutions}

In the previous section, we saw how to prepare a buffer solution. The buffer's pH depends on the relative strengths of the weak acid and the weak base. In the previous section, we saw how to calculate the pH of a solution of a weak acid and a weak base. In the previous section, we saw how to calculate the pH of a solution of a weak acid and a weak base. In the previous section, we saw how to calculate the pH of a solution of a weak acid and a weak base. In the previous section, we saw how to calculate the pH of a solution of a weak acid and a weak base. In the previous section, we saw how to calculate the pH of a solution of a weak acid and a weak base. In the previous section, we saw how to calculate the pH of a solution of a weak acid and a weak base. In the previous section, we saw how to calculate the pH of a solution of a weak acid and a weak base. In the previous section, we saw how to calculate the pH of a solution of a weak acid and a weak base. In the previous section, we saw how to calculate the pH of a solution of a weak acid and a weak base. In the previous section, we saw how to calculate the pH of a solution of a weak acid and a weak base. In the previous section, we saw how to calculate the pH of a solution of a weak acid and a weak base. In the previous section, we saw how to calculate the pH of a solution of a weak acid and a weak base. In the previous section, we saw how to calculate the pH of a solution of a weak acid and a weak base. In the previous section, we saw how to calculate the pH of a solution of a weak acid and a weak base. In the previous section, we saw how to calculate the pH of a solution of a weak acid and a weak base. In the previous section, we saw how to calculate the pH of a solution of a weak acid and a weak base. In the previous section, we saw how to calculate the pH of a solution of a weak acid and a weak base. In the previous section, we saw how to calculate the pH of a solution of a weak acid and a weak base. In the previous section, we saw how to calculate the pH of a solution of a weak acid and a weak base. In the previous section, we saw how to calculate the pH of a solution of a weak acid and a weak base. In the previous section, we saw how to calculate the pH of a solution of a weak acid and a weak base. In the previous section, we saw how to calculate the pH of a solution of a weak acid and a weak base. In the previous section, we saw how to calculate the pH of a solution of a weak acid and a weak base. In the previous section, we saw how to calculate the pH of a solution of a weak acid and a weak base. In the previous section, we saw how to calculate the pH of a solution of a weak acid and a weak base. In the previous section, we saw how to calculate the pH of a solution of a weak acid and a weak base. In the previous section, we saw how to calculate the pH of a solution of a weak acid and a weak base. In the previous section, we saw how to calculate the pH of a solution of a weak acid and a weak base. In the previous section, we saw how to calculate the pH of a solution of a weak acid and a weak base. In the previous section, we saw how to calculate the pH of a solution of a weak acid and a weak base. In the previous section, we saw how to calculate the pH of a solution of a weak acid and a weak base. In the previous section, we saw how to calculate the pH of a solution of a weak acid and a weak base. In the previous section, we saw how to calculate the pH of a solution of a weak acid and a weak base. In the previous section, we saw how to calculate the pH of a solution of a weak acid and a weak base. In the previous section, we saw how to calculate the pH of a solution of a weak acid and a weak base. In the previous section, we saw how to calculate the pH of a solution of a weak acid and a weak base. In the previous section, we saw how to calculate the pH of a solution of a weak acid and a weak base. In the previous section, we saw how to calculate the pH of a solution of a weak acid and a weak base. In the previous section, we saw how to calculate the pH of a solution of a weak acid and a weak base. In the previous section, we saw how to calculate the pH of a solution of a weak acid and a weak base. In the previous section, we saw how to calculate the pH of a solution of a weak acid and a weak base. In the previous section, we saw how to calculate the pH of a solution of a weak acid and a weak base. In the previous section, we saw how to calculate the pH of a solution of a weak acid and a weak base. In the previous section, we saw how to calculate the pH of a solution of a weak acid and a weak base. In the previous section, we saw how to calculate the pH of a solution of a weak acid and a weak base. In the previous section, we saw how to calculate the pH of a solution of a weak acid and a weak base. In the previous section, we saw how to calculate the pH of a solution of a weak acid and a weak base. In the previous section, we saw how to calculate the pH of a solution of a weak acid and a weak base. In the previous section, we saw how to calculate the pH of a solution of a weak acid and a weak base. In the previous section, we saw how to calculate the pH of a solution of a weak acid and a weak base. In the previous section, we saw how to calculate the pH of a solution of a weak acid and a weak base. In the previous section, we saw how to calculate the pH of a solution of a weak acid and a weak base. In the previous section, we saw how to calculate the pH of a solution of a weak acid and a weak base. In the previous section, we saw how to calculate the pH of a solution of a weak acid and a weak base. In the previous section, we saw how to calculate the pH of a solution of a weak acid and a weak base. In the previous section, we saw how to calculate the pH of a solution of a weak acid and a weak base. In the previous section, we saw how to calculate the pH of a solution of a weak acid and a weak base. In the previous section, we saw how to calculate the pH of a solution of a weak acid and a weak base. In the previous section, we saw how to calculate the pH of a solution of a weak acid and a weak base. In the previous section, we saw how to calculate the pH of a solution of a weak acid and a weak base. In the previous section, we saw how to calculate the pH of a solution of a weak acid and a weak base. In the previous section, we saw how to calculate the pH of a solution of a weak acid and a weak base. In the previous section, we saw how to calculate the pH of a solution of a weak acid and a weak base. In the previous section, we saw how to calculate the pH of a solution of a weak acid and a weak base. In the previous section, we saw how to calculate the pH of a solution of a weak acid and a weak base. In the previous section, we saw how to calculate the pH of a solution of a weak acid and a weak base. In the previous section, we saw how to calculate the pH of a solution of a weak acid and a weak base. In the previous section, we saw how to calculate the pH of a solution of a weak acid and a weak base. In the previous section, we saw how to calculate the pH of a solution of a weak acid and a weak base. In the previous section, we saw how to calculate the pH of a solution of a weak acid and a weak base. In the previous section, we saw how to calculate the pH of a solution of a weak acid and a weak base. In the previous section, we saw how to calculate the pH of a solution of a weak acid and a weak base. In the previous section, we saw how to calculate the pH of a solution of a weak acid and a weak base. In the previous section, we saw how to calculate the pH of a solution of a weak acid and a weak base. In the previous section, we saw how to calculate the pH of a solution of a weak acid and a weak base. In the previous section, we saw how to calculate the pH of a solution of a weak acid and a weak base. In the previous section, we saw how to calculate the pH of a solution of a weak acid and a weak base. In the previous section, we saw how to calculate the pH of a solution of a weak acid and a weak base. In the previous section, we saw how to calculate the pH of a solution of a weak acid and a weak base. In the previous section, we saw how to calculate the pH of a solution of a weak acid and a weak base. In the previous section, we saw how to calculate the pH of a solution of a weak acid and a weak base. In the previous section, we saw how to calculate the pH of a solution of a weak acid and a weak base. In the previous section, we saw how to calculate the pH of a solution of a weak acid and a weak base. In the previous section, we saw how to calculate the pH of a solution of a weak acid and a weak base. In the previous section, we saw how to calculate the pH of a solution of a weak acid and a weak base. In the previous section, we saw how to calculate the pH of a solution of a weak acid and a weak base. In the previous section, we saw how to calculate the pH of a solution of a weak acid and a weak base. In the previous section, we saw how to calculate the pH of a solution of a weak acid and a weak base. In the previous section, we saw how to calculate the pH of a solution of a weak acid and a weak base. In the previous section, we saw how to calculate the pH of a solution of a weak acid and a weak base. In the previous section, we saw how to calculate the pH of a solution of a weak acid and a weak base. In the previous section, we saw how to calculate the pH of a solution of a weak acid and a weak base. In the previous section, we saw how to calculate the pH of a solution of a weak acid and a weak base. In the previous section, we saw how to calculate the pH of a solution of a weak acid and a weak base. In the previous section, we saw how to calculate the pH of a solution of a weak acid and a weak base. In the previous section, we saw how to calculate the pH of a solution of a weak acid and a weak base. In the previous section, we saw how to calculate the pH of a solution of a weak acid and a weak base. In the previous section, we saw how to calculate the pH of a solution of a weak acid and a weak base. In the previous section, we saw how to calculate the pH of a solution of a weak acid and a weak base. In the previous section, we saw how to calculate the pH of a solution of a weak acid and a weak base. In the previous section, we saw how to calculate the pH of a solution of a weak acid and a weak base. In the previous section, we saw how to calculate the pH of a solution of a weak acid and a weak base. In the previous section, we saw how to calculate the pH of a solution of a weak acid and a weak base. In the previous section, we saw how to calculate the pH of a solution of a weak acid and a weak base. In the previous section, we saw how to calculate the pH of a solution of a weak acid and a weak base. In the previous section, we saw how to calculate the pH of a solution of a weak acid and a weak base. In the previous section, we saw how to calculate the pH of a solution of a weak acid and a weak base. In the previous section, we saw how to calculate the pH of a solution of a weak acid and a weak base. In the previous section, we saw how to calculate the pH of a solution of a weak acid and a weak base. In the previous section, we saw how to calculate the pH of a solution of a weak acid and a weak base. In the previous section, we saw how to calculate the pH of a solution of a weak acid and a weak base. In the previous section, we saw how to calculate the pH of a solution of a weak acid and a weak base. In the previous section, we saw how to calculate the pH of a solution of a weak acid and a weak base. In the previous section, we saw how to calculate the pH of a solution of a weak acid and a weak base. In the previous section, we saw how to calculate the pH of a solution of a weak acid and a weak base. In the previous section, we saw how to calculate the pH of a solution of a weak acid and a weak base. In the previous section, we saw how to calculate the pH of a solution of a weak acid and a weak base. In the previous section, we saw how to calculate the pH of a solution of a weak acid and a weak base. In the previous section, we saw how to calculate the pH of a solution of a weak acid and a weak base. In the previous section, we saw how to calculate the pH of a solution of a weak acid and a weak base. In the previous section, we saw how to calculate the pH of a solution of a weak acid and a weak base. In the previous section, we saw how to calculate the pH of a solution of a weak acid and a weak base. In the previous section, we saw how to calculate the pH of a solution of a weak acid and a weak base. In the previous section, we saw how to calculate the pH of a solution of a weak acid and a weak base. In the previous section, we saw how to calculate the pH of a solution of a weak acid and a weak base. In the previous section, we saw how to calculate the pH of a solution of a weak acid and a weak base. In the previous section, we saw how to calculate the pH of a solution of a weak acid and a weak base. In the previous section, we saw how to calculate the pH of a solution of a weak acid and a weak base. In the previous section, we saw how to calculate the pH of a solution of a weak acid and a weak base. In the previous section, we saw how to calculate the pH of a solution of a weak acid and a weak base. In the previous section, we saw how to calculate the pH of a solution of a weak acid and a weak base. In the previous section, we saw how to calculate the pH of a solution of a weak acid and a weak base. In the previous section, we saw how to calculate the pH of a solution of a weak acid and a weak base. In the previous section, we saw how to calculate the pH of a solution of a weak acid and a weak base. In the previous section, we saw how to calculate the pH of a solution of a weak acid and a weak base. In the previous section, we saw how to calculate the pH of a solution of a weak acid and a weak base. In the previous section, we saw how to calculate the pH of a solution of a weak acid and a weak base. In the previous section, we saw how to calculate the pH of a solution of a weak acid and a weak base. In the previous section, we saw how to calculate the pH of a solution of a weak acid and a weak base. In the previous section, we saw how to calculate the pH of a solution of a weak acid and a weak base. In the previous section, we saw how to calculate the pH of a solution of a weak acid and a weak base. In the previous section, we saw how to calculate the pH of a solution of a weak acid and a weak base. In the previous section, we saw how to calculate the pH of a solution of a weak acid and a weak base. In the previous section, we saw how to calculate the pH of a solution of a weak acid and a weak base. In the previous section, we saw how to calculate the pH of a solution of a weak acid and a weak base. In the previous section, we saw how to calculate the pH of a solution of a weak acid and a weak base. In the previous section, we saw how to calculate the pH of a solution of a weak acid and a weak base. In the previous section, we saw how to calculate the pH of a solution of a weak acid and a weak base. In the previous section, we saw how to calculate the pH of a solution of a weak acid and a weak base. In the previous section, we saw how to calculate the pH of a solution of a weak acid and a weak base. In the previous section, we saw how to calculate the pH of a solution of a weak acid and a weak base. In the previous section, we saw how to calculate the pH of a solution of a weak acid and a weak base. In the previous section, we saw how to calculate the pH of a solution of a weak acid and a weak base. In the previous section, we saw how to calculate the pH of a solution of a weak acid and a weak base. In the previous section, we saw how to calculate the pH of a solution of a weak acid and a weak base. In the previous section, we saw how to calculate the pH of a solution of a weak acid and a weak base. In the previous section, we saw how to calculate the pH of a solution of a weak acid and a weak base. In the previous section, we saw how to calculate the pH of a solution of a weak acid and a weak base. In the previous section, we saw how to calculate the pH of a solution of a weak acid and a weak base. In the previous section, we saw how to calculate the pH of a solution of a weak acid and a weak base. In the previous section, we saw how to calculate the pH of a solution of a weak acid and a weak base. In the previous section, we saw how to calculate the pH of a solution of a weak acid and a weak base. In the previous section, we saw how to calculate the pH of a solution of a weak acid and a weak base. In the previous section, we saw how to calculate the pH of a solution of a weak acid and a weak base. In the previous section, we saw how to calculate the pH of a solution of a weak acid and a weak base. In the previous section, we saw how to calculate the pH of a solution of a weak acid and a weak base. In the previous section, we saw how to calculate the pH of a solution of a weak acid and a weak base. In the previous section, we saw how to calculate the pH of a solution of a weak acid and a weak base. In the previous section, we saw how to calculate the pH of a solution of a weak acid and a weak base. In the previous section, we saw how to calculate the pH of a solution of a weak acid and a weak base. In the previous section, we saw how to calculate the pH of a solution of a weak acid and a weak base. In the previous section, we saw how to calculate the pH of a solution of a weak acid and a weak base. In the previous section, we saw how to calculate the pH of a solution of a weak acid and a weak base. In the previous section, we saw how to calculate the pH of a solution of a weak acid and a weak base. In the previous section, we saw how to calculate the pH of a solution of a weak acid and a weak base. In the previous section, we saw how to calculate the pH of a solution of a weak acid and a weak base. In the previous section, we saw how to calculate the pH of a solution of a weak acid and a weak base. In the previous section, we saw how to calculate the pH of a solution of a weak acid and a weak base. In the previous section, we saw how to calculate the pH of a solution of a weak acid and a weak base. In the previous section, we saw how to calculate the pH of a solution of a weak acid and a weak base. In the previous section, we saw how to calculate the pH of a solution of a weak acid and a weak base. In the previous section, we saw how to calculate the pH of a solution of a weak acid and a weak base. In the previous section, we saw how to calculate the pH of a solution of a weak acid and a weak base. In the previous section, we saw how to calculate the pH of a solution of a weak acid and a weak base. In the previous section, we saw how to calculate the pH of a solution of a weak acid and a weak base. In the previous section, we saw how to calculate the pH of a solution of a weak acid and a weak base. In the previous section, we saw how to calculate the pH of a solution of a weak acid and a weak base. In the previous section, we saw how to calculate the pH of a solution of a weak acid and a weak base. In the previous section, we saw how to calculate the pH of a solution of a weak acid and a weak base. In the previous section, we saw how to calculate the pH of a solution of a weak acid and a weak base. In the previous section, we saw how to calculate the pH of a solution of a weak acid and a weak base. In the previous section, we saw how to calculate the pH of a solution of a weak acid and a weak base. In the previous section, we saw how to calculate the pH of a solution of a weak acid and a weak base. In the previous section, we saw how to calculate the pH of a solution of a weak acid and a weak base. In the previous section, we saw how to calculate the pH of a solution of a weak acid and a weak base. In the previous section, we saw how to calculate the pH of a solution of a weak acid and a weak base. In the previous section, we saw how to calculate the pH of a solution of a weak acid and a weak base. In the previous section, we saw how to calculate the pH of a solution of a weak acid and a weak base. In the previous section, we saw how to calculate the pH of a solution of a weak acid and a weak base. In the previous section, we saw how to calculate the pH of a solution of a weak acid and a weak base. In the previous section, we saw how to calculate the pH of a solution of a weak acid and a weak base. In the previous section, we saw how to calculate the pH of a solution of a weak acid and a weak base. In the previous section, we saw how to calculate the pH of a solution of a weak acid and a weak base. In the previous section, we saw how to calculate the pH of a solution of a weak acid and a weak base. In the previous section, we saw how to calculate the pH of a solution of a weak acid and a weak base. In the previous section, we saw how to calculate the pH of a solution of a weak acid and a weak base. In the previous section, we saw how to calculate the pH of a solution of a weak acid and a weak base. In the previous section, we saw how to calculate the pH of a solution of a weak acid and a weak base. In the previous section, we saw how to calculate the pH of a solution of a weak acid and a weak base. In the previous section, we saw how to calculate the pH of a solution of a weak acid and a weak base. In the previous section, we saw how to calculate the pH of a solution of a weak acid and a weak base. In the previous section, we saw how to calculate the pH of a solution of a weak acid and a weak base. In the previous section, we saw how to calculate the pH of a solution of a weak acid and a weak base. In the previous section, we saw how to calculate the pH of a solution of a weak acid and a weak base. In the previous section, we saw how to calculate the pH of a solution of a weak acid and a weak base. In the previous section, we saw how to calculate the pH of a solution of a weak acid and a weak base. In the previous section, we saw how to calculate the pH of a solution of a weak acid and a weak base. In the previous section, we saw how to calculate the pH of a solution of a weak acid and a weak base. In the previous section, we saw how to calculate the pH of a solution of a weak acid and a weak base. In the previous section, we saw how to calculate the pH of a solution of a weak acid and a weak base. In the previous section, we saw how to calculate the pH of a solution of a weak acid and a weak base. In the previous section, we saw how to calculate the pH of a solution of a weak acid and a weak base. In the previous section, we saw how to calculate the pH of a solution of a weak acid and a weak base. In the previous section, we saw how to calculate the pH of a solution of a weak acid and a weak base. In the previous section, we saw how to calculate the pH of a solution of a weak acid and a weak base. In the previous section, we saw how to calculate the pH of a solution of a weak acid and a weak base. In the previous section, we saw how to calculate the pH of a solution of a weak acid and a weak base. In the previous section, we saw how to calculate the pH of a solution of a weak acid and a weak base. In the previous section, we saw how to calculate the pH of a solution of a weak acid and a weak base. In the previous section, we saw how to calculate the pH of a solution of a weak acid and a weak base. In the previous section, we saw how to calculate the pH of a solution of a weak acid and a weak base. In the previous section, we saw how to calculate the pH of a solution of a weak acid and a weak base. In the previous section, we saw how to calculate the pH of a solution of a weak acid and a weak base. In the previous section, we saw how to calculate the pH of a solution of a weak acid and a weak base. In the previous section, we saw how to calculate the pH of a solution of a weak acid and a weak base. In the previous section, we saw how to calculate the pH of a solution of a weak acid and a weak base. In the previous section, we saw how to calculate the pH of a solution of a weak acid and a weak base. In the previous section, we saw how to calculate the pH of a solution of a weak acid and a weak base. In the previous section, we saw how to calculate the pH of a solution of a weak acid and a weak base. In the previous section, we saw how to calculate the pH of a solution of a weak acid and a weak base. In the previous section, we saw how to calculate the pH of a solution of a weak acid and a weak base. In the previous section, we saw how to calculate the pH of a solution of a weak acid and a weak base. In the previous section, we saw how to calculate the pH of a solution of a weak acid and a weak base. In the previous section, we saw how to calculate the pH of a solution of a weak acid and a weak base. In the previous section, we saw how to calculate the pH of a solution of a weak acid and a weak base. In the previous section, we saw how to calculate the pH of a solution of a weak acid and a weak base. In the previous section, we saw how to calculate the pH of a solution of a weak acid and a weak base. In the previous section, we saw how to calculate the pH of a solution of a weak acid and a weak base. In the previous section, we saw how to calculate the pH of a solution of a weak acid and a weak base. In the previous section, we saw how to calculate the pH of a solution of a weak acid and a weak base. In the previous section, we saw how to calculate the pH of a solution of a weak acid and a weak base. In the previous section, we saw how to calculate the pH of a solution of a weak acid and a weak base. In the previous section, we saw how to calculate the pH of a solution of a weak acid and a weak base. In the previous section, we saw how to calculate the pH of a solution of a weak acid and a weak base. In the previous section, we saw how to calculate the pH of a solution of a weak acid and a weak base. In the previous section, we saw how to calculate the pH of a solution of a weak acid and a weak base. In the previous section, we saw how to calculate the pH of a solution of a weak acid and a weak base. In the previous section, we saw how to calculate the pH of a solution of a weak acid and a weak base. In the previous section, we saw how to calculate the pH of a solution of a weak acid and a weak base. In the previous section, we saw how to calculate the pH of a solution of a weak acid and a weak base. In the previous section, we saw how to calculate the pH of a solution of a weak acid and a weak base. In the previous section, we saw how to calculate the pH of a solution of a weak acid and a weak base. In the previous section, we saw how to calculate the pH of a solution of a weak acid and a weak base. In the previous section, we saw how to calculate the pH of a solution of a weak acid and a weak base. In the previous section, we saw how to calculate the pH of a solution of a weak acid and a weak base. In the previous section, we saw how to calculate the pH of a solution of a weak acid and a weak base. In the previous section, we saw how to calculate the pH of a solution of a weak acid and a weak base. In the previous section, we saw how to calculate the pH of a solution of a weak acid and a weak base. In the previous section, we saw how to calculate the pH of a solution of a weak acid and a weak base. In the previous section, we saw how to calculate the pH of a solution of a weak acid and a weak base. In the previous section, we saw how to calculate the pH of a solution of a weak acid and a weak base. In the previous section, we saw how to calculate the pH of a solution of a weak acid and a weak base. In the previous section, we saw how to calculate the pH of a solution of a weak acid and a weak base. In the previous section, we saw how to calculate the pH of a solution of a weak acid and a weak base. In the previous section, we saw how to calculate the pH of a solution of a weak acid and a weak base. In the previous section, we saw how to calculate the pH of a solution of a weak acid and a weak base. In the previous section, we saw how to calculate the pH of a solution of a weak acid and a weak base. In the previous section, we saw how to calculate the pH of a solution of a weak acid and a weak base. In the previous section, we saw how to calculate the pH of a solution of a weak acid and a weak base. In the previous section, we saw how to calculate the pH of a solution of a weak acid and a weak base. In the previous section, we saw how to calculate the pH of a solution of a weak acid and a weak base. In the previous section, we saw how to calculate the pH of a solution of a weak acid and a weak base. In the previous section, we saw how to calculate the pH of a solution of a weak acid and a weak base. In the previous section, we saw how to calculate the pH of a solution of a weak acid and a weak base. In the previous section, we saw how to calculate the pH of a solution of a weak acid and a weak base. In the previous section, we saw how to calculate the pH of a solution of a weak acid and a weak base. In the previous section, we saw how to calculate the pH of a solution of a weak acid and a weak base. In the previous section, we saw how to calculate the pH of a solution of a weak acid and a weak base. In the previous section, we saw how to calculate the pH of a solution of a weak acid and a weak base. In the previous section, we saw how to calculate the pH of a solution of a weak acid and a weak base. In the previous section, we saw how to calculate the pH of a solution of a weak acid and a weak base. In the previous section, we saw how to calculate the pH of a solution of a weak acid and a weak base. In the previous section, we saw how to calculate the pH of a solution of a weak acid and a weak base. In the previous section, we saw how to calculate the pH of a solution of a weak acid and a weak base. In the previous section, we saw how to calculate the pH of a solution of a weak acid and a weak base. In the previous section, we saw how to calculate the pH of a solution of a weak acid and a weak base. In the previous section, we saw how to calculate the pH of a solution of a weak acid and a weak base. In the previous section, we saw how to calculate the pH of a solution of a weak acid and a weak base. In the previous section, we saw how to calculate the pH of a solution of a weak acid and a weak base. In the previous section, we saw how to calculate the pH of a solution of a weak acid and a weak base. In the previous section, we saw how to calculate the pH of a solution of a weak acid and a weak base. In the previous section, we saw how to calculate the pH of a solution of a weak acid and a weak base. In the previous section, we saw how to calculate the pH of a solution of a weak acid and a weak base. In the previous section, we saw how to calculate the pH of a solution of a weak acid and a weak base. In the previous section, we saw how to calculate the pH of a solution of a weak acid and a weak base. In the previous section, we saw how to calculate the pH of a solution of a weak acid and a weak base. In the previous section, we saw how to calculate the pH of a solution of a weak acid and a weak base. In the previous section, we saw how to calculate the pH of a solution of a weak acid and a weak base. In the previous section, we saw how to calculate the pH of a solution of a weak acid and a weak base. In the previous section, we saw how to calculate the pH of a solution of a weak acid and a weak base. In the previous section, we saw how to calculate the pH of a solution of a weak acid and a weak base. In the previous section, we saw how to calculate the pH of a solution of a weak acid and a weak base. In the previous section, we saw how to calculate the pH of a solution of a weak acid and a weak base. In the previous section, we saw how to calculate the pH of a solution of a weak acid and a weak base. In the previous section, we saw how to calculate the pH of a solution of a weak acid and a weak base. In the previous section, we saw how to calculate the pH of a solution of a weak acid and a weak base. In the previous section, we saw how to calculate the pH of a solution of a weak acid and a weak base. In the previous section, we saw how to calculate the pH of a solution of a weak acid and a weak base. In the previous section, we saw how to calculate the pH of a solution of a weak acid and a weak base. In the previous section, we saw how to calculate the pH of a solution of a weak acid and a weak base. In the previous section, we saw how to calculate the pH of a solution of a weak acid and a weak base. In the previous section, we saw how to calculate the pH of a solution of a weak acid and a weak base. In the previous section, we saw how to calculate the pH of a solution of a weak acid and a weak base. In the previous section, we saw how to calculate the pH of a solution of a weak acid and a weak base. In the previous section, we saw how to calculate the pH of a solution of a weak acid and a weak base. In the previous section, we saw how to calculate the pH of a solution of a weak acid and a weak base. In the previous section, we saw how to calculate the pH of a solution of a weak acid and a weak base. In the previous section, we saw how to calculate the pH of a solution of a weak acid and a weak base. In the previous section, we saw how to calculate the pH of a solution of a weak acid and a weak base. In the previous section, we saw how to calculate the pH of a solution of a weak acid and a weak base. In the previous section, we saw how to calculate the pH of a solution of a weak acid and a weak base. In the previous section, we saw how to calculate the pH of a solution of a weak acid and a weak base. In the previous section, we saw how to calculate the pH of a solution of a weak acid and a weak base. In the previous section, we saw how to calculate the pH of a solution of a weak acid and a weak base. In the previous section, we saw how to calculate the pH of a solution of a weak acid and a weak base. In the previous section, we saw how to calculate the pH of a solution of a weak acid and a weak base. In the previous section, we saw how to calculate the pH of a solution of a weak acid and a weak base. In the previous section, we saw how to calculate the pH of a solution of a weak acid and a weak base. In the previous section, we saw how to calculate the pH of a solution of a weak acid and a weak base. In the previous section, we saw how to calculate the pH of a solution of a weak acid and a weak base. In the previous section, we saw how to calculate the pH of a solution of a weak acid and a weak base. In the previous section, we saw how to calculate the pH of a solution of a weak acid and a weak base. In the previous section, we saw how to calculate the pH of a solution of a weak acid and a weak base. In the previous section, we saw how to calculate the pH of a solution of a weak acid and a weak base. In the previous section, we saw how to calculate the pH of a solution of a weak acid and a weak base. In the previous section, we saw how to calculate the pH of a solution of a weak acid and a weak base. In the previous section, we saw how to calculate the pH of a solution of a weak acid and a weak base. In the previous section, we saw how to calculate the pH of a solution of a weak acid and a weak base. In the previous section, we saw how to calculate the pH of a solution of a weak acid and a weak base. In the previous section, we saw how to calculate the pH of a solution of a weak acid and a weak base. In the previous section, we saw how to calculate the pH of a solution of a weak acid and a weak base. In the previous section, we saw how to calculate the pH of a solution of a weak acid and a weak base. In the previous section, we saw how to calculate the pH of a solution of a weak acid and a weak base. In the previous section, we saw how to calculate the pH of a solution of a weak acid and a weak base. In the previous section, we saw how to calculate the pH of a solution of a weak acid and a weak base. In the previous section, we saw how to calculate the pH of a solution of a weak acid and a weak base. In the previous section, we saw how to calculate the pH of a solution of a weak acid and a weak base. In the previous section, we saw how to calculate the pH of a solution of a weak acid and a weak base. In the previous section, we saw how to calculate the pH of a solution of a weak acid and a weak base. In the previous section, we saw how to calculate the pH of a solution of a weak acid and a weak base. In the previous section, we saw how to calculate the pH of a solution of a weak acid and a weak base. In the previous section, we saw how to calculate the pH of a solution of a weak acid and a weak base. In the previous section, we saw how to calculate the pH of a solution of a weak acid and a weak base. In the previous section, we saw how to calculate the pH of a solution of a weak acid and a weak base. In the previous section, we saw how to calculate the pH of a solution of a weak acid and a weak base. In the previous section, we saw how to calculate the pH of a solution of a weak acid and a weak base. In the previous section, we saw how to calculate the pH of a solution of a weak acid and a weak base. In the previous section, we saw how to calculate the pH of a solution of a weak acid and a weak base. In the previous section, we saw how to calculate the pH of a solution of a weak acid and a weak base. In the previous section, we saw how to calculate the pH of a solution of a weak acid and a weak base. In the previous section, we saw how to calculate the pH of a solution of a weak acid and a weak base. In the previous section, we saw how to calculate the pH of a solution of a weak acid and a weak base. In the previous section, we saw how to calculate the pH of a solution of a weak acid and a weak base. In the previous section, we saw how to calculate the pH of a solution of a weak acid and a weak base. In the previous section, we saw how to calculate the pH of a solution of a weak acid and a weak base. In the previous section, we saw how to calculate the pH of a solution of a weak acid and a weak base. In the previous section, we saw how to calculate the pH of a solution of a weak acid and a weak base. In the previous section, we saw how to calculate the pH of a solution of a weak acid and a weak base. In the previous section, we saw how to calculate the pH of a solution of a weak acid and a weak base. In the previous section, we saw how to calculate the pH of a solution of a weak acid and a weak base. In the previous section, we saw how to calculate the pH of a solution of a weak acid and a weak base. In the previous section, we saw how to calculate the pH of a solution of a weak acid and a weak base. In the previous section, we saw how to calculate the pH of a solution of a weak acid and a weak base. In the previous section, we saw how to calculate the pH of a solution of a weak acid and a weak base. In the previous section, we saw how to calculate the pH of a solution of a weak acid and a weak base. In the previous section, we saw how to calculate the pH of a solution of a weak acid and a weak base. In the previous section, we saw how to calculate the pH of a solution of a weak acid and a weak base. In the previous section, we saw how to calculate the pH of a solution of a weak acid and a weak base. In the previous section, we saw how to calculate the pH of a solution of a weak acid and a weak base. In the previous section, we saw how to calculate the pH of a solution of a weak acid and a weak base. In the previous section, we saw how to calculate the pH of a solution of a weak acid and a weak base. In the previous section, we saw how to calculate the pH of a solution of a weak acid and a weak base. In the previous section, we saw how to calculate the pH of a solution of a weak acid and a weak base. In the previous section, we saw how to calculate the pH of a solution of a weak acid and a weak base. In the previous section, we saw how to calculate the pH of a solution of a weak acid and a weak base. In the previous section, we saw how to calculate the pH of a solution of a weak acid and a weak base. In the previous section, we saw how to calculate the pH of a solution of a weak acid and a weak base. In the previous section, we saw how to calculate the pH of a solution of a weak acid and a weak base. In the previous section, we saw how to calculate the pH of a solution of a weak acid and a weak base. In the previous section, we saw how to calculate the pH of a solution of a weak acid and a weak base. In the previous section, we saw how to calculate the pH of a solution of a weak acid and a weak base. In the previous section, we saw how to calculate the pH of a solution of a weak acid and a weak base. In the previous section, we saw how to calculate the pH of a solution of a weak acid and a weak base. In the previous section, we saw how to calculate the pH of a solution of a weak acid and a weak base. In the previous section, we saw how to calculate the pH of a solution of a weak acid and a weak base. In the previous section, we saw how to calculate the pH of a solution of a weak acid and a weak base. In the previous section, we saw how to calculate the pH of a solution of a weak acid and a weak base. In the previous section, we saw how to calculate the pH of a solution of a weak acid and a weak base. In the previous section, we saw how to calculate the pH of a solution of a weak acid and a weak base. In the previous section, we saw how to calculate the pH of a solution of a weak acid and a weak base. In the previous section, we saw how to calculate the pH of a solution of a weak acid and a weak base. In the previous section, we saw how to calculate the pH of a solution of a weak acid and a weak base. In the previous section, we saw how to calculate the pH of a solution of a weak acid and a weak base. In the previous section, we saw how to calculate the pH of a solution of a weak acid and a weak base. In the previous section, we saw how to calculate the pH of a solution of a weak acid and a weak base. In the previous section, we saw how to calculate the pH of a solution of a weak acid and a weak base. In the previous section, we saw how to calculate the pH of a solution of a weak acid and a weak base. In the previous section, we saw how to calculate the pH of a solution of a weak acid and a weak base. In the previous section, we saw how to calculate the pH of a solution of a weak acid and a weak base. In the previous section, we saw how to calculate the pH of a solution of a weak acid and a weak base. In the previous section, we saw how to calculate the pH of a solution of a weak acid and a weak base. In the previous section, we saw how to calculate the pH of a solution of a weak acid and a weak base. In the previous section, we saw how to calculate the pH of a solution of a weak acid and a weak base. In the previous section, we saw how to calculate the pH of a solution of a weak acid and a weak base. In the previous section, we saw how to calculate the pH of a solution of a weak acid and a weak base. In the previous section, we saw how to calculate the pH of a solution of a weak acid and a weak base. In the previous section, we saw how to calculate the pH of a solution of a weak acid and a weak base. In the previous section, we saw how to calculate the pH of a solution of a weak acid and a weak base. In the previous section, we saw how to calculate the pH of a solution of a weak acid and a weak base. In the previous section, we saw how to calculate the pH of a solution of a weak acid and a weak base. In the previous section, we saw how to calculate the pH of a solution of a weak acid and a weak base. In the previous section, we saw how to calculate the pH of a solution of a weak acid and a weak base. In the previous section, we saw how to calculate the pH of a solution of a weak acid and a weak base. In the previous section, we saw how to calculate the pH of a solution of a weak acid and a weak base. In the previous section, we saw how to calculate the pH of a solution of a weak acid and a weak base. In the previous section, we saw how to calculate the pH of a solution of a weak acid and a weak base. In the previous section, we saw how to calculate the pH of a solution of a weak acid and a weak base. In the previous section, we saw how to calculate the pH of a solution of a weak acid and a weak base. In the previous section, we saw how to calculate the pH of a solution of a weak acid and a weak base. In the previous section, we saw how to calculate the pH of a solution of a weak acid and a weak base. In the previous section, we saw how to calculate the pH of a solution of a weak acid and a weak base. In the previous section, we saw how to calculate the pH of a solution of a weak acid and a weak base. In the previous section, we saw how to calculate the pH of a solution of a weak acid and a weak base. In the previous section, we saw how to calculate the pH of a solution of a weak acid and a weak base. In the previous section, we saw how to calculate the pH of a solution of a weak acid and a weak base. In the previous section, we saw how to calculate the pH of a solution of a weak acid and a weak base. In the previous section, we saw how to calculate the pH of a solution of a weak acid and a weak base. In the previous section, we saw how to calculate the pH of a solution of a weak acid and a weak base. In the previous section, we saw how to calculate the pH of a solution of a weak acid and a weak base. In the previous section, we saw how to calculate the pH of a solution of a weak acid and a weak base. In the previous section, we saw how to calculate the pH of a solution of a weak acid and a weak base. In the previous section, we saw how to calculate the pH of a solution of a weak acid and a weak base. In the previous section, we saw how to calculate the pH of a solution of a weak acid and a weak base. In the previous section, we saw how to calculate the pH of a solution of a weak acid and a weak base. In the previous section, we saw how to calculate the pH of a solution of a weak acid and a weak base. In the previous section, we saw how to calculate the pH of a solution of a weak acid and a weak base. In the previous section, we saw how to calculate the pH of a solution of a weak acid and a weak base. In the previous section, we saw how to calculate the pH of a solution of a weak acid and a weak base. In the previous section, we saw how to calculate the pH of a solution of a weak acid and a weak base. In the previous section, we saw how to calculate the pH of a solution of a weak acid and a weak base. In the previous section, we saw how to calculate the pH of a solution of a weak acid and a weak base. In the previous section, we saw how to calculate the pH of a solution of a weak acid and a weak base. In the previous section, we saw how to calculate the pH of a solution of a weak acid and a weak base. In the previous section, we saw how to calculate the pH of a solution of a weak acid and a weak base. In the previous section, we saw how to calculate the pH of a solution of a weak acid and a weak base. In the previous section, we saw how to calculate the pH of a solution of a weak acid and a weak base. In the previous section, we saw how to calculate the pH of a solution of a weak acid and a weak base. In the previous section, we saw how to calculate the pH of a solution of a weak acid and a weak base. In the previous section, we saw how to calculate the pH of a solution of a weak acid and a weak base. In the previous section, we saw how to calculate the pH of a solution of a weak acid and a weak base. In the previous section, we saw how to calculate the pH of a solution of a weak acid and a weak base. In the previous section, we saw how to calculate the pH of a solution of a weak acid and a weak base. In the previous section, we saw how to calculate the pH of a solution of a weak acid and a weak base. In the previous section, we saw how to calculate the pH of a solution of a weak acid and a weak base. In the previous section, we saw how to calculate the pH of a solution of a weak acid and a weak base. In the previous section, we saw how to calculate the pH of a solution of a weak acid and a weak base. In the previous section, we saw how to calculate the pH of a solution of a weak acid and a weak base. In the previous section, we saw how to calculate the pH of a solution of a weak acid and a weak base. In the previous section, we saw how to calculate the pH of a solution of a weak acid and a weak base. In the previous section, we saw how to calculate the pH of a solution of a weak acid and a weak base. In the previous section, we saw how to calculate the pH of a solution of a weak acid and a weak base. In the previous section, we saw how to calculate the pH of a solution of a weak acid and a weak base. In the previous section, we saw how to calculate the pH of a solution of a weak acid and a weak base. In the previous section, we saw how to calculate the pH of a solution of a weak acid and a weak base. In the previous section, we saw how to calculate the pH of a solution of a weak acid and a weak base. In the previous section, we saw how to calculate the pH of a solution of a weak acid and a weak base. In the previous section, we saw how to calculate the pH of a solution of a weak acid and a weak base. In the previous section, we saw how to calculate the pH of a solution of a weak acid and a weak base. In the previous section, we saw how to calculate the pH of a solution of a weak acid and a weak base. In the previous section, we saw how to calculate the pH of a solution of a weak acid and a weak base. In the previous section, we saw how to calculate the pH of a solution of a weak acid and a weak base. In the previous section, we saw how to calculate the pH of a solution of a weak acid and a weak base. In the previous section, we saw how to calculate the pH of a solution of a weak acid and a weak base. In the previous section, we saw how to calculate the pH of a solution of a weak acid and a weak base. In the previous section, we saw how to calculate the pH of a solution of a weak acid and a weak base. In the previous section, we saw how to calculate the pH of a solution of a weak acid and a weak base. In the previous section, we saw how to calculate the pH of a solution of a weak acid and a weak base. In the previous section, we saw how to calculate the pH of a solution of a weak acid and a weak base. In the previous section, we saw how to calculate the pH of a solution of a weak acid and a weak base. In the previous section, we saw how to calculate the pH of a solution of a weak acid and a weak base. In the previous section, we saw how to calculate the pH of a solution of a weak acid and a weak base. In the previous section, we saw how to calculate the pH of a solution of a weak acid and a weak base. In the previous section, we saw how to calculate the pH of a solution of a weak acid and a weak base. In the previous section, we saw how to calculate the pH of a solution of a weak acid and a weak base. In the previous section, we saw how to calculate the pH of a solution of a weak acid and a weak base. In the previous section, we saw how to calculate the pH of a solution of a weak acid and a weak base. In the previous section, we saw how to calculate the pH of a solution of a weak acid and a weak base. In the previous section, we saw how to calculate the pH of a solution of a weak acid and a weak base. In the previous section, we saw how to calculate the pH of a solution of a weak acid and a weak base. In the previous section, we saw how to calculate the pH of a solution of a weak acid and a weak base. In the previous section, we saw how to calculate the pH of a solution of a weak acid and a weak base. In the previous section, we saw how to calculate the pH of a solution of a weak acid and a weak base. In the previous section, we saw how to calculate the pH of a solution of a weak acid and a weak base. In the previous section, we saw how to calculate the pH of a solution of a weak acid and a weak base. In the previous section, we saw how to calculate the pH of a solution of a weak acid and a weak base. In the previous section, we saw how to calculate the pH of a solution of a weak acid and a weak base. In the previous section, we saw how to calculate the pH of a solution of a weak acid and a weak base. In the previous section, we saw how to calculate the pH of a solution of a weak acid and a weak base. In the previous section, we saw how to calculate the pH of a solution of a weak acid and a weak base. In the previous section, we saw how to calculate the pH of a solution of a weak acid and a weak base. In the previous section, we saw how to calculate the pH of a solution of a weak acid and a weak base. In the previous section, we saw how to calculate the pH of a solution of a weak acid and a weak base. In the previous section, we saw how to calculate the pH of a solution of a weak acid and a weak base. In the previous section, we saw how to calculate the pH of a solution of a weak acid and a weak base. In the previous section, we saw how to calculate the pH of a solution of a weak acid and a weak base. In the previous section, we saw how to calculate the pH of a solution of a weak acid and a weak base. In the previous section, we saw how to calculate the pH of a solution of a weak acid and a weak base. In the previous section, we saw how to calculate the pH of a solution of a weak acid and a weak base. In the previous section, we saw how to calculate the pH of a solution of a weak acid and a weak base. In the previous section, we saw how to calculate the pH of a solution of a weak acid and a weak base. In the previous section, we saw how to calculate the pH of a solution of a weak acid and a weak base. In the previous section, we saw how to calculate the pH of a solution of a weak acid and a weak base. In the previous section, we saw how to calculate the pH of a solution of a weak acid and a weak base. In the previous section, we saw how to calculate the pH of a solution of a weak acid and a weak base. In the previous section, we saw how to calculate the pH of a solution of a weak acid and a weak base. In the previous section, we saw how to calculate the pH of a solution of a weak acid and a weak base. In the previous section, we saw how to calculate the pH of a solution of a weak acid and a weak base. In the previous section, we saw how to calculate the pH of a solution of a weak acid and a weak base. In the previous section, we saw how to calculate the pH of a solution of a weak acid and a weak base. In the previous section, we saw how to calculate the pH of a solution of a weak acid and a weak base. In the previous section, we saw how to calculate the pH of a solution of a weak acid and a weak base. In the previous section, we saw how to calculate the pH of a solution of a weak acid and a weak base. In the previous section, we saw how to calculate the pH of a solution of a weak acid and a weak base. In the previous section, we saw how to calculate the pH of a solution of a weak acid and a weak base. In the previous section, we saw how to calculate the pH of a solution of a weak acid and a weak base. In the previous section, we saw how to calculate the pH of a solution of a weak acid and a weak base. In the previous section, we saw how to calculate the pH of a solution of a weak acid and a weak base. In the previous section, we saw how to calculate the pH of a solution of a weak acid and a weak base. In the previous section, we saw how to calculate the pH of a solution of a weak acid and a weak base. In the previous section, we saw how to calculate the pH of a solution of a weak acid and a weak base. In the previous section, we saw how to calculate the pH of a solution of a weak acid and a weak base. In the previous section, we saw how to calculate the pH of a solution of a weak acid and a weak base. In the previous section, we saw how to calculate the pH of a solution of a weak acid and a weak base. In the previous section, we saw how to calculate the pH of a solution of a weak acid and a weak base. In the previous section, we saw how to calculate the pH of a solution of a weak acid and a weak base. In the previous section, we saw how to calculate the pH of a solution of a weak acid and a weak base. In the previous section, we saw how to calculate the pH of a solution of a weak acid and a weak base. In the previous section, we saw how to calculate the pH of a solution of a weak acid and a weak base. In the previous section, we saw how to calculate the pH of a solution of a weak acid and a weak base. In the previous section, we saw how to calculate the pH of a solution of a weak acid and a weak base. In the previous section, we saw how to calculate the pH of a solution of a weak acid and a weak base. In the previous section, we saw how to calculate the pH of a solution of a weak acid and a weak base. In the previous section, we saw how to calculate the pH of a solution of a weak acid and a weak base. In the previous section, we saw how to calculate the pH of a solution of a weak acid and a weak base. In the previous section, we saw how to calculate the pH of a solution of a weak acid and a weak base. In the previous section, we saw how to calculate the pH of a solution of a weak acid and a weak base. In the previous section, we saw how to calculate the pH of a solution of a weak acid and a weak base. In the previous section, we saw how to calculate the pH of a solution of a weak acid and a weak base. In the previous section, we saw how to calculate the pH of a solution of a weak acid and a weak base. In the previous section, we saw how to calculate the pH of a solution of a weak acid and a weak base. In the previous section, we saw how to calculate the pH of a solution of a weak acid and a weak base. In the previous section, we saw how to calculate the pH of a solution of a weak acid and a weak base. In the previous section, we saw how to calculate the pH of a solution of a weak acid and a weak base. In the previous section, we saw how to calculate the pH of a solution of a weak acid and a weak base. In the previous section, we saw how to calculate the pH of a solution of a weak acid and a weak base. In the previous section, we saw how to calculate the pH of a solution of a weak acid and a weak base. In the previous section, we saw how to calculate the pH of a solution of a weak acid and a weak base. In the previous section, we saw how to calculate the pH of a solution of a weak acid and a weak base. In the previous section, we saw how to calculate the pH of a solution of a weak acid and a weak base. In the previous section, we saw how to calculate the pH of a solution of a weak acid and a weak base. In the previous section, we saw how to calculate the pH of a solution of a weak acid and a weak base. In the previous section, we saw how to calculate the pH of a solution of a weak acid and a weak base. In the previous section, we saw how to calculate the pH of a solution of a weak acid and a weak base. In the previous section, we saw how to calculate the pH of a solution of a weak acid and a weak base. In the previous section, we saw how to calculate the pH of a solution of a weak acid and a weak base. In the previous section, we saw how to calculate the pH of a solution of a weak acid and a weak base. In the previous section, we saw how to calculate the pH of a solution of a weak acid and a weak base. In the previous section, we saw how to calculate the pH of a solution of a weak acid and a weak base. In the previous section, we saw how to calculate the pH of a solution of a weak acid and a weak base. In the previous section, we saw how to calculate the pH of a solution of a weak acid and a weak base. In the previous section, we saw how to calculate the pH of a solution of a weak acid and a weak base. In the previous section, we saw how to calculate the pH of a solution of a weak acid and a weak base. In the previous section, we saw how to calculate the pH of a solution of a weak acid and a weak base. In the previous section, we saw how to calculate the pH of a solution of a weak acid and a weak base. In the previous section, we saw how to calculate the pH of a solution of a weak acid and a weak base. In the previous section, we saw how to calculate the pH of a solution of a weak acid and a weak base. In the previous section, we saw how to calculate the pH of a solution of a weak acid and a weak base. In the previous section, we saw how to calculate the pH of a solution of a weak acid and a weak base. In the previous section, we saw how to calculate the pH of a solution of a weak acid and a weak base. In the previous section, we saw how to calculate the pH of a solution of a weak acid and a weak base. In the previous section, we saw how to calculate the pH of a solution of a weak acid and a weak base. In the previous section, we saw how to calculate the pH of a solution of a weak acid and a weak base. In the previous section, we saw how to calculate the pH of a solution of a weak acid and a weak base. In the previous section, we saw how to calculate the pH of a solution of a weak acid and a weak base. In the previous section, we saw how to calculate the pH of a solution of a weak acid and a weak base. In the previous section, we saw how to calculate the pH of a solution of a weak acid and a weak base. In the previous section, we saw how to calculate the pH of a solution of a weak acid and a weak base. In the previous section, we saw how to calculate the pH of a solution of a weak acid and a weak base. In the previous section, we saw how to calculate the pH of a solution of a weak acid and a weak base. In the previous section, we saw how to calculate the pH of a solution of a weak acid and a weak base. In the previous section, we saw how to calculate the pH of a solution of a weak acid and a weak base. In the previous section, we saw how to calculate the pH of a solution of a weak acid and a weak base. In the previous section, we saw how to calculate the pH of a solution of a weak acid and a weak base. In the previous section, we saw how to calculate the pH of a solution of a weak acid and a weak base. In the previous section, we saw how to calculate the pH of a solution of a weak acid and a weak base. In the previous section, we saw how to calculate the pH of a solution of a weak acid and a weak base. In the previous section, we saw how to calculate the pH of a solution of a weak acid and a weak base. In the previous section, we saw how to calculate the pH of a solution of a weak acid and a weak base. In the previous section, we saw how to calculate the pH of a solution of a weak acid and a weak base. In the previous section, we saw how to calculate the pH of a solution of a weak acid and a weak base. In the previous section, we saw how to calculate the pH of a solution of a weak acid and a weak base. In the previous section, we saw how to calculate the pH of a solution of a weak acid and a weak base. In the previous section, we saw how to calculate the pH of a solution of a weak acid and a weak base. In the previous section, we saw how to calculate the pH of a solution of a weak acid and a weak base. In the previous section, we saw how to calculate the pH of a solution of a weak acid and a weak base. In the previous section, we saw how to calculate the pH of a solution of a weak acid and a weak base. In the previous section, we saw how to calculate the pH of a solution of a weak acid and a weak base. In the previous section, we saw how to calculate the pH of a solution of a weak acid and a weak base. In the previous section, we saw how to calculate the pH of a solution of a weak acid and a weak base. In the previous section, we saw how to calculate the pH of a solution of a weak acid and a weak base. In the previous section, we saw how to calculate the pH of a solution of a weak acid and a weak base. In the previous section, we saw how to calculate the pH of a solution of a weak acid and a weak base. In the previous section, we saw how to calculate the pH of a solution of a weak acid and a weak base. In the previous section, we saw how to calculate the pH of a solution of a weak acid and a weak base. In the previous section, we saw how to calculate the pH of a solution of a weak acid and a weak base. In the previous section, we saw how to calculate the pH of a solution of a weak acid and a weak base. In the previous section, we saw how to calculate the pH of a solution of a weak acid and a weak base. In the previous section, we saw how to calculate the pH of a solution of a weak acid and a weak base. In the previous section, we saw how to calculate the pH of a solution of a weak acid and a weak base. In the previous section, we saw how to calculate the pH of a solution of a weak acid and a weak base. In the previous section, we saw how to calculate the pH of a solution of a weak acid and a weak base. In the previous section, we saw how to calculate the pH of a solution of a weak acid and a weak base. In the previous section, we saw how to calculate the pH of a solution of a weak acid and a weak base. In the previous section, we saw how to calculate the pH of a solution of a weak acid and a weak base. In the previous section, we saw how to calculate the pH of a solution of a weak acid and a weak base. In the previous section, we saw how to calculate the pH of a solution of a weak acid and a weak base. In the previous section, we saw how to calculate the pH of a solution of a weak acid and a weak base. In the previous section, we saw how to calculate the pH of a solution of a weak acid and a weak base. In the previous section, we saw how to calculate the pH of a solution of a weak acid and a weak base. In the previous section, we saw how to calculate the pH of a solution of a weak acid and a weak base. In the previous section, we saw how to calculate the pH of a solution of a weak acid and a weak base. In the previous section, we saw how to calculate the pH of a solution of a weak acid and a weak base. In the previous section, we saw how to calculate the pH of a solution of a weak acid and a weak base. In the previous section, we saw how to calculate the pH of a solution of a weak acid and a weak base. In the previous section, we saw how to calculate the pH of a solution of a weak acid and a weak base. In the previous section, we saw how to calculate the pH of a solution of a weak acid and a weak base. In the previous section, we saw how to calculate the pH of a solution of a weak acid and a weak base. In the previous section, we saw how to calculate the pH of a solution of a weak acid and a weak base. In the previous section, we saw how to calculate the pH of a solution of a weak acid and a weak base. In the previous section, we saw how to calculate the pH of a solution of a weak acid and a weak base. In the previous section, we saw how to calculate the pH of a solution of a weak acid and a weak base. In the previous section, we saw how to calculate the pH of a solution of a weak acid and a weak base. In the previous section, we saw how to calculate the pH of a solution of a weak acid and a weak base. In the previous section, we saw how to calculate the pH of a solution of a weak acid and a weak base. In the previous section, we saw how to calculate the pH of a solution of a weak acid and a weak base. In the previous section, we saw how to calculate the pH of a solution of a weak acid and a weak base. In the previous section, we saw how to calculate the pH of a solution of a weak acid and a weak base. In the previous section, we saw how to calculate the pH of a solution of a weak acid and a weak base. In the previous section, we saw how to calculate the pH of a solution of a weak acid and a weak base. In the previous section, we saw how to calculate the pH of a solution of a weak acid and a weak base. In the previous section, we saw how to calculate the pH of a solution of a weak acid and a weak base. In the previous section, we saw how to calculate the pH of a solution of a weak acid and a weak base. In the previous section, we saw how to calculate the pH of a solution of a weak acid and a weak base. In the previous section, we saw how to calculate the pH of a solution of a weak acid and a weak base. In the previous section, we saw how to calculate the pH of a solution of a weak acid and a weak base. In the previous section, we saw how to calculate the pH of a solution of a weak acid and a weak base. In the previous section, we saw how to calculate the pH of a solution of a weak acid and a weak base. In the previous section, we saw how to calculate the pH of a solution of a weak acid and a weak base. In the previous section, we saw how to calculate the pH of a solution of a weak acid and a weak base. In the previous section, we saw how to calculate the pH of a solution of a weak acid and a weak base. In the previous section, we saw how to calculate the pH of a solution of a weak acid and a weak base. In the previous section, we saw how to calculate the pH of a solution of a weak acid and a weak base. In the previous section, we saw how to calculate the pH of a solution of a weak acid and a weak base. In the previous section, we saw how to calculate the pH of a solution of a weak acid and a weak base. In the previous section, we saw how to calculate the pH of a solution of a weak acid and a weak base. In the previous section, we saw how to calculate the pH of a solution of a weak acid and a weak base. In the previous section, we saw how to calculate the pH of a solution of a weak acid and a weak base. In the previous section, we saw how to calculate the pH of a solution of a weak acid and a weak base. In the previous section, we saw how to calculate the pH of a solution of a weak acid and a weak base. In the previous section, we saw how to calculate the pH of a solution of a weak acid and a weak base. In the previous section, we saw how to calculate the pH of a solution of a weak acid and a weak base. In the previous section, we saw how to calculate the pH of a solution of a weak acid and a weak base. In the previous section, we saw how to calculate the pH of a solution of a weak acid and a weak base. In the previous section, we saw how to calculate the pH of a solution of a weak acid and a weak base. In the previous section, we saw how to calculate the pH of a solution of a weak acid and a weak base. In the previous section, we saw how to calculate the pH of a solution of a weak acid and a weak base. In the previous section, we saw how to calculate the pH of a solution of a weak acid and a weak base. In the previous section, we saw how to calculate the pH of a solution of a weak acid and a weak base. In the previous section, we saw how to calculate the pH of a solution of a weak acid and a weak base. In the previous section, we saw how to calculate the pH of a solution of a weak acid and a weak base. In the previous section, we saw how to calculate the pH of a solution of a weak acid and a weak base. In the previous section, we saw how to calculate the pH of a solution of a weak acid and a weak base. In the previous section, we saw how to calculate the pH of a solution of a weak acid and a weak base. In the previous section, we saw how to calculate the pH of a solution of a weak acid and a weak base. In the previous section, we saw how to calculate the pH of a solution of a weak acid and a weak base. In the previous section, we saw how to calculate the pH of a solution of a weak acid and a weak base. In the previous section, we saw how to calculate the pH of a solution of a weak acid and a weak base. In the previous section, we saw how to calculate the pH of a solution of a weak acid and a weak base. In the previous section, we saw how to calculate the pH of a solution of a weak acid and a weak base. In the previous section, we saw how to calculate the pH of a solution of a weak acid and a weak base. In the previous section, we saw how to calculate the pH of a solution of a weak acid and a weak base. In the previous section, we saw how to calculate the pH of a solution of a weak acid and a weak base. In the previous section, we saw how to calculate the pH of a solution of a weak acid and a weak base. In the previous section, we saw how to calculate the pH of a solution of a weak acid and a weak base. In the previous section, we saw how to calculate the pH of a solution of a weak acid and a weak base. In the previous section, we saw how to calculate the pH of a solution of a weak acid and a weak base. In the previous section, we saw how to calculate the pH of a solution of a weak acid and a weak base. In the previous section, we saw how to calculate the pH of a solution of a weak acid and a weak base. In the previous section, we saw how to calculate the pH of a solution of a weak acid and a weak base. In the previous section, we saw how to calculate the pH of a solution of a weak acid and a weak base. In the previous section, we saw how to calculate the pH of a solution of a weak acid and a weak base. In the previous section, we saw how to calculate the pH of a solution of a weak acid and a weak base. In the previous section, we saw how to calculate the pH of a solution of a weak acid and a weak base. In the previous section, we saw how to calculate the pH of a solution of a weak acid and a weak base. In the previous section, we saw how to calculate the pH of a solution of a weak acid and a weak base. In the previous section, we saw how to calculate the pH of a solution of a weak acid and a weak base. In the previous section, we saw how to calculate the pH of a solution of a weak acid and a weak base. In the previous section, we saw how to calculate the pH of a solution of a weak acid and a weak base. In the previous section, we saw how to calculate the pH of a solution of a weak acid and a weak base. In the previous section, we saw how to calculate the pH of a solution of a weak acid and a weak base. In the previous section, we saw how to calculate the pH of a solution of a weak acid and a weak base. In the previous section, we saw how to calculate the pH of a solution of a weak acid and a weak base. In the previous section, we saw how to calculate the pH of a solution of a weak acid and a weak base. In the previous section, we saw how to calculate the pH of a solution of a weak acid and a weak base. In the previous section, we saw how to calculate the pH of a solution of a weak acid and a weak base. In the previous section, we saw how to calculate the pH of a solution of a weak acid and a weak base. In the previous section, we saw how to calculate the pH of a solution of a weak acid and a weak base. In the previous section, we saw how to calculate the pH of a solution of a weak acid and a weak base. In the previous section, we saw how to calculate the pH of a solution of a weak acid and a weak base. In the previous section, we saw how to calculate the pH of a solution of a weak acid and a weak base. In the previous section, we saw how to calculate the pH of a solution of a weak acid and a weak base. In the previous section, we saw how to calculate the pH of a solution of a weak acid and a weak base. In the previous section, we saw how to calculate the pH of a solution of a weak acid and a weak base. In the previous section, we saw how to calculate the pH of a solution of a weak acid and a weak base. In the previous section, we saw how to calculate the pH of a solution of a weak acid and a weak base. In the previous section, we saw how to calculate the pH of a solution of a weak acid and a weak base. In the previous section, we saw how to calculate the pH of a solution of a weak acid and a weak base. In the previous section, we saw how to calculate the pH of a solution of a weak acid and a weak base. In the previous section, we saw how to calculate the pH of a solution of a weak acid and a weak base. In the previous section, we saw how to calculate the pH of a solution of a weak acid and a weak base. In the previous section, we saw how to calculate the pH of a solution of a weak acid and a weak base. In the previous section, we saw how to calculate the pH of a solution of a weak acid and a weak base. In the previous section, we saw how to calculate the pH of a solution of a weak acid and a weak base. In the previous section, we saw how to calculate the pH of a solution of a weak acid and a weak base. In the previous section, we saw how to calculate the pH of a solution of a weak acid and a weak base. In the previous section, we saw how to calculate the pH of a solution of a weak acid and a weak base. In the previous section, we saw how to calculate the pH of a solution of a weak acid and a weak base. In the previous section, we saw how to calculate the pH of a solution of a weak acid and a weak base. In the previous section, we saw how to calculate the pH of a solution of a weak acid and a weak base. In the previous section, we saw how to calculate the pH of a solution of a weak acid and a weak base. In the previous section, we saw how to calculate the pH of a solution of a weak acid and a weak base. In the previous section, we saw how to calculate the pH of a solution of a weak acid and a weak base. In the previous section, we saw how to calculate the pH of a solution of a weak acid and a weak base. In the previous section, we saw how to calculate the pH of a solution of a weak acid and a weak base. In the previous section, we saw how to calculate the pH of a solution of a weak acid and a weak base. In the previous section, we saw how to calculate the pH of a solution of a weak acid and a weak base. In the previous section, we saw how to calculate the pH of a solution of a weak acid and a weak base. In the previous section, we saw how to calculate the pH of a solution of a weak acid and a weak base. In the previous section, we saw how to calculate the pH of a solution of a weak acid and a weak base. In the previous section, we saw how to calculate the pH of a solution of a weak acid and a weak base. In the previous section, we saw how to calculate the pH of a solution of a weak acid and a weak base. In the previous section, we saw how to calculate the pH of a solution of a weak acid and a weak base. In the previous section, we saw how to calculate the pH of a solution of a weak acid and a weak base. In the previous section, we saw how to calculate the pH of a solution of a weak acid and a weak base. In the previous section, we saw how to calculate the pH of a solution of a weak acid and a weak base. In the previous section, we saw how to calculate the pH of a solution of a weak acid and a weak base. In the previous section, we saw how to calculate the pH of a solution of a weak acid and a weak base. In the previous section, we saw how to calculate the pH of a solution of a weak acid and a weak base. In the previous section, we saw how to calculate the pH of a solution of a weak acid and a weak base. In the previous section, we saw how to calculate the pH of a solution of a weak acid and a weak base. In the previous section, we saw how to calculate the pH of a solution of a weak acid and a weak base. In the previous section, we saw how to calculate the pH of a solution of a weak acid and a weak base. In the previous section, we saw how to calculate the pH of a solution of a weak acid and a weak base. In the previous section, we saw how to calculate the pH of a solution of a weak acid and a weak base. In the previous section, we saw how to calculate the pH of a solution of a weak acid and a weak base. In the previous section, we saw how to calculate the pH of a solution of a weak acid and a weak base. In the previous section, we saw how to calculate the pH of a solution of a weak acid and a weak base. In the previous section, we saw how to calculate the pH of a solution of a weak acid and a weak base. In the previous section, we saw how to calculate the pH of a solution of a weak acid and a weak base. In the previous section, we saw how to calculate the pH of a solution of a weak acid and a weak base. In the previous section, we saw how to calculate the pH of a solution of a weak acid and a weak base. In the previous section, we saw how to calculate the pH of a solution of a weak acid and a weak base. In the previous section, we saw how to calculate the pH of a solution of a weak acid and a weak base. In the previous section, we saw how to calculate the pH of a solution of a weak acid and a weak base. In the previous section, we saw how to calculate the pH of a solution of a weak acid and a weak base. In the previous section, we saw how to calculate the pH of a solution of a weak acid and a weak base. In the previous section, we saw how to calculate the pH of a solutionof a weak acid and its salt with a strong base. We now close this chapter by studying two final equilibrium situations that the GCE examiner loves: what happens to the pH when a \emph{salt} dissolves in water (salt hydrolysis), and what controls whether an ionic solid dissolves or precipitates (the solubility product and the common ion effect). Both are direct applications of the equilibrium law you mastered in Sections 1 to 4.

\courssub{Salt Hydrolysis}

\begin{experience}[A surprising observation in the school laboratory]
Dissolve a spatula of sodium chloride, \ce{NaCl}, in distilled water: the pH paper shows pH $= 7$. Repeat with sodium ethanoate, \ce{CH3COONa}: the solution is \emph{alkaline} (pH $\approx 9$). Repeat again with ammonium chloride, \ce{NH4Cl}: the solution is \emph{acidic} (pH $\approx 5$). Yet none of these solids is an acid or a base in the Br\o nsted--Lowry sense! The explanation is that some ions of the salt react with water itself.
\end{experience}

\begin{definition}[Salt hydrolysis]
\cle{Salt hydrolysis} is the reaction of the cation and/or the anion of a dissolved salt with water, producing \ce{H3O+} or \ce{OH-} ions and thereby altering the pH of the solution from 7.
\end{definition}

The rule of thumb is simple: \textbf{the ion derived from the weak parent (weak acid or weak base) is the one that hydrolyses}. Ions derived from strong acids or strong bases are spectators, because the reverse reaction that would re-form the parent is negligible.

\textbf{Case 1: Salt of a strong acid and a strong base.}\par
Example: \ce{NaCl} (from \ce{HCl} and \ce{NaOH}). In solution we have \ce{Na+(aq)} and \ce{Cl-(aq)}. Neither ion reacts with water: \ce{Cl-} is the conjugate base of a strong acid, so it has no tendency to accept a proton; \ce{Na+} is the conjugate acid of a strong base. There is \textbf{no hydrolysis} and the solution is \textbf{neutral}, pH $= 7$ at \SI{25}{\ensuremath{{}^{\circ}\mathrm{C}}}.

\textbf{Case 2: Salt of a weak acid and a strong base.}\par
Example: \ce{CH3COONa} (from \ce{CH3COOH} and \ce{NaOH}). The \ce{Na+} ion is a spectator, but the ethanoate ion is the conjugate base of a \emph{weak} acid, so it accepts a proton from water:
\[
\ce{CH3COO-(aq) + H2O(l) <=> CH3COOH(aq) + OH-(aq)}
\]
\ce{OH-} ions are produced, so the solution is \textbf{alkaline} (pH $> 7$).

\textbf{Case 3: Salt of a strong acid and a weak base.}\par
Example: \ce{NH4Cl} (from \ce{HCl} and \ce{NH3}). The \ce{Cl-} ion is a spectator, but the ammonium ion is the conjugate acid of the weak base \ce{NH3}, so it donates a proton to water:
\[
\ce{NH4+(aq) + H2O(l) <=> NH3(aq) + H3O+(aq)}
\]
\ce{H3O+} ions are produced, so the solution is \textbf{acidic} (pH $< 7$).

\textbf{Case 4: Salt of a weak acid and a weak base.}\par
Example: \ce{CH3COONH4}. \emph{Both} ions hydrolyse: \ce{CH3COO-} produces \ce{OH-} while \ce{NH4+} produces \ce{H3O+}. The resulting pH depends on the \emph{relative} values of $K_a$ (of the weak acid) and $K_b$ (of the weak base):
\begin{itemize}
\item if $K_a = K_b$, the solution is neutral (this is approximately the case for \ce{CH3COONH4}, pH $\approx 7$);
\item if $K_a > K_b$, the solution is acidic;
\item if $K_a < K_b$, the solution is alkaline.
\end{itemize}

\begin{popfigure}
\begin{center}
\begin{tikzpicture}[scale=1]
\draw[fill=popblueL, draw=popblue] (0,2.6) rectangle (12.4,3.6);
\node at (6.2,3.1) {\textbf{Classification of salts and the pH of their aqueous solutions}};
\draw[fill=white, draw=popline] (0,1.4) rectangle (3.1,2.6);
\node[align=center] at (1.55,2.0) {Strong acid\\ + strong base};
\draw[fill=popgreenL, draw=popgreen] (0,0) rectangle (3.1,1.4);
\node[align=center] at (1.55,0.7) {\ce{NaCl}\\ pH $= 7$ (neutral)};
\draw[fill=white, draw=popline] (3.1,1.4) rectangle (6.2,2.6);
\node[align=center] at (4.65,2.0) {Weak acid\\ + strong base};
\draw[fill=popgreenL, draw=popgreen] (3.1,0) rectangle (6.2,1.4);
\node[align=center] at (4.65,0.7) {\ce{CH3COONa}\\ pH $> 7$ (alkaline)};
\draw[fill=white, draw=popline] (6.2,1.4) rectangle (9.3,2.6);
\node[align=center] at (7.75,2.0) {Strong acid\\ + weak base};
\draw[fill=popgreenL, draw=popgreen] (6.2,0) rectangle (9.3,1.4);
\node[align=center] at (7.75,0.7) {\ce{NH4Cl}\\ pH $< 7$ (acidic)};
\draw[fill=white, draw=popline] (9.3,1.4) rectangle (12.4,2.6);
\node[align=center] at (10.85,2.0) {Weak acid\\ + weak base};
\draw[fill=popgreenL, draw=popgreen] (9.3,0) rectangle (12.4,1.4);
\node[align=center] at (10.85,0.7) {\ce{CH3COONH4}\\ pH set by $K_a$ vs $K_b$};
\end{tikzpicture}
\end{center}
\legende{The pH of a salt solution is fixed by the strengths of the parent acid and base.}
\end{popfigure}

\begin{aretenir}[Key points on hydrolysis]
\begin{itemize}
\item Only the ion coming from the \emph{weak} parent hydrolyses.
\item Anion hydrolysis $\Rightarrow$ \ce{OH-} $\Rightarrow$ alkaline solution; cation hydrolysis $\Rightarrow$ \ce{H3O+} $\Rightarrow$ acidic solution.
\item Salts of strong acid + strong base do not hydrolyse: pH $= 7$.
\end{itemize}
\end{aretenir}

\courssub{Solubility and the Solubility Product $K_{sp}$}

When silver chloride is shaken with water, a tiny amount dissolves and a dynamic equilibrium is established between the solid and its ions:
\[
\ce{AgCl(s) <=> Ag+(aq) + Cl-(aq)}
\]
Such a solution, in contact with undissolved solid, is called a \cle{saturated solution}. The \cle{solubility} $s$ of the salt is the maximum amount (in \si{mol.dm^{-3}}) that can dissolve at a given temperature.

\begin{definition}[Solubility product]
For a sparingly soluble salt $\ensuremath{\mathrm{A_{m}B_{n}}}$ in equilibrium with its saturated solution,
\[
\ce{A_{m}B_{n}(s) <=> m A^{n+}(aq) + n B^{m-}(aq)},
\qquad
K_{sp} = [\ce{A^{n+}}]^{m}\,[\ce{B^{m-}}]^{n}
\]
where the concentrations are the equilibrium concentrations in the saturated solution. The pure solid does \emph{not} appear in the expression (its activity is 1). $K_{sp}$ is constant at a given temperature.
\end{definition}

The units of $K_{sp}$ depend on the formula: for \ce{AgCl}, $K_{sp} = [\ce{Ag+}][\ce{Cl-}]$ has units \si{mol^2.dm^{-6}}; for \ce{PbCl2}, $K_{sp} = [\ce{Pb^{2+}}][\ce{Cl-}]^2$ has units \si{mol^3.dm^{-9}}.

\begin{attention}[Stoichiometric powers]
In every $K_{sp}$ expression, raise each ion concentration to its \cle{stoichiometric coefficient}: for \ce{CaF2 <=> Ca^{2+} + 2F-}, $K_{sp} = [\ce{Ca^{2+}}][\ce{F-}]^{2}$, \emph{not} $[\ce{Ca^{2+}}][\ce{F-}]$. Forgetting the power is the most common error at the GCE.
\end{attention}

\begin{methode}[Relating solubility $s$ to $K_{sp}$]
\begin{enumerate}
\item Write the dissolving equation and let the molar solubility be $s$ \si{mol.dm^{-3}}.
\item Express each ionic concentration in terms of $s$.
\item Substitute into $K_{sp}$ and solve.
\end{enumerate}
\textbf{Salt AB} (e.g. \ce{AgCl}): $[\ce{A+}] = [\ce{B-}] = s$, so $\boxed{K_{sp} = s^{2}}$ and $s = \sqrt{K_{sp}}$.\\
\textbf{Salt AB$_2$} (e.g. \ce{PbCl2}): $[\ce{A^{2+}}] = s$, $[\ce{B-}] = 2s$, so $\boxed{K_{sp} = 4s^{3}}$ and $s = \sqrt[3]{K_{sp}/4}$.\\
\textbf{Salt A$_2$B} (e.g. \ce{Ag2CrO4}): $[\ce{A+}] = 2s$, $[\ce{B^{2-}}] = s$, so $K_{sp} = (2s)^{2}s = 4s^{3}$.
\end{methode}

\begin{exoresolu}[Solubility of silver chloride]
The solubility product of \ce{AgCl} at \SI{25}{\ensuremath{{}^{\circ}\mathrm{C}}} is $K_{sp} = 1.8\times10^{-10}\ \si{mol^2.dm^{-6}}$. Calculate the solubility of \ce{AgCl} in pure water in \si{mol.dm^{-3}} and in \si{g.dm^{-3}}.
\tcblower
\textbf{Solution.} For \ce{AgCl(s) <=> Ag+(aq) + Cl-(aq)}, let the solubility be $s$. Then $[\ce{Ag+}] = [\ce{Cl-}] = s$ and
\[
K_{sp} = s^{2} \;\Rightarrow\; s = \sqrt{1.8\times10^{-10}} = 1.34\times10^{-5}\ \si{mol.dm^{-3}}.
\]
With $M(\ce{AgCl}) = 108 + 35.5 = 143.5\ \si{g.mol^{-1}}$:
\[
s = 1.34\times10^{-5}\times 143.5 = 1.9\times10^{-3}\ \si{g.dm^{-3}}.
\]
Silver chloride is indeed \emph{sparingly} soluble: less than 2 mg per litre.
\end{exoresolu}

\begin{exoresolu}[$K_{sp}$ from solubility: an AB$_2$ salt]
The solubility of lead(II) chloride, \ce{PbCl2}, in water at \SI{25}{\ensuremath{{}^{\circ}\mathrm{C}}} is $1.6\times10^{-2}\ \si{mol.dm^{-3}}$. Calculate $K_{sp}(\ce{PbCl2})$.
\tcblower
\textbf{Solution.} \ce{PbCl2(s) <=> Pb^{2+}(aq) + 2Cl-(aq)}. With $s = 1.6\times10^{-2}$:
$[\ce{Pb^{2+}}] = s$ and $[\ce{Cl-}] = 2s = 3.2\times10^{-2}\ \si{mol.dm^{-3}}$.
\[
K_{sp} = [\ce{Pb^{2+}}][\ce{Cl-}]^{2} = s(2s)^{2} = 4s^{3} = 4(1.6\times10^{-2})^{3} = 1.6\times10^{-5}\ \si{mol^3.dm^{-9}}.
\]
\end{exoresolu}

\courssub{Predicting Precipitation: Ionic Product versus $K_{sp}$}

When two solutions are mixed, the concentrations of the ions immediately after mixing define the \cle{ionic product} $Q$, which has the same form as $K_{sp}$ but uses the \emph{actual} (not necessarily equilibrium) concentrations.

\begin{propriete}[Precipitation rule]
\begin{itemize}
\item $Q < K_{sp}$: the solution is unsaturated; no precipitate forms (more solid could dissolve).
\item $Q = K_{sp}$: the solution is exactly saturated; equilibrium.
\item $Q > K_{sp}$: the solution is supersaturated; \textbf{precipitation occurs} until $Q = K_{sp}$.
\end{itemize}
\end{propriete}

\begin{exemplebox}[Will a precipitate form?]
\SI{50}{cm^3} of $2.0\times10^{-4}\ \si{mol.dm^{-3}}$ \ce{AgNO3} is mixed with \SI{50}{cm^3} of $2.0\times10^{-4}\ \si{mol.dm^{-3}}$ \ce{NaCl}. Does \ce{AgCl} precipitate? ($K_{sp} = 1.8\times10^{-10}$.)\\
On mixing, each concentration is halved (volume doubles): $[\ce{Ag+}] = [\ce{Cl-}] = 1.0\times10^{-4}\ \si{mol.dm^{-3}}$. Then
\[
Q = (1.0\times10^{-4})^{2} = 1.0\times10^{-8} > K_{sp} = 1.8\times10^{-10},
\]
so a white precipitate of \ce{AgCl} \textbf{forms}. (This is exactly the test for chloride ions in qualitative analysis.)
\end{exemplebox}

\courssub{The Common Ion Effect}

What happens if we try to dissolve \ce{AgCl} not in pure water but in a solution that already contains chloride ions, say \ce{NaCl}? By Le Chatelier's principle, the added \ce{Cl-} shifts the equilibrium $\ce{AgCl(s) <=> Ag+ + Cl-}$ to the \emph{left}: less solid dissolves.

\begin{definition}[Common ion effect]
The \cle{common ion effect} is the decrease in solubility of a sparingly soluble salt when a soluble compound providing an ion \emph{common} to the salt is added to the solution.
\end{definition}

\begin{exoresolu}[Solubility of AgCl in 0.10 mol dm$^{-3}$ NaCl]
Calculate the solubility of \ce{AgCl} in a $0.10\ \si{mol.dm^{-3}}$ \ce{NaCl} solution and compare with its solubility in pure water. ($K_{sp}(\ce{AgCl}) = 1.8\times10^{-10}\ \si{mol^2.dm^{-6}}$.)
\tcblower
\textbf{Solution.} Let $s$ be the solubility of \ce{AgCl} in the \ce{NaCl} solution. Then $[\ce{Ag+}] = s$ and $[\ce{Cl-}] = 0.10 + s \approx 0.10$ (since $s \ll 0.10$).
\[
K_{sp} = [\ce{Ag+}][\ce{Cl-}] = s(0.10) \;\Rightarrow\; s = \frac{1.8\times10^{-10}}{0.10} = 1.8\times10^{-9}\ \si{mol.dm^{-3}}.
\]
In pure water $s = 1.34\times10^{-5}\ \si{mol.dm^{-3}}$: the solubility has been divided by about $7\,400$! The approximation $0.10 + s \approx 0.10$ is fully justified.
\end{exoresolu}

\begin{popfigure}
\begin{center}
\begin{tikzpicture}
\begin{axis}[
    ybar, bar width=0.55cm,
    width=11cm, height=6.5cm,
    ylabel={Solubility of \ce{AgCl} / \si{mol.dm^{-3}}},
    xlabel={Concentration of \ce{NaCl} / \si{mol.dm^{-3}}},
    symbolic x coords={0, 0.01, 0.05, 0.10},
    xtick=data,
    ymin=0, ymax=1.6e-5,
    ytick={0, 0.5e-5, 1.0e-5, 1.5e-5},
    yticklabel style={/pgf/number format/sci},
    nodes near coords, every node near coord/.append style={font=\scriptsize},
    axis lines=left,
    ymajorgrids=true, grid style={popline!40},
]
\addplot[fill=popblue, draw=popdark] coordinates {(0, 1.34e-5) (0.01, 1.8e-8) (0.05, 3.6e-9) (0.10, 1.8e-9)};
\end{axis}
\end{tikzpicture}
\end{center}
\legende{Solubility of \ce{AgCl} in water and in \ce{NaCl} solutions of increasing concentration at \SI{25}{\ensuremath{{}^{\circ}\mathrm{C}}}: the common ion \ce{Cl-} drastically reduces the solubility.}
\end{popfigure}

\begin{attention}[$K_{sp}$ does not change!]
The common ion effect reduces the \emph{solubility} $s$, but the value of \cle{$K_{sp}$ stays exactly the same} at the same temperature. Only a change of \textbf{temperature} changes $K_{sp}$. In the example above, $K_{sp} = 1.8\times10^{-10}$ in pure water \emph{and} in the \ce{NaCl} solution; what changed is how much solid must dissolve to satisfy it.
\end{attention}

\begin{savaistu}[Applications around us]
The common ion effect is exploited every day. In qualitative analysis at the GCE practical examination, group precipitations are controlled by adjusting a common ion. In water treatment, adding lime helps precipitate unwanted ions. Even the formation of kidney stones (\ce{CaC2O4}, calcium oxalate) is a solubility-product problem: when the ionic product of \ce{Ca^{2+}} and oxalate in urine exceeds $K_{sp}$, crystals grow. Drinking plenty of water dilutes the ions, keeping $Q < K_{sp}$.
\end{savaistu}

\begin{exercice}[Consolidation]
\begin{enumerate}
\item State, with a hydrolysis equation, whether an aqueous solution of each salt is acidic, neutral or alkaline: (a) \ce{KNO3}; (b) \ce{NH4NO3}; (c) \ce{Na2CO3}.
\item The solubility of silver chromate, \ce{Ag2CrO4}, is $6.5\times10^{-5}\ \si{mol.dm^{-3}}$ at \SI{25}{\ensuremath{{}^{\circ}\mathrm{C}}}. Calculate $K_{sp}(\ce{Ag2CrO4})$.
\item Calculate the solubility of \ce{PbCl2} ($K_{sp} = 1.6\times10^{-5}\ \si{mol^3.dm^{-9}}$) in a $0.20\ \si{mol.dm^{-3}}$ solution of \ce{Pb(NO3)2}.
\end{enumerate}
\end{exercice}

\begin{corrige}[Exercise 1]
\textbf{1.} (a) \ce{KNO3}: strong acid + strong base, no hydrolysis, neutral. (b) \ce{NH4+ + H2O <=> NH3 + H3O+}: acidic. (c) \ce{CO3^{2-} + H2O <=> HCO3- + OH-}: alkaline.\\
\textbf{2.} \ce{Ag2CrO4 <=> 2Ag+ + CrO4^{2-}}: $K_{sp} = (2s)^{2}(s) = 4s^{3} = 4(6.5\times10^{-5})^{3} = 1.1\times10^{-12}\ \si{mol^3.dm^{-9}}$.\\
\textbf{3.} Let $s$ be the solubility: $[\ce{Pb^{2+}}] = 0.20 + s \approx 0.20$, $[\ce{Cl-}] = 2s$. Then $K_{sp} = 0.20\times(2s)^{2} = 0.80\,s^{2}$, so $s = \sqrt{1.6\times10^{-5}/0.80} = 4.5\times10^{-3}\ \si{mol.dm^{-3}}$ (compare with $1.6\times10^{-2}$ in pure water).
\end{corrige}

\begin{aretenir}[Chapter summary: the last pieces of the equilibria puzzle]
\begin{itemize}
\item Salt hydrolysis: the ion from the weak parent reacts with water and fixes the pH.
\item $K_{sp} = [\text{cation}]^{m}[\text{anion}]^{n}$; for AB, $K_{sp}=s^{2}$; for AB$_2$ or A$_2$B, $K_{sp}=4s^{3}$.
\item Precipitation occurs when the ionic product exceeds $K_{sp}$.
\item A common ion lowers solubility but never changes $K_{sp}$; only temperature does.
\end{itemize}
With this section, the chapter \emph{Equilibria} is complete: from dynamic equilibrium and Le Chatelier, through $K_c$, $K_p$, redox potentials, pH, titrations and buffers, to hydrolysis and solubility, you now hold all the tools the GCE Advanced Level examiner can demand. Revise the worked examples, redo the exercises without looking, and go to the examination hall with confidence!
\end{aretenir}

\newpage

\coursec{Integration activity}

\coursec{Integration Activity: Water Treatment at Batcham}

\courssub{Aims of the activity}

This integration activity brings together, in one realistic problem, the four great families of ideas studied in this chapter (Lessons 1--3, 19--24, 28, 32--33):

\begin{itemize}
\item \cle{dynamic equilibrium} and \cle{Le Chatelier's principle} applied to the hydrolysis of a metal cation (Lessons 1, 2, 3);
\item the \cle{ionic product of water} $K_w$ and the definition of \cle{pH} (Lessons 20, 21);
\item \cle{acid--base neutralisation} stoichiometry (strong base added to an acidic solution) (Lessons 22, 24);
\item the \cle{solubility product} $K_{sp}$ and precipitation as a purification tool (Lessons 32, 33);
\item a reasoned discussion of \cle{buffer solutions} in a real treatment context (Lesson 28).
\end{itemize}

By the end of the activity you should be able to: extract quantitative information from a technical data sheet; perform pH, neutralisation and $K_{sp}$ calculations with correct units and significant figures; justify a treatment decision using equilibrium arguments; and present a short, well-organised technical report, as a junior chemist of a water company (for example \emph{Camwater}) would do.

\courssub{Situation}

\begin{experience}[The Village Water Treatment Challenge]
The village of \textbf{Batcham}, near Bafoussam in the West Region, is supplied by a borehole drilled into iron-rich lateritic ground. Since the last rainy season, families complain that the water tastes sour and that washed clothes come out with \textbf{reddish-brown stains}. The village water committee sends a sample to the regional laboratory. You are the junior chemist who receives the analysis report:

\begin{center}
\begin{tabularx}{\linewidth}{|l|X|}
\hline
\rowcolor{popblueL} \textbf{Parameter} & \textbf{Measured value} \\
\hline
pH (pH meter, \SI{25}{\ensuremath{{}^{\circ}\mathrm{C}}}) & $5.6$ \\
\hline
Total dissolved iron, as \ce{Fe^3+} & $c_0 = 2.0\times10^{-4}\ \mathrm{mol\,L^{-1}}$ \\
\hline
Temperature of all measurements & \SI{25}{\ensuremath{{}^{\circ}\mathrm{C}}} \\
\hline
\end{tabularx}
\end{center}

\textbf{Data provided by the laboratory (all at \SI{25}{\ensuremath{{}^{\circ}\mathrm{C}}}):}
\begin{itemize}
\item Ionic product of water: $K_w = 1.0\times10^{-14}\ \mathrm{mol^2\,L^{-2}}$.
\item Solubility product of iron(III) hydroxide: $K_{sp}\big(\ce{Fe(OH)3}\big) = 1.0\times10^{-38}\ \mathrm{mol^4\,L^{-4}}$.
\item Iron(III) ions hydrolyse in water; the first step is
\[
\ce{Fe^3+(aq) + H2O(l) <=> Fe(OH)^{2+}(aq) + H3O+(aq)}
\]
\item The treatment reagent available locally is \textbf{slaked lime}, \ce{Ca(OH)2} ($M = 74\ \mathrm{g\,mol^{-1}}$), assumed to dissolve and dissociate completely:
\[
\ce{Ca(OH)2(s) -> Ca^{2+}(aq) + 2OH^{-}(aq)}
\]
\item The borehole supplies a daily tank of volume $V = 10\,000\ \mathrm{L}$ ($10\ \mathrm{m^3}$).
\item Health guideline: the iron concentration in drinking water should fall below $1.0\times10^{-6}\ \mathrm{mol\,L^{-1}}$, and the pH should be brought close to $7$.
\end{itemize}

Your team must design the treatment and write a \textbf{one-page treatment report} containing all calculations and justifications.
\end{experience}

\begin{popfigure}
\begin{tikzpicture}[scale=0.9, every node/.style={font=\small}]
% Hydrolysis equilibrium
\node[draw, rounded corners, fill=popblueL!20, text width=4.5cm, align=center] (hyd) at (0,0) {
\textbf{Hydrolysis of \ce{Fe^3+}} \\
\ce{Fe^3+(aq) + H2O(l) <=> Fe(OH)^{2+}(aq) + H3O+(aq)} \\
\emph{Acidic hydrolysis} \\
Produces \ce{H3O+} $\rightarrow$ pH < 7
};
% Precipitation
\node[draw, rounded corners, fill=popgreenL!20, text width=4.5cm, align=center] (ppt) at (10,0) {
\textbf{Precipitation of \ce{Fe(OH)3}} \\
\ce{Fe^3+(aq) + 3OH^{-}(aq) <=> Fe(OH)3(s)} \\
$K_{sp} = [\ce{Fe^3+}][\ce{OH-}]^3$ \\
High \ce{OH-} $\rightarrow$ shifts left $\rightarrow$ solid \ce{Fe(OH)3}
};
% Lime addition
\node[draw, rounded corners, fill=poporange!20, text width=4cm, align=center] (lime) at (5,-3.5) {
\textbf{Slaked lime addition} \\
\ce{Ca(OH)2(s) -> Ca^{2+}(aq) + 2OH^{-}(aq)} \\
\ce{OH-} neutralises \ce{H3O+} \\
\ce{OH-} raises pH $\rightarrow$ precipitates \ce{Fe(OH)3}
};
\draw[->, thick, popblue] (hyd.east) -- node[above] {add \ce{OH-}} (ppt.west);
\draw[->, thick, poporange] (lime.north west) -- (hyd.south west);
\draw[->, thick, poporange] (lime.north east) -- (ppt.south east);
\end{tikzpicture}
\legende{Equilibria involved in the Batcham water treatment: hydrolysis of \ce{Fe^3+} (acidic), precipitation of \ce{Fe(OH)3} (favoured by high pH), and the role of slaked lime.}
\end{popfigure}

\courssub{Tasks}

\begin{enumerate}
\item \textbf{Characterising the acidity.} Using the measured pH and $K_w$, calculate $[\ce{H3O+}]$ and $[\ce{OH-}]$ in the borehole water. Is the water acidic, neutral or basic? Justify.

\item \textbf{Why is the water acidic?} The acidity is partly due to dissolved \ce{Fe^3+} salts. Using the hydrolysis equation given and \cle{Le Chatelier's principle}, explain why a solution of \ce{Fe^3+} is acidic, and predict what happens to the hydrolysis equilibrium when \ce{OH-} ions are added. Name the type of reaction by which a cation such as \ce{Fe^3+} produces \ce{H3O+} in water.

\item \textbf{Dosing the lime -- free acidity neutralisation.} The committee first considers neutralising only the free acidity. Calculate the amount (in moles) of \ce{OH-} needed to neutralise the free \ce{H3O+} in the daily tank, the corresponding amount of \ce{Ca(OH)2}, and the mass of slaked lime required. Write the neutralisation equation. Why, in practice, must \emph{much more} lime than this be added (think of the hydrolysis equilibrium of Task 2 and the precipitation of iron)?

\item \textbf{Removing the iron.} Write the dissolution equilibrium of \ce{Fe(OH)3} and the expression of its $K_{sp}$. Show that as lime raises the pH, \ce{Fe(OH)3} precipitates. Calculate the residual concentration $[\ce{Fe^3+}]$ when the treated water reaches pH $= 7.0$. Compare with the health guideline and conclude. What mass of iron (as \ce{Fe}) is removed per day from the $10\,000\ \mathrm{L}$ tank? ($M(\ce{Fe}) = 56\ \mathrm{g\,mol^{-1}}$.)

\item \textbf{Total lime requirement.} Calculate the total amount of \ce{Ca(OH)2} needed per day to (i) neutralise the free \ce{H3O+}, (ii) precipitate essentially all \ce{Fe^3+} as \ce{Fe(OH)3}, and (iii) provide a slight excess to reach pH 7. Give the total mass in kg.

\item \textbf{Would a buffer help?} A trainee suggests adding a buffer to ``fix the pH at 7 permanently''. Using what you know about the composition and functioning of buffer solutions, explain why installing a buffer is \emph{undesirable} here (think of the quantities of reagents needed, the cost, and the purpose of a buffer versus the purpose of a treatment).

\item \textbf{Report.} Assemble your results into a one-page treatment report: diagnosis, proposed treatment, quantities of reagent, expected final quality of the water, and one precaution for safe operation.
\end{enumerate}

\courssub{Model answer}

\textbf{Task 1 — Ion concentrations.}

By definition $\mathrm{pH} = -\log_{10}[\ce{H3O+}]$, so
\[
[\ce{H3O+}] = 10^{-\mathrm{pH}} = 10^{-5.6} = 2.5\times10^{-6}\ \mathrm{mol\,L^{-1}} \quad (\text{2 significant figures}).
\]
Since $K_w = [\ce{H3O+}][\ce{OH-}] = 1.0\times10^{-14}\ \mathrm{mol^2\,L^{-2}}$ at \SI{25}{\ensuremath{{}^{\circ}\mathrm{C}}},
\[
[\ce{OH-}] = \frac{K_w}{[\ce{H3O+}]} = \frac{1.0\times10^{-14}}{2.5\times10^{-6}} = 4.0\times10^{-9}\ \mathrm{mol\,L^{-1}}.
\]
Because $[\ce{H3O+}] > [\ce{OH-}]$ (and pH $< 7$), the water is \textbf{acidic}.

\begin{aretenir}[pH and $K_w$ at \SI{25}{\ensuremath{{}^{\circ}\mathrm{C}}}]
$\mathrm{pH} = -\log[\ce{H3O+}]$, $\mathrm{pOH} = -\log[\ce{OH-}]$, $\mathrm{pH} + \mathrm{pOH} = 14.00$.
Acidic: $\mathrm{pH} < 7$ ($[\ce{H3O+}] > [\ce{OH-}]$); Neutral: $\mathrm{pH} = 7$; Basic: $\mathrm{pH} > 7$.
\end{aretenir}

\textbf{Task 2 — Origin of the acidity.}

The \ce{Fe^3+} ion is small and highly charged (high charge density); it strongly polarises the \ce{O-H} bonds in the water molecules coordinated to it in the hydration sphere, facilitating the loss of a proton:
\[
\ce{Fe^3+(aq) + H2O(l) <=> Fe(OH)^{2+}(aq) + H3O+(aq)} \qquad (\text{acidic hydrolysis})
\]
This is \cle{salt hydrolysis} (specifically, acidic hydrolysis of the cation of a weak base, \ce{Fe(OH)3}). In Br\o nsted--Lowry terms, the hydrated \ce{Fe^3+} ion acts as an \cle{acid} (proton donor). The equilibrium produces \ce{H3O+}, lowering the pH.

When \ce{OH-} is added (from lime), it reacts with the free \ce{H3O+}:
\[
\ce{H3O+(aq) + OH^{-}(aq) -> 2H2O(l)}
\]
By \cle{Le Chatelier's principle}, removing the product \ce{H3O+} shifts the hydrolysis equilibrium to the \textbf{right}, causing more \ce{Fe^3+} to hydrolyse and release additional \ce{H3O+}. This continues until most \ce{Fe^3+} is converted to \ce{Fe(OH)^{2+}} and eventually to solid \ce{Fe(OH)3} (see Task 4). Hence the \emph{total} acidity (free \ce{H3O+} + hydrolysable \ce{Fe^3+}) is much larger than the free \ce{H3O+} alone.

\begin{attention}[Common mistake]
Do not confuse \emph{free acidity} (measured by pH) with \emph{total acidity} (including the \ce{H3O+} that will be released as hydrolysable cations react with added base). The titration curve of a \ce{Fe^3+} solution with a strong base shows a large buffer-like region due to this hydrolysis.
\end{attention}

\textbf{Task 3 — Lime dose for free acidity only.}

Neutralisation equation:
\[
\ce{H3O+(aq) + OH^{-}(aq) -> 2H2O(l)} \qquad \text{mole ratio } 1:1.
\]
Volume of daily tank: $V = 10\,000\ \mathrm{L}$.
\[
n(\ce{H3O+})_{\text{free}} = [\ce{H3O+}] \times V = 2.5\times10^{-6}\ \mathrm{mol\,L^{-1}} \times 10\,000\ \mathrm{L} = 2.5\times10^{-2}\ \mathrm{mol}.
\]
Hence $n(\ce{OH-})_{\text{neutralisation}} = 2.5\times10^{-2}\ \mathrm{mol}$.
Each mole of \ce{Ca(OH)2} yields $2\ \mathrm{mol}\ \ce{OH-}$:
\[
n\big(\ce{Ca(OH)2}\big)_{\text{free}} = \frac{2.5\times10^{-2}}{2} = 1.25\times10^{-2}\ \mathrm{mol}.
\]
Mass of slaked lime:
\[
m = n \times M = 1.25\times10^{-2}\ \mathrm{mol} \times 74\ \mathrm{g\,mol^{-1}} = 0.925\ \mathrm{g} \approx \mathbf{0.93\ g}.
\]

\textbf{Why much more lime is needed in practice:}
\begin{enumerate}
\item The hydrolysis equilibrium (Task 2) shifts right as free \ce{H3O+} is consumed, releasing more \ce{H3O+} from \ce{Fe^3+}. This continues until the \ce{Fe^3+} is largely precipitated.
\item The precipitation of \ce{Fe(OH)3} itself consumes \ce{OH-} ions (3 mol \ce{OH-} per mol \ce{Fe^3+}).
\item A slight excess is needed to reach and maintain pH $\approx 7$.
\end{enumerate}
The 0.93 g calculated above covers \emph{only} the free \ce{H3O+}; the dominant consumption is for iron precipitation (Task 5).

\textbf{Task 4 — Precipitation of iron.}

Dissolution equilibrium of iron(III) hydroxide:
\[
\ce{Fe(OH)3(s) <=> Fe^{3+}(aq) + 3OH^{-}(aq)}
\]
Solubility product expression:
\[
K_{sp} = [\ce{Fe^3+}][\ce{OH-}]^3 = 1.0\times10^{-38}\ \mathrm{mol^4\,L^{-4}} \quad (\text{at \SI{25}{\ensuremath{{}^{\circ}\mathrm{C}}}}).
\]
Adding lime increases $[\ce{OH-}]$; the ionic product $Q = [\ce{Fe^3+}][\ce{OH-}]^3$ quickly exceeds $K_{sp}$, so the equilibrium shifts to the \textbf{left} and \textbf{\ce{Fe(OH)3} precipitates} as a reddish-brown solid (the stains on clothes).

At target pH $= 7.0$:
\[
[\ce{H3O+}] = 10^{-7.0} = 1.0\times10^{-7}\ \mathrm{mol\,L^{-1}}, \qquad
[\ce{OH-}] = \frac{K_w}{[\ce{H3O+}]} = \frac{1.0\times10^{-14}}{1.0\times10^{-7}} = 1.0\times10^{-7}\ \mathrm{mol\,L^{-1}}.
\]
Residual \ce{Fe^3+} concentration:
\[
[\ce{Fe^3+}]_{\text{residual}} = \frac{K_{sp}}{[\ce{OH-}]^3}
= \frac{1.0\times10^{-38}}{(1.0\times10^{-7})^3}
= 1.0\times10^{-17}\ \mathrm{mol\,L^{-1}}.
\]
Comparison with guideline ($1.0\times10^{-6}\ \mathrm{mol\,L^{-1}}$):
\[
1.0\times10^{-17} \ll 1.0\times10^{-6} \quad \Rightarrow \quad \textbf{the iron is effectively completely removed.}
\]

Mass of iron removed per day (assuming initial $c_0 = 2.0\times10^{-4}\ \mathrm{mol\,L^{-1}}$ is fully precipitated):
\[
n(\ce{Fe})_{\text{removed}} \approx c_0 V = 2.0\times10^{-4}\ \mathrm{mol\,L^{-1}} \times 10\,000\ \mathrm{L} = 2.0\ \mathrm{mol}.
\]
\[
m(\ce{Fe}) = n \times M(\ce{Fe}) = 2.0\ \mathrm{mol} \times 56\ \mathrm{g\,mol^{-1}} = \mathbf{112\ g\ per\ day}.
\]
(Recovered as \ce{Fe(OH)3} sludge; molar mass \ce{Fe(OH)3} $\approx 107\ \mathrm{g\,mol^{-1}}$, so $\approx 214\ \mathrm{g}$ sludge/day.)

\textbf{Task 5 — Total lime requirement.}

\begin{enumerate}
\item \textbf{Neutralise free \ce{H3O+}:} $n(\ce{OH-})_1 = 2.5\times10^{-2}\ \mathrm{mol}$.
\item \textbf{Precipitate \ce{Fe^3+}:} $\ce{Fe^3+ + 3OH- -> Fe(OH)3(s)}$.
$n(\ce{Fe^3+})_{\text{initial}} = 2.0\ \mathrm{mol}$.
$n(\ce{OH-})_2 = 3 \times 2.0 = 6.0\ \mathrm{mol}$.
\item \textbf{Excess to reach pH 7:} At pH 7, $[\ce{OH-}] = 1.0\times10^{-7}\ \mathrm{M}$ in $10^4\ \mathrm{L}$ $\Rightarrow$ $n(\ce{OH-})_{\text{excess}} = 1.0\times10^{-3}\ \mathrm{mol}$ (negligible).
\end{enumerate}
Total $\ce{OH-}$ required $\approx 6.025\ \mathrm{mol}$.
\[
n\big(\ce{Ca(OH)2}\big)_{\text{total}} = \frac{6.025}{2} = 3.0125\ \mathrm{mol}.
\]
\[
m\big(\ce{Ca(OH)2}\big)_{\text{total}} = 3.0125 \times 74 = 222.9\ \mathrm{g} \approx \mathbf{0.223\ kg\ per\ day}.
\]
(Add $\sim 5\%$ excess for incomplete mixing/kinetics $\rightarrow$ \textbf{$\approx 0.23\ \mathrm{kg/day}$}.)

\begin{attention}[Key stoichiometry]
$\ce{Ca(OH)2} \rightarrow \ce{Ca^2+} + \mathbf{2}\ \ce{OH-}$.
$\ce{Fe^3+} + \mathbf{3}\ \ce{OH-} \rightarrow \ce{Fe(OH)3(s)}$.
The cube in $K_{sp} = [\ce{Fe^3+}][\ce{OH-}]^3$ means a small increase in $[\ce{OH-}]$ drastically reduces $[\ce{Fe^3+}]$.
\end{attention}

\textbf{Task 6 — Why no buffer.}

A buffer solution (e.g., \ce{CH3COOH}/\ce{CH3COO-} or \ce{NH4+}/\ce{NH3}) contains a weak acid/base conjugate pair in comparable concentrations. It \emph{resists} pH change upon addition of \emph{small} amounts of acid or base, but only while both components remain present.

Here the goal is the \emph{opposite}: to \textbf{permanently change} the pH from 5.6 to 7 and to \textbf{keep iron precipitated}. Buffering the water at pH 7 would require:
\begin{itemize}
\item Dissolving large quantities of organic reagents (or \ce{NH3}/\ce{NH4+} salts) in $10\ \mathrm{m^3}$ of drinking water \emph{every day} — costly and impractical.
\item Adding unnecessary dissolved substances to water meant for human consumption.
\item Fighting against the very adjustment we want: a buffer would consume the added \ce{OH-} to maintain pH, preventing the pH rise needed to precipitate \ce{Fe(OH)3}.
\end{itemize}
A one-off, controlled dose of cheap slaked lime achieves the treatment; a buffer would be counterproductive and uneconomical.

\textbf{Task 7 — Report (summary).}

\begin{center}
\begin{tabularx}{\linewidth}{|l|X|}
\hline
\rowcolor{popblueL} \textbf{BATCHAM WATER TREATMENT REPORT} & \textbf{Date: [Today]} \\
\hline
\textbf{Diagnosis} & Borehole water: pH 5.6 (acidic), $[\ce{Fe^3+}] = 2.0\times10^{-4}\ \mathrm{M}$. Acidity due to \ce{Fe^3+} hydrolysis. Reddish stains = \ce{Fe(OH)3} precipitating on exposure to air. \\
\hline
\textbf{Proposed treatment} & Daily dose of slaked lime \ce{Ca(OH)2} to $10\,000\ \mathrm{L}$ tank. Stir, allow \ce{Fe(OH)3} to precipitate (15--30 min), filter/decant. \\
\hline
\textbf{Reagent quantity} & \textbf{0.23 kg \ce{Ca(OH)2} per day} (covers free acidity, complete \ce{Fe^3+} precipitation, and pH 7 target). \\
\hline
\textbf{Expected final quality} & pH $\approx 7.0$; $[\ce{Fe^3+}] \approx 1\times10^{-17}\ \mathrm{M}$ (guideline: $< 1\times10^{-6}\ \mathrm{M}$); iron removed $\approx 112\ \mathrm{g/day}$ as \ce{Fe(OH)3} sludge. \\
\hline
\textbf{Precaution} & Add lime in small portions while monitoring pH with a meter; avoid overshooting above pH 8 (causes \ce{CaCO3} scaling, bitter taste). Wear gloves/goggles — \ce{Ca(OH)2} is caustic. \\
\hline
\end{tabularx}
\end{center}

\begin{attention}[Common mistakes to avoid in the examination]
\begin{itemize}
\item Confusing \emph{free acidity} ($[\ce{H3O+}]$ from pH) with \emph{total acidity} (including hydrolysable \ce{Fe^3+}). The latter determines the true lime demand.
\item Forgetting the $1:2$ stoichiometry of \ce{Ca(OH)2} to \ce{OH-}, or the $1:3$ stoichiometry of \ce{Fe^3+} to \ce{OH-} in precipitation.
\item Miswriting $K_{sp}$ expression: it is $[\ce{Fe^3+}][\ce{OH-}]^3$ (cube!), not $[\ce{Fe^3+}][\ce{OH-}]$.
\item Using $K_w$ or $K_{sp}$ values at the wrong temperature; always state \SI{25}{\ensuremath{{}^{\circ}\mathrm{C}}}.
\item Suggesting a buffer for a process that requires a \emph{permanent} pH shift and precipitation.
\item Omitting units or giving answers to unjustified significant figures (pH 5.6 $\rightarrow$ 2 s.f. in concentrations).
\end{itemize}
\end{attention}

\begin{savaistu}[Cameroon context]
Lateritic soils in the West Region (Bafoussam, Dschang, Bafang) are rich in iron oxides. Boreholes often yield water with high \ce{Fe^2+}/\ce{Fe^3+} and low pH. The standard Camwater treatment is aeration (oxidises \ce{Fe^2+} to \ce{Fe^3+}) followed by liming and filtration. The \ce{Fe(OH)3} sludge is a disposal issue — it can be dried and used as a pigment or in brick-making.
\end{savaistu}

\newpage

\coursec{Methods and worked examples}

\courssub{Method 1 — Writing the Equilibrium Law and Calculating $K_c$}

\begin{methode}[Calculating $K_c$ from equilibrium data]
\begin{enumerate}
\item \textbf{Balance the equation} and write it with state symbols.
\item \textbf{Write the equilibrium law}: for $\ce{aA + bB <=> cC + dD}$,
\[K_c=\dfrac{[\ce{C}]^c\,[\ce{D}]^d}{[\ce{A}]^a\,[\ce{B}]^b}\]
Products on top, reactants below, each concentration raised to its \cle{stoichiometric coefficient}.
\item \textbf{Build a ``start--change--equilibrium'' (ICE) table} in moles (or directly in concentrations if the volume is given): initial amounts, change (in the ratio of the coefficients), equilibrium amounts.
\item \textbf{Convert moles to concentrations}: $c=\dfrac{n}{V}$. Never forget the volume of the container!
\item \textbf{Substitute} into the expression, keeping units ($\mathrm{mol\,dm^{-3}}$) at each step, and work out the \textbf{units of $K_c$} by cancelling: units $=(\mathrm{mol\,dm^{-3}})^{\Delta n}$ where $\Delta n = (c+d)-(a+b)$.
\item Quote the answer to 2--3 significant figures.
\end{enumerate}
\textbf{Reflexes:} pure solids and pure liquids never appear in the expression; if $K_c$ is asked for the \emph{reverse} reaction, take $\dfrac{1}{K_c}$; if the equation is doubled, square $K_c$.
\end{methode}

\begin{exoresolu}[Esterification equilibrium]
\SI{3.0}{mol} of ethanoic acid and \SI{2.0}{mol} of ethanol are mixed and allowed to reach equilibrium at \SI{298}{K} in a total volume of \SI{2.0}{dm^3}:
\[\ce{CH3COOH(l) + C2H5OH(l) <=> CH3COOC2H5(l) + H2O(l)}\]
At equilibrium, \SI{1.2}{mol} of ethyl ethanoate is present. Calculate $K_c$ and state its units.
\tcblower
\textbf{Step 1 -- ICE table (moles):} since 1.2 mol of ester has formed, 1.2 mol of each reactant has been consumed (coefficients all equal to 1).
\begin{center}
\begin{tabular}{|l|c|c|c|c|}
\hline
\rowcolor{popblueL} & \ce{CH3COOH} & \ce{C2H5OH} & \ce{CH3COOC2H5} & \ce{H2O} \\
\hline
Start / mol & 3.0 & 2.0 & 0 & 0 \\
\hline
Change / mol & $-1.2$ & $-1.2$ & $+1.2$ & $+1.2$ \\
\hline
Equilibrium / mol & 1.8 & 0.8 & 1.2 & 1.2 \\
\hline
\end{tabular}
\end{center}
\textbf{Step 2 -- Concentrations} ($V=\SI{2.0}{dm^3}$, using $c=n/V$):
\[[\ce{CH3COOH}]=0.90,\quad [\ce{C2H5OH}]=0.40,\quad [\ce{CH3COOC2H5}]=[\ce{H2O}]=0.60\ \mathrm{mol\,dm^{-3}}\]
\textbf{Step 3 -- Substitute:}
\[K_c=\dfrac{[\ce{CH3COOC2H5}][\ce{H2O}]}{[\ce{CH3COOH}][\ce{C2H5OH}]}=\dfrac{(0.60)(0.60)}{(0.90)(0.40)}=\dfrac{0.36}{0.36}=1.0\]
\textbf{Units:} $\Delta n = 2-2=0$, so $K_c$ has \textbf{no units}. \checkmark\ Note how the volume $V$ cancels here (equal numbers of moles on both sides) — a useful check on your working.
\end{exoresolu}

\courssub{Method 2 — Calculating $K_p$ from Partial Pressures}

\begin{methode}[$K_p$ and partial pressures]
\begin{enumerate}
\item Write the equation and the expression $K_p=\dfrac{p_{\text{products}}^{\nu}}{p_{\text{reactants}}^{\nu}}$ (gaseous species only).
\item Find the \textbf{mole fraction} of each gas: $x_i=\dfrac{n_i}{n_{\text{total}}}$.
\item Compute each \textbf{partial pressure}: $p_i=x_i\,P_{\text{total}}$.
\item Substitute into $K_p$; units $=(\text{pressure unit})^{\Delta n_g}$, where $\Delta n_g$ counts gaseous moles only.
\item Link with $K_c$ if needed: $K_p=K_c(RT)^{\Delta n_g}$.
\end{enumerate}
\textbf{Reflex:} if the \emph{total pressure at equilibrium} is given together with the amount dissociated, first complete the ICE table in moles, then get mole fractions, then partial pressures.
\end{methode}

\begin{exoresolu}[Dissociation of dinitrogen tetroxide]
\SI{1.00}{mol} of \ce{N2O4} is placed in a vessel at \SI{350}{K}. At equilibrium the total pressure is \SI{2.0}{atm} and \SI{0.20}{mol} of \ce{N2O4} has dissociated:
\[\ce{N2O4(g) <=> 2NO2(g)}\]
Calculate $K_p$ at this temperature and give its units.
\tcblower
\textbf{Step 1 -- ICE table (moles):} each mole of \ce{N2O4} that dissociates gives \emph{two} moles of \ce{NO2}.
\begin{center}
\begin{tabular}{|l|c|c|}
\hline
\rowcolor{popblueL} & \ce{N2O4} & \ce{NO2} \\
\hline
Start / mol & 1.00 & 0 \\
\hline
Change / mol & $-0.20$ & $+0.40$ \\
\hline
Equilibrium / mol & 0.80 & 0.40 \\
\hline
\end{tabular}
\end{center}
Total moles at equilibrium: $n_{\text{tot}}=0.80+0.40=1.20\ \mathrm{mol}$.
\textbf{Step 2 -- Partial pressures} ($p_i=x_iP_{\text{total}}$):
\[p_{\ce{N2O4}}=\frac{0.80}{1.20}\times 2.0=1.33\ \mathrm{atm};\qquad p_{\ce{NO2}}=\frac{0.40}{1.20}\times 2.0=0.667\ \mathrm{atm}\]
\textbf{Step 3 -- $K_p$:}
\[K_p=\dfrac{p_{\ce{NO2}}^{2}}{p_{\ce{N2O4}}}=\dfrac{(0.667)^2}{1.33}=\dfrac{0.444}{1.33}=0.33\ \mathrm{atm}\]
\textbf{Units:} $\Delta n_g=2-1=+1$, so $K_p$ is in \textbf{atm} (or Pa if pressures are in Pa). \checkmark
\end{exoresolu}

\courssub{Method 3 — Heterogeneous Equilibria and the Effect of Temperature on $K$}

\begin{methode}[Heterogeneous equilibria and van 't Hoff reasoning]
\begin{enumerate}
\item \textbf{Omit pure solids and pure liquids} from the equilibrium expression — their concentrations are constant and are absorbed into the value of $K$.
\item For a thermal decomposition such as $\ce{CaCO3(s) <=> CaO(s) + CO2(g)}$, $K_p=p_{\ce{CO2}}$ only.
\item \textbf{Temperature effect:} $K$ changes \emph{only} with temperature. For an \cle{endothermic} forward reaction ($\Delta H>0$), $K$ \emph{increases} as $T$ rises; for an \cle{exothermic} one ($\Delta H<0$), $K$ decreases as $T$ rises.
\item Justify with Le Chatelier's principle: raising $T$ shifts the equilibrium in the direction that absorbs heat.
\end{enumerate}
\textbf{Reflex:} catalysts, pressure and concentration changes alter the \emph{position} of equilibrium but \emph{never} the value of $K$.
\end{methode}

\begin{exoresolu}[Decomposition of calcium carbonate]
Limestone from a quarry near Figuil is heated in a closed kiln:
\[\ce{CaCO3(s) <=> CaO(s) + CO2(g)} \qquad \Delta H^{\theta}=+178\ \mathrm{kJ\,mol^{-1}}\]
(a) Write the expression for $K_p$. (b) At \SI{1100}{K} the pressure of \ce{CO2} at equilibrium is \SI{0.50}{atm}. State $K_p$. (c) Explain the effect on $K_p$ of (i) raising the temperature, (ii) adding more \ce{CaCO3}, (iii) pumping away \ce{CO2}.
\tcblower
\textbf{(a)} Both solids are omitted:
\[K_p=p_{\ce{CO2}}\]
\textbf{(b)} $K_p=0.50$ atm.
\textbf{(c)} (i) The forward reaction is endothermic ($\Delta H>0$). Raising $T$ shifts the equilibrium to the right (the direction that absorbs heat), so $p_{\ce{CO2}}$ rises and \textbf{$K_p$ increases}.
(ii) Adding \ce{CaCO3(s)}: a pure solid has a constant concentration term; the position of equilibrium and $K_p$ are \textbf{unchanged}.
(iii) Removing \ce{CO2} momentarily lowers $p_{\ce{CO2}}$; by Le Chatelier the system decomposes more \ce{CaCO3} until $p_{\ce{CO2}}$ returns to $0.50$ atm. The \emph{position} shifts right, but \textbf{$K_p$ is unchanged} because the temperature is constant. \checkmark
\end{exoresolu}

\courssub{Method 4 — Balancing Redox Equations (Ion–Electron Method)}

\begin{methode}[Half-equation method in acid solution]
\begin{enumerate}
\item Assign \textbf{oxidation numbers} to identify what is oxidised (oxidation number rises) and what is reduced (oxidation number falls).
\item Write the two \textbf{half-equations}: balance the element being oxidised/reduced, then balance O with \ce{H2O}, then H with \ce{H+}, then charge with electrons \ce{e-}.
\item Multiply the half-equations so that the \textbf{numbers of electrons are equal}.
\item \textbf{Add} the half-equations and cancel electrons (and any \ce{H+} or \ce{H2O} appearing on both sides).
\item \textbf{Check}: atoms balanced \emph{and} total charge balanced on both sides.
\end{enumerate}
\textbf{Reflex:} in alkaline solution, neutralise each \ce{H+} by adding \ce{OH-} to both sides, then combine \ce{H+ + OH-} into \ce{H2O}.
\end{methode}

\begin{exoresolu}[Permanganate and iron(II)]
Acidified potassium manganate(VII), \ce{KMnO4}, oxidises iron(II) ions to iron(III). Given
\[\ce{MnO4-(aq) + 8H+(aq) + 5e- -> Mn^{2+}(aq) + 4H2O(l)}\]
construct the balanced ionic equation for the titration reaction used in redox analysis.
\tcblower
\textbf{Step 1 -- Oxidation numbers:} Mn falls from $+7$ to $+2$ (reduction, gain of 5 \ce{e-}); Fe rises from $+2$ to $+3$ (oxidation, loss of 1 \ce{e-}).
\textbf{Step 2 -- Oxidation half-equation:}
\[\ce{Fe^{2+}(aq) -> Fe^{3+}(aq) + e-}\]
\textbf{Step 3 -- Equalise electrons:} multiply the iron half-equation by 5:
\[\ce{5Fe^{2+}(aq) -> 5Fe^{3+}(aq) + 5e-}\]
\textbf{Step 4 -- Add} and cancel the 5 electrons:
\[\ce{MnO4-(aq) + 8H+(aq) + 5Fe^{2+}(aq) -> Mn^{2+}(aq) + 5Fe^{3+}(aq) + 4H2O(l)}\]
\textbf{Step 5 -- Check:} Mn: 1/1; O: 4/4; H: 8/8; Fe: 5/5. Charge: left $=(-1)+(+8)+5(+2)=+17$; right $=(+2)+5(+3)=+17$. \checkmark\ Balanced. This is why \SI{1}{mol} of \ce{MnO4-} reacts with exactly \SI{5}{mol} of \ce{Fe^{2+}} in titrations — the basis of redox titration calculations.
\end{exoresolu}

\courssub{Method 5 — Cell emf and Predicting Spontaneous Redox Reactions}

\begin{methode}[Using standard electrode potentials]
\begin{enumerate}
\item Write both half-cells as \textbf{reductions} with their $E^{\theta}$ values.
\item The \textbf{more positive} $E^{\theta}$ system is reduced (cathode, positive pole); the other is reversed and acts as the oxidation half-cell (anode, negative pole).
\item Cell emf: $E^{\theta}_{\text{cell}}=E^{\theta}_{\text{cathode}}-E^{\theta}_{\text{anode}}$ (right-hand electrode minus left-hand electrode in the conventional cell diagram).
\item If $E^{\theta}_{\text{cell}}>0$, the reaction is \textbf{feasible (spontaneous)} under standard conditions.
\item \textbf{Never multiply} $E^{\theta}$ when multiplying a half-equation — it is an intensive quantity (it does not depend on amounts).
\end{enumerate}
\end{methode}

\begin{exoresolu}[Copper–silver cell]
Given $E^{\theta}(\ce{Cu^{2+}/Cu})=+0.34\ \mathrm{V}$ and $E^{\theta}(\ce{Ag+/Ag})=+0.80\ \mathrm{V}$:
(a) Write the cell diagram for the standard cell. (b) Calculate $E^{\theta}_{\text{cell}}$. (c) Write the overall cell reaction. (d) State the direction of electron flow in the external circuit.
\tcblower
\textbf{(a)} The copper electrode (less positive $E^{\theta}$) is the negative pole, written on the left by convention:
\[\ce{Cu(s) | Cu^{2+}(aq, 1 mol dm^{-3}) || Ag+(aq, 1 mol dm^{-3}) | Ag(s)}\]
\textbf{(b)}
\[E^{\theta}_{\text{cell}}=E^{\theta}_{\text{right}}-E^{\theta}_{\text{left}}=(+0.80)-(+0.34)=+0.46\ \mathrm{V}\]
\textbf{(c)} Silver is reduced (cathode): $\ce{Ag+ + e- -> Ag}$; copper is oxidised (anode): $\ce{Cu -> Cu^{2+} + 2e-}$. Doubling the silver half-equation (without touching $E^{\theta}$!):
\[\ce{Cu(s) + 2Ag+(aq) -> Cu^{2+}(aq) + 2Ag(s)}\]
\textbf{(d)} Electrons flow through the external circuit from the \textbf{copper electrode (anode, negative pole)} to the \textbf{silver electrode (cathode, positive pole)}. Since $E^{\theta}_{\text{cell}}>0$, the reaction is spontaneous under standard conditions. \checkmark
\end{exoresolu}

\courssub{Method 6 — pH of Strong and Weak Acids}

\begin{methode}[pH calculations]
\begin{enumerate}
\item \textbf{Strong monoprotic acid} (full dissociation): $[\ce{H3O+}]=c_a$, so $\mathrm{pH}=-\log c_a$.
\item \textbf{Strong base}: $[\ce{OH-}]=c_b$, then $[\ce{H3O+}]=\dfrac{K_w}{[\ce{OH-}]}$ with $K_w=1.0\times10^{-14}\ \mathrm{mol^2\,dm^{-6}}$ at \SI{298}{K}; equivalently $\mathrm{pH}=14-\mathrm{pOH}$.
\item \textbf{Weak acid}: set up $K_a=\dfrac{[\ce{H3O+}]^2}{c_a-[\ce{H3O+}]}\approx\dfrac{[\ce{H3O+}]^2}{c_a}$, hence $[\ce{H3O+}]\approx\sqrt{K_a\,c_a}$ and $\mathrm{pH}=\tfrac12(\mathrm{p}K_a-\log c_a)$.
\item \textbf{Check the approximation}: the degree of dissociation must be less than about $5\%$; otherwise solve the quadratic equation.
\item \textbf{Weak base}: same pattern with $K_b$: $[\ce{OH-}]\approx\sqrt{K_b\,c_b}$, then convert through $K_w$.
\end{enumerate}
\end{methode}

\begin{exoresolu}[Comparing a strong and a weak acid]
(a) Calculate the pH of \SI{0.010}{mol.dm^{-3}} hydrochloric acid. (b) Calculate the pH of \SI{0.010}{mol.dm^{-3}} ethanoic acid, $K_a=1.8\times10^{-5}\ \mathrm{mol\,dm^{-3}}$. (c) Comment on the difference.
\tcblower
\textbf{(a)} \ce{HCl} is a strong acid, fully ionised in water:
\[[\ce{H3O+}]=0.010\ \mathrm{mol\,dm^{-3}}\quad\Rightarrow\quad \mathrm{pH}=-\log(0.010)=2.00\]
\textbf{(b)} Ethanoic acid is weak: $\ce{CH3COOH + H2O <=> CH3COO- + H3O+}$
\[K_a=\frac{[\ce{H3O+}][\ce{CH3COO-}]}{[\ce{CH3COOH}]}\approx\frac{[\ce{H3O+}]^2}{c_a}\]
\[[\ce{H3O+}]=\sqrt{K_a\,c_a}=\sqrt{1.8\times10^{-5}\times0.010}=\sqrt{1.8\times10^{-7}}=4.24\times10^{-4}\ \mathrm{mol\,dm^{-3}}\]
\[\mathrm{pH}=-\log(4.24\times10^{-4})=3.37\]
\textbf{Check:} degree of dissociation $=\dfrac{4.24\times10^{-4}}{0.010}\times100\%=4.2\%<5\%$ — the approximation is valid. \checkmark
\textbf{(c)} At the \emph{same concentration}, the weak acid has a higher pH (3.37 against 2.00) because it is only partially ionised: $[\ce{H3O+}]$ is about 24 times smaller than in the strong acid.
\end{exoresolu}

\courssub{Method 7 — pH of Buffer Solutions}

\begin{methode}[Buffer pH (Henderson--Hasselbalch)]
\begin{enumerate}
\item \textbf{Identify the buffer pair}: weak acid + its conjugate base (acidic buffer) or weak base + its conjugate acid (alkaline buffer).
\item Acidic buffer: $\mathrm{pH}=\mathrm{p}K_a+\log\dfrac{[\text{salt}]}{[\text{acid}]}$.
\item Alkaline buffer: $\mathrm{pOH}=\mathrm{p}K_b+\log\dfrac{[\text{salt}]}{[\text{base}]}$, then $\mathrm{pH}=14-\mathrm{pOH}$.
\item Use the \textbf{mole ratio} directly when both species share the same solution volume — the volumes cancel.
\item A buffer works best when $[\text{salt}]\approx[\text{acid}]$, i.e. when $\mathrm{pH}\approx\mathrm{p}K_a$.
\end{enumerate}
\end{methode}

\begin{exoresolu}[Ethanoate buffer]
A buffer is prepared by dissolving \SI{0.10}{mol} of ethanoic acid and \SI{0.15}{mol} of sodium ethanoate in water to make \SI{1.0}{dm^3} of solution. Given $\mathrm{p}K_a(\ce{CH3COOH})=4.74$:
(a) Calculate the pH of the buffer. (b) Explain, with equations, how the buffer resists pH change when a small amount of \ce{HCl} is added.
\tcblower
\textbf{(a)} Both species are in the same volume, so the mole ratio may be used directly:
\[\mathrm{pH}=\mathrm{p}K_a+\log\frac{[\ce{CH3COO-}]}{[\ce{CH3COOH}]}=4.74+\log\frac{0.15}{0.10}=4.74+\log(1.5)=4.74+0.18=4.92\]
\textbf{(b)} The added \ce{H+} is removed by the reservoir of ethanoate ions:
\[\ce{CH3COO-(aq) + H+(aq) -> CH3COOH(aq)}\]
The ratio $\dfrac{[\ce{CH3COO-}]}{[\ce{CH3COOH}]}$ changes only slightly (a little salt is converted into acid, but both are present in large amounts), so $\log(\text{ratio})$ — and hence the pH — barely changes. If a base were added instead, the acid reservoir would neutralise it: $\ce{CH3COOH + OH- -> CH3COO- + H2O}$. \checkmark
\end{exoresolu}

\courssub{Method 8 — Solubility Product and the Common Ion Effect}

\begin{methode}[$K_{sp}$ calculations]
\begin{enumerate}
\item Write the \textbf{dissolution equation}, e.g. $\ce{AgCl(s) <=> Ag+(aq) + Cl-(aq)}$.
\item Write $K_{sp}$ as the product of the ion concentrations raised to their coefficients (the solid is omitted).
\item If the \textbf{molar solubility} is $s$, express each ion concentration in terms of $s$ (e.g. for $\ce{CaF2}$: $[\ce{Ca^{2+}}]=s$, $[\ce{F-}]=2s$, so $K_{sp}=4s^3$), then substitute.
\item \textbf{Common ion effect}: in a solution already containing a common ion of concentration $c$, write $K_{sp}$ with $(c+s)$ for that ion and approximate $c+s\approx c$, since $s$ becomes very small.
\item Precipitation occurs when the \textbf{ionic product} $Q$ exceeds $K_{sp}$.
\end{enumerate}
\end{methode}

\begin{exoresolu}[Solubility of silver chloride]
$K_{sp}(\ce{AgCl})=1.8\times10^{-10}\ \mathrm{mol^2\,dm^{-6}}$ at \SI{298}{K}.
(a) Calculate the solubility of \ce{AgCl} in pure water in $\mathrm{mol\,dm^{-3}}$.
(b) Calculate its solubility in \SI{0.10}{mol.dm^{-3}} \ce{NaCl} solution, and comment.
\tcblower
\textbf{(a)} $\ce{AgCl(s) <=> Ag+(aq) + Cl-(aq)}$; let the molar solubility be $s$, so $[\ce{Ag+}]=[\ce{Cl-}]=s$:
\[K_{sp}=[\ce{Ag+}][\ce{Cl-}]=s^2\quad\Rightarrow\quad s=\sqrt{1.8\times10^{-10}}=1.34\times10^{-5}\ \mathrm{mol\,dm^{-3}}\]
\textbf{(b)} In \SI{0.10}{mol.dm^{-3}} \ce{NaCl}, the chloride ion is the common ion: $[\ce{Cl-}]=0.10+s'\approx0.10\ \mathrm{mol\,dm^{-3}}$. Then:
\[K_{sp}=s'(0.10)\quad\Rightarrow\quad s'=\frac{1.8\times10^{-10}}{0.10}=1.8\times10^{-9}\ \mathrm{mol\,dm^{-3}}\]
\textbf{Comment:} the solubility falls from $1.34\times10^{-5}$ to $1.8\times10^{-9}\ \mathrm{mol\,dm^{-3}}$ — about $7\,000$ times smaller. The added \ce{Cl-} shifts the dissolution equilibrium to the left (Le Chatelier's principle). This is exactly why, in qualitative analysis, an excess of the precipitating reagent is used to precipitate an ion almost completely. \checkmark
\end{exoresolu}

\begin{attention}[Frequent traps in equilibrium calculations]
\begin{itemize}
\item Using \textbf{moles instead of concentrations} in $K_c$ — always divide by the volume of the container.
\item Multiplying $E^{\theta}$ by stoichiometric coefficients — electrode potentials are intensive quantities and never change with scaling of the equation.
\item Forgetting that $K$ changes \textbf{only with temperature}; a catalyst speeds up the attainment of equilibrium but leaves both $K$ and the equilibrium position unchanged.
\item For weak acids, wrongly using $\mathrm{pH}=-\log c_a$ (that formula applies only to \emph{strong} monoprotic acids).
\item In $K_{sp}$ for salts of formula \ce{A2B} or \ce{AB2}, forgetting the factor 2 and the cube: $K_{sp}=(2s)^2(s)=4s^3$.
\item Omitting pure solids and pure liquids is correct — but do not omit \emph{gases} or \emph{aqueous species} from $K_p$ or $K_c$ expressions.
\end{itemize}
\end{attention}

\newpage

\coursec{Exercises}

\courssub{I check my knowledge}

\begin{exercice}[Multiple choice questions]
For each question, choose the single correct answer and justify your choice briefly.
\begin{enumerate}
\item For the equilibrium \ce{N2(g) + 3H2(g) <=> 2NH3(g)}, $\Delta H^{\theta} = -92\ \mathrm{kJ\,mol^{-1}}$, raising the temperature at constant pressure:
\begin{enumerate}
\item increases both the yield of \ce{NH3} and the value of $K_p$;
\item decreases the yield of \ce{NH3} and decreases $K_p$;
\item decreases the yield of \ce{NH3} but increases $K_p$;
\item has no effect because a catalyst is present.
\end{enumerate}
\item The units of $K_c$ for the equilibrium \ce{CH3COOH(l) + C2H5OH(l) <=> CH3COOC2H5(l) + H2O(l)} are:
\begin{enumerate}
\item $\mathrm{mol\,dm^{-3}}$;
\item $\mathrm{mol^{-1}\,dm^{3}}$;
\item $\mathrm{mol^{2}\,dm^{-6}}$;
\item no units.
\end{enumerate}
\item The pH of a $0.010\ \mathrm{mol\,dm^{-3}}$ solution of \ce{NaOH} at $25\ ^{\circ}\mathrm{C}$ ($K_w = 1.0\times 10^{-14}\ \mathrm{mol^{2}\,dm^{-6}}$) is:
\begin{enumerate}
\item $2.0$; \item $7.0$; \item $10.0$; \item $12.0$.
\end{enumerate}
\item In the cell \ce{Zn(s)|Zn^{2+}(aq)||Cu^{2+}(aq)|Cu(s)}, the electrode at which oxidation occurs is:
\begin{enumerate}
\item the copper electrode (cathode);
\item the zinc electrode (anode);
\item the salt bridge;
\item neither electrode.
\end{enumerate}
\end{enumerate}
\end{exercice}

\begin{exercice}[True or false]
State whether each statement is true or false, and justify your answer in one or two sentences.
\begin{enumerate}
\item At dynamic equilibrium, the forward and reverse reactions have both stopped.
\item Adding a catalyst to a system at equilibrium increases the equilibrium yield of the products.
\item The standard hydrogen electrode is assigned a standard electrode potential of $0.00\ \mathrm{V}$ under standard conditions.
\item A solution of \ce{NH4Cl} in pure water has a pH greater than $7$ at $25\ ^{\circ}\mathrm{C}$.
\item The solubility of \ce{AgCl} in $0.10\ \mathrm{mol\,dm^{-3}}$ \ce{NaCl} is smaller than its solubility in pure water.
\end{enumerate}
\end{exercice}

\begin{exercice}[Definitions]
Define each of the following terms precisely, giving one example in each case:
\begin{enumerate}
\item dynamic equilibrium;
\item Le Chatelier's principle;
\item standard electrode potential, $E^{\theta}$;
\item a Br\o nsted--Lowry acid and a Lewis base;
\item a buffer solution;
\item the solubility product, $K_{sp}$.
\end{enumerate}
\end{exercice}

\begin{exercice}[Direct calculations]
\begin{enumerate}
\item Calculate the pH at $25\ ^{\circ}\mathrm{C}$ of: (i) $0.050\ \mathrm{mol\,dm^{-3}}$ \ce{HCl}; (ii) $0.050\ \mathrm{mol\,dm^{-3}}$ \ce{NaOH}.
\item For \ce{PCl5(g) <=> PCl3(g) + Cl2(g)}, the equilibrium mixture contains $[\ce{PCl5}] = 0.10$, $[\ce{PCl3}] = 0.40$ and $[\ce{Cl2}] = 0.40\ \mathrm{mol\,dm^{-3}}$. Calculate $K_c$ and state its units.
\item Given $E^{\theta}(\ce{Zn^{2+}/Zn}) = -0.76\ \mathrm{V}$ and $E^{\theta}(\ce{Cu^{2+}/Cu}) = +0.34\ \mathrm{V}$, calculate the standard e.m.f. of the Daniell cell and identify the positive electrode.
\end{enumerate}
\end{exercice}

\courssub{I practise}

\begin{exercice}[Decomposition of phosphorus pentachloride]
$1.00\ \mathrm{mol}$ of \ce{PCl5} is placed in a closed vessel of volume $2.00\ \mathrm{dm^{3}}$ and heated to a temperature $T$. At equilibrium, $0.20\ \mathrm{mol}$ of \ce{PCl5} remains.
\[\ce{PCl5(g) <=> PCl3(g) + Cl2(g)}\]
\begin{enumerate}
\item Construct an ICE (initial--change--equilibrium) table in moles.
\item Deduce the equilibrium concentrations of the three species.
\item Calculate $K_c$ at temperature $T$ and state its units.
\item State and explain the effect on the equilibrium position of decreasing the volume of the vessel at constant temperature.
\end{enumerate}
\end{exercice}

\begin{exercice}[Redox titration of iron(II)]
\begin{enumerate}
\item Using the half-equation method, balance the equation for the reaction between acidified potassium manganate(VII) and iron(II) sulphate:
\[\ce{MnO4^{-}(aq) + Fe^{2+}(aq) + H^{+}(aq) -> Mn^{2+}(aq) + Fe^{3+}(aq) + H2O(l)}\]
\item Calculate the volume of $0.0200\ \mathrm{mol\,dm^{-3}}$ \ce{KMnO4} solution required to oxidise exactly the \ce{Fe^{2+}} contained in a sample of $1.52\ \mathrm{g}$ of \ce{FeSO4.7H2O} dissolved in dilute sulphuric acid. ($M(\ce{FeSO4.7H2O}) = 278\ \mathrm{g\,mol^{-1}}$)
\item State the colour change observed at the end-point and explain why no indicator is needed.
\end{enumerate}
\end{exercice}

\begin{exercice}[pH of a buffer solution]
A buffer solution is prepared by dissolving ethanoic acid and sodium ethanoate in water so that $[\ce{CH3COOH}] = 0.10\ \mathrm{mol\,dm^{-3}}$ and $[\ce{CH3COONa}] = 0.15\ \mathrm{mol\,dm^{-3}}$. Given $K_a(\ce{CH3COOH}) = 1.8\times 10^{-5}\ \mathrm{mol\,dm^{-3}}$ at $25\ ^{\circ}\mathrm{C}$:
\begin{enumerate}
\item Write the expression for $K_a$ of ethanoic acid and derive the buffer formula $\mathrm{pH} = \mathrm{p}K_a + \log\dfrac{[\text{salt}]}{[\text{acid}]}$.
\item Calculate the pH of the buffer.
\item $0.010\ \mathrm{mol}$ of \ce{NaOH} is added to $1.00\ \mathrm{dm^{3}}$ of this buffer (assume no volume change). Write the equation of the reaction that occurs and calculate the new pH.
\item Comment on the pH change compared with adding the same amount of \ce{NaOH} to $1.00\ \mathrm{dm^{3}}$ of pure water.
\end{enumerate}
\end{exercice}

\begin{exercice}[Solubility and common ion effect]
Given at $25\ ^{\circ}\mathrm{C}$: $K_{sp}(\ce{AgCl}) = 1.8\times 10^{-10}\ \mathrm{mol^{2}\,dm^{-6}}$ and $K_{sp}(\ce{Mg(OH)2}) = 1.8\times 10^{-11}\ \mathrm{mol^{3}\,dm^{-9}}$.
\begin{enumerate}
\item Calculate the solubility, in $\mathrm{mol\,dm^{-3}}$, of \ce{AgCl} and of \ce{Mg(OH)2} in pure water.
\item Calculate the solubility of \ce{AgCl} in a $0.10\ \mathrm{mol\,dm^{-3}}$ \ce{NaCl} solution, and explain the result using Le Chatelier's principle.
\item $50.0\ \mathrm{cm^{3}}$ of $2.0\times 10^{-4}\ \mathrm{mol\,dm^{-3}}$ \ce{AgNO3} is mixed with $50.0\ \mathrm{cm^{3}}$ of $2.0\times 10^{-4}\ \mathrm{mol\,dm^{-3}}$ \ce{NaCl}. Determine, by calculation, whether a precipitate of \ce{AgCl} forms.
\end{enumerate}
\end{exercice}

\courssub{I prepare for the GCE}

\begin{exercice}[The Haber process]
Ammonia is manufactured industrially by the Haber process:
\[\ce{N2(g) + 3H2(g) <=> 2NH3(g)} \qquad \Delta H^{\theta} = -92\ \mathrm{kJ\,mol^{-1}}\]
\begin{enumerate}
\item Write the expression for $K_p$ for this equilibrium and state its units. \hfill (3 marks)
\item State and explain, using Le Chatelier's principle, the effect on the equilibrium yield of ammonia of: (i) increasing the pressure; (ii) increasing the temperature. \hfill (4 marks)
\item State the effect, if any, of (i) increasing the pressure and (ii) increasing the temperature on the \emph{value} of $K_p$. Justify your answers. \hfill (3 marks)
\item Explain why an iron catalyst is used, and state its effect on the equilibrium position and on the rate at which equilibrium is attained. \hfill (3 marks)
\item Typical industrial conditions are about $450\ ^{\circ}\mathrm{C}$, $200\ \mathrm{atm}$ and an iron catalyst. Explain why these conditions represent a compromise between yield and rate. \hfill (3 marks)
\item At a certain temperature, an equilibrium mixture contains \ce{N2}, \ce{H2} and \ce{NH3} at partial pressures $20\ \mathrm{atm}$, $60\ \mathrm{atm}$ and $40\ \mathrm{atm}$ respectively. Calculate $K_p$ at this temperature and give its units. \hfill (4 marks)
\end{enumerate}
\hfill (Total: 20 marks)
\end{exercice}

\begin{exercice}[Electrode potentials and electrochemical cells]
Use the following standard electrode potentials at $25\ ^{\circ}\mathrm{C}$:
\begin{center}
\begin{tabular}{|l|c|}
\hline
\rowcolor{popblueL} Half-cell & $E^{\theta}\ /\ \mathrm{V}$ \\
\hline
\ce{Zn^{2+}(aq) + 2e- <=> Zn(s)} & $-0.76$ \\
\ce{Cu^{2+}(aq) + 2e- <=> Cu(s)} & $+0.34$ \\
\ce{Fe^{3+}(aq) + e- <=> Fe^{2+}(aq)} & $+0.77$ \\
\ce{Ag^{+}(aq) + e- <=> Ag(s)} & $+0.80$ \\
\hline
\end{tabular}
\end{center}
\begin{enumerate}
\item Describe, with the aid of a labelled diagram, how the standard electrode potential of the \ce{Cu^{2+}/Cu} half-cell is measured using the standard hydrogen electrode. State the standard conditions of measurement. \hfill (5 marks)
\item Draw a labelled diagram of the Daniell cell, showing the electrodes, the solutions, the salt bridge and the direction of electron flow in the external circuit. Write the conventional cell notation. \hfill (4 marks)
\item Calculate the standard e.m.f. of the Daniell cell and write the overall cell reaction. \hfill (3 marks)
\item Predict, with a calculation to support your answer, whether copper metal can displace silver from a solution of silver nitrate. Write the equation for any reaction that occurs. \hfill (4 marks)
\item State and explain, using the Nernst equation qualitatively, how the electrode potential of the \ce{Cu^{2+}/Cu} electrode changes when the concentration of \ce{Cu^{2+}} ions is decreased. \hfill (2 marks)
\item Give two reasons why a salt bridge containing \ce{KNO3} is preferred to one containing \ce{KCl} when a solution of \ce{AgNO3} is involved. \hfill (2 marks)
\end{enumerate}
\hfill (Total: 20 marks)
\end{exercice}

\begin{exercice}[Acid--base equilibria and titration curves]
Ethanoic acid is a weak acid with $K_a = 1.8\times 10^{-5}\ \mathrm{mol\,dm^{-3}}$ at $25\ ^{\circ}\mathrm{C}$ ($K_w = 1.0\times 10^{-14}\ \mathrm{mol^{2}\,dm^{-6}}$).
\begin{enumerate}
\item Define a weak acid according to the Br\o nsted--Lowry theory and write the equation for the ionisation of ethanoic acid in water. \hfill (2 marks)
\item Calculate the pH of a $0.100\ \mathrm{mol\,dm^{-3}}$ solution of \ce{CH3COOH}, stating any approximation you make. \hfill (4 marks)
\item $25.0\ \mathrm{cm^{3}}$ of this $0.100\ \mathrm{mol\,dm^{-3}}$ \ce{CH3COOH} solution is titrated with $0.100\ \mathrm{mol\,dm^{-3}}$ \ce{NaOH}. Sketch and fully label the titration curve (pH against volume of \ce{NaOH} added), marking: the initial pH, the half-equivalence point (where $\mathrm{pH} = \mathrm{p}K_a$), the equivalence point and the buffer region. \hfill (5 marks)
\item Calculate the volume of \ce{NaOH} needed to reach the equivalence point, and explain why the pH at the equivalence point is greater than $7$. \hfill (3 marks)
\item Choosing between methyl orange (range $3.1$--$4.4$) and phenolphthalein (range $8.3$--$10.0$), state which indicator is suitable for this titration and justify your choice. \hfill (2 marks)
\item Sodium carbonate is titrated with dilute hydrochloric acid. Write the equations of the two successive reactions that occur, and sketch the shape of the titration curve showing the two equivalence points. \hfill (4 marks)
\end{enumerate}
\hfill (Total: 20 marks)
\end{exercice}

\newpage

\coursec{Solutions to the exercises}

\begin{corrige}[Exercise 1 — Multiple choice questions]
\begin{enumerate}
\item \textbf{Answer: (b) decreases the yield of \ce{NH3} and decreases $K_p$.}
\par Justification: The reaction \ce{N2(g) + 3H2(g) <=> 2NH3(g)} is exothermic ($\Delta H^\theta = -92\ \mathrm{kJ\,mol^{-1}}$). According to Le Chatelier's principle, increasing the temperature shifts the equilibrium to the left (endothermic direction), decreasing the yield of \ce{NH3}. The equilibrium constant $K_p$ for an exothermic reaction decreases with increasing temperature (van't Hoff equation: $\frac{d\ln K_p}{dT} = \frac{\Delta H^\theta}{RT^2} < 0$).
\item \textbf{Answer: (d) no units.}
\par Justification: For the reaction \ce{CH3COOH(l) + C2H5OH(l) <=> CH3COOC2H5(l) + H2O(l)}, all species are pure liquids. Their molar concentrations (density/molar mass) are constant at constant temperature and are incorporated into the equilibrium constant. The thermodynamic equilibrium constant $K$ (and the concentration quotient $K_c$ for this homogeneous liquid system) is therefore dimensionless.
\item \textbf{Answer: (d) $12.0$.}
\par Justification: \ce{NaOH} is a strong base, fully dissociated: $[\ce{OH-}] = 0.010\ \mathrm{mol\,dm^{-3}}$. $\mathrm{pOH} = -\log_{10}(0.010) = 2.00$. At $25\ ^\circ\mathrm{C}$, $K_w = 1.0\times 10^{-14}$, so $\mathrm{pH} + \mathrm{pOH} = 14.00$, giving $\mathrm{pH} = 14.00 - 2.00 = 12.00$.
\item \textbf{Answer: (b) the zinc electrode (anode).}
\par Justification: In the cell notation \ce{Zn(s)|Zn^{2+}(aq)||Cu^{2+}(aq)|Cu(s)}, the left-hand electrode is the anode where oxidation occurs: \ce{Zn(s) -> Zn^{2+}(aq) + 2e-}. The right-hand electrode is the cathode where reduction occurs: \ce{Cu^{2+}(aq) + 2e- -> Cu(s)}.
\end{enumerate}
\end{corrige}

\begin{corrige}[Exercise 2 — True or false]
\begin{enumerate}
\item \textbf{False.} At dynamic equilibrium, the forward and reverse reactions continue to occur at equal rates; they have not stopped. The concentrations of reactants and products remain constant because the rates are equal, not because the reactions have ceased.
\item \textbf{False.} A catalyst increases the rate of both the forward and reverse reactions equally by providing an alternative reaction pathway with lower activation energy. It does not change the equilibrium constant or the equilibrium yield; it only helps the system reach equilibrium faster.
\item \textbf{True.} By international convention (IUPAC), the standard hydrogen electrode (SHE) \ce{Pt|H2(g, 1 bar)|H+(aq, 1 mol dm^{-3})} is assigned a standard electrode potential of exactly $0.00\ \mathrm{V}$ at $298\ \mathrm{K}$ ($25\ ^\circ\mathrm{C}$).
\item \textbf{False.} \ce{NH4Cl} is the salt of a weak base (\ce{NH3}, $K_b = 1.8\times 10^{-5}$) and a strong acid (\ce{HCl}). In water, \ce{NH4+} hydrolyses: \ce{NH4+(aq) + H2O(l) <=> NH3(aq) + H3O+(aq)}, producing an acidic solution with $\mathrm{pH} < 7$ at $25\ ^\circ\mathrm{C}$.
\item \textbf{True.} The common ion effect: adding \ce{NaCl} provides a common \ce{Cl-} ion, shifting the equilibrium \ce{AgCl(s) <=> Ag+(aq) + Cl-(aq)} to the left (Le Chatelier's principle), decreasing the solubility of \ce{AgCl} compared to its solubility in pure water.
\end{enumerate}
\end{corrige}

\begin{corrige}[Exercise 3 — Definitions]
\begin{enumerate}
\item \textbf{Dynamic equilibrium:} A state in a reversible reaction occurring in a closed system at constant temperature where the rate of the forward reaction equals the rate of the reverse reaction, so the concentrations (or partial pressures) of reactants and products remain constant over time. \emph{Example:} \ce{N2(g) + 3H2(g) <=> 2NH3(g)} in a closed vessel at constant temperature and pressure.
\item \textbf{Le Chatelier's principle:} If a system at equilibrium is subjected to a change in concentration, temperature, pressure (or volume), the system will shift its equilibrium position to counteract the effect of the change. \emph{Example:} Increasing the total pressure on \ce{N2(g) + 3H2(g) <=> 2NH3(g)} (4 mol gas $\to$ 2 mol gas) shifts the equilibrium to the right, increasing the yield of \ce{NH3}.
\item \textbf{Standard electrode potential, $E^\theta$:} The potential difference (electromotive force) of a half-cell relative to the standard hydrogen electrode (SHE) when all species are in their standard states: temperature $298\ \mathrm{K}$ ($25\ ^\circ\mathrm{C}$), pressure $1\ \mathrm{bar}$ for gases, concentration $1\ \mathrm{mol\,dm^{-3}}$ for solutions. \emph{Example:} $E^\theta(\ce{Cu^{2+}/Cu}) = +0.34\ \mathrm{V}$.
\item \textbf{Brønsted–Lowry acid:} A proton (\ce{H+}) donor. \emph{Example:} \ce{HCl} donates a proton to water: \ce{HCl + H2O -> H3O+ + Cl-}. \textbf{Lewis base:} An electron-pair donor. \emph{Example:} \ce{NH3} donates its lone pair to \ce{H+} to form \ce{NH4+}.
\item \textbf{Buffer solution:} A solution that resists changes in pH when small amounts of acid or base are added, or upon dilution. It typically consists of a weak acid and its conjugate base (acidic buffer, e.g., \ce{CH3COOH/CH3COONa}) or a weak base and its conjugate acid (alkaline buffer, e.g., \ce{NH3/NH4Cl}).
\item \textbf{Solubility product, $K_{\mathrm{sp}}$:} The equilibrium constant for the dissolution of a sparingly soluble ionic solid in water, expressed as the product of the equilibrium concentrations of the ions raised to the power of their stoichiometric coefficients in the balanced dissolution equation. \emph{Example:} For \ce{AgCl(s) <=> Ag+(aq) + Cl-(aq)}, $K_{\mathrm{sp}} = [\ce{Ag+}][\ce{Cl-}]$; for \ce{Mg(OH)2(s) <=> Mg^{2+}(aq) + 2OH-(aq)}, $K_{\mathrm{sp}} = [\ce{Mg^{2+}}][\ce{OH-}]^2$.
\end{enumerate}
\end{corrige}

\begin{corrige}[Exercise 4 — Direct calculations]
\begin{enumerate}
\item \textbf{(i) $0.050\ \mathrm{mol\,dm^{-3}}$ \ce{HCl}:} Strong acid, fully dissociated: $[\ce{H+}] = 0.050\ \mathrm{mol\,dm^{-3}}$. $\mathrm{pH} = -\log_{10}(0.050) = 1.30$.
\par \textbf{(ii) $0.050\ \mathrm{mol\,dm^{-3}}$ \ce{NaOH}:} Strong base, fully dissociated: $[\ce{OH-}] = 0.050\ \mathrm{mol\,dm^{-3}}$. $\mathrm{pOH} = -\log_{10}(0.050) = 1.30$. At $25\ ^\circ\mathrm{C}$, $\mathrm{pH} = 14.00 - \mathrm{pOH} = 12.70$.
\item For the equilibrium \ce{PCl5(g) <=> PCl3(g) + Cl2(g)} at a certain temperature:
\[
K_c = \frac{[\ce{PCl3}][\ce{Cl2}]}{[\ce{PCl5}]} = \frac{(0.40)(0.40)}{0.10} = 1.6
\]
Units: $\frac{(\mathrm{mol\,dm^{-3}})(\mathrm{mol\,dm^{-3}})}{\mathrm{mol\,dm^{-3}}} = \mathrm{mol\,dm^{-3}}$.
\item Daniell cell: \ce{Zn(s)|Zn^{2+}(aq)||Cu^{2+}(aq)|Cu(s)}.
Standard e.m.f. $E^\theta_{\text{cell}} = E^\theta_{\text{cathode}} - E^\theta_{\text{anode}} = E^\theta(\ce{Cu^{2+}/Cu}) - E^\theta(\ce{Zn^{2+}/Zn}) = 0.34 - (-0.76) = +1.10\ \mathrm{V}$.
Positive electrode (cathode): copper electrode (where reduction \ce{Cu^{2+} + 2e- -> Cu} occurs).
\end{enumerate}
\end{corrige}

\begin{corrige}[Exercise 5 — Decomposition of phosphorus pentachloride]
\begin{enumerate}
\item \textbf{ICE table (amounts in moles):}
\begin{center}
\begin{tabular}{|l|ccc|}
\hline
 & \ce{PCl5(g)} & \ce{PCl3(g)} & \ce{Cl2(g)} \\
\hline
Initial & 1.00 & 0 & 0 \\
Change & -0.80 & +0.80 & +0.80 \\
Equilibrium & 0.20 & 0.80 & 0.80 \\
\hline
\end{tabular}
\end{center}
\item Volume of vessel = $2.00\ \mathrm{dm^3}$.
Equilibrium concentrations:
\[
[\ce{PCl5}] = \frac{0.20}{2.00} = 0.10\ \mathrm{mol\,dm^{-3}},\quad
[\ce{PCl3}] = \frac{0.80}{2.00} = 0.40\ \mathrm{mol\,dm^{-3}},\quad
[\ce{Cl2}] = \frac{0.80}{2.00} = 0.40\ \mathrm{mol\,dm^{-3}}.
\]
\item $K_c = \frac{[\ce{PCl3}][\ce{Cl2}]}{[\ce{PCl5}]} = \frac{(0.40)(0.40)}{0.10} = 1.6\ \mathrm{mol\,dm^{-3}}$.
\item Decreasing the volume at constant temperature increases the total pressure. According to Le Chatelier's principle, the equilibrium shifts to the side with fewer moles of gas to counteract the pressure increase. Here, the left side has 1 mole of gas, the right side has 2 moles of gas. Therefore, the equilibrium shifts to the \textbf{left} (towards \ce{PCl5}), decreasing the yield of \ce{PCl3} and \ce{Cl2}. The value of $K_c$ remains unchanged (temperature constant).
\end{enumerate}
\end{corrige}

\begin{corrige}[Exercise 6 — Redox titration of iron(II) with permanganate]
\begin{enumerate}
\item \textbf{Half-equations:}
Oxidation: \ce{Fe^{2+}(aq) -> Fe^{3+}(aq) + e-} \quad ($\times 5$)
Reduction: \ce{MnO4^{-}(aq) + 8H^{+}(aq) + 5e- -> Mn^{2+}(aq) + 4H2O(l)}
\par \textbf{Overall balanced ionic equation:}
\[
\ce{MnO4^{-}(aq) + 5Fe^{2+}(aq) + 8H^{+}(aq) -> Mn^{2+}(aq) + 5Fe^{3+}(aq) + 4H2O(l)}
\]
\item Molar mass of \ce{FeSO4.7H2O} = $56 + 32 + 4\times16 + 7\times(2+16) = 278\ \mathrm{g\,mol^{-1}}$.
Amount of \ce{FeSO4.7H2O} = $\frac{1.52\ \mathrm{g}}{278\ \mathrm{g\,mol^{-1}}} = 5.468\times 10^{-3}\ \mathrm{mol}$.
Amount of \ce{Fe^{2+}} = $5.468\times 10^{-3}\ \mathrm{mol}$.
Stoichiometry: $1\ \mathrm{mol}\ \ce{MnO4-} \equiv 5\ \mathrm{mol}\ \ce{Fe^{2+}}$.
Amount of \ce{MnO4-} required = $\frac{5.468\times 10^{-3}}{5} = 1.0936\times 10^{-3}\ \mathrm{mol}$.
Volume of $0.0200\ \mathrm{mol\,dm^{-3}}\ \ce{KMnO4}$ = $\frac{1.0936\times 10^{-3}\ \mathrm{mol}}{0.0200\ \mathrm{mol\,dm^{-3}}} = 0.05468\ \mathrm{dm^3} = \mathbf{54.7\ \mathrm{cm^3}}$.
\item \textbf{Colour change at end-point:} The intense purple colour of \ce{MnO4-} disappears (solution becomes colourless or very pale pink due to \ce{Mn^{2+}}). The first permanent faint pink colour (due to a tiny excess of \ce{MnO4-}) signals the equivalence point.
\par \textbf{Indicator:} No external indicator is needed; \ce{MnO4-} is intensely coloured and acts as its own indicator (self-indicating).
\end{enumerate}
\end{corrige}

\begin{corrige}[Exercise 7 — pH of a buffer solution]
\begin{enumerate}
\item \textbf{$K_a$ expression for \ce{CH3COOH}:} $K_a = \frac{[\ce{H+}][\ce{CH3COO-}]}{[\ce{CH3COOH}]}$.
\par \textbf{Derivation of Henderson–Hasselbalch equation:}
\[
[\ce{H+}] = K_a \frac{[\ce{CH3COOH}]}{[\ce{CH3COO-}]}
\]
Taking $-\log_{10}$:
\[
\mathrm{pH} = \mathrm{p}K_a + \log_{10}\frac{[\ce{CH3COO-}]}{[\ce{CH3COOH}]}
\]
Since \ce{CH3COONa} is a strong electrolyte, it fully dissociates: $[\ce{CH3COO-}] \approx [\text{salt}]$. The weak acid is only slightly dissociated, so $[\ce{CH3COOH}] \approx [\text{acid}]$.
\item $\mathrm{p}K_a = -\log_{10}(1.8\times 10^{-5}) = 4.7447$.
\[
\mathrm{pH} = 4.7447 + \log_{10}\left(\frac{0.15}{0.10}\right) = 4.7447 + \log_{10}(1.5) = 4.7447 + 0.1761 = \mathbf{4.92}
\]
\item \textbf{Reaction with added \ce{NaOH}:} \ce{CH3COOH(aq) + OH^{-}(aq) -> CH3COO^{-}(aq) + H2O(l)}.
Amount of \ce{NaOH} added = $0.010\ \mathrm{mol}$ in $1.00\ \mathrm{dm^3}$.
New amounts: acid = $0.10 - 0.010 = 0.090\ \mathrm{mol}$; salt = $0.15 + 0.010 = 0.16\ \mathrm{mol}$.
Concentrations (volume $\approx$ constant): $[\text{acid}] = 0.090\ \mathrm{M}$, $[\text{salt}] = 0.16\ \mathrm{M}$.
\[
\mathrm{pH} = 4.7447 + \log_{10}\left(\frac{0.16}{0.090}\right) = 4.7447 + \log_{10}(1.7778) = 4.7447 + 0.2499 = \mathbf{4.99}
\]
\item Adding $0.010\ \mathrm{mol}\ \ce{NaOH}$ to $1.00\ \mathrm{dm^3}$ pure water gives $[\ce{OH-}] = 0.010\ \mathrm{M}$, $\mathrm{pOH}=2.00$, $\mathrm{pH}=12.00$. The buffer pH changes only from $4.92$ to $4.99$ ($\Delta\mathrm{pH}=0.07$), demonstrating its capacity to resist pH change effectively.
\end{enumerate}
\end{corrige}

\begin{corrige}[Exercise 8 — Solubility and common ion effect]
\begin{enumerate}
\item \textbf{\ce{AgCl}:} \ce{AgCl(s) <=> Ag+(aq) + Cl-(aq)}. Let $s$ = molar solubility in $\mathrm{mol\,dm^{-3}}$.
$K_{\mathrm{sp}} = s^2 = 1.8\times 10^{-10} \Rightarrow s = \sqrt{1.8\times 10^{-10}} = \mathbf{1.34\times 10^{-5}\ \mathrm{mol\,dm^{-3}}}$.
\par \textbf{\ce{Mg(OH)2}:} \ce{Mg(OH)2(s) <=> Mg^{2+}(aq) + 2OH-(aq)}. Let $s$ = molar solubility.
$K_{\mathrm{sp}} = [\ce{Mg^{2+}}][\ce{OH-}]^2 = s(2s)^2 = 4s^3 = 1.8\times 10^{-11}$.
$s^3 = 4.5\times 10^{-12} \Rightarrow s = \sqrt[3]{4.5\times 10^{-12}} = \mathbf{1.65\times 10^{-4}\ \mathrm{mol\,dm^{-3}}}$.
\item In $0.10\ \mathrm{M}\ \ce{NaCl}$, $[\ce{Cl-}] \approx 0.10\ \mathrm{M}$ (from complete dissociation of \ce{NaCl}; contribution from \ce{AgCl} negligible).
Let $s'$ = solubility in $0.10\ \mathrm{M}\ \ce{NaCl}$.
$K_{\mathrm{sp}} = [\ce{Ag+}][\ce{Cl-}] = s' \times 0.10 = 1.8\times 10^{-10} \Rightarrow s' = \mathbf{1.8\times 10^{-9}\ \mathrm{mol\,dm^{-3}}}$.
\par \textbf{Explanation:} The common ion \ce{Cl-} from \ce{NaCl} increases $[\ce{Cl-}]$, shifting the equilibrium \ce{AgCl(s) <=> Ag+ + Cl-} to the left (Le Chatelier's principle), decreasing the solubility of \ce{AgCl} by a factor of $\approx 7500$.
\item Mixing equal volumes halves all concentrations:
$[\ce{Ag+}] = \frac{2.0\times 10^{-4}}{2} = 1.0\times 10^{-4}\ \mathrm{M}$.
$[\ce{Cl-}] = \frac{2.0\times 10^{-4}}{2} = 1.0\times 10^{-4}\ \mathrm{M}$.
Ionic product $Q = [\ce{Ag+}][\ce{Cl-}] = (1.0\times 10^{-4})^2 = 1.0\times 10^{-8}$.
Since $Q = 1.0\times 10^{-8} > K_{\mathrm{sp}} = 1.8\times 10^{-10}$, the solution is supersaturated and a precipitate of \ce{AgCl} \textbf{will form} until $Q = K_{\mathrm{sp}}$.
\end{enumerate}
\end{corrige}

\begin{corrige}[Exercise 9 — The Haber process]
\textbf{Mark scheme and examiner expectations:}
\begin{enumerate}
\item \textbf{Expression for $K_p$:} $K_p = \frac{(p_{\ce{NH3}})^2}{(p_{\ce{N2}})(p_{\ce{H2}})^3}$.
\par \textbf{Units:} Pressure$^{-2}$ (e.g., $\mathrm{atm^{-2}}$ or $\mathrm{Pa^{-2}}$). \hfill (3 marks: 1 for expression, 1 for correct powers, 1 for units)
\item \textbf{(i) Increasing pressure:} Equilibrium shifts to the side with fewer moles of gas (4 moles reactants $\to$ 2 moles products). Yield of \ce{NH3} \textbf{increases}. \hfill (2 marks: shift direction + reason)
\par \textbf{(ii) Increasing temperature:} Reaction is exothermic ($\Delta H^\theta = -92\ \mathrm{kJ\,mol^{-1}}$). Equilibrium shifts to the left (endothermic direction). Yield of \ce{NH3} \textbf{decreases}. \hfill (2 marks: shift direction + reason)
\item \textbf{(i) Effect on $K_p$ of increasing pressure:} \textbf{No effect.} $K_p$ depends only on temperature. \hfill (1 mark)
\par \textbf{(ii) Effect on $K_p$ of increasing temperature:} \textbf{Decreases $K_p$.} For an exothermic reaction ($\Delta H^\theta < 0$), the van't Hoff equation shows $\frac{d\ln K_p}{dT} < 0$, so $K_p$ decreases as temperature increases. \hfill (2 marks: direction + justification)
\item \textbf{Iron catalyst:} Used to increase the rate of reaction (lowers activation energy). \textbf{Effect on equilibrium position:} None (catalyst does not change $K_p$ or equilibrium composition). \textbf{Effect on rate:} Increases the rate at which equilibrium is attained, allowing lower temperatures to be used for a given rate. \hfill (3 marks: 1 for purpose, 1 for equilibrium position, 1 for rate)
\item \textbf{Compromise conditions (industrial practice):} 
- High pressure favours yield but is expensive/dangerous; $200\text{--}250\ \mathrm{atm}$ is a practical limit.
- Low temperature favours yield (exothermic) but slows rate; $400\text{--}450\ ^\circ\mathrm{C}$ gives a reasonable rate with acceptable yield ($\approx 15\text{--}20\%$).
- Iron catalyst (with promoters \ce{K2O}, \ce{Al2O3}) allows operation at lower temperature for same rate, improving yield.
\par \emph{Examiner expects:} Mention of yield vs rate trade-off, cost/safety of high pressure, catalyst role. \hfill (3 marks)
\item \textbf{Calculation of $K_p$ at equilibrium:}
$p_{\ce{N2}} = 20\ \mathrm{atm}$, $p_{\ce{H2}} = 60\ \mathrm{atm}$, $p_{\ce{NH3}} = 40\ \mathrm{atm}$.
\[
K_p = \frac{(40)^2}{(20)(60)^3} = \frac{1600}{20 \times 216\,000} = \frac{1600}{4\,320\,000} = \mathbf{3.70\times 10^{-4}\ \mathrm{atm^{-2}}}
\]
\hfill (4 marks: 1 for substitution, 1 for arithmetic, 1 for correct value, 1 for units)
\end{enumerate}
\emph{Total: 20 marks.}
\end{corrige}

\begin{corrige}[Exercise 10 — Electrode potentials and electrochemical cells]
\textbf{Mark scheme and examiner expectations:}
\begin{enumerate}
\item \textbf{Diagram description for measuring $E^\theta(\ce{Cu^{2+}/Cu})$:} 
- SHE: Platinum electrode in $1\ \mathrm{mol\,dm^{-3}}\ \ce{H+}$ (e.g., \ce{H2SO4}) with \ce{H2} gas at $1\ \mathrm{bar}$ bubbling over it.
- \ce{Cu^{2+}/Cu} half-cell: Copper rod in $1\ \mathrm{mol\,dm^{-3}}\ \ce{Cu^{2+}}$ (e.g., \ce{CuSO4}).
- Salt bridge (e.g., saturated \ce{KNO3} in agar gel) connecting the two solutions to complete the circuit and minimise liquid junction potential.
- High-resistance voltmeter (or potentiometer) connecting the two electrodes.
\par \textbf{Standard conditions:} $298\ \mathrm{K}$ ($25\ ^\circ\mathrm{C}$), $1\ \mathrm{bar}$ for gases, $1\ \mathrm{mol\,dm^{-3}}$ for all solutions. \hfill (5 marks: 2 for diagram labels, 2 for description, 1 for conditions)
\item \textbf{Daniell cell diagram:}
- Left half-cell (anode): Zn electrode in \ce{Zn^{2+}} (oxidation: \ce{Zn -> Zn^{2+} + 2e-}).
- Right half-cell (cathode): Cu electrode in \ce{Cu^{2+}} (reduction: \ce{Cu^{2+} + 2e- -> Cu}).
- Salt bridge (e.g., \ce{KNO3}) connecting solutions.
- External circuit with voltmeter; electrons flow from Zn (negative) to Cu (positive).
\par \textbf{Cell notation (IUPAC):} \ce{Zn(s)|Zn^{2+}(aq, 1 M)||Cu^{2+}(aq, 1 M)|Cu(s)}. \hfill (4 marks: 2 for diagram, 1 for electron flow, 1 for notation)
\item \textbf{Standard e.m.f.:} $E^\theta_{\text{cell}} = E^\theta_{\text{cathode}} - E^\theta_{\text{anode}} = E^\theta(\ce{Cu^{2+}/Cu}) - E^\theta(\ce{Zn^{2+}/Zn}) = 0.34 - (-0.76) = \mathbf{+1.10\ \mathrm{V}}$.
\par \textbf{Overall cell reaction:} \ce{Zn(s) + Cu^{2+}(aq) -> Zn^{2+}(aq) + Cu(s)}. \hfill (3 marks: 1 for calculation, 1 for sign, 1 for equation)
\item \textbf{Can Cu displace Ag from \ce{Ag+} solution?} 
Compare standard reduction potentials: $E^\theta(\ce{Ag+/Ag}) = +0.80\ \mathrm{V}$, $E^\theta(\ce{Cu^{2+}/Cu}) = +0.34\ \mathrm{V}$.
For the reaction \ce{Cu(s) + 2Ag+(aq) -> Cu^{2+}(aq) + 2Ag(s)}:
$E^\theta_{\text{cell}} = E^\theta(\ce{Ag+/Ag}) - E^\theta(\ce{Cu^{2+}/Cu}) = 0.80 - 0.34 = +0.46\ \mathrm{V} > 0$.
Since $E^\theta_{\text{cell}} > 0$, the reaction is thermodynamically spontaneous. \textbf{Yes, copper displaces silver.}
\par \textbf{Equation:} \ce{Cu(s) + 2Ag+(aq) -> Cu^{2+}(aq) + 2Ag(s)}. \hfill (4 marks: 1 for comparison, 1 for $E^\theta_{\text{cell}}$ calculation, 1 for conclusion, 1 for equation)
\item \textbf{Effect of decreasing $[\ce{Cu^{2+}}]$ on $E(\ce{Cu^{2+}/Cu})$:}
Nernst equation at $298\ \mathrm{K}$: $E = E^\theta - \frac{0.059}{2}\log\frac{1}{[\ce{Cu^{2+}}]} = E^\theta + \frac{0.059}{2}\log[\ce{Cu^{2+}}]$.
Decreasing $[\ce{Cu^{2+}}]$ makes $\log[\ce{Cu^{2+}}]$ more negative, so $E$ becomes \textbf{less positive (decreases)}. \hfill (2 marks: 1 for direction, 1 for explanation)
\item \textbf{Reasons for using \ce{KNO3} (not \ce{KCl}) as salt bridge with \ce{AgNO3} half-cell:}
1. \ce{Cl-} from \ce{KCl} would react with \ce{Ag+} in the half-cell or salt bridge to form insoluble \ce{AgCl(s)}, blocking the bridge and altering the electrode potential.
2. \ce{K+} and \ce{NO3-} have very similar ionic mobilities, minimising the liquid junction potential at the salt bridge/solution interfaces. \hfill (2 marks: 1 each)
\end{enumerate}
\emph{Total: 20 marks.}
\end{corrige}

\begin{corrige}[Exercise 11 — Acid–base equilibria and titration curves]
\textbf{Mark scheme and examiner expectations:}
\begin{enumerate}
\item \textbf{Weak acid (Brønsted–Lowry):} A substance that partially donates protons to water, establishing an equilibrium with its conjugate base. \hfill (1 mark)
\par \textbf{Ionisation equation:} \ce{CH3COOH(aq) + H2O(l) <=> CH3COO-(aq) + H3O+(aq)} (or simplified: \ce{CH3COOH(aq) <=> H+(aq) + CH3COO-(aq)}). \hfill (1 mark)
\item \textbf{Calculation of pH of $0.100\ \mathrm{M}\ \ce{CH3COOH}$ ($K_a = 1.8\times 10^{-5}$):}
Approximation: $[\ce{H+}] \ll c$, so $c - [\ce{H+}] \approx c$.
$K_a = \frac{[\ce{H+}]^2}{c} \Rightarrow [\ce{H+}] = \sqrt{K_a c} = \sqrt{(1.8\times 10^{-5})(0.100)} = \sqrt{1.8\times 10^{-6}} = 1.34\times 10^{-3}\ \mathrm{mol\,dm^{-3}}$.
$\mathrm{pH} = -\log_{10}(1.34\times 10^{-3}) = \mathbf{2.87}$.
\par \textbf{Justification of approximation:} $\frac{[\ce{H+}]}{c} = \frac{1.34\times 10^{-3}}{0.100} = 1.34\% < 5\%$, so approximation is valid. \hfill (4 marks: 1 for approximation statement, 1 for formula, 1 for $[\ce{H+}]$, 1 for pH)
\item \textbf{Titration curve sketch (\ce{25.0 cm^3} $0.100\ \mathrm{M}\ \ce{CH3COOH}$ vs $0.100\ \mathrm{M}\ \ce{NaOH}$):} 
- Axes: pH (vertical, 0--14) vs Volume \ce{NaOH} added (horizontal, 0--50 cm3).
- Initial point ($V=0$): $\mathrm{pH} = 2.87$.
- Buffer region: gradual rise.
- Half-equivalence point: $V = 12.5\ \mathrm{cm^3}$, $\mathrm{pH} = \mathrm{p}K_a = 4.74$.
- Equivalence point: $V = 25.0\ \mathrm{cm^3}$, $\mathrm{pH} \approx 8.7$ (vertical jump $\approx$ pH 7--10).
- After equivalence: gradual rise approaching pH of excess \ce{NaOH} (e.g., pH $\approx$ 12.6 at $V=50\ \mathrm{cm^3}$).
\par \emph{Examiner expects:} Correct S-shape with asymmetric vertical jump, labelled axes, initial pH, half-equivalence point, equivalence point volume and pH. \hfill (5 marks)
\item \textbf{Equivalence volume:} Moles \ce{CH3COOH} = $0.0250\ \mathrm{dm^3} \times 0.100\ \mathrm{mol\,dm^{-3}} = 0.00250\ \mathrm{mol}$. $1:1$ stoichiometry $\Rightarrow$ moles \ce{NaOH} = $0.00250\ \mathrm{mol}$. Volume = $\frac{0.00250\ \mathrm{mol}}{0.100\ \mathrm{mol\,dm^{-3}}} = 0.0250\ \mathrm{dm^3} = \mathbf{25.0\ \mathrm{cm^3}}$. \hfill (1 mark)
\par \textbf{pH > 7 at equivalence:} Salt formed is \ce{CH3COONa}, the conjugate base of a weak acid. \ce{CH3COO-} hydrolyses: \ce{CH3COO-(aq) + H2O(l) <=> CH3COOH(aq) + OH-(aq)}, making solution alkaline. \hfill (2 marks)
\item \textbf{Suitable indicator:} Phenolphthalein (transition range 8.3--10.0, colourless to pink). \textbf{Justification:} The pH jump at equivalence (approx 7--10) falls entirely within phenolphthalein's range. Methyl orange (3.1--4.4) changes colour in the buffer region, well before equivalence, giving a large titration error. \hfill (2 marks)
\item \textbf{Sodium carbonate (\ce{Na2CO3}) titration with \ce{HCl} (diprotic base):}
\ce{CO3^{2-}(aq) + H+(aq) -> HCO3-(aq)} \quad (1st equivalence point, $\mathrm{pH} \approx 8.3$)
\ce{HCO3-(aq) + H+(aq) -> H2CO3(aq) -> CO2(g) + H2O(l)} \quad (2nd equivalence point, $\mathrm{pH} \approx 3.7$)
\par \textbf{Curve sketch:} Two distinct steps (two buffer regions). First vertical jump at $V_1$ (1st eq. pt., pH $\approx$ 8.3), second vertical jump at $V_2 = 2V_1$ (2nd eq. pt., pH $\approx$ 3.7). \hfill (4 marks: 2 for equations, 2 for sketch with two equivalence points labelled)
\end{enumerate}
\emph{Total: 20 marks.}
\end{corrige}

\newpage

\coursec{Summary sheet}

\begin{popfigure}
\begin{center}\tcbox[colback=popblue,colframe=popblue,coltext=white,arc=9pt,boxsep=4pt,left=6pt,right=6pt]{\bfseries\small Chemical Equilibria}\end{center}
\vspace{-2pt}
\begin{tcbraster}[raster columns=3,raster equal height=rows,raster column skip=6pt,raster row skip=6pt,enhanced,blankest]
\begin{tcolorbox}[enhanced,colback=poporange,colframe=poporange,coltext=white,boxrule=0.9pt,arc=3pt,left=4pt,right=4pt,top=3pt,bottom=3pt,boxsep=1pt,halign=left]\bfseries\scriptsize Dynamic equilibrium \& Le Chatelier\end{tcolorbox}
\begin{tcolorbox}[enhanced,colback=popgreen,colframe=popgreen,coltext=white,boxrule=0.9pt,arc=3pt,left=4pt,right=4pt,top=3pt,bottom=3pt,boxsep=1pt,halign=left]\bfseries\scriptsize Equilibrium law: $K_c$, $K_p$\end{tcolorbox}
\begin{tcolorbox}[enhanced,colback=poppurple,colframe=poppurple,coltext=white,boxrule=0.9pt,arc=3pt,left=4pt,right=4pt,top=3pt,bottom=3pt,boxsep=1pt,halign=left]\bfseries\scriptsize Heterogeneous equilibria; $K$ vs $T$\end{tcolorbox}
\begin{tcolorbox}[enhanced,colback=poppink,colframe=poppink,coltext=white,boxrule=0.9pt,arc=3pt,left=4pt,right=4pt,top=3pt,bottom=3pt,boxsep=1pt,halign=left]\bfseries\scriptsize Redox \& electrode potentials $E^\theta$\end{tcolorbox}
\begin{tcolorbox}[enhanced,colback=popgold,colframe=popgold,coltext=white,boxrule=0.9pt,arc=3pt,left=4pt,right=4pt,top=3pt,bottom=3pt,boxsep=1pt,halign=left]\bfseries\scriptsize Acids, bases, pH, $K_w$\end{tcolorbox}
\begin{tcolorbox}[enhanced,colback=popmuted,colframe=popmuted,coltext=white,boxrule=0.9pt,arc=3pt,left=4pt,right=4pt,top=3pt,bottom=3pt,boxsep=1pt,halign=left]\bfseries\scriptsize Indicators, titration curves, buffers\end{tcolorbox}
\begin{tcolorbox}[enhanced,colback=popdark,colframe=popdark,coltext=white,boxrule=0.9pt,arc=3pt,left=4pt,right=4pt,top=3pt,bottom=3pt,boxsep=1pt,halign=left]\bfseries\scriptsize Hydrolysis, $K_{sp}$, common ion\end{tcolorbox}
\end{tcbraster}

\legende{Mind map of the chapter \emph{Equilibria}: seven interconnected themes, from gaseous and heterogeneous equilibria through redox systems to acid--base equilibria in aqueous solution.}
\end{popfigure}

\begin{bilan}
\textbf{Key definitions}
\begin{itemize}
\item \cle{Reversible reaction}: a reaction that can proceed in both the forward and the reverse directions, denoted by $\rightleftharpoons$; it never goes to completion.
\item \cle{Dynamic equilibrium}: the state reached in a \emph{closed} system when the forward and reverse reactions proceed at \emph{equal rates}; macroscopic properties (concentration, colour, pressure) are constant, but particles keep reacting in both directions.
\item \cle{Le Chatelier's principle}: when a constraint (change of concentration, pressure or temperature) is imposed on a system at equilibrium, the system adjusts itself so as to oppose (annul the effect of) the constraint.
\item \cle{Equilibrium constant}: for $\ce{aA + bB <=> cC + dD}$,
\[ K_c=\dfrac{[\ce{C}]^c[\ce{D}]^d}{[\ce{A}]^a[\ce{B}]^b} \qquad\text{and}\qquad K_p=\dfrac{p_{\ce{C}}^{c}\,p_{\ce{D}}^{d}}{p_{\ce{A}}^{a}\,p_{\ce{B}}^{b}}, \]
related by $K_p=K_c(RT)^{\Delta n}$, where $\Delta n=(c+d)-(a+b)$ is the change in the number of moles of \emph{gas}.
\item \cle{Standard electrode potential} $E^\theta$: the emf of a half-cell measured against the standard hydrogen electrode (SHE) under standard conditions: temperature $\SI{298}{K}$, ion concentration $\SI{1}{mol.dm^{-3}}$, gas pressure $\SI{1}{atm}$ ($\SI{101}{kPa}$).
\item \cle{pH}: $\mathrm{pH}=-\log_{10}[\ce{H3O+}]$; the ionic product of water is $K_w=[\ce{H3O+}][\ce{OH-}]=1.0\times10^{-14}\ \mathrm{mol^2\,dm^{-6}}$ at $\SI{298}{K}$, so that $\mathrm{pH}+\mathrm{pOH}=14$ at this temperature.
\item \cle{Buffer solution}: a solution that resists a change of pH on addition of small amounts of acid or base; it consists of a weak acid and one of its salts (acidic buffer) or a weak base and one of its salts (alkaline buffer).
\item \cle{Solubility product} $K_{sp}$: the product of the concentrations of the ions of a sparingly soluble salt, each raised to its stoichiometric power, in a \emph{saturated} solution at a given temperature.
\end{itemize}

\textbf{Laws and formulas to know}
\begin{itemize}
\item Reaction quotient $Q$ (same expression as $K$, at any instant): if $Q<K$ the reaction proceeds in the forward direction; if $Q>K$ it proceeds in the reverse direction; if $Q=K$ the system is at equilibrium.
\item Effect of temperature on $K$: $K$ increases with temperature for an \emph{endothermic} forward reaction and decreases for an \emph{exothermic} one. Quantitatively (van 't Hoff equation):
\[ \ln\dfrac{K_2}{K_1}=-\dfrac{\Delta H^\theta}{R}\left(\dfrac{1}{T_2}-\dfrac{1}{T_1}\right). \]
A catalyst changes \emph{neither} $K$ \emph{nor} the position of equilibrium; it only increases the rate at which equilibrium is attained.
\item Cell emf: $E^\theta_{\text{cell}}=E^\theta_{\text{cathode}}-E^\theta_{\text{anode}}$ (right-hand electrode minus left-hand electrode in the conventional cell diagram); a reaction is feasible (spontaneous) under standard conditions if $E^\theta_{\text{cell}}>0$.
\item Strong acid (concentration $c$, fully ionised): $\mathrm{pH}=-\log_{10} c$. Strong base at $\SI{298}{K}$: $\mathrm{pH}=14+\log_{10} c$.
\item Weak acid: $K_a=\dfrac{[\ce{H3O+}][\ce{A-}]}{[\ce{HA}]}=\dfrac{[\ce{H3O+}]^2}{c-[\ce{H3O+}]}\approx\dfrac{[\ce{H3O+}]^2}{c}$, giving $\mathrm{pH}=\tfrac12\left(\mathrm{p}K_a-\log_{10} c\right)$. Weak base analogously with $K_b$, and $\mathrm{p}K_a+\mathrm{p}K_b=14$ for a conjugate pair at $\SI{298}{K}$.
\item Acidic buffer (Henderson--Hasselbalch): $\mathrm{pH}=\mathrm{p}K_a+\log_{10}\dfrac{[\text{salt}]}{[\text{acid}]}$; alkaline buffer: $\mathrm{pOH}=\mathrm{p}K_b+\log_{10}\dfrac{[\text{salt}]}{[\text{base}]}$.
\item Salt hydrolysis: a salt of a weak acid and a strong base gives $\mathrm{pH}>7$; a salt of a weak base and a strong acid gives $\mathrm{pH}<7$; a salt of a strong acid and a strong base gives $\mathrm{pH}=7$ (at $\SI{298}{K}$).
\item $K_{sp}$ for $\ce{A_mB_n}$ of molar solubility $s$: $K_{sp}=(ms)^m(ns)^n=m^m n^n\, s^{m+n}$. The common ion effect \emph{decreases} the solubility of a sparingly soluble salt.
\end{itemize}

\textbf{Methods}
\begin{itemize}
\item $K_c$/$K_p$ calculations: write the balanced equation with state symbols; set up an ICE table (Initial moles, Change, Equilibrium moles); convert to concentrations (divide by the volume $V$) or to partial pressures ($p_i=x_i\,P_{\text{total}}$, where $x_i$ is the mole fraction); substitute into the expression; state the units from $\Delta n$.
\item Balancing redox equations (ion--electron method): split into oxidation and reduction half-equations; balance atoms other than O and H; balance O with \ce{H2O} and H with \ce{H+}; balance charge with electrons; multiply so that electrons cancel; add the half-equations and simplify.
\item Titration curves: identify the four regions (initial pH, buffer region where relevant, equivalence point, excess of titrant); the vertical jump brackets the equivalence point; choose an indicator whose colour-change range falls \emph{within} the vertical jump.
\end{itemize}

\textbf{Common mistakes}
\begin{itemize}
\item Including pure solids or pure liquids in $K_c$/$K_p$ expressions: their ``concentration'' is constant and is incorporated into $K$.
\item Forgetting the units of $K_c$/$K_p$, which depend on $\Delta n$; $K$ is unitless only when $\Delta n=0$.
\item Claiming that a catalyst shifts the position of equilibrium: it affects only the \emph{rate} of attainment, never $K$ nor the position.
\item Writing $E^\theta_{\text{cell}}=E^\theta_{\text{anode}}-E^\theta_{\text{cathode}}$ (wrong sign order), or multiplying $E^\theta$ values by stoichiometric coefficients ($E^\theta$ is an intensive quantity).
\item Using $\mathrm{pH}=-\log_{10} c$ for a \emph{weak} acid, or forgetting that dilution of a weak acid increases its degree of ionisation (Ostwald's dilution law).
\item Assuming the equivalence point of a weak acid--strong base titration is at pH 7: it lies \emph{above} 7 because the anion of the weak acid is hydrolysed.
\item Confusing solubility $s$ (in $\mathrm{mol\,dm^{-3}}$) with $K_{sp}$, and forgetting to raise each ion concentration to its stoichiometric power.
\end{itemize}

\textbf{GCE tips}
\begin{itemize}
\item Always write the \emph{balanced equation with state symbols} before any equilibrium calculation; examiners award marks for a clearly laid-out ICE table.
\item Quote Le Chatelier's principle precisely: state the imposed change, the direction of the shift, and \emph{why} (which side has fewer moles of gas; which direction is endothermic).
\item In electrode-potential questions, state the standard conditions, draw the cell diagram with the salt bridge, and remember that the half-cell with the more negative $E^\theta$ is the anode (oxidation).
\item For buffers and $K_{sp}$, state the approximation used (e.g.\ $[\ce{H3O+}]\ll c$) and check that it is valid.
\item Give pH to 2 decimal places, and $K$ values with sensible significant figures and correct units.
\end{itemize}
\end{bilan}

\findecours

\end{document}
